We defined the Wiener process as a Gaussian process as follows:
\begin{defn}{Wiener Process}{}
    The \textit{Wiener process} $(W_t, t \geq 0)$ is a zero-mean Gaussian process with continuous sample paths and covariance function
    \begin{equation}
    \gamma_{s,t} = \min\{s, t\} =: s \wedge t, \quad s, t \geq 0. \notag
    \label{eq:WienerProcessCovarianceFunction}
\end{equation}
\end{defn}

A Wiener process is a \Levy \; process with \textit{continuous} sample paths, as shown in the following result, which can also serve as an equivalent definition of the Wiener process.

\begin{thm}{Wiener Process as a \Levy \; Process}{}
$W = (W_t)_{t \geq 0}$ is a Wiener process if and only if the following three properties hold:
\begin{enumerate}
    \item \textit{(Independent increments):} For any $t_1 < t_2 \leq t_3 < t_4$, the increments $W_{t_4} - W_{t_3}$ and $W_{t_2} - W_{t_1}$ are independent.
    \item \textit{(Gaussian stationarity):} For all $t, u \geq 0$, $W_{t+u} - W_t \sim \mathcal{N}(0, u)$.
    \item \textit{(Continuity of paths):} $W$ has continuous paths, with $W_0 = 0$.
\end{enumerate}
\end{thm}

The following algorithm, \cref{alg:SimulatingWienerProcess} uses the same ideas as in the previous proof to simulate the Wiener process at specific time points $t_1,\dots,t_n$

\begin{algorithm}[H]
    \caption{Simulating a Wiener Process}
    \label{alg:SimulatingWienerProcess}
    \begin{algorithmic}[1]
        \State Let $0 = t_0 < t_1 < \dots < t_n$ be the times at which the process needs to be simulated.
        \State Simulate $Z_1, \dots, Z_n \iidsim \mathcal{N}(0, 1)$.
        \State Return $W_0 \coloneqq 0$ and $W_{t_k} \coloneqq \sum_{i=1}^k Z_i \sqrt{t_i - t_{i-1}}, \quad k = 1, \dots, n$.
    \end{algorithmic}
\end{algorithm}

\begin{figure}[H]
    \centering
    %% Creator: Matplotlib, PGF backend
%%
%% To include the figure in your LaTeX document, write
%%   \input{<filename>.pgf}
%%
%% Make sure the required packages are loaded in your preamble
%%   \usepackage{pgf}
%%
%% Also ensure that all the required font packages are loaded; for instance,
%% the lmodern package is sometimes necessary when using math font.
%%   \usepackage{lmodern}
%%
%% Figures using additional raster images can only be included by \input if
%% they are in the same directory as the main LaTeX file. For loading figures
%% from other directories you can use the `import` package
%%   \usepackage{import}
%%
%% and then include the figures with
%%   \import{<path to file>}{<filename>.pgf}
%%
%% Matplotlib used the following preamble
%%   \def\mathdefault#1{#1}
%%   \everymath=\expandafter{\the\everymath\displaystyle}
%%   \IfFileExists{scrextend.sty}{
%%     \usepackage[fontsize=10.000000pt]{scrextend}
%%   }{
%%     \renewcommand{\normalsize}{\fontsize{10.000000}{12.000000}\selectfont}
%%     \normalsize
%%   }
%%   
%%   \makeatletter\@ifpackageloaded{underscore}{}{\usepackage[strings]{underscore}}\makeatother
%%
\begingroup%
\makeatletter%
\begin{pgfpicture}%
\pgfpathrectangle{\pgfpointorigin}{\pgfqpoint{5.739353in}{4.455080in}}%
\pgfusepath{use as bounding box, clip}%
\begin{pgfscope}%
\pgfsetbuttcap%
\pgfsetmiterjoin%
\definecolor{currentfill}{rgb}{1.000000,1.000000,1.000000}%
\pgfsetfillcolor{currentfill}%
\pgfsetlinewidth{0.000000pt}%
\definecolor{currentstroke}{rgb}{1.000000,1.000000,1.000000}%
\pgfsetstrokecolor{currentstroke}%
\pgfsetdash{}{0pt}%
\pgfpathmoveto{\pgfqpoint{0.000000in}{0.000000in}}%
\pgfpathlineto{\pgfqpoint{5.739353in}{0.000000in}}%
\pgfpathlineto{\pgfqpoint{5.739353in}{4.455080in}}%
\pgfpathlineto{\pgfqpoint{0.000000in}{4.455080in}}%
\pgfpathlineto{\pgfqpoint{0.000000in}{0.000000in}}%
\pgfpathclose%
\pgfusepath{fill}%
\end{pgfscope}%
\begin{pgfscope}%
\pgfsetbuttcap%
\pgfsetmiterjoin%
\definecolor{currentfill}{rgb}{1.000000,1.000000,1.000000}%
\pgfsetfillcolor{currentfill}%
\pgfsetlinewidth{0.000000pt}%
\definecolor{currentstroke}{rgb}{0.000000,0.000000,0.000000}%
\pgfsetstrokecolor{currentstroke}%
\pgfsetstrokeopacity{0.000000}%
\pgfsetdash{}{0pt}%
\pgfpathmoveto{\pgfqpoint{0.563118in}{0.401080in}}%
\pgfpathlineto{\pgfqpoint{5.600618in}{0.401080in}}%
\pgfpathlineto{\pgfqpoint{5.600618in}{4.405080in}}%
\pgfpathlineto{\pgfqpoint{0.563118in}{4.405080in}}%
\pgfpathlineto{\pgfqpoint{0.563118in}{0.401080in}}%
\pgfpathclose%
\pgfusepath{fill}%
\end{pgfscope}%
\begin{pgfscope}%
\pgfsetbuttcap%
\pgfsetroundjoin%
\definecolor{currentfill}{rgb}{0.000000,0.000000,0.000000}%
\pgfsetfillcolor{currentfill}%
\pgfsetlinewidth{0.501875pt}%
\definecolor{currentstroke}{rgb}{0.000000,0.000000,0.000000}%
\pgfsetstrokecolor{currentstroke}%
\pgfsetdash{}{0pt}%
\pgfsys@defobject{currentmarker}{\pgfqpoint{0.000000in}{0.000000in}}{\pgfqpoint{0.000000in}{0.041667in}}{%
\pgfpathmoveto{\pgfqpoint{0.000000in}{0.000000in}}%
\pgfpathlineto{\pgfqpoint{0.000000in}{0.041667in}}%
\pgfusepath{stroke,fill}%
}%
\begin{pgfscope}%
\pgfsys@transformshift{0.563118in}{0.401080in}%
\pgfsys@useobject{currentmarker}{}%
\end{pgfscope}%
\end{pgfscope}%
\begin{pgfscope}%
\pgfsetbuttcap%
\pgfsetroundjoin%
\definecolor{currentfill}{rgb}{0.000000,0.000000,0.000000}%
\pgfsetfillcolor{currentfill}%
\pgfsetlinewidth{0.501875pt}%
\definecolor{currentstroke}{rgb}{0.000000,0.000000,0.000000}%
\pgfsetstrokecolor{currentstroke}%
\pgfsetdash{}{0pt}%
\pgfsys@defobject{currentmarker}{\pgfqpoint{0.000000in}{-0.041667in}}{\pgfqpoint{0.000000in}{0.000000in}}{%
\pgfpathmoveto{\pgfqpoint{0.000000in}{0.000000in}}%
\pgfpathlineto{\pgfqpoint{0.000000in}{-0.041667in}}%
\pgfusepath{stroke,fill}%
}%
\begin{pgfscope}%
\pgfsys@transformshift{0.563118in}{4.405080in}%
\pgfsys@useobject{currentmarker}{}%
\end{pgfscope}%
\end{pgfscope}%
\begin{pgfscope}%
\definecolor{textcolor}{rgb}{0.000000,0.000000,0.000000}%
\pgfsetstrokecolor{textcolor}%
\pgfsetfillcolor{textcolor}%
\pgftext[x=0.563118in,y=0.352469in,,top]{\color{textcolor}{\rmfamily\fontsize{10.000000}{12.000000}\selectfont\catcode`\^=\active\def^{\ifmmode\sp\else\^{}\fi}\catcode`\%=\active\def%{\%}$\mathdefault{0.0}$}}%
\end{pgfscope}%
\begin{pgfscope}%
\pgfsetbuttcap%
\pgfsetroundjoin%
\definecolor{currentfill}{rgb}{0.000000,0.000000,0.000000}%
\pgfsetfillcolor{currentfill}%
\pgfsetlinewidth{0.501875pt}%
\definecolor{currentstroke}{rgb}{0.000000,0.000000,0.000000}%
\pgfsetstrokecolor{currentstroke}%
\pgfsetdash{}{0pt}%
\pgfsys@defobject{currentmarker}{\pgfqpoint{0.000000in}{0.000000in}}{\pgfqpoint{0.000000in}{0.041667in}}{%
\pgfpathmoveto{\pgfqpoint{0.000000in}{0.000000in}}%
\pgfpathlineto{\pgfqpoint{0.000000in}{0.041667in}}%
\pgfusepath{stroke,fill}%
}%
\begin{pgfscope}%
\pgfsys@transformshift{1.570618in}{0.401080in}%
\pgfsys@useobject{currentmarker}{}%
\end{pgfscope}%
\end{pgfscope}%
\begin{pgfscope}%
\pgfsetbuttcap%
\pgfsetroundjoin%
\definecolor{currentfill}{rgb}{0.000000,0.000000,0.000000}%
\pgfsetfillcolor{currentfill}%
\pgfsetlinewidth{0.501875pt}%
\definecolor{currentstroke}{rgb}{0.000000,0.000000,0.000000}%
\pgfsetstrokecolor{currentstroke}%
\pgfsetdash{}{0pt}%
\pgfsys@defobject{currentmarker}{\pgfqpoint{0.000000in}{-0.041667in}}{\pgfqpoint{0.000000in}{0.000000in}}{%
\pgfpathmoveto{\pgfqpoint{0.000000in}{0.000000in}}%
\pgfpathlineto{\pgfqpoint{0.000000in}{-0.041667in}}%
\pgfusepath{stroke,fill}%
}%
\begin{pgfscope}%
\pgfsys@transformshift{1.570618in}{4.405080in}%
\pgfsys@useobject{currentmarker}{}%
\end{pgfscope}%
\end{pgfscope}%
\begin{pgfscope}%
\definecolor{textcolor}{rgb}{0.000000,0.000000,0.000000}%
\pgfsetstrokecolor{textcolor}%
\pgfsetfillcolor{textcolor}%
\pgftext[x=1.570618in,y=0.352469in,,top]{\color{textcolor}{\rmfamily\fontsize{10.000000}{12.000000}\selectfont\catcode`\^=\active\def^{\ifmmode\sp\else\^{}\fi}\catcode`\%=\active\def%{\%}$\mathdefault{0.2}$}}%
\end{pgfscope}%
\begin{pgfscope}%
\pgfsetbuttcap%
\pgfsetroundjoin%
\definecolor{currentfill}{rgb}{0.000000,0.000000,0.000000}%
\pgfsetfillcolor{currentfill}%
\pgfsetlinewidth{0.501875pt}%
\definecolor{currentstroke}{rgb}{0.000000,0.000000,0.000000}%
\pgfsetstrokecolor{currentstroke}%
\pgfsetdash{}{0pt}%
\pgfsys@defobject{currentmarker}{\pgfqpoint{0.000000in}{0.000000in}}{\pgfqpoint{0.000000in}{0.041667in}}{%
\pgfpathmoveto{\pgfqpoint{0.000000in}{0.000000in}}%
\pgfpathlineto{\pgfqpoint{0.000000in}{0.041667in}}%
\pgfusepath{stroke,fill}%
}%
\begin{pgfscope}%
\pgfsys@transformshift{2.578118in}{0.401080in}%
\pgfsys@useobject{currentmarker}{}%
\end{pgfscope}%
\end{pgfscope}%
\begin{pgfscope}%
\pgfsetbuttcap%
\pgfsetroundjoin%
\definecolor{currentfill}{rgb}{0.000000,0.000000,0.000000}%
\pgfsetfillcolor{currentfill}%
\pgfsetlinewidth{0.501875pt}%
\definecolor{currentstroke}{rgb}{0.000000,0.000000,0.000000}%
\pgfsetstrokecolor{currentstroke}%
\pgfsetdash{}{0pt}%
\pgfsys@defobject{currentmarker}{\pgfqpoint{0.000000in}{-0.041667in}}{\pgfqpoint{0.000000in}{0.000000in}}{%
\pgfpathmoveto{\pgfqpoint{0.000000in}{0.000000in}}%
\pgfpathlineto{\pgfqpoint{0.000000in}{-0.041667in}}%
\pgfusepath{stroke,fill}%
}%
\begin{pgfscope}%
\pgfsys@transformshift{2.578118in}{4.405080in}%
\pgfsys@useobject{currentmarker}{}%
\end{pgfscope}%
\end{pgfscope}%
\begin{pgfscope}%
\definecolor{textcolor}{rgb}{0.000000,0.000000,0.000000}%
\pgfsetstrokecolor{textcolor}%
\pgfsetfillcolor{textcolor}%
\pgftext[x=2.578118in,y=0.352469in,,top]{\color{textcolor}{\rmfamily\fontsize{10.000000}{12.000000}\selectfont\catcode`\^=\active\def^{\ifmmode\sp\else\^{}\fi}\catcode`\%=\active\def%{\%}$\mathdefault{0.4}$}}%
\end{pgfscope}%
\begin{pgfscope}%
\pgfsetbuttcap%
\pgfsetroundjoin%
\definecolor{currentfill}{rgb}{0.000000,0.000000,0.000000}%
\pgfsetfillcolor{currentfill}%
\pgfsetlinewidth{0.501875pt}%
\definecolor{currentstroke}{rgb}{0.000000,0.000000,0.000000}%
\pgfsetstrokecolor{currentstroke}%
\pgfsetdash{}{0pt}%
\pgfsys@defobject{currentmarker}{\pgfqpoint{0.000000in}{0.000000in}}{\pgfqpoint{0.000000in}{0.041667in}}{%
\pgfpathmoveto{\pgfqpoint{0.000000in}{0.000000in}}%
\pgfpathlineto{\pgfqpoint{0.000000in}{0.041667in}}%
\pgfusepath{stroke,fill}%
}%
\begin{pgfscope}%
\pgfsys@transformshift{3.585618in}{0.401080in}%
\pgfsys@useobject{currentmarker}{}%
\end{pgfscope}%
\end{pgfscope}%
\begin{pgfscope}%
\pgfsetbuttcap%
\pgfsetroundjoin%
\definecolor{currentfill}{rgb}{0.000000,0.000000,0.000000}%
\pgfsetfillcolor{currentfill}%
\pgfsetlinewidth{0.501875pt}%
\definecolor{currentstroke}{rgb}{0.000000,0.000000,0.000000}%
\pgfsetstrokecolor{currentstroke}%
\pgfsetdash{}{0pt}%
\pgfsys@defobject{currentmarker}{\pgfqpoint{0.000000in}{-0.041667in}}{\pgfqpoint{0.000000in}{0.000000in}}{%
\pgfpathmoveto{\pgfqpoint{0.000000in}{0.000000in}}%
\pgfpathlineto{\pgfqpoint{0.000000in}{-0.041667in}}%
\pgfusepath{stroke,fill}%
}%
\begin{pgfscope}%
\pgfsys@transformshift{3.585618in}{4.405080in}%
\pgfsys@useobject{currentmarker}{}%
\end{pgfscope}%
\end{pgfscope}%
\begin{pgfscope}%
\definecolor{textcolor}{rgb}{0.000000,0.000000,0.000000}%
\pgfsetstrokecolor{textcolor}%
\pgfsetfillcolor{textcolor}%
\pgftext[x=3.585618in,y=0.352469in,,top]{\color{textcolor}{\rmfamily\fontsize{10.000000}{12.000000}\selectfont\catcode`\^=\active\def^{\ifmmode\sp\else\^{}\fi}\catcode`\%=\active\def%{\%}$\mathdefault{0.6}$}}%
\end{pgfscope}%
\begin{pgfscope}%
\pgfsetbuttcap%
\pgfsetroundjoin%
\definecolor{currentfill}{rgb}{0.000000,0.000000,0.000000}%
\pgfsetfillcolor{currentfill}%
\pgfsetlinewidth{0.501875pt}%
\definecolor{currentstroke}{rgb}{0.000000,0.000000,0.000000}%
\pgfsetstrokecolor{currentstroke}%
\pgfsetdash{}{0pt}%
\pgfsys@defobject{currentmarker}{\pgfqpoint{0.000000in}{0.000000in}}{\pgfqpoint{0.000000in}{0.041667in}}{%
\pgfpathmoveto{\pgfqpoint{0.000000in}{0.000000in}}%
\pgfpathlineto{\pgfqpoint{0.000000in}{0.041667in}}%
\pgfusepath{stroke,fill}%
}%
\begin{pgfscope}%
\pgfsys@transformshift{4.593118in}{0.401080in}%
\pgfsys@useobject{currentmarker}{}%
\end{pgfscope}%
\end{pgfscope}%
\begin{pgfscope}%
\pgfsetbuttcap%
\pgfsetroundjoin%
\definecolor{currentfill}{rgb}{0.000000,0.000000,0.000000}%
\pgfsetfillcolor{currentfill}%
\pgfsetlinewidth{0.501875pt}%
\definecolor{currentstroke}{rgb}{0.000000,0.000000,0.000000}%
\pgfsetstrokecolor{currentstroke}%
\pgfsetdash{}{0pt}%
\pgfsys@defobject{currentmarker}{\pgfqpoint{0.000000in}{-0.041667in}}{\pgfqpoint{0.000000in}{0.000000in}}{%
\pgfpathmoveto{\pgfqpoint{0.000000in}{0.000000in}}%
\pgfpathlineto{\pgfqpoint{0.000000in}{-0.041667in}}%
\pgfusepath{stroke,fill}%
}%
\begin{pgfscope}%
\pgfsys@transformshift{4.593118in}{4.405080in}%
\pgfsys@useobject{currentmarker}{}%
\end{pgfscope}%
\end{pgfscope}%
\begin{pgfscope}%
\definecolor{textcolor}{rgb}{0.000000,0.000000,0.000000}%
\pgfsetstrokecolor{textcolor}%
\pgfsetfillcolor{textcolor}%
\pgftext[x=4.593118in,y=0.352469in,,top]{\color{textcolor}{\rmfamily\fontsize{10.000000}{12.000000}\selectfont\catcode`\^=\active\def^{\ifmmode\sp\else\^{}\fi}\catcode`\%=\active\def%{\%}$\mathdefault{0.8}$}}%
\end{pgfscope}%
\begin{pgfscope}%
\pgfsetbuttcap%
\pgfsetroundjoin%
\definecolor{currentfill}{rgb}{0.000000,0.000000,0.000000}%
\pgfsetfillcolor{currentfill}%
\pgfsetlinewidth{0.501875pt}%
\definecolor{currentstroke}{rgb}{0.000000,0.000000,0.000000}%
\pgfsetstrokecolor{currentstroke}%
\pgfsetdash{}{0pt}%
\pgfsys@defobject{currentmarker}{\pgfqpoint{0.000000in}{0.000000in}}{\pgfqpoint{0.000000in}{0.041667in}}{%
\pgfpathmoveto{\pgfqpoint{0.000000in}{0.000000in}}%
\pgfpathlineto{\pgfqpoint{0.000000in}{0.041667in}}%
\pgfusepath{stroke,fill}%
}%
\begin{pgfscope}%
\pgfsys@transformshift{5.600618in}{0.401080in}%
\pgfsys@useobject{currentmarker}{}%
\end{pgfscope}%
\end{pgfscope}%
\begin{pgfscope}%
\pgfsetbuttcap%
\pgfsetroundjoin%
\definecolor{currentfill}{rgb}{0.000000,0.000000,0.000000}%
\pgfsetfillcolor{currentfill}%
\pgfsetlinewidth{0.501875pt}%
\definecolor{currentstroke}{rgb}{0.000000,0.000000,0.000000}%
\pgfsetstrokecolor{currentstroke}%
\pgfsetdash{}{0pt}%
\pgfsys@defobject{currentmarker}{\pgfqpoint{0.000000in}{-0.041667in}}{\pgfqpoint{0.000000in}{0.000000in}}{%
\pgfpathmoveto{\pgfqpoint{0.000000in}{0.000000in}}%
\pgfpathlineto{\pgfqpoint{0.000000in}{-0.041667in}}%
\pgfusepath{stroke,fill}%
}%
\begin{pgfscope}%
\pgfsys@transformshift{5.600618in}{4.405080in}%
\pgfsys@useobject{currentmarker}{}%
\end{pgfscope}%
\end{pgfscope}%
\begin{pgfscope}%
\definecolor{textcolor}{rgb}{0.000000,0.000000,0.000000}%
\pgfsetstrokecolor{textcolor}%
\pgfsetfillcolor{textcolor}%
\pgftext[x=5.600618in,y=0.352469in,,top]{\color{textcolor}{\rmfamily\fontsize{10.000000}{12.000000}\selectfont\catcode`\^=\active\def^{\ifmmode\sp\else\^{}\fi}\catcode`\%=\active\def%{\%}$\mathdefault{1.0}$}}%
\end{pgfscope}%
\begin{pgfscope}%
\pgfsetbuttcap%
\pgfsetroundjoin%
\definecolor{currentfill}{rgb}{0.000000,0.000000,0.000000}%
\pgfsetfillcolor{currentfill}%
\pgfsetlinewidth{0.501875pt}%
\definecolor{currentstroke}{rgb}{0.000000,0.000000,0.000000}%
\pgfsetstrokecolor{currentstroke}%
\pgfsetdash{}{0pt}%
\pgfsys@defobject{currentmarker}{\pgfqpoint{0.000000in}{0.000000in}}{\pgfqpoint{0.000000in}{0.020833in}}{%
\pgfpathmoveto{\pgfqpoint{0.000000in}{0.000000in}}%
\pgfpathlineto{\pgfqpoint{0.000000in}{0.020833in}}%
\pgfusepath{stroke,fill}%
}%
\begin{pgfscope}%
\pgfsys@transformshift{0.814993in}{0.401080in}%
\pgfsys@useobject{currentmarker}{}%
\end{pgfscope}%
\end{pgfscope}%
\begin{pgfscope}%
\pgfsetbuttcap%
\pgfsetroundjoin%
\definecolor{currentfill}{rgb}{0.000000,0.000000,0.000000}%
\pgfsetfillcolor{currentfill}%
\pgfsetlinewidth{0.501875pt}%
\definecolor{currentstroke}{rgb}{0.000000,0.000000,0.000000}%
\pgfsetstrokecolor{currentstroke}%
\pgfsetdash{}{0pt}%
\pgfsys@defobject{currentmarker}{\pgfqpoint{0.000000in}{-0.020833in}}{\pgfqpoint{0.000000in}{0.000000in}}{%
\pgfpathmoveto{\pgfqpoint{0.000000in}{0.000000in}}%
\pgfpathlineto{\pgfqpoint{0.000000in}{-0.020833in}}%
\pgfusepath{stroke,fill}%
}%
\begin{pgfscope}%
\pgfsys@transformshift{0.814993in}{4.405080in}%
\pgfsys@useobject{currentmarker}{}%
\end{pgfscope}%
\end{pgfscope}%
\begin{pgfscope}%
\pgfsetbuttcap%
\pgfsetroundjoin%
\definecolor{currentfill}{rgb}{0.000000,0.000000,0.000000}%
\pgfsetfillcolor{currentfill}%
\pgfsetlinewidth{0.501875pt}%
\definecolor{currentstroke}{rgb}{0.000000,0.000000,0.000000}%
\pgfsetstrokecolor{currentstroke}%
\pgfsetdash{}{0pt}%
\pgfsys@defobject{currentmarker}{\pgfqpoint{0.000000in}{0.000000in}}{\pgfqpoint{0.000000in}{0.020833in}}{%
\pgfpathmoveto{\pgfqpoint{0.000000in}{0.000000in}}%
\pgfpathlineto{\pgfqpoint{0.000000in}{0.020833in}}%
\pgfusepath{stroke,fill}%
}%
\begin{pgfscope}%
\pgfsys@transformshift{1.066868in}{0.401080in}%
\pgfsys@useobject{currentmarker}{}%
\end{pgfscope}%
\end{pgfscope}%
\begin{pgfscope}%
\pgfsetbuttcap%
\pgfsetroundjoin%
\definecolor{currentfill}{rgb}{0.000000,0.000000,0.000000}%
\pgfsetfillcolor{currentfill}%
\pgfsetlinewidth{0.501875pt}%
\definecolor{currentstroke}{rgb}{0.000000,0.000000,0.000000}%
\pgfsetstrokecolor{currentstroke}%
\pgfsetdash{}{0pt}%
\pgfsys@defobject{currentmarker}{\pgfqpoint{0.000000in}{-0.020833in}}{\pgfqpoint{0.000000in}{0.000000in}}{%
\pgfpathmoveto{\pgfqpoint{0.000000in}{0.000000in}}%
\pgfpathlineto{\pgfqpoint{0.000000in}{-0.020833in}}%
\pgfusepath{stroke,fill}%
}%
\begin{pgfscope}%
\pgfsys@transformshift{1.066868in}{4.405080in}%
\pgfsys@useobject{currentmarker}{}%
\end{pgfscope}%
\end{pgfscope}%
\begin{pgfscope}%
\pgfsetbuttcap%
\pgfsetroundjoin%
\definecolor{currentfill}{rgb}{0.000000,0.000000,0.000000}%
\pgfsetfillcolor{currentfill}%
\pgfsetlinewidth{0.501875pt}%
\definecolor{currentstroke}{rgb}{0.000000,0.000000,0.000000}%
\pgfsetstrokecolor{currentstroke}%
\pgfsetdash{}{0pt}%
\pgfsys@defobject{currentmarker}{\pgfqpoint{0.000000in}{0.000000in}}{\pgfqpoint{0.000000in}{0.020833in}}{%
\pgfpathmoveto{\pgfqpoint{0.000000in}{0.000000in}}%
\pgfpathlineto{\pgfqpoint{0.000000in}{0.020833in}}%
\pgfusepath{stroke,fill}%
}%
\begin{pgfscope}%
\pgfsys@transformshift{1.318743in}{0.401080in}%
\pgfsys@useobject{currentmarker}{}%
\end{pgfscope}%
\end{pgfscope}%
\begin{pgfscope}%
\pgfsetbuttcap%
\pgfsetroundjoin%
\definecolor{currentfill}{rgb}{0.000000,0.000000,0.000000}%
\pgfsetfillcolor{currentfill}%
\pgfsetlinewidth{0.501875pt}%
\definecolor{currentstroke}{rgb}{0.000000,0.000000,0.000000}%
\pgfsetstrokecolor{currentstroke}%
\pgfsetdash{}{0pt}%
\pgfsys@defobject{currentmarker}{\pgfqpoint{0.000000in}{-0.020833in}}{\pgfqpoint{0.000000in}{0.000000in}}{%
\pgfpathmoveto{\pgfqpoint{0.000000in}{0.000000in}}%
\pgfpathlineto{\pgfqpoint{0.000000in}{-0.020833in}}%
\pgfusepath{stroke,fill}%
}%
\begin{pgfscope}%
\pgfsys@transformshift{1.318743in}{4.405080in}%
\pgfsys@useobject{currentmarker}{}%
\end{pgfscope}%
\end{pgfscope}%
\begin{pgfscope}%
\pgfsetbuttcap%
\pgfsetroundjoin%
\definecolor{currentfill}{rgb}{0.000000,0.000000,0.000000}%
\pgfsetfillcolor{currentfill}%
\pgfsetlinewidth{0.501875pt}%
\definecolor{currentstroke}{rgb}{0.000000,0.000000,0.000000}%
\pgfsetstrokecolor{currentstroke}%
\pgfsetdash{}{0pt}%
\pgfsys@defobject{currentmarker}{\pgfqpoint{0.000000in}{0.000000in}}{\pgfqpoint{0.000000in}{0.020833in}}{%
\pgfpathmoveto{\pgfqpoint{0.000000in}{0.000000in}}%
\pgfpathlineto{\pgfqpoint{0.000000in}{0.020833in}}%
\pgfusepath{stroke,fill}%
}%
\begin{pgfscope}%
\pgfsys@transformshift{1.822493in}{0.401080in}%
\pgfsys@useobject{currentmarker}{}%
\end{pgfscope}%
\end{pgfscope}%
\begin{pgfscope}%
\pgfsetbuttcap%
\pgfsetroundjoin%
\definecolor{currentfill}{rgb}{0.000000,0.000000,0.000000}%
\pgfsetfillcolor{currentfill}%
\pgfsetlinewidth{0.501875pt}%
\definecolor{currentstroke}{rgb}{0.000000,0.000000,0.000000}%
\pgfsetstrokecolor{currentstroke}%
\pgfsetdash{}{0pt}%
\pgfsys@defobject{currentmarker}{\pgfqpoint{0.000000in}{-0.020833in}}{\pgfqpoint{0.000000in}{0.000000in}}{%
\pgfpathmoveto{\pgfqpoint{0.000000in}{0.000000in}}%
\pgfpathlineto{\pgfqpoint{0.000000in}{-0.020833in}}%
\pgfusepath{stroke,fill}%
}%
\begin{pgfscope}%
\pgfsys@transformshift{1.822493in}{4.405080in}%
\pgfsys@useobject{currentmarker}{}%
\end{pgfscope}%
\end{pgfscope}%
\begin{pgfscope}%
\pgfsetbuttcap%
\pgfsetroundjoin%
\definecolor{currentfill}{rgb}{0.000000,0.000000,0.000000}%
\pgfsetfillcolor{currentfill}%
\pgfsetlinewidth{0.501875pt}%
\definecolor{currentstroke}{rgb}{0.000000,0.000000,0.000000}%
\pgfsetstrokecolor{currentstroke}%
\pgfsetdash{}{0pt}%
\pgfsys@defobject{currentmarker}{\pgfqpoint{0.000000in}{0.000000in}}{\pgfqpoint{0.000000in}{0.020833in}}{%
\pgfpathmoveto{\pgfqpoint{0.000000in}{0.000000in}}%
\pgfpathlineto{\pgfqpoint{0.000000in}{0.020833in}}%
\pgfusepath{stroke,fill}%
}%
\begin{pgfscope}%
\pgfsys@transformshift{2.074368in}{0.401080in}%
\pgfsys@useobject{currentmarker}{}%
\end{pgfscope}%
\end{pgfscope}%
\begin{pgfscope}%
\pgfsetbuttcap%
\pgfsetroundjoin%
\definecolor{currentfill}{rgb}{0.000000,0.000000,0.000000}%
\pgfsetfillcolor{currentfill}%
\pgfsetlinewidth{0.501875pt}%
\definecolor{currentstroke}{rgb}{0.000000,0.000000,0.000000}%
\pgfsetstrokecolor{currentstroke}%
\pgfsetdash{}{0pt}%
\pgfsys@defobject{currentmarker}{\pgfqpoint{0.000000in}{-0.020833in}}{\pgfqpoint{0.000000in}{0.000000in}}{%
\pgfpathmoveto{\pgfqpoint{0.000000in}{0.000000in}}%
\pgfpathlineto{\pgfqpoint{0.000000in}{-0.020833in}}%
\pgfusepath{stroke,fill}%
}%
\begin{pgfscope}%
\pgfsys@transformshift{2.074368in}{4.405080in}%
\pgfsys@useobject{currentmarker}{}%
\end{pgfscope}%
\end{pgfscope}%
\begin{pgfscope}%
\pgfsetbuttcap%
\pgfsetroundjoin%
\definecolor{currentfill}{rgb}{0.000000,0.000000,0.000000}%
\pgfsetfillcolor{currentfill}%
\pgfsetlinewidth{0.501875pt}%
\definecolor{currentstroke}{rgb}{0.000000,0.000000,0.000000}%
\pgfsetstrokecolor{currentstroke}%
\pgfsetdash{}{0pt}%
\pgfsys@defobject{currentmarker}{\pgfqpoint{0.000000in}{0.000000in}}{\pgfqpoint{0.000000in}{0.020833in}}{%
\pgfpathmoveto{\pgfqpoint{0.000000in}{0.000000in}}%
\pgfpathlineto{\pgfqpoint{0.000000in}{0.020833in}}%
\pgfusepath{stroke,fill}%
}%
\begin{pgfscope}%
\pgfsys@transformshift{2.326243in}{0.401080in}%
\pgfsys@useobject{currentmarker}{}%
\end{pgfscope}%
\end{pgfscope}%
\begin{pgfscope}%
\pgfsetbuttcap%
\pgfsetroundjoin%
\definecolor{currentfill}{rgb}{0.000000,0.000000,0.000000}%
\pgfsetfillcolor{currentfill}%
\pgfsetlinewidth{0.501875pt}%
\definecolor{currentstroke}{rgb}{0.000000,0.000000,0.000000}%
\pgfsetstrokecolor{currentstroke}%
\pgfsetdash{}{0pt}%
\pgfsys@defobject{currentmarker}{\pgfqpoint{0.000000in}{-0.020833in}}{\pgfqpoint{0.000000in}{0.000000in}}{%
\pgfpathmoveto{\pgfqpoint{0.000000in}{0.000000in}}%
\pgfpathlineto{\pgfqpoint{0.000000in}{-0.020833in}}%
\pgfusepath{stroke,fill}%
}%
\begin{pgfscope}%
\pgfsys@transformshift{2.326243in}{4.405080in}%
\pgfsys@useobject{currentmarker}{}%
\end{pgfscope}%
\end{pgfscope}%
\begin{pgfscope}%
\pgfsetbuttcap%
\pgfsetroundjoin%
\definecolor{currentfill}{rgb}{0.000000,0.000000,0.000000}%
\pgfsetfillcolor{currentfill}%
\pgfsetlinewidth{0.501875pt}%
\definecolor{currentstroke}{rgb}{0.000000,0.000000,0.000000}%
\pgfsetstrokecolor{currentstroke}%
\pgfsetdash{}{0pt}%
\pgfsys@defobject{currentmarker}{\pgfqpoint{0.000000in}{0.000000in}}{\pgfqpoint{0.000000in}{0.020833in}}{%
\pgfpathmoveto{\pgfqpoint{0.000000in}{0.000000in}}%
\pgfpathlineto{\pgfqpoint{0.000000in}{0.020833in}}%
\pgfusepath{stroke,fill}%
}%
\begin{pgfscope}%
\pgfsys@transformshift{2.829993in}{0.401080in}%
\pgfsys@useobject{currentmarker}{}%
\end{pgfscope}%
\end{pgfscope}%
\begin{pgfscope}%
\pgfsetbuttcap%
\pgfsetroundjoin%
\definecolor{currentfill}{rgb}{0.000000,0.000000,0.000000}%
\pgfsetfillcolor{currentfill}%
\pgfsetlinewidth{0.501875pt}%
\definecolor{currentstroke}{rgb}{0.000000,0.000000,0.000000}%
\pgfsetstrokecolor{currentstroke}%
\pgfsetdash{}{0pt}%
\pgfsys@defobject{currentmarker}{\pgfqpoint{0.000000in}{-0.020833in}}{\pgfqpoint{0.000000in}{0.000000in}}{%
\pgfpathmoveto{\pgfqpoint{0.000000in}{0.000000in}}%
\pgfpathlineto{\pgfqpoint{0.000000in}{-0.020833in}}%
\pgfusepath{stroke,fill}%
}%
\begin{pgfscope}%
\pgfsys@transformshift{2.829993in}{4.405080in}%
\pgfsys@useobject{currentmarker}{}%
\end{pgfscope}%
\end{pgfscope}%
\begin{pgfscope}%
\pgfsetbuttcap%
\pgfsetroundjoin%
\definecolor{currentfill}{rgb}{0.000000,0.000000,0.000000}%
\pgfsetfillcolor{currentfill}%
\pgfsetlinewidth{0.501875pt}%
\definecolor{currentstroke}{rgb}{0.000000,0.000000,0.000000}%
\pgfsetstrokecolor{currentstroke}%
\pgfsetdash{}{0pt}%
\pgfsys@defobject{currentmarker}{\pgfqpoint{0.000000in}{0.000000in}}{\pgfqpoint{0.000000in}{0.020833in}}{%
\pgfpathmoveto{\pgfqpoint{0.000000in}{0.000000in}}%
\pgfpathlineto{\pgfqpoint{0.000000in}{0.020833in}}%
\pgfusepath{stroke,fill}%
}%
\begin{pgfscope}%
\pgfsys@transformshift{3.081868in}{0.401080in}%
\pgfsys@useobject{currentmarker}{}%
\end{pgfscope}%
\end{pgfscope}%
\begin{pgfscope}%
\pgfsetbuttcap%
\pgfsetroundjoin%
\definecolor{currentfill}{rgb}{0.000000,0.000000,0.000000}%
\pgfsetfillcolor{currentfill}%
\pgfsetlinewidth{0.501875pt}%
\definecolor{currentstroke}{rgb}{0.000000,0.000000,0.000000}%
\pgfsetstrokecolor{currentstroke}%
\pgfsetdash{}{0pt}%
\pgfsys@defobject{currentmarker}{\pgfqpoint{0.000000in}{-0.020833in}}{\pgfqpoint{0.000000in}{0.000000in}}{%
\pgfpathmoveto{\pgfqpoint{0.000000in}{0.000000in}}%
\pgfpathlineto{\pgfqpoint{0.000000in}{-0.020833in}}%
\pgfusepath{stroke,fill}%
}%
\begin{pgfscope}%
\pgfsys@transformshift{3.081868in}{4.405080in}%
\pgfsys@useobject{currentmarker}{}%
\end{pgfscope}%
\end{pgfscope}%
\begin{pgfscope}%
\pgfsetbuttcap%
\pgfsetroundjoin%
\definecolor{currentfill}{rgb}{0.000000,0.000000,0.000000}%
\pgfsetfillcolor{currentfill}%
\pgfsetlinewidth{0.501875pt}%
\definecolor{currentstroke}{rgb}{0.000000,0.000000,0.000000}%
\pgfsetstrokecolor{currentstroke}%
\pgfsetdash{}{0pt}%
\pgfsys@defobject{currentmarker}{\pgfqpoint{0.000000in}{0.000000in}}{\pgfqpoint{0.000000in}{0.020833in}}{%
\pgfpathmoveto{\pgfqpoint{0.000000in}{0.000000in}}%
\pgfpathlineto{\pgfqpoint{0.000000in}{0.020833in}}%
\pgfusepath{stroke,fill}%
}%
\begin{pgfscope}%
\pgfsys@transformshift{3.333743in}{0.401080in}%
\pgfsys@useobject{currentmarker}{}%
\end{pgfscope}%
\end{pgfscope}%
\begin{pgfscope}%
\pgfsetbuttcap%
\pgfsetroundjoin%
\definecolor{currentfill}{rgb}{0.000000,0.000000,0.000000}%
\pgfsetfillcolor{currentfill}%
\pgfsetlinewidth{0.501875pt}%
\definecolor{currentstroke}{rgb}{0.000000,0.000000,0.000000}%
\pgfsetstrokecolor{currentstroke}%
\pgfsetdash{}{0pt}%
\pgfsys@defobject{currentmarker}{\pgfqpoint{0.000000in}{-0.020833in}}{\pgfqpoint{0.000000in}{0.000000in}}{%
\pgfpathmoveto{\pgfqpoint{0.000000in}{0.000000in}}%
\pgfpathlineto{\pgfqpoint{0.000000in}{-0.020833in}}%
\pgfusepath{stroke,fill}%
}%
\begin{pgfscope}%
\pgfsys@transformshift{3.333743in}{4.405080in}%
\pgfsys@useobject{currentmarker}{}%
\end{pgfscope}%
\end{pgfscope}%
\begin{pgfscope}%
\pgfsetbuttcap%
\pgfsetroundjoin%
\definecolor{currentfill}{rgb}{0.000000,0.000000,0.000000}%
\pgfsetfillcolor{currentfill}%
\pgfsetlinewidth{0.501875pt}%
\definecolor{currentstroke}{rgb}{0.000000,0.000000,0.000000}%
\pgfsetstrokecolor{currentstroke}%
\pgfsetdash{}{0pt}%
\pgfsys@defobject{currentmarker}{\pgfqpoint{0.000000in}{0.000000in}}{\pgfqpoint{0.000000in}{0.020833in}}{%
\pgfpathmoveto{\pgfqpoint{0.000000in}{0.000000in}}%
\pgfpathlineto{\pgfqpoint{0.000000in}{0.020833in}}%
\pgfusepath{stroke,fill}%
}%
\begin{pgfscope}%
\pgfsys@transformshift{3.837493in}{0.401080in}%
\pgfsys@useobject{currentmarker}{}%
\end{pgfscope}%
\end{pgfscope}%
\begin{pgfscope}%
\pgfsetbuttcap%
\pgfsetroundjoin%
\definecolor{currentfill}{rgb}{0.000000,0.000000,0.000000}%
\pgfsetfillcolor{currentfill}%
\pgfsetlinewidth{0.501875pt}%
\definecolor{currentstroke}{rgb}{0.000000,0.000000,0.000000}%
\pgfsetstrokecolor{currentstroke}%
\pgfsetdash{}{0pt}%
\pgfsys@defobject{currentmarker}{\pgfqpoint{0.000000in}{-0.020833in}}{\pgfqpoint{0.000000in}{0.000000in}}{%
\pgfpathmoveto{\pgfqpoint{0.000000in}{0.000000in}}%
\pgfpathlineto{\pgfqpoint{0.000000in}{-0.020833in}}%
\pgfusepath{stroke,fill}%
}%
\begin{pgfscope}%
\pgfsys@transformshift{3.837493in}{4.405080in}%
\pgfsys@useobject{currentmarker}{}%
\end{pgfscope}%
\end{pgfscope}%
\begin{pgfscope}%
\pgfsetbuttcap%
\pgfsetroundjoin%
\definecolor{currentfill}{rgb}{0.000000,0.000000,0.000000}%
\pgfsetfillcolor{currentfill}%
\pgfsetlinewidth{0.501875pt}%
\definecolor{currentstroke}{rgb}{0.000000,0.000000,0.000000}%
\pgfsetstrokecolor{currentstroke}%
\pgfsetdash{}{0pt}%
\pgfsys@defobject{currentmarker}{\pgfqpoint{0.000000in}{0.000000in}}{\pgfqpoint{0.000000in}{0.020833in}}{%
\pgfpathmoveto{\pgfqpoint{0.000000in}{0.000000in}}%
\pgfpathlineto{\pgfqpoint{0.000000in}{0.020833in}}%
\pgfusepath{stroke,fill}%
}%
\begin{pgfscope}%
\pgfsys@transformshift{4.089368in}{0.401080in}%
\pgfsys@useobject{currentmarker}{}%
\end{pgfscope}%
\end{pgfscope}%
\begin{pgfscope}%
\pgfsetbuttcap%
\pgfsetroundjoin%
\definecolor{currentfill}{rgb}{0.000000,0.000000,0.000000}%
\pgfsetfillcolor{currentfill}%
\pgfsetlinewidth{0.501875pt}%
\definecolor{currentstroke}{rgb}{0.000000,0.000000,0.000000}%
\pgfsetstrokecolor{currentstroke}%
\pgfsetdash{}{0pt}%
\pgfsys@defobject{currentmarker}{\pgfqpoint{0.000000in}{-0.020833in}}{\pgfqpoint{0.000000in}{0.000000in}}{%
\pgfpathmoveto{\pgfqpoint{0.000000in}{0.000000in}}%
\pgfpathlineto{\pgfqpoint{0.000000in}{-0.020833in}}%
\pgfusepath{stroke,fill}%
}%
\begin{pgfscope}%
\pgfsys@transformshift{4.089368in}{4.405080in}%
\pgfsys@useobject{currentmarker}{}%
\end{pgfscope}%
\end{pgfscope}%
\begin{pgfscope}%
\pgfsetbuttcap%
\pgfsetroundjoin%
\definecolor{currentfill}{rgb}{0.000000,0.000000,0.000000}%
\pgfsetfillcolor{currentfill}%
\pgfsetlinewidth{0.501875pt}%
\definecolor{currentstroke}{rgb}{0.000000,0.000000,0.000000}%
\pgfsetstrokecolor{currentstroke}%
\pgfsetdash{}{0pt}%
\pgfsys@defobject{currentmarker}{\pgfqpoint{0.000000in}{0.000000in}}{\pgfqpoint{0.000000in}{0.020833in}}{%
\pgfpathmoveto{\pgfqpoint{0.000000in}{0.000000in}}%
\pgfpathlineto{\pgfqpoint{0.000000in}{0.020833in}}%
\pgfusepath{stroke,fill}%
}%
\begin{pgfscope}%
\pgfsys@transformshift{4.341243in}{0.401080in}%
\pgfsys@useobject{currentmarker}{}%
\end{pgfscope}%
\end{pgfscope}%
\begin{pgfscope}%
\pgfsetbuttcap%
\pgfsetroundjoin%
\definecolor{currentfill}{rgb}{0.000000,0.000000,0.000000}%
\pgfsetfillcolor{currentfill}%
\pgfsetlinewidth{0.501875pt}%
\definecolor{currentstroke}{rgb}{0.000000,0.000000,0.000000}%
\pgfsetstrokecolor{currentstroke}%
\pgfsetdash{}{0pt}%
\pgfsys@defobject{currentmarker}{\pgfqpoint{0.000000in}{-0.020833in}}{\pgfqpoint{0.000000in}{0.000000in}}{%
\pgfpathmoveto{\pgfqpoint{0.000000in}{0.000000in}}%
\pgfpathlineto{\pgfqpoint{0.000000in}{-0.020833in}}%
\pgfusepath{stroke,fill}%
}%
\begin{pgfscope}%
\pgfsys@transformshift{4.341243in}{4.405080in}%
\pgfsys@useobject{currentmarker}{}%
\end{pgfscope}%
\end{pgfscope}%
\begin{pgfscope}%
\pgfsetbuttcap%
\pgfsetroundjoin%
\definecolor{currentfill}{rgb}{0.000000,0.000000,0.000000}%
\pgfsetfillcolor{currentfill}%
\pgfsetlinewidth{0.501875pt}%
\definecolor{currentstroke}{rgb}{0.000000,0.000000,0.000000}%
\pgfsetstrokecolor{currentstroke}%
\pgfsetdash{}{0pt}%
\pgfsys@defobject{currentmarker}{\pgfqpoint{0.000000in}{0.000000in}}{\pgfqpoint{0.000000in}{0.020833in}}{%
\pgfpathmoveto{\pgfqpoint{0.000000in}{0.000000in}}%
\pgfpathlineto{\pgfqpoint{0.000000in}{0.020833in}}%
\pgfusepath{stroke,fill}%
}%
\begin{pgfscope}%
\pgfsys@transformshift{4.844993in}{0.401080in}%
\pgfsys@useobject{currentmarker}{}%
\end{pgfscope}%
\end{pgfscope}%
\begin{pgfscope}%
\pgfsetbuttcap%
\pgfsetroundjoin%
\definecolor{currentfill}{rgb}{0.000000,0.000000,0.000000}%
\pgfsetfillcolor{currentfill}%
\pgfsetlinewidth{0.501875pt}%
\definecolor{currentstroke}{rgb}{0.000000,0.000000,0.000000}%
\pgfsetstrokecolor{currentstroke}%
\pgfsetdash{}{0pt}%
\pgfsys@defobject{currentmarker}{\pgfqpoint{0.000000in}{-0.020833in}}{\pgfqpoint{0.000000in}{0.000000in}}{%
\pgfpathmoveto{\pgfqpoint{0.000000in}{0.000000in}}%
\pgfpathlineto{\pgfqpoint{0.000000in}{-0.020833in}}%
\pgfusepath{stroke,fill}%
}%
\begin{pgfscope}%
\pgfsys@transformshift{4.844993in}{4.405080in}%
\pgfsys@useobject{currentmarker}{}%
\end{pgfscope}%
\end{pgfscope}%
\begin{pgfscope}%
\pgfsetbuttcap%
\pgfsetroundjoin%
\definecolor{currentfill}{rgb}{0.000000,0.000000,0.000000}%
\pgfsetfillcolor{currentfill}%
\pgfsetlinewidth{0.501875pt}%
\definecolor{currentstroke}{rgb}{0.000000,0.000000,0.000000}%
\pgfsetstrokecolor{currentstroke}%
\pgfsetdash{}{0pt}%
\pgfsys@defobject{currentmarker}{\pgfqpoint{0.000000in}{0.000000in}}{\pgfqpoint{0.000000in}{0.020833in}}{%
\pgfpathmoveto{\pgfqpoint{0.000000in}{0.000000in}}%
\pgfpathlineto{\pgfqpoint{0.000000in}{0.020833in}}%
\pgfusepath{stroke,fill}%
}%
\begin{pgfscope}%
\pgfsys@transformshift{5.096868in}{0.401080in}%
\pgfsys@useobject{currentmarker}{}%
\end{pgfscope}%
\end{pgfscope}%
\begin{pgfscope}%
\pgfsetbuttcap%
\pgfsetroundjoin%
\definecolor{currentfill}{rgb}{0.000000,0.000000,0.000000}%
\pgfsetfillcolor{currentfill}%
\pgfsetlinewidth{0.501875pt}%
\definecolor{currentstroke}{rgb}{0.000000,0.000000,0.000000}%
\pgfsetstrokecolor{currentstroke}%
\pgfsetdash{}{0pt}%
\pgfsys@defobject{currentmarker}{\pgfqpoint{0.000000in}{-0.020833in}}{\pgfqpoint{0.000000in}{0.000000in}}{%
\pgfpathmoveto{\pgfqpoint{0.000000in}{0.000000in}}%
\pgfpathlineto{\pgfqpoint{0.000000in}{-0.020833in}}%
\pgfusepath{stroke,fill}%
}%
\begin{pgfscope}%
\pgfsys@transformshift{5.096868in}{4.405080in}%
\pgfsys@useobject{currentmarker}{}%
\end{pgfscope}%
\end{pgfscope}%
\begin{pgfscope}%
\pgfsetbuttcap%
\pgfsetroundjoin%
\definecolor{currentfill}{rgb}{0.000000,0.000000,0.000000}%
\pgfsetfillcolor{currentfill}%
\pgfsetlinewidth{0.501875pt}%
\definecolor{currentstroke}{rgb}{0.000000,0.000000,0.000000}%
\pgfsetstrokecolor{currentstroke}%
\pgfsetdash{}{0pt}%
\pgfsys@defobject{currentmarker}{\pgfqpoint{0.000000in}{0.000000in}}{\pgfqpoint{0.000000in}{0.020833in}}{%
\pgfpathmoveto{\pgfqpoint{0.000000in}{0.000000in}}%
\pgfpathlineto{\pgfqpoint{0.000000in}{0.020833in}}%
\pgfusepath{stroke,fill}%
}%
\begin{pgfscope}%
\pgfsys@transformshift{5.348743in}{0.401080in}%
\pgfsys@useobject{currentmarker}{}%
\end{pgfscope}%
\end{pgfscope}%
\begin{pgfscope}%
\pgfsetbuttcap%
\pgfsetroundjoin%
\definecolor{currentfill}{rgb}{0.000000,0.000000,0.000000}%
\pgfsetfillcolor{currentfill}%
\pgfsetlinewidth{0.501875pt}%
\definecolor{currentstroke}{rgb}{0.000000,0.000000,0.000000}%
\pgfsetstrokecolor{currentstroke}%
\pgfsetdash{}{0pt}%
\pgfsys@defobject{currentmarker}{\pgfqpoint{0.000000in}{-0.020833in}}{\pgfqpoint{0.000000in}{0.000000in}}{%
\pgfpathmoveto{\pgfqpoint{0.000000in}{0.000000in}}%
\pgfpathlineto{\pgfqpoint{0.000000in}{-0.020833in}}%
\pgfusepath{stroke,fill}%
}%
\begin{pgfscope}%
\pgfsys@transformshift{5.348743in}{4.405080in}%
\pgfsys@useobject{currentmarker}{}%
\end{pgfscope}%
\end{pgfscope}%
\begin{pgfscope}%
\definecolor{textcolor}{rgb}{0.000000,0.000000,0.000000}%
\pgfsetstrokecolor{textcolor}%
\pgfsetfillcolor{textcolor}%
\pgftext[x=3.081868in,y=0.173457in,,top]{\color{textcolor}{\rmfamily\fontsize{10.000000}{12.000000}\selectfont\catcode`\^=\active\def^{\ifmmode\sp\else\^{}\fi}\catcode`\%=\active\def%{\%}$t$}}%
\end{pgfscope}%
\begin{pgfscope}%
\pgfsetbuttcap%
\pgfsetroundjoin%
\definecolor{currentfill}{rgb}{0.000000,0.000000,0.000000}%
\pgfsetfillcolor{currentfill}%
\pgfsetlinewidth{0.501875pt}%
\definecolor{currentstroke}{rgb}{0.000000,0.000000,0.000000}%
\pgfsetstrokecolor{currentstroke}%
\pgfsetdash{}{0pt}%
\pgfsys@defobject{currentmarker}{\pgfqpoint{0.000000in}{0.000000in}}{\pgfqpoint{0.041667in}{0.000000in}}{%
\pgfpathmoveto{\pgfqpoint{0.000000in}{0.000000in}}%
\pgfpathlineto{\pgfqpoint{0.041667in}{0.000000in}}%
\pgfusepath{stroke,fill}%
}%
\begin{pgfscope}%
\pgfsys@transformshift{0.563118in}{0.719023in}%
\pgfsys@useobject{currentmarker}{}%
\end{pgfscope}%
\end{pgfscope}%
\begin{pgfscope}%
\pgfsetbuttcap%
\pgfsetroundjoin%
\definecolor{currentfill}{rgb}{0.000000,0.000000,0.000000}%
\pgfsetfillcolor{currentfill}%
\pgfsetlinewidth{0.501875pt}%
\definecolor{currentstroke}{rgb}{0.000000,0.000000,0.000000}%
\pgfsetstrokecolor{currentstroke}%
\pgfsetdash{}{0pt}%
\pgfsys@defobject{currentmarker}{\pgfqpoint{-0.041667in}{0.000000in}}{\pgfqpoint{-0.000000in}{0.000000in}}{%
\pgfpathmoveto{\pgfqpoint{-0.000000in}{0.000000in}}%
\pgfpathlineto{\pgfqpoint{-0.041667in}{0.000000in}}%
\pgfusepath{stroke,fill}%
}%
\begin{pgfscope}%
\pgfsys@transformshift{5.600618in}{0.719023in}%
\pgfsys@useobject{currentmarker}{}%
\end{pgfscope}%
\end{pgfscope}%
\begin{pgfscope}%
\definecolor{textcolor}{rgb}{0.000000,0.000000,0.000000}%
\pgfsetstrokecolor{textcolor}%
\pgfsetfillcolor{textcolor}%
\pgftext[x=0.229012in, y=0.670798in, left, base]{\color{textcolor}{\rmfamily\fontsize{10.000000}{12.000000}\selectfont\catcode`\^=\active\def^{\ifmmode\sp\else\^{}\fi}\catcode`\%=\active\def%{\%}$\mathdefault{-1.0}$}}%
\end{pgfscope}%
\begin{pgfscope}%
\pgfsetbuttcap%
\pgfsetroundjoin%
\definecolor{currentfill}{rgb}{0.000000,0.000000,0.000000}%
\pgfsetfillcolor{currentfill}%
\pgfsetlinewidth{0.501875pt}%
\definecolor{currentstroke}{rgb}{0.000000,0.000000,0.000000}%
\pgfsetstrokecolor{currentstroke}%
\pgfsetdash{}{0pt}%
\pgfsys@defobject{currentmarker}{\pgfqpoint{0.000000in}{0.000000in}}{\pgfqpoint{0.041667in}{0.000000in}}{%
\pgfpathmoveto{\pgfqpoint{0.000000in}{0.000000in}}%
\pgfpathlineto{\pgfqpoint{0.041667in}{0.000000in}}%
\pgfusepath{stroke,fill}%
}%
\begin{pgfscope}%
\pgfsys@transformshift{0.563118in}{1.354001in}%
\pgfsys@useobject{currentmarker}{}%
\end{pgfscope}%
\end{pgfscope}%
\begin{pgfscope}%
\pgfsetbuttcap%
\pgfsetroundjoin%
\definecolor{currentfill}{rgb}{0.000000,0.000000,0.000000}%
\pgfsetfillcolor{currentfill}%
\pgfsetlinewidth{0.501875pt}%
\definecolor{currentstroke}{rgb}{0.000000,0.000000,0.000000}%
\pgfsetstrokecolor{currentstroke}%
\pgfsetdash{}{0pt}%
\pgfsys@defobject{currentmarker}{\pgfqpoint{-0.041667in}{0.000000in}}{\pgfqpoint{-0.000000in}{0.000000in}}{%
\pgfpathmoveto{\pgfqpoint{-0.000000in}{0.000000in}}%
\pgfpathlineto{\pgfqpoint{-0.041667in}{0.000000in}}%
\pgfusepath{stroke,fill}%
}%
\begin{pgfscope}%
\pgfsys@transformshift{5.600618in}{1.354001in}%
\pgfsys@useobject{currentmarker}{}%
\end{pgfscope}%
\end{pgfscope}%
\begin{pgfscope}%
\definecolor{textcolor}{rgb}{0.000000,0.000000,0.000000}%
\pgfsetstrokecolor{textcolor}%
\pgfsetfillcolor{textcolor}%
\pgftext[x=0.229012in, y=1.305776in, left, base]{\color{textcolor}{\rmfamily\fontsize{10.000000}{12.000000}\selectfont\catcode`\^=\active\def^{\ifmmode\sp\else\^{}\fi}\catcode`\%=\active\def%{\%}$\mathdefault{-0.5}$}}%
\end{pgfscope}%
\begin{pgfscope}%
\pgfsetbuttcap%
\pgfsetroundjoin%
\definecolor{currentfill}{rgb}{0.000000,0.000000,0.000000}%
\pgfsetfillcolor{currentfill}%
\pgfsetlinewidth{0.501875pt}%
\definecolor{currentstroke}{rgb}{0.000000,0.000000,0.000000}%
\pgfsetstrokecolor{currentstroke}%
\pgfsetdash{}{0pt}%
\pgfsys@defobject{currentmarker}{\pgfqpoint{0.000000in}{0.000000in}}{\pgfqpoint{0.041667in}{0.000000in}}{%
\pgfpathmoveto{\pgfqpoint{0.000000in}{0.000000in}}%
\pgfpathlineto{\pgfqpoint{0.041667in}{0.000000in}}%
\pgfusepath{stroke,fill}%
}%
\begin{pgfscope}%
\pgfsys@transformshift{0.563118in}{1.988980in}%
\pgfsys@useobject{currentmarker}{}%
\end{pgfscope}%
\end{pgfscope}%
\begin{pgfscope}%
\pgfsetbuttcap%
\pgfsetroundjoin%
\definecolor{currentfill}{rgb}{0.000000,0.000000,0.000000}%
\pgfsetfillcolor{currentfill}%
\pgfsetlinewidth{0.501875pt}%
\definecolor{currentstroke}{rgb}{0.000000,0.000000,0.000000}%
\pgfsetstrokecolor{currentstroke}%
\pgfsetdash{}{0pt}%
\pgfsys@defobject{currentmarker}{\pgfqpoint{-0.041667in}{0.000000in}}{\pgfqpoint{-0.000000in}{0.000000in}}{%
\pgfpathmoveto{\pgfqpoint{-0.000000in}{0.000000in}}%
\pgfpathlineto{\pgfqpoint{-0.041667in}{0.000000in}}%
\pgfusepath{stroke,fill}%
}%
\begin{pgfscope}%
\pgfsys@transformshift{5.600618in}{1.988980in}%
\pgfsys@useobject{currentmarker}{}%
\end{pgfscope}%
\end{pgfscope}%
\begin{pgfscope}%
\definecolor{textcolor}{rgb}{0.000000,0.000000,0.000000}%
\pgfsetstrokecolor{textcolor}%
\pgfsetfillcolor{textcolor}%
\pgftext[x=0.337037in, y=1.940755in, left, base]{\color{textcolor}{\rmfamily\fontsize{10.000000}{12.000000}\selectfont\catcode`\^=\active\def^{\ifmmode\sp\else\^{}\fi}\catcode`\%=\active\def%{\%}$\mathdefault{0.0}$}}%
\end{pgfscope}%
\begin{pgfscope}%
\pgfsetbuttcap%
\pgfsetroundjoin%
\definecolor{currentfill}{rgb}{0.000000,0.000000,0.000000}%
\pgfsetfillcolor{currentfill}%
\pgfsetlinewidth{0.501875pt}%
\definecolor{currentstroke}{rgb}{0.000000,0.000000,0.000000}%
\pgfsetstrokecolor{currentstroke}%
\pgfsetdash{}{0pt}%
\pgfsys@defobject{currentmarker}{\pgfqpoint{0.000000in}{0.000000in}}{\pgfqpoint{0.041667in}{0.000000in}}{%
\pgfpathmoveto{\pgfqpoint{0.000000in}{0.000000in}}%
\pgfpathlineto{\pgfqpoint{0.041667in}{0.000000in}}%
\pgfusepath{stroke,fill}%
}%
\begin{pgfscope}%
\pgfsys@transformshift{0.563118in}{2.623959in}%
\pgfsys@useobject{currentmarker}{}%
\end{pgfscope}%
\end{pgfscope}%
\begin{pgfscope}%
\pgfsetbuttcap%
\pgfsetroundjoin%
\definecolor{currentfill}{rgb}{0.000000,0.000000,0.000000}%
\pgfsetfillcolor{currentfill}%
\pgfsetlinewidth{0.501875pt}%
\definecolor{currentstroke}{rgb}{0.000000,0.000000,0.000000}%
\pgfsetstrokecolor{currentstroke}%
\pgfsetdash{}{0pt}%
\pgfsys@defobject{currentmarker}{\pgfqpoint{-0.041667in}{0.000000in}}{\pgfqpoint{-0.000000in}{0.000000in}}{%
\pgfpathmoveto{\pgfqpoint{-0.000000in}{0.000000in}}%
\pgfpathlineto{\pgfqpoint{-0.041667in}{0.000000in}}%
\pgfusepath{stroke,fill}%
}%
\begin{pgfscope}%
\pgfsys@transformshift{5.600618in}{2.623959in}%
\pgfsys@useobject{currentmarker}{}%
\end{pgfscope}%
\end{pgfscope}%
\begin{pgfscope}%
\definecolor{textcolor}{rgb}{0.000000,0.000000,0.000000}%
\pgfsetstrokecolor{textcolor}%
\pgfsetfillcolor{textcolor}%
\pgftext[x=0.337037in, y=2.575734in, left, base]{\color{textcolor}{\rmfamily\fontsize{10.000000}{12.000000}\selectfont\catcode`\^=\active\def^{\ifmmode\sp\else\^{}\fi}\catcode`\%=\active\def%{\%}$\mathdefault{0.5}$}}%
\end{pgfscope}%
\begin{pgfscope}%
\pgfsetbuttcap%
\pgfsetroundjoin%
\definecolor{currentfill}{rgb}{0.000000,0.000000,0.000000}%
\pgfsetfillcolor{currentfill}%
\pgfsetlinewidth{0.501875pt}%
\definecolor{currentstroke}{rgb}{0.000000,0.000000,0.000000}%
\pgfsetstrokecolor{currentstroke}%
\pgfsetdash{}{0pt}%
\pgfsys@defobject{currentmarker}{\pgfqpoint{0.000000in}{0.000000in}}{\pgfqpoint{0.041667in}{0.000000in}}{%
\pgfpathmoveto{\pgfqpoint{0.000000in}{0.000000in}}%
\pgfpathlineto{\pgfqpoint{0.041667in}{0.000000in}}%
\pgfusepath{stroke,fill}%
}%
\begin{pgfscope}%
\pgfsys@transformshift{0.563118in}{3.258938in}%
\pgfsys@useobject{currentmarker}{}%
\end{pgfscope}%
\end{pgfscope}%
\begin{pgfscope}%
\pgfsetbuttcap%
\pgfsetroundjoin%
\definecolor{currentfill}{rgb}{0.000000,0.000000,0.000000}%
\pgfsetfillcolor{currentfill}%
\pgfsetlinewidth{0.501875pt}%
\definecolor{currentstroke}{rgb}{0.000000,0.000000,0.000000}%
\pgfsetstrokecolor{currentstroke}%
\pgfsetdash{}{0pt}%
\pgfsys@defobject{currentmarker}{\pgfqpoint{-0.041667in}{0.000000in}}{\pgfqpoint{-0.000000in}{0.000000in}}{%
\pgfpathmoveto{\pgfqpoint{-0.000000in}{0.000000in}}%
\pgfpathlineto{\pgfqpoint{-0.041667in}{0.000000in}}%
\pgfusepath{stroke,fill}%
}%
\begin{pgfscope}%
\pgfsys@transformshift{5.600618in}{3.258938in}%
\pgfsys@useobject{currentmarker}{}%
\end{pgfscope}%
\end{pgfscope}%
\begin{pgfscope}%
\definecolor{textcolor}{rgb}{0.000000,0.000000,0.000000}%
\pgfsetstrokecolor{textcolor}%
\pgfsetfillcolor{textcolor}%
\pgftext[x=0.337037in, y=3.210712in, left, base]{\color{textcolor}{\rmfamily\fontsize{10.000000}{12.000000}\selectfont\catcode`\^=\active\def^{\ifmmode\sp\else\^{}\fi}\catcode`\%=\active\def%{\%}$\mathdefault{1.0}$}}%
\end{pgfscope}%
\begin{pgfscope}%
\pgfsetbuttcap%
\pgfsetroundjoin%
\definecolor{currentfill}{rgb}{0.000000,0.000000,0.000000}%
\pgfsetfillcolor{currentfill}%
\pgfsetlinewidth{0.501875pt}%
\definecolor{currentstroke}{rgb}{0.000000,0.000000,0.000000}%
\pgfsetstrokecolor{currentstroke}%
\pgfsetdash{}{0pt}%
\pgfsys@defobject{currentmarker}{\pgfqpoint{0.000000in}{0.000000in}}{\pgfqpoint{0.041667in}{0.000000in}}{%
\pgfpathmoveto{\pgfqpoint{0.000000in}{0.000000in}}%
\pgfpathlineto{\pgfqpoint{0.041667in}{0.000000in}}%
\pgfusepath{stroke,fill}%
}%
\begin{pgfscope}%
\pgfsys@transformshift{0.563118in}{3.893916in}%
\pgfsys@useobject{currentmarker}{}%
\end{pgfscope}%
\end{pgfscope}%
\begin{pgfscope}%
\pgfsetbuttcap%
\pgfsetroundjoin%
\definecolor{currentfill}{rgb}{0.000000,0.000000,0.000000}%
\pgfsetfillcolor{currentfill}%
\pgfsetlinewidth{0.501875pt}%
\definecolor{currentstroke}{rgb}{0.000000,0.000000,0.000000}%
\pgfsetstrokecolor{currentstroke}%
\pgfsetdash{}{0pt}%
\pgfsys@defobject{currentmarker}{\pgfqpoint{-0.041667in}{0.000000in}}{\pgfqpoint{-0.000000in}{0.000000in}}{%
\pgfpathmoveto{\pgfqpoint{-0.000000in}{0.000000in}}%
\pgfpathlineto{\pgfqpoint{-0.041667in}{0.000000in}}%
\pgfusepath{stroke,fill}%
}%
\begin{pgfscope}%
\pgfsys@transformshift{5.600618in}{3.893916in}%
\pgfsys@useobject{currentmarker}{}%
\end{pgfscope}%
\end{pgfscope}%
\begin{pgfscope}%
\definecolor{textcolor}{rgb}{0.000000,0.000000,0.000000}%
\pgfsetstrokecolor{textcolor}%
\pgfsetfillcolor{textcolor}%
\pgftext[x=0.337037in, y=3.845691in, left, base]{\color{textcolor}{\rmfamily\fontsize{10.000000}{12.000000}\selectfont\catcode`\^=\active\def^{\ifmmode\sp\else\^{}\fi}\catcode`\%=\active\def%{\%}$\mathdefault{1.5}$}}%
\end{pgfscope}%
\begin{pgfscope}%
\pgfsetbuttcap%
\pgfsetroundjoin%
\definecolor{currentfill}{rgb}{0.000000,0.000000,0.000000}%
\pgfsetfillcolor{currentfill}%
\pgfsetlinewidth{0.501875pt}%
\definecolor{currentstroke}{rgb}{0.000000,0.000000,0.000000}%
\pgfsetstrokecolor{currentstroke}%
\pgfsetdash{}{0pt}%
\pgfsys@defobject{currentmarker}{\pgfqpoint{0.000000in}{0.000000in}}{\pgfqpoint{0.020833in}{0.000000in}}{%
\pgfpathmoveto{\pgfqpoint{0.000000in}{0.000000in}}%
\pgfpathlineto{\pgfqpoint{0.020833in}{0.000000in}}%
\pgfusepath{stroke,fill}%
}%
\begin{pgfscope}%
\pgfsys@transformshift{0.563118in}{0.465031in}%
\pgfsys@useobject{currentmarker}{}%
\end{pgfscope}%
\end{pgfscope}%
\begin{pgfscope}%
\pgfsetbuttcap%
\pgfsetroundjoin%
\definecolor{currentfill}{rgb}{0.000000,0.000000,0.000000}%
\pgfsetfillcolor{currentfill}%
\pgfsetlinewidth{0.501875pt}%
\definecolor{currentstroke}{rgb}{0.000000,0.000000,0.000000}%
\pgfsetstrokecolor{currentstroke}%
\pgfsetdash{}{0pt}%
\pgfsys@defobject{currentmarker}{\pgfqpoint{-0.020833in}{0.000000in}}{\pgfqpoint{-0.000000in}{0.000000in}}{%
\pgfpathmoveto{\pgfqpoint{-0.000000in}{0.000000in}}%
\pgfpathlineto{\pgfqpoint{-0.020833in}{0.000000in}}%
\pgfusepath{stroke,fill}%
}%
\begin{pgfscope}%
\pgfsys@transformshift{5.600618in}{0.465031in}%
\pgfsys@useobject{currentmarker}{}%
\end{pgfscope}%
\end{pgfscope}%
\begin{pgfscope}%
\pgfsetbuttcap%
\pgfsetroundjoin%
\definecolor{currentfill}{rgb}{0.000000,0.000000,0.000000}%
\pgfsetfillcolor{currentfill}%
\pgfsetlinewidth{0.501875pt}%
\definecolor{currentstroke}{rgb}{0.000000,0.000000,0.000000}%
\pgfsetstrokecolor{currentstroke}%
\pgfsetdash{}{0pt}%
\pgfsys@defobject{currentmarker}{\pgfqpoint{0.000000in}{0.000000in}}{\pgfqpoint{0.020833in}{0.000000in}}{%
\pgfpathmoveto{\pgfqpoint{0.000000in}{0.000000in}}%
\pgfpathlineto{\pgfqpoint{0.020833in}{0.000000in}}%
\pgfusepath{stroke,fill}%
}%
\begin{pgfscope}%
\pgfsys@transformshift{0.563118in}{0.592027in}%
\pgfsys@useobject{currentmarker}{}%
\end{pgfscope}%
\end{pgfscope}%
\begin{pgfscope}%
\pgfsetbuttcap%
\pgfsetroundjoin%
\definecolor{currentfill}{rgb}{0.000000,0.000000,0.000000}%
\pgfsetfillcolor{currentfill}%
\pgfsetlinewidth{0.501875pt}%
\definecolor{currentstroke}{rgb}{0.000000,0.000000,0.000000}%
\pgfsetstrokecolor{currentstroke}%
\pgfsetdash{}{0pt}%
\pgfsys@defobject{currentmarker}{\pgfqpoint{-0.020833in}{0.000000in}}{\pgfqpoint{-0.000000in}{0.000000in}}{%
\pgfpathmoveto{\pgfqpoint{-0.000000in}{0.000000in}}%
\pgfpathlineto{\pgfqpoint{-0.020833in}{0.000000in}}%
\pgfusepath{stroke,fill}%
}%
\begin{pgfscope}%
\pgfsys@transformshift{5.600618in}{0.592027in}%
\pgfsys@useobject{currentmarker}{}%
\end{pgfscope}%
\end{pgfscope}%
\begin{pgfscope}%
\pgfsetbuttcap%
\pgfsetroundjoin%
\definecolor{currentfill}{rgb}{0.000000,0.000000,0.000000}%
\pgfsetfillcolor{currentfill}%
\pgfsetlinewidth{0.501875pt}%
\definecolor{currentstroke}{rgb}{0.000000,0.000000,0.000000}%
\pgfsetstrokecolor{currentstroke}%
\pgfsetdash{}{0pt}%
\pgfsys@defobject{currentmarker}{\pgfqpoint{0.000000in}{0.000000in}}{\pgfqpoint{0.020833in}{0.000000in}}{%
\pgfpathmoveto{\pgfqpoint{0.000000in}{0.000000in}}%
\pgfpathlineto{\pgfqpoint{0.020833in}{0.000000in}}%
\pgfusepath{stroke,fill}%
}%
\begin{pgfscope}%
\pgfsys@transformshift{0.563118in}{0.846019in}%
\pgfsys@useobject{currentmarker}{}%
\end{pgfscope}%
\end{pgfscope}%
\begin{pgfscope}%
\pgfsetbuttcap%
\pgfsetroundjoin%
\definecolor{currentfill}{rgb}{0.000000,0.000000,0.000000}%
\pgfsetfillcolor{currentfill}%
\pgfsetlinewidth{0.501875pt}%
\definecolor{currentstroke}{rgb}{0.000000,0.000000,0.000000}%
\pgfsetstrokecolor{currentstroke}%
\pgfsetdash{}{0pt}%
\pgfsys@defobject{currentmarker}{\pgfqpoint{-0.020833in}{0.000000in}}{\pgfqpoint{-0.000000in}{0.000000in}}{%
\pgfpathmoveto{\pgfqpoint{-0.000000in}{0.000000in}}%
\pgfpathlineto{\pgfqpoint{-0.020833in}{0.000000in}}%
\pgfusepath{stroke,fill}%
}%
\begin{pgfscope}%
\pgfsys@transformshift{5.600618in}{0.846019in}%
\pgfsys@useobject{currentmarker}{}%
\end{pgfscope}%
\end{pgfscope}%
\begin{pgfscope}%
\pgfsetbuttcap%
\pgfsetroundjoin%
\definecolor{currentfill}{rgb}{0.000000,0.000000,0.000000}%
\pgfsetfillcolor{currentfill}%
\pgfsetlinewidth{0.501875pt}%
\definecolor{currentstroke}{rgb}{0.000000,0.000000,0.000000}%
\pgfsetstrokecolor{currentstroke}%
\pgfsetdash{}{0pt}%
\pgfsys@defobject{currentmarker}{\pgfqpoint{0.000000in}{0.000000in}}{\pgfqpoint{0.020833in}{0.000000in}}{%
\pgfpathmoveto{\pgfqpoint{0.000000in}{0.000000in}}%
\pgfpathlineto{\pgfqpoint{0.020833in}{0.000000in}}%
\pgfusepath{stroke,fill}%
}%
\begin{pgfscope}%
\pgfsys@transformshift{0.563118in}{0.973014in}%
\pgfsys@useobject{currentmarker}{}%
\end{pgfscope}%
\end{pgfscope}%
\begin{pgfscope}%
\pgfsetbuttcap%
\pgfsetroundjoin%
\definecolor{currentfill}{rgb}{0.000000,0.000000,0.000000}%
\pgfsetfillcolor{currentfill}%
\pgfsetlinewidth{0.501875pt}%
\definecolor{currentstroke}{rgb}{0.000000,0.000000,0.000000}%
\pgfsetstrokecolor{currentstroke}%
\pgfsetdash{}{0pt}%
\pgfsys@defobject{currentmarker}{\pgfqpoint{-0.020833in}{0.000000in}}{\pgfqpoint{-0.000000in}{0.000000in}}{%
\pgfpathmoveto{\pgfqpoint{-0.000000in}{0.000000in}}%
\pgfpathlineto{\pgfqpoint{-0.020833in}{0.000000in}}%
\pgfusepath{stroke,fill}%
}%
\begin{pgfscope}%
\pgfsys@transformshift{5.600618in}{0.973014in}%
\pgfsys@useobject{currentmarker}{}%
\end{pgfscope}%
\end{pgfscope}%
\begin{pgfscope}%
\pgfsetbuttcap%
\pgfsetroundjoin%
\definecolor{currentfill}{rgb}{0.000000,0.000000,0.000000}%
\pgfsetfillcolor{currentfill}%
\pgfsetlinewidth{0.501875pt}%
\definecolor{currentstroke}{rgb}{0.000000,0.000000,0.000000}%
\pgfsetstrokecolor{currentstroke}%
\pgfsetdash{}{0pt}%
\pgfsys@defobject{currentmarker}{\pgfqpoint{0.000000in}{0.000000in}}{\pgfqpoint{0.020833in}{0.000000in}}{%
\pgfpathmoveto{\pgfqpoint{0.000000in}{0.000000in}}%
\pgfpathlineto{\pgfqpoint{0.020833in}{0.000000in}}%
\pgfusepath{stroke,fill}%
}%
\begin{pgfscope}%
\pgfsys@transformshift{0.563118in}{1.100010in}%
\pgfsys@useobject{currentmarker}{}%
\end{pgfscope}%
\end{pgfscope}%
\begin{pgfscope}%
\pgfsetbuttcap%
\pgfsetroundjoin%
\definecolor{currentfill}{rgb}{0.000000,0.000000,0.000000}%
\pgfsetfillcolor{currentfill}%
\pgfsetlinewidth{0.501875pt}%
\definecolor{currentstroke}{rgb}{0.000000,0.000000,0.000000}%
\pgfsetstrokecolor{currentstroke}%
\pgfsetdash{}{0pt}%
\pgfsys@defobject{currentmarker}{\pgfqpoint{-0.020833in}{0.000000in}}{\pgfqpoint{-0.000000in}{0.000000in}}{%
\pgfpathmoveto{\pgfqpoint{-0.000000in}{0.000000in}}%
\pgfpathlineto{\pgfqpoint{-0.020833in}{0.000000in}}%
\pgfusepath{stroke,fill}%
}%
\begin{pgfscope}%
\pgfsys@transformshift{5.600618in}{1.100010in}%
\pgfsys@useobject{currentmarker}{}%
\end{pgfscope}%
\end{pgfscope}%
\begin{pgfscope}%
\pgfsetbuttcap%
\pgfsetroundjoin%
\definecolor{currentfill}{rgb}{0.000000,0.000000,0.000000}%
\pgfsetfillcolor{currentfill}%
\pgfsetlinewidth{0.501875pt}%
\definecolor{currentstroke}{rgb}{0.000000,0.000000,0.000000}%
\pgfsetstrokecolor{currentstroke}%
\pgfsetdash{}{0pt}%
\pgfsys@defobject{currentmarker}{\pgfqpoint{0.000000in}{0.000000in}}{\pgfqpoint{0.020833in}{0.000000in}}{%
\pgfpathmoveto{\pgfqpoint{0.000000in}{0.000000in}}%
\pgfpathlineto{\pgfqpoint{0.020833in}{0.000000in}}%
\pgfusepath{stroke,fill}%
}%
\begin{pgfscope}%
\pgfsys@transformshift{0.563118in}{1.227006in}%
\pgfsys@useobject{currentmarker}{}%
\end{pgfscope}%
\end{pgfscope}%
\begin{pgfscope}%
\pgfsetbuttcap%
\pgfsetroundjoin%
\definecolor{currentfill}{rgb}{0.000000,0.000000,0.000000}%
\pgfsetfillcolor{currentfill}%
\pgfsetlinewidth{0.501875pt}%
\definecolor{currentstroke}{rgb}{0.000000,0.000000,0.000000}%
\pgfsetstrokecolor{currentstroke}%
\pgfsetdash{}{0pt}%
\pgfsys@defobject{currentmarker}{\pgfqpoint{-0.020833in}{0.000000in}}{\pgfqpoint{-0.000000in}{0.000000in}}{%
\pgfpathmoveto{\pgfqpoint{-0.000000in}{0.000000in}}%
\pgfpathlineto{\pgfqpoint{-0.020833in}{0.000000in}}%
\pgfusepath{stroke,fill}%
}%
\begin{pgfscope}%
\pgfsys@transformshift{5.600618in}{1.227006in}%
\pgfsys@useobject{currentmarker}{}%
\end{pgfscope}%
\end{pgfscope}%
\begin{pgfscope}%
\pgfsetbuttcap%
\pgfsetroundjoin%
\definecolor{currentfill}{rgb}{0.000000,0.000000,0.000000}%
\pgfsetfillcolor{currentfill}%
\pgfsetlinewidth{0.501875pt}%
\definecolor{currentstroke}{rgb}{0.000000,0.000000,0.000000}%
\pgfsetstrokecolor{currentstroke}%
\pgfsetdash{}{0pt}%
\pgfsys@defobject{currentmarker}{\pgfqpoint{0.000000in}{0.000000in}}{\pgfqpoint{0.020833in}{0.000000in}}{%
\pgfpathmoveto{\pgfqpoint{0.000000in}{0.000000in}}%
\pgfpathlineto{\pgfqpoint{0.020833in}{0.000000in}}%
\pgfusepath{stroke,fill}%
}%
\begin{pgfscope}%
\pgfsys@transformshift{0.563118in}{1.480997in}%
\pgfsys@useobject{currentmarker}{}%
\end{pgfscope}%
\end{pgfscope}%
\begin{pgfscope}%
\pgfsetbuttcap%
\pgfsetroundjoin%
\definecolor{currentfill}{rgb}{0.000000,0.000000,0.000000}%
\pgfsetfillcolor{currentfill}%
\pgfsetlinewidth{0.501875pt}%
\definecolor{currentstroke}{rgb}{0.000000,0.000000,0.000000}%
\pgfsetstrokecolor{currentstroke}%
\pgfsetdash{}{0pt}%
\pgfsys@defobject{currentmarker}{\pgfqpoint{-0.020833in}{0.000000in}}{\pgfqpoint{-0.000000in}{0.000000in}}{%
\pgfpathmoveto{\pgfqpoint{-0.000000in}{0.000000in}}%
\pgfpathlineto{\pgfqpoint{-0.020833in}{0.000000in}}%
\pgfusepath{stroke,fill}%
}%
\begin{pgfscope}%
\pgfsys@transformshift{5.600618in}{1.480997in}%
\pgfsys@useobject{currentmarker}{}%
\end{pgfscope}%
\end{pgfscope}%
\begin{pgfscope}%
\pgfsetbuttcap%
\pgfsetroundjoin%
\definecolor{currentfill}{rgb}{0.000000,0.000000,0.000000}%
\pgfsetfillcolor{currentfill}%
\pgfsetlinewidth{0.501875pt}%
\definecolor{currentstroke}{rgb}{0.000000,0.000000,0.000000}%
\pgfsetstrokecolor{currentstroke}%
\pgfsetdash{}{0pt}%
\pgfsys@defobject{currentmarker}{\pgfqpoint{0.000000in}{0.000000in}}{\pgfqpoint{0.020833in}{0.000000in}}{%
\pgfpathmoveto{\pgfqpoint{0.000000in}{0.000000in}}%
\pgfpathlineto{\pgfqpoint{0.020833in}{0.000000in}}%
\pgfusepath{stroke,fill}%
}%
\begin{pgfscope}%
\pgfsys@transformshift{0.563118in}{1.607993in}%
\pgfsys@useobject{currentmarker}{}%
\end{pgfscope}%
\end{pgfscope}%
\begin{pgfscope}%
\pgfsetbuttcap%
\pgfsetroundjoin%
\definecolor{currentfill}{rgb}{0.000000,0.000000,0.000000}%
\pgfsetfillcolor{currentfill}%
\pgfsetlinewidth{0.501875pt}%
\definecolor{currentstroke}{rgb}{0.000000,0.000000,0.000000}%
\pgfsetstrokecolor{currentstroke}%
\pgfsetdash{}{0pt}%
\pgfsys@defobject{currentmarker}{\pgfqpoint{-0.020833in}{0.000000in}}{\pgfqpoint{-0.000000in}{0.000000in}}{%
\pgfpathmoveto{\pgfqpoint{-0.000000in}{0.000000in}}%
\pgfpathlineto{\pgfqpoint{-0.020833in}{0.000000in}}%
\pgfusepath{stroke,fill}%
}%
\begin{pgfscope}%
\pgfsys@transformshift{5.600618in}{1.607993in}%
\pgfsys@useobject{currentmarker}{}%
\end{pgfscope}%
\end{pgfscope}%
\begin{pgfscope}%
\pgfsetbuttcap%
\pgfsetroundjoin%
\definecolor{currentfill}{rgb}{0.000000,0.000000,0.000000}%
\pgfsetfillcolor{currentfill}%
\pgfsetlinewidth{0.501875pt}%
\definecolor{currentstroke}{rgb}{0.000000,0.000000,0.000000}%
\pgfsetstrokecolor{currentstroke}%
\pgfsetdash{}{0pt}%
\pgfsys@defobject{currentmarker}{\pgfqpoint{0.000000in}{0.000000in}}{\pgfqpoint{0.020833in}{0.000000in}}{%
\pgfpathmoveto{\pgfqpoint{0.000000in}{0.000000in}}%
\pgfpathlineto{\pgfqpoint{0.020833in}{0.000000in}}%
\pgfusepath{stroke,fill}%
}%
\begin{pgfscope}%
\pgfsys@transformshift{0.563118in}{1.734989in}%
\pgfsys@useobject{currentmarker}{}%
\end{pgfscope}%
\end{pgfscope}%
\begin{pgfscope}%
\pgfsetbuttcap%
\pgfsetroundjoin%
\definecolor{currentfill}{rgb}{0.000000,0.000000,0.000000}%
\pgfsetfillcolor{currentfill}%
\pgfsetlinewidth{0.501875pt}%
\definecolor{currentstroke}{rgb}{0.000000,0.000000,0.000000}%
\pgfsetstrokecolor{currentstroke}%
\pgfsetdash{}{0pt}%
\pgfsys@defobject{currentmarker}{\pgfqpoint{-0.020833in}{0.000000in}}{\pgfqpoint{-0.000000in}{0.000000in}}{%
\pgfpathmoveto{\pgfqpoint{-0.000000in}{0.000000in}}%
\pgfpathlineto{\pgfqpoint{-0.020833in}{0.000000in}}%
\pgfusepath{stroke,fill}%
}%
\begin{pgfscope}%
\pgfsys@transformshift{5.600618in}{1.734989in}%
\pgfsys@useobject{currentmarker}{}%
\end{pgfscope}%
\end{pgfscope}%
\begin{pgfscope}%
\pgfsetbuttcap%
\pgfsetroundjoin%
\definecolor{currentfill}{rgb}{0.000000,0.000000,0.000000}%
\pgfsetfillcolor{currentfill}%
\pgfsetlinewidth{0.501875pt}%
\definecolor{currentstroke}{rgb}{0.000000,0.000000,0.000000}%
\pgfsetstrokecolor{currentstroke}%
\pgfsetdash{}{0pt}%
\pgfsys@defobject{currentmarker}{\pgfqpoint{0.000000in}{0.000000in}}{\pgfqpoint{0.020833in}{0.000000in}}{%
\pgfpathmoveto{\pgfqpoint{0.000000in}{0.000000in}}%
\pgfpathlineto{\pgfqpoint{0.020833in}{0.000000in}}%
\pgfusepath{stroke,fill}%
}%
\begin{pgfscope}%
\pgfsys@transformshift{0.563118in}{1.861984in}%
\pgfsys@useobject{currentmarker}{}%
\end{pgfscope}%
\end{pgfscope}%
\begin{pgfscope}%
\pgfsetbuttcap%
\pgfsetroundjoin%
\definecolor{currentfill}{rgb}{0.000000,0.000000,0.000000}%
\pgfsetfillcolor{currentfill}%
\pgfsetlinewidth{0.501875pt}%
\definecolor{currentstroke}{rgb}{0.000000,0.000000,0.000000}%
\pgfsetstrokecolor{currentstroke}%
\pgfsetdash{}{0pt}%
\pgfsys@defobject{currentmarker}{\pgfqpoint{-0.020833in}{0.000000in}}{\pgfqpoint{-0.000000in}{0.000000in}}{%
\pgfpathmoveto{\pgfqpoint{-0.000000in}{0.000000in}}%
\pgfpathlineto{\pgfqpoint{-0.020833in}{0.000000in}}%
\pgfusepath{stroke,fill}%
}%
\begin{pgfscope}%
\pgfsys@transformshift{5.600618in}{1.861984in}%
\pgfsys@useobject{currentmarker}{}%
\end{pgfscope}%
\end{pgfscope}%
\begin{pgfscope}%
\pgfsetbuttcap%
\pgfsetroundjoin%
\definecolor{currentfill}{rgb}{0.000000,0.000000,0.000000}%
\pgfsetfillcolor{currentfill}%
\pgfsetlinewidth{0.501875pt}%
\definecolor{currentstroke}{rgb}{0.000000,0.000000,0.000000}%
\pgfsetstrokecolor{currentstroke}%
\pgfsetdash{}{0pt}%
\pgfsys@defobject{currentmarker}{\pgfqpoint{0.000000in}{0.000000in}}{\pgfqpoint{0.020833in}{0.000000in}}{%
\pgfpathmoveto{\pgfqpoint{0.000000in}{0.000000in}}%
\pgfpathlineto{\pgfqpoint{0.020833in}{0.000000in}}%
\pgfusepath{stroke,fill}%
}%
\begin{pgfscope}%
\pgfsys@transformshift{0.563118in}{2.115976in}%
\pgfsys@useobject{currentmarker}{}%
\end{pgfscope}%
\end{pgfscope}%
\begin{pgfscope}%
\pgfsetbuttcap%
\pgfsetroundjoin%
\definecolor{currentfill}{rgb}{0.000000,0.000000,0.000000}%
\pgfsetfillcolor{currentfill}%
\pgfsetlinewidth{0.501875pt}%
\definecolor{currentstroke}{rgb}{0.000000,0.000000,0.000000}%
\pgfsetstrokecolor{currentstroke}%
\pgfsetdash{}{0pt}%
\pgfsys@defobject{currentmarker}{\pgfqpoint{-0.020833in}{0.000000in}}{\pgfqpoint{-0.000000in}{0.000000in}}{%
\pgfpathmoveto{\pgfqpoint{-0.000000in}{0.000000in}}%
\pgfpathlineto{\pgfqpoint{-0.020833in}{0.000000in}}%
\pgfusepath{stroke,fill}%
}%
\begin{pgfscope}%
\pgfsys@transformshift{5.600618in}{2.115976in}%
\pgfsys@useobject{currentmarker}{}%
\end{pgfscope}%
\end{pgfscope}%
\begin{pgfscope}%
\pgfsetbuttcap%
\pgfsetroundjoin%
\definecolor{currentfill}{rgb}{0.000000,0.000000,0.000000}%
\pgfsetfillcolor{currentfill}%
\pgfsetlinewidth{0.501875pt}%
\definecolor{currentstroke}{rgb}{0.000000,0.000000,0.000000}%
\pgfsetstrokecolor{currentstroke}%
\pgfsetdash{}{0pt}%
\pgfsys@defobject{currentmarker}{\pgfqpoint{0.000000in}{0.000000in}}{\pgfqpoint{0.020833in}{0.000000in}}{%
\pgfpathmoveto{\pgfqpoint{0.000000in}{0.000000in}}%
\pgfpathlineto{\pgfqpoint{0.020833in}{0.000000in}}%
\pgfusepath{stroke,fill}%
}%
\begin{pgfscope}%
\pgfsys@transformshift{0.563118in}{2.242972in}%
\pgfsys@useobject{currentmarker}{}%
\end{pgfscope}%
\end{pgfscope}%
\begin{pgfscope}%
\pgfsetbuttcap%
\pgfsetroundjoin%
\definecolor{currentfill}{rgb}{0.000000,0.000000,0.000000}%
\pgfsetfillcolor{currentfill}%
\pgfsetlinewidth{0.501875pt}%
\definecolor{currentstroke}{rgb}{0.000000,0.000000,0.000000}%
\pgfsetstrokecolor{currentstroke}%
\pgfsetdash{}{0pt}%
\pgfsys@defobject{currentmarker}{\pgfqpoint{-0.020833in}{0.000000in}}{\pgfqpoint{-0.000000in}{0.000000in}}{%
\pgfpathmoveto{\pgfqpoint{-0.000000in}{0.000000in}}%
\pgfpathlineto{\pgfqpoint{-0.020833in}{0.000000in}}%
\pgfusepath{stroke,fill}%
}%
\begin{pgfscope}%
\pgfsys@transformshift{5.600618in}{2.242972in}%
\pgfsys@useobject{currentmarker}{}%
\end{pgfscope}%
\end{pgfscope}%
\begin{pgfscope}%
\pgfsetbuttcap%
\pgfsetroundjoin%
\definecolor{currentfill}{rgb}{0.000000,0.000000,0.000000}%
\pgfsetfillcolor{currentfill}%
\pgfsetlinewidth{0.501875pt}%
\definecolor{currentstroke}{rgb}{0.000000,0.000000,0.000000}%
\pgfsetstrokecolor{currentstroke}%
\pgfsetdash{}{0pt}%
\pgfsys@defobject{currentmarker}{\pgfqpoint{0.000000in}{0.000000in}}{\pgfqpoint{0.020833in}{0.000000in}}{%
\pgfpathmoveto{\pgfqpoint{0.000000in}{0.000000in}}%
\pgfpathlineto{\pgfqpoint{0.020833in}{0.000000in}}%
\pgfusepath{stroke,fill}%
}%
\begin{pgfscope}%
\pgfsys@transformshift{0.563118in}{2.369967in}%
\pgfsys@useobject{currentmarker}{}%
\end{pgfscope}%
\end{pgfscope}%
\begin{pgfscope}%
\pgfsetbuttcap%
\pgfsetroundjoin%
\definecolor{currentfill}{rgb}{0.000000,0.000000,0.000000}%
\pgfsetfillcolor{currentfill}%
\pgfsetlinewidth{0.501875pt}%
\definecolor{currentstroke}{rgb}{0.000000,0.000000,0.000000}%
\pgfsetstrokecolor{currentstroke}%
\pgfsetdash{}{0pt}%
\pgfsys@defobject{currentmarker}{\pgfqpoint{-0.020833in}{0.000000in}}{\pgfqpoint{-0.000000in}{0.000000in}}{%
\pgfpathmoveto{\pgfqpoint{-0.000000in}{0.000000in}}%
\pgfpathlineto{\pgfqpoint{-0.020833in}{0.000000in}}%
\pgfusepath{stroke,fill}%
}%
\begin{pgfscope}%
\pgfsys@transformshift{5.600618in}{2.369967in}%
\pgfsys@useobject{currentmarker}{}%
\end{pgfscope}%
\end{pgfscope}%
\begin{pgfscope}%
\pgfsetbuttcap%
\pgfsetroundjoin%
\definecolor{currentfill}{rgb}{0.000000,0.000000,0.000000}%
\pgfsetfillcolor{currentfill}%
\pgfsetlinewidth{0.501875pt}%
\definecolor{currentstroke}{rgb}{0.000000,0.000000,0.000000}%
\pgfsetstrokecolor{currentstroke}%
\pgfsetdash{}{0pt}%
\pgfsys@defobject{currentmarker}{\pgfqpoint{0.000000in}{0.000000in}}{\pgfqpoint{0.020833in}{0.000000in}}{%
\pgfpathmoveto{\pgfqpoint{0.000000in}{0.000000in}}%
\pgfpathlineto{\pgfqpoint{0.020833in}{0.000000in}}%
\pgfusepath{stroke,fill}%
}%
\begin{pgfscope}%
\pgfsys@transformshift{0.563118in}{2.496963in}%
\pgfsys@useobject{currentmarker}{}%
\end{pgfscope}%
\end{pgfscope}%
\begin{pgfscope}%
\pgfsetbuttcap%
\pgfsetroundjoin%
\definecolor{currentfill}{rgb}{0.000000,0.000000,0.000000}%
\pgfsetfillcolor{currentfill}%
\pgfsetlinewidth{0.501875pt}%
\definecolor{currentstroke}{rgb}{0.000000,0.000000,0.000000}%
\pgfsetstrokecolor{currentstroke}%
\pgfsetdash{}{0pt}%
\pgfsys@defobject{currentmarker}{\pgfqpoint{-0.020833in}{0.000000in}}{\pgfqpoint{-0.000000in}{0.000000in}}{%
\pgfpathmoveto{\pgfqpoint{-0.000000in}{0.000000in}}%
\pgfpathlineto{\pgfqpoint{-0.020833in}{0.000000in}}%
\pgfusepath{stroke,fill}%
}%
\begin{pgfscope}%
\pgfsys@transformshift{5.600618in}{2.496963in}%
\pgfsys@useobject{currentmarker}{}%
\end{pgfscope}%
\end{pgfscope}%
\begin{pgfscope}%
\pgfsetbuttcap%
\pgfsetroundjoin%
\definecolor{currentfill}{rgb}{0.000000,0.000000,0.000000}%
\pgfsetfillcolor{currentfill}%
\pgfsetlinewidth{0.501875pt}%
\definecolor{currentstroke}{rgb}{0.000000,0.000000,0.000000}%
\pgfsetstrokecolor{currentstroke}%
\pgfsetdash{}{0pt}%
\pgfsys@defobject{currentmarker}{\pgfqpoint{0.000000in}{0.000000in}}{\pgfqpoint{0.020833in}{0.000000in}}{%
\pgfpathmoveto{\pgfqpoint{0.000000in}{0.000000in}}%
\pgfpathlineto{\pgfqpoint{0.020833in}{0.000000in}}%
\pgfusepath{stroke,fill}%
}%
\begin{pgfscope}%
\pgfsys@transformshift{0.563118in}{2.750955in}%
\pgfsys@useobject{currentmarker}{}%
\end{pgfscope}%
\end{pgfscope}%
\begin{pgfscope}%
\pgfsetbuttcap%
\pgfsetroundjoin%
\definecolor{currentfill}{rgb}{0.000000,0.000000,0.000000}%
\pgfsetfillcolor{currentfill}%
\pgfsetlinewidth{0.501875pt}%
\definecolor{currentstroke}{rgb}{0.000000,0.000000,0.000000}%
\pgfsetstrokecolor{currentstroke}%
\pgfsetdash{}{0pt}%
\pgfsys@defobject{currentmarker}{\pgfqpoint{-0.020833in}{0.000000in}}{\pgfqpoint{-0.000000in}{0.000000in}}{%
\pgfpathmoveto{\pgfqpoint{-0.000000in}{0.000000in}}%
\pgfpathlineto{\pgfqpoint{-0.020833in}{0.000000in}}%
\pgfusepath{stroke,fill}%
}%
\begin{pgfscope}%
\pgfsys@transformshift{5.600618in}{2.750955in}%
\pgfsys@useobject{currentmarker}{}%
\end{pgfscope}%
\end{pgfscope}%
\begin{pgfscope}%
\pgfsetbuttcap%
\pgfsetroundjoin%
\definecolor{currentfill}{rgb}{0.000000,0.000000,0.000000}%
\pgfsetfillcolor{currentfill}%
\pgfsetlinewidth{0.501875pt}%
\definecolor{currentstroke}{rgb}{0.000000,0.000000,0.000000}%
\pgfsetstrokecolor{currentstroke}%
\pgfsetdash{}{0pt}%
\pgfsys@defobject{currentmarker}{\pgfqpoint{0.000000in}{0.000000in}}{\pgfqpoint{0.020833in}{0.000000in}}{%
\pgfpathmoveto{\pgfqpoint{0.000000in}{0.000000in}}%
\pgfpathlineto{\pgfqpoint{0.020833in}{0.000000in}}%
\pgfusepath{stroke,fill}%
}%
\begin{pgfscope}%
\pgfsys@transformshift{0.563118in}{2.877950in}%
\pgfsys@useobject{currentmarker}{}%
\end{pgfscope}%
\end{pgfscope}%
\begin{pgfscope}%
\pgfsetbuttcap%
\pgfsetroundjoin%
\definecolor{currentfill}{rgb}{0.000000,0.000000,0.000000}%
\pgfsetfillcolor{currentfill}%
\pgfsetlinewidth{0.501875pt}%
\definecolor{currentstroke}{rgb}{0.000000,0.000000,0.000000}%
\pgfsetstrokecolor{currentstroke}%
\pgfsetdash{}{0pt}%
\pgfsys@defobject{currentmarker}{\pgfqpoint{-0.020833in}{0.000000in}}{\pgfqpoint{-0.000000in}{0.000000in}}{%
\pgfpathmoveto{\pgfqpoint{-0.000000in}{0.000000in}}%
\pgfpathlineto{\pgfqpoint{-0.020833in}{0.000000in}}%
\pgfusepath{stroke,fill}%
}%
\begin{pgfscope}%
\pgfsys@transformshift{5.600618in}{2.877950in}%
\pgfsys@useobject{currentmarker}{}%
\end{pgfscope}%
\end{pgfscope}%
\begin{pgfscope}%
\pgfsetbuttcap%
\pgfsetroundjoin%
\definecolor{currentfill}{rgb}{0.000000,0.000000,0.000000}%
\pgfsetfillcolor{currentfill}%
\pgfsetlinewidth{0.501875pt}%
\definecolor{currentstroke}{rgb}{0.000000,0.000000,0.000000}%
\pgfsetstrokecolor{currentstroke}%
\pgfsetdash{}{0pt}%
\pgfsys@defobject{currentmarker}{\pgfqpoint{0.000000in}{0.000000in}}{\pgfqpoint{0.020833in}{0.000000in}}{%
\pgfpathmoveto{\pgfqpoint{0.000000in}{0.000000in}}%
\pgfpathlineto{\pgfqpoint{0.020833in}{0.000000in}}%
\pgfusepath{stroke,fill}%
}%
\begin{pgfscope}%
\pgfsys@transformshift{0.563118in}{3.004946in}%
\pgfsys@useobject{currentmarker}{}%
\end{pgfscope}%
\end{pgfscope}%
\begin{pgfscope}%
\pgfsetbuttcap%
\pgfsetroundjoin%
\definecolor{currentfill}{rgb}{0.000000,0.000000,0.000000}%
\pgfsetfillcolor{currentfill}%
\pgfsetlinewidth{0.501875pt}%
\definecolor{currentstroke}{rgb}{0.000000,0.000000,0.000000}%
\pgfsetstrokecolor{currentstroke}%
\pgfsetdash{}{0pt}%
\pgfsys@defobject{currentmarker}{\pgfqpoint{-0.020833in}{0.000000in}}{\pgfqpoint{-0.000000in}{0.000000in}}{%
\pgfpathmoveto{\pgfqpoint{-0.000000in}{0.000000in}}%
\pgfpathlineto{\pgfqpoint{-0.020833in}{0.000000in}}%
\pgfusepath{stroke,fill}%
}%
\begin{pgfscope}%
\pgfsys@transformshift{5.600618in}{3.004946in}%
\pgfsys@useobject{currentmarker}{}%
\end{pgfscope}%
\end{pgfscope}%
\begin{pgfscope}%
\pgfsetbuttcap%
\pgfsetroundjoin%
\definecolor{currentfill}{rgb}{0.000000,0.000000,0.000000}%
\pgfsetfillcolor{currentfill}%
\pgfsetlinewidth{0.501875pt}%
\definecolor{currentstroke}{rgb}{0.000000,0.000000,0.000000}%
\pgfsetstrokecolor{currentstroke}%
\pgfsetdash{}{0pt}%
\pgfsys@defobject{currentmarker}{\pgfqpoint{0.000000in}{0.000000in}}{\pgfqpoint{0.020833in}{0.000000in}}{%
\pgfpathmoveto{\pgfqpoint{0.000000in}{0.000000in}}%
\pgfpathlineto{\pgfqpoint{0.020833in}{0.000000in}}%
\pgfusepath{stroke,fill}%
}%
\begin{pgfscope}%
\pgfsys@transformshift{0.563118in}{3.131942in}%
\pgfsys@useobject{currentmarker}{}%
\end{pgfscope}%
\end{pgfscope}%
\begin{pgfscope}%
\pgfsetbuttcap%
\pgfsetroundjoin%
\definecolor{currentfill}{rgb}{0.000000,0.000000,0.000000}%
\pgfsetfillcolor{currentfill}%
\pgfsetlinewidth{0.501875pt}%
\definecolor{currentstroke}{rgb}{0.000000,0.000000,0.000000}%
\pgfsetstrokecolor{currentstroke}%
\pgfsetdash{}{0pt}%
\pgfsys@defobject{currentmarker}{\pgfqpoint{-0.020833in}{0.000000in}}{\pgfqpoint{-0.000000in}{0.000000in}}{%
\pgfpathmoveto{\pgfqpoint{-0.000000in}{0.000000in}}%
\pgfpathlineto{\pgfqpoint{-0.020833in}{0.000000in}}%
\pgfusepath{stroke,fill}%
}%
\begin{pgfscope}%
\pgfsys@transformshift{5.600618in}{3.131942in}%
\pgfsys@useobject{currentmarker}{}%
\end{pgfscope}%
\end{pgfscope}%
\begin{pgfscope}%
\pgfsetbuttcap%
\pgfsetroundjoin%
\definecolor{currentfill}{rgb}{0.000000,0.000000,0.000000}%
\pgfsetfillcolor{currentfill}%
\pgfsetlinewidth{0.501875pt}%
\definecolor{currentstroke}{rgb}{0.000000,0.000000,0.000000}%
\pgfsetstrokecolor{currentstroke}%
\pgfsetdash{}{0pt}%
\pgfsys@defobject{currentmarker}{\pgfqpoint{0.000000in}{0.000000in}}{\pgfqpoint{0.020833in}{0.000000in}}{%
\pgfpathmoveto{\pgfqpoint{0.000000in}{0.000000in}}%
\pgfpathlineto{\pgfqpoint{0.020833in}{0.000000in}}%
\pgfusepath{stroke,fill}%
}%
\begin{pgfscope}%
\pgfsys@transformshift{0.563118in}{3.385933in}%
\pgfsys@useobject{currentmarker}{}%
\end{pgfscope}%
\end{pgfscope}%
\begin{pgfscope}%
\pgfsetbuttcap%
\pgfsetroundjoin%
\definecolor{currentfill}{rgb}{0.000000,0.000000,0.000000}%
\pgfsetfillcolor{currentfill}%
\pgfsetlinewidth{0.501875pt}%
\definecolor{currentstroke}{rgb}{0.000000,0.000000,0.000000}%
\pgfsetstrokecolor{currentstroke}%
\pgfsetdash{}{0pt}%
\pgfsys@defobject{currentmarker}{\pgfqpoint{-0.020833in}{0.000000in}}{\pgfqpoint{-0.000000in}{0.000000in}}{%
\pgfpathmoveto{\pgfqpoint{-0.000000in}{0.000000in}}%
\pgfpathlineto{\pgfqpoint{-0.020833in}{0.000000in}}%
\pgfusepath{stroke,fill}%
}%
\begin{pgfscope}%
\pgfsys@transformshift{5.600618in}{3.385933in}%
\pgfsys@useobject{currentmarker}{}%
\end{pgfscope}%
\end{pgfscope}%
\begin{pgfscope}%
\pgfsetbuttcap%
\pgfsetroundjoin%
\definecolor{currentfill}{rgb}{0.000000,0.000000,0.000000}%
\pgfsetfillcolor{currentfill}%
\pgfsetlinewidth{0.501875pt}%
\definecolor{currentstroke}{rgb}{0.000000,0.000000,0.000000}%
\pgfsetstrokecolor{currentstroke}%
\pgfsetdash{}{0pt}%
\pgfsys@defobject{currentmarker}{\pgfqpoint{0.000000in}{0.000000in}}{\pgfqpoint{0.020833in}{0.000000in}}{%
\pgfpathmoveto{\pgfqpoint{0.000000in}{0.000000in}}%
\pgfpathlineto{\pgfqpoint{0.020833in}{0.000000in}}%
\pgfusepath{stroke,fill}%
}%
\begin{pgfscope}%
\pgfsys@transformshift{0.563118in}{3.512929in}%
\pgfsys@useobject{currentmarker}{}%
\end{pgfscope}%
\end{pgfscope}%
\begin{pgfscope}%
\pgfsetbuttcap%
\pgfsetroundjoin%
\definecolor{currentfill}{rgb}{0.000000,0.000000,0.000000}%
\pgfsetfillcolor{currentfill}%
\pgfsetlinewidth{0.501875pt}%
\definecolor{currentstroke}{rgb}{0.000000,0.000000,0.000000}%
\pgfsetstrokecolor{currentstroke}%
\pgfsetdash{}{0pt}%
\pgfsys@defobject{currentmarker}{\pgfqpoint{-0.020833in}{0.000000in}}{\pgfqpoint{-0.000000in}{0.000000in}}{%
\pgfpathmoveto{\pgfqpoint{-0.000000in}{0.000000in}}%
\pgfpathlineto{\pgfqpoint{-0.020833in}{0.000000in}}%
\pgfusepath{stroke,fill}%
}%
\begin{pgfscope}%
\pgfsys@transformshift{5.600618in}{3.512929in}%
\pgfsys@useobject{currentmarker}{}%
\end{pgfscope}%
\end{pgfscope}%
\begin{pgfscope}%
\pgfsetbuttcap%
\pgfsetroundjoin%
\definecolor{currentfill}{rgb}{0.000000,0.000000,0.000000}%
\pgfsetfillcolor{currentfill}%
\pgfsetlinewidth{0.501875pt}%
\definecolor{currentstroke}{rgb}{0.000000,0.000000,0.000000}%
\pgfsetstrokecolor{currentstroke}%
\pgfsetdash{}{0pt}%
\pgfsys@defobject{currentmarker}{\pgfqpoint{0.000000in}{0.000000in}}{\pgfqpoint{0.020833in}{0.000000in}}{%
\pgfpathmoveto{\pgfqpoint{0.000000in}{0.000000in}}%
\pgfpathlineto{\pgfqpoint{0.020833in}{0.000000in}}%
\pgfusepath{stroke,fill}%
}%
\begin{pgfscope}%
\pgfsys@transformshift{0.563118in}{3.639925in}%
\pgfsys@useobject{currentmarker}{}%
\end{pgfscope}%
\end{pgfscope}%
\begin{pgfscope}%
\pgfsetbuttcap%
\pgfsetroundjoin%
\definecolor{currentfill}{rgb}{0.000000,0.000000,0.000000}%
\pgfsetfillcolor{currentfill}%
\pgfsetlinewidth{0.501875pt}%
\definecolor{currentstroke}{rgb}{0.000000,0.000000,0.000000}%
\pgfsetstrokecolor{currentstroke}%
\pgfsetdash{}{0pt}%
\pgfsys@defobject{currentmarker}{\pgfqpoint{-0.020833in}{0.000000in}}{\pgfqpoint{-0.000000in}{0.000000in}}{%
\pgfpathmoveto{\pgfqpoint{-0.000000in}{0.000000in}}%
\pgfpathlineto{\pgfqpoint{-0.020833in}{0.000000in}}%
\pgfusepath{stroke,fill}%
}%
\begin{pgfscope}%
\pgfsys@transformshift{5.600618in}{3.639925in}%
\pgfsys@useobject{currentmarker}{}%
\end{pgfscope}%
\end{pgfscope}%
\begin{pgfscope}%
\pgfsetbuttcap%
\pgfsetroundjoin%
\definecolor{currentfill}{rgb}{0.000000,0.000000,0.000000}%
\pgfsetfillcolor{currentfill}%
\pgfsetlinewidth{0.501875pt}%
\definecolor{currentstroke}{rgb}{0.000000,0.000000,0.000000}%
\pgfsetstrokecolor{currentstroke}%
\pgfsetdash{}{0pt}%
\pgfsys@defobject{currentmarker}{\pgfqpoint{0.000000in}{0.000000in}}{\pgfqpoint{0.020833in}{0.000000in}}{%
\pgfpathmoveto{\pgfqpoint{0.000000in}{0.000000in}}%
\pgfpathlineto{\pgfqpoint{0.020833in}{0.000000in}}%
\pgfusepath{stroke,fill}%
}%
\begin{pgfscope}%
\pgfsys@transformshift{0.563118in}{3.766920in}%
\pgfsys@useobject{currentmarker}{}%
\end{pgfscope}%
\end{pgfscope}%
\begin{pgfscope}%
\pgfsetbuttcap%
\pgfsetroundjoin%
\definecolor{currentfill}{rgb}{0.000000,0.000000,0.000000}%
\pgfsetfillcolor{currentfill}%
\pgfsetlinewidth{0.501875pt}%
\definecolor{currentstroke}{rgb}{0.000000,0.000000,0.000000}%
\pgfsetstrokecolor{currentstroke}%
\pgfsetdash{}{0pt}%
\pgfsys@defobject{currentmarker}{\pgfqpoint{-0.020833in}{0.000000in}}{\pgfqpoint{-0.000000in}{0.000000in}}{%
\pgfpathmoveto{\pgfqpoint{-0.000000in}{0.000000in}}%
\pgfpathlineto{\pgfqpoint{-0.020833in}{0.000000in}}%
\pgfusepath{stroke,fill}%
}%
\begin{pgfscope}%
\pgfsys@transformshift{5.600618in}{3.766920in}%
\pgfsys@useobject{currentmarker}{}%
\end{pgfscope}%
\end{pgfscope}%
\begin{pgfscope}%
\pgfsetbuttcap%
\pgfsetroundjoin%
\definecolor{currentfill}{rgb}{0.000000,0.000000,0.000000}%
\pgfsetfillcolor{currentfill}%
\pgfsetlinewidth{0.501875pt}%
\definecolor{currentstroke}{rgb}{0.000000,0.000000,0.000000}%
\pgfsetstrokecolor{currentstroke}%
\pgfsetdash{}{0pt}%
\pgfsys@defobject{currentmarker}{\pgfqpoint{0.000000in}{0.000000in}}{\pgfqpoint{0.020833in}{0.000000in}}{%
\pgfpathmoveto{\pgfqpoint{0.000000in}{0.000000in}}%
\pgfpathlineto{\pgfqpoint{0.020833in}{0.000000in}}%
\pgfusepath{stroke,fill}%
}%
\begin{pgfscope}%
\pgfsys@transformshift{0.563118in}{4.020912in}%
\pgfsys@useobject{currentmarker}{}%
\end{pgfscope}%
\end{pgfscope}%
\begin{pgfscope}%
\pgfsetbuttcap%
\pgfsetroundjoin%
\definecolor{currentfill}{rgb}{0.000000,0.000000,0.000000}%
\pgfsetfillcolor{currentfill}%
\pgfsetlinewidth{0.501875pt}%
\definecolor{currentstroke}{rgb}{0.000000,0.000000,0.000000}%
\pgfsetstrokecolor{currentstroke}%
\pgfsetdash{}{0pt}%
\pgfsys@defobject{currentmarker}{\pgfqpoint{-0.020833in}{0.000000in}}{\pgfqpoint{-0.000000in}{0.000000in}}{%
\pgfpathmoveto{\pgfqpoint{-0.000000in}{0.000000in}}%
\pgfpathlineto{\pgfqpoint{-0.020833in}{0.000000in}}%
\pgfusepath{stroke,fill}%
}%
\begin{pgfscope}%
\pgfsys@transformshift{5.600618in}{4.020912in}%
\pgfsys@useobject{currentmarker}{}%
\end{pgfscope}%
\end{pgfscope}%
\begin{pgfscope}%
\pgfsetbuttcap%
\pgfsetroundjoin%
\definecolor{currentfill}{rgb}{0.000000,0.000000,0.000000}%
\pgfsetfillcolor{currentfill}%
\pgfsetlinewidth{0.501875pt}%
\definecolor{currentstroke}{rgb}{0.000000,0.000000,0.000000}%
\pgfsetstrokecolor{currentstroke}%
\pgfsetdash{}{0pt}%
\pgfsys@defobject{currentmarker}{\pgfqpoint{0.000000in}{0.000000in}}{\pgfqpoint{0.020833in}{0.000000in}}{%
\pgfpathmoveto{\pgfqpoint{0.000000in}{0.000000in}}%
\pgfpathlineto{\pgfqpoint{0.020833in}{0.000000in}}%
\pgfusepath{stroke,fill}%
}%
\begin{pgfscope}%
\pgfsys@transformshift{0.563118in}{4.147908in}%
\pgfsys@useobject{currentmarker}{}%
\end{pgfscope}%
\end{pgfscope}%
\begin{pgfscope}%
\pgfsetbuttcap%
\pgfsetroundjoin%
\definecolor{currentfill}{rgb}{0.000000,0.000000,0.000000}%
\pgfsetfillcolor{currentfill}%
\pgfsetlinewidth{0.501875pt}%
\definecolor{currentstroke}{rgb}{0.000000,0.000000,0.000000}%
\pgfsetstrokecolor{currentstroke}%
\pgfsetdash{}{0pt}%
\pgfsys@defobject{currentmarker}{\pgfqpoint{-0.020833in}{0.000000in}}{\pgfqpoint{-0.000000in}{0.000000in}}{%
\pgfpathmoveto{\pgfqpoint{-0.000000in}{0.000000in}}%
\pgfpathlineto{\pgfqpoint{-0.020833in}{0.000000in}}%
\pgfusepath{stroke,fill}%
}%
\begin{pgfscope}%
\pgfsys@transformshift{5.600618in}{4.147908in}%
\pgfsys@useobject{currentmarker}{}%
\end{pgfscope}%
\end{pgfscope}%
\begin{pgfscope}%
\pgfsetbuttcap%
\pgfsetroundjoin%
\definecolor{currentfill}{rgb}{0.000000,0.000000,0.000000}%
\pgfsetfillcolor{currentfill}%
\pgfsetlinewidth{0.501875pt}%
\definecolor{currentstroke}{rgb}{0.000000,0.000000,0.000000}%
\pgfsetstrokecolor{currentstroke}%
\pgfsetdash{}{0pt}%
\pgfsys@defobject{currentmarker}{\pgfqpoint{0.000000in}{0.000000in}}{\pgfqpoint{0.020833in}{0.000000in}}{%
\pgfpathmoveto{\pgfqpoint{0.000000in}{0.000000in}}%
\pgfpathlineto{\pgfqpoint{0.020833in}{0.000000in}}%
\pgfusepath{stroke,fill}%
}%
\begin{pgfscope}%
\pgfsys@transformshift{0.563118in}{4.274903in}%
\pgfsys@useobject{currentmarker}{}%
\end{pgfscope}%
\end{pgfscope}%
\begin{pgfscope}%
\pgfsetbuttcap%
\pgfsetroundjoin%
\definecolor{currentfill}{rgb}{0.000000,0.000000,0.000000}%
\pgfsetfillcolor{currentfill}%
\pgfsetlinewidth{0.501875pt}%
\definecolor{currentstroke}{rgb}{0.000000,0.000000,0.000000}%
\pgfsetstrokecolor{currentstroke}%
\pgfsetdash{}{0pt}%
\pgfsys@defobject{currentmarker}{\pgfqpoint{-0.020833in}{0.000000in}}{\pgfqpoint{-0.000000in}{0.000000in}}{%
\pgfpathmoveto{\pgfqpoint{-0.000000in}{0.000000in}}%
\pgfpathlineto{\pgfqpoint{-0.020833in}{0.000000in}}%
\pgfusepath{stroke,fill}%
}%
\begin{pgfscope}%
\pgfsys@transformshift{5.600618in}{4.274903in}%
\pgfsys@useobject{currentmarker}{}%
\end{pgfscope}%
\end{pgfscope}%
\begin{pgfscope}%
\pgfsetbuttcap%
\pgfsetroundjoin%
\definecolor{currentfill}{rgb}{0.000000,0.000000,0.000000}%
\pgfsetfillcolor{currentfill}%
\pgfsetlinewidth{0.501875pt}%
\definecolor{currentstroke}{rgb}{0.000000,0.000000,0.000000}%
\pgfsetstrokecolor{currentstroke}%
\pgfsetdash{}{0pt}%
\pgfsys@defobject{currentmarker}{\pgfqpoint{0.000000in}{0.000000in}}{\pgfqpoint{0.020833in}{0.000000in}}{%
\pgfpathmoveto{\pgfqpoint{0.000000in}{0.000000in}}%
\pgfpathlineto{\pgfqpoint{0.020833in}{0.000000in}}%
\pgfusepath{stroke,fill}%
}%
\begin{pgfscope}%
\pgfsys@transformshift{0.563118in}{4.401899in}%
\pgfsys@useobject{currentmarker}{}%
\end{pgfscope}%
\end{pgfscope}%
\begin{pgfscope}%
\pgfsetbuttcap%
\pgfsetroundjoin%
\definecolor{currentfill}{rgb}{0.000000,0.000000,0.000000}%
\pgfsetfillcolor{currentfill}%
\pgfsetlinewidth{0.501875pt}%
\definecolor{currentstroke}{rgb}{0.000000,0.000000,0.000000}%
\pgfsetstrokecolor{currentstroke}%
\pgfsetdash{}{0pt}%
\pgfsys@defobject{currentmarker}{\pgfqpoint{-0.020833in}{0.000000in}}{\pgfqpoint{-0.000000in}{0.000000in}}{%
\pgfpathmoveto{\pgfqpoint{-0.000000in}{0.000000in}}%
\pgfpathlineto{\pgfqpoint{-0.020833in}{0.000000in}}%
\pgfusepath{stroke,fill}%
}%
\begin{pgfscope}%
\pgfsys@transformshift{5.600618in}{4.401899in}%
\pgfsys@useobject{currentmarker}{}%
\end{pgfscope}%
\end{pgfscope}%
\begin{pgfscope}%
\definecolor{textcolor}{rgb}{0.000000,0.000000,0.000000}%
\pgfsetstrokecolor{textcolor}%
\pgfsetfillcolor{textcolor}%
\pgftext[x=0.173457in,y=2.403080in,,bottom,rotate=90.000000]{\color{textcolor}{\rmfamily\fontsize{10.000000}{12.000000}\selectfont\catcode`\^=\active\def^{\ifmmode\sp\else\^{}\fi}\catcode`\%=\active\def%{\%}$W_t$}}%
\end{pgfscope}%
\begin{pgfscope}%
\pgfpathrectangle{\pgfqpoint{0.563118in}{0.401080in}}{\pgfqpoint{5.037500in}{4.004000in}}%
\pgfusepath{clip}%
\pgfsetrectcap%
\pgfsetroundjoin%
\pgfsetlinewidth{0.803000pt}%
\definecolor{currentstroke}{rgb}{0.000000,0.000000,0.000000}%
\pgfsetstrokecolor{currentstroke}%
\pgfsetdash{}{0pt}%
\pgfpathmoveto{\pgfqpoint{0.563118in}{1.988980in}}%
\pgfpathlineto{\pgfqpoint{0.565637in}{2.000227in}}%
\pgfpathlineto{\pgfqpoint{0.568156in}{2.032790in}}%
\pgfpathlineto{\pgfqpoint{0.570674in}{2.013115in}}%
\pgfpathlineto{\pgfqpoint{0.573193in}{2.014755in}}%
\pgfpathlineto{\pgfqpoint{0.575712in}{2.012791in}}%
\pgfpathlineto{\pgfqpoint{0.578231in}{2.075383in}}%
\pgfpathlineto{\pgfqpoint{0.580749in}{2.114559in}}%
\pgfpathlineto{\pgfqpoint{0.583268in}{2.139475in}}%
\pgfpathlineto{\pgfqpoint{0.585787in}{2.178059in}}%
\pgfpathlineto{\pgfqpoint{0.588306in}{2.136320in}}%
\pgfpathlineto{\pgfqpoint{0.590824in}{2.133905in}}%
\pgfpathlineto{\pgfqpoint{0.593343in}{2.144136in}}%
\pgfpathlineto{\pgfqpoint{0.595862in}{2.148664in}}%
\pgfpathlineto{\pgfqpoint{0.598381in}{2.131411in}}%
\pgfpathlineto{\pgfqpoint{0.600899in}{2.129993in}}%
\pgfpathlineto{\pgfqpoint{0.603418in}{2.084396in}}%
\pgfpathlineto{\pgfqpoint{0.605937in}{2.070252in}}%
\pgfpathlineto{\pgfqpoint{0.608456in}{2.081137in}}%
\pgfpathlineto{\pgfqpoint{0.610974in}{2.083093in}}%
\pgfpathlineto{\pgfqpoint{0.613493in}{2.103588in}}%
\pgfpathlineto{\pgfqpoint{0.616012in}{2.135889in}}%
\pgfpathlineto{\pgfqpoint{0.618531in}{2.113625in}}%
\pgfpathlineto{\pgfqpoint{0.621049in}{2.121157in}}%
\pgfpathlineto{\pgfqpoint{0.623568in}{2.130937in}}%
\pgfpathlineto{\pgfqpoint{0.626087in}{2.171198in}}%
\pgfpathlineto{\pgfqpoint{0.628606in}{2.188230in}}%
\pgfpathlineto{\pgfqpoint{0.631124in}{2.194980in}}%
\pgfpathlineto{\pgfqpoint{0.633643in}{2.187389in}}%
\pgfpathlineto{\pgfqpoint{0.636162in}{2.172988in}}%
\pgfpathlineto{\pgfqpoint{0.638681in}{2.212776in}}%
\pgfpathlineto{\pgfqpoint{0.641199in}{2.218493in}}%
\pgfpathlineto{\pgfqpoint{0.643718in}{2.195635in}}%
\pgfpathlineto{\pgfqpoint{0.646237in}{2.208672in}}%
\pgfpathlineto{\pgfqpoint{0.648756in}{2.198609in}}%
\pgfpathlineto{\pgfqpoint{0.651274in}{2.191666in}}%
\pgfpathlineto{\pgfqpoint{0.653793in}{2.181537in}}%
\pgfpathlineto{\pgfqpoint{0.656312in}{2.179332in}}%
\pgfpathlineto{\pgfqpoint{0.658831in}{2.169088in}}%
\pgfpathlineto{\pgfqpoint{0.661349in}{2.206227in}}%
\pgfpathlineto{\pgfqpoint{0.663868in}{2.177495in}}%
\pgfpathlineto{\pgfqpoint{0.666387in}{2.247743in}}%
\pgfpathlineto{\pgfqpoint{0.668906in}{2.259979in}}%
\pgfpathlineto{\pgfqpoint{0.671424in}{2.289792in}}%
\pgfpathlineto{\pgfqpoint{0.673943in}{2.283404in}}%
\pgfpathlineto{\pgfqpoint{0.676462in}{2.277751in}}%
\pgfpathlineto{\pgfqpoint{0.678981in}{2.309668in}}%
\pgfpathlineto{\pgfqpoint{0.681499in}{2.243528in}}%
\pgfpathlineto{\pgfqpoint{0.684018in}{2.222351in}}%
\pgfpathlineto{\pgfqpoint{0.686537in}{2.178205in}}%
\pgfpathlineto{\pgfqpoint{0.689056in}{2.185801in}}%
\pgfpathlineto{\pgfqpoint{0.691574in}{2.154049in}}%
\pgfpathlineto{\pgfqpoint{0.694093in}{2.149775in}}%
\pgfpathlineto{\pgfqpoint{0.696612in}{2.076559in}}%
\pgfpathlineto{\pgfqpoint{0.699131in}{2.067402in}}%
\pgfpathlineto{\pgfqpoint{0.701649in}{2.038861in}}%
\pgfpathlineto{\pgfqpoint{0.704168in}{2.028612in}}%
\pgfpathlineto{\pgfqpoint{0.706687in}{2.058643in}}%
\pgfpathlineto{\pgfqpoint{0.709206in}{2.107423in}}%
\pgfpathlineto{\pgfqpoint{0.711724in}{2.117376in}}%
\pgfpathlineto{\pgfqpoint{0.714243in}{2.124021in}}%
\pgfpathlineto{\pgfqpoint{0.716762in}{2.073560in}}%
\pgfpathlineto{\pgfqpoint{0.719281in}{2.025239in}}%
\pgfpathlineto{\pgfqpoint{0.721799in}{2.055863in}}%
\pgfpathlineto{\pgfqpoint{0.724318in}{2.054779in}}%
\pgfpathlineto{\pgfqpoint{0.726837in}{2.005266in}}%
\pgfpathlineto{\pgfqpoint{0.729356in}{1.963336in}}%
\pgfpathlineto{\pgfqpoint{0.731874in}{1.993429in}}%
\pgfpathlineto{\pgfqpoint{0.734393in}{1.979822in}}%
\pgfpathlineto{\pgfqpoint{0.736912in}{2.011835in}}%
\pgfpathlineto{\pgfqpoint{0.739431in}{2.027475in}}%
\pgfpathlineto{\pgfqpoint{0.741949in}{1.980200in}}%
\pgfpathlineto{\pgfqpoint{0.744468in}{2.012227in}}%
\pgfpathlineto{\pgfqpoint{0.746987in}{2.030368in}}%
\pgfpathlineto{\pgfqpoint{0.749506in}{2.050314in}}%
\pgfpathlineto{\pgfqpoint{0.752024in}{2.064558in}}%
\pgfpathlineto{\pgfqpoint{0.754543in}{2.021068in}}%
\pgfpathlineto{\pgfqpoint{0.757062in}{1.971181in}}%
\pgfpathlineto{\pgfqpoint{0.759581in}{1.991432in}}%
\pgfpathlineto{\pgfqpoint{0.762099in}{1.978058in}}%
\pgfpathlineto{\pgfqpoint{0.764618in}{1.973215in}}%
\pgfpathlineto{\pgfqpoint{0.767137in}{1.996892in}}%
\pgfpathlineto{\pgfqpoint{0.769656in}{2.036005in}}%
\pgfpathlineto{\pgfqpoint{0.772174in}{2.049915in}}%
\pgfpathlineto{\pgfqpoint{0.774693in}{2.073084in}}%
\pgfpathlineto{\pgfqpoint{0.777212in}{2.083404in}}%
\pgfpathlineto{\pgfqpoint{0.779731in}{2.057921in}}%
\pgfpathlineto{\pgfqpoint{0.782249in}{2.029150in}}%
\pgfpathlineto{\pgfqpoint{0.784768in}{2.071628in}}%
\pgfpathlineto{\pgfqpoint{0.787287in}{2.094081in}}%
\pgfpathlineto{\pgfqpoint{0.789806in}{2.067017in}}%
\pgfpathlineto{\pgfqpoint{0.792324in}{2.112724in}}%
\pgfpathlineto{\pgfqpoint{0.794843in}{2.056292in}}%
\pgfpathlineto{\pgfqpoint{0.797362in}{2.005992in}}%
\pgfpathlineto{\pgfqpoint{0.799881in}{2.007136in}}%
\pgfpathlineto{\pgfqpoint{0.802399in}{1.978418in}}%
\pgfpathlineto{\pgfqpoint{0.804918in}{1.955867in}}%
\pgfpathlineto{\pgfqpoint{0.807437in}{1.956565in}}%
\pgfpathlineto{\pgfqpoint{0.809956in}{1.933148in}}%
\pgfpathlineto{\pgfqpoint{0.812474in}{1.935510in}}%
\pgfpathlineto{\pgfqpoint{0.814993in}{1.970497in}}%
\pgfpathlineto{\pgfqpoint{0.817512in}{2.026767in}}%
\pgfpathlineto{\pgfqpoint{0.820031in}{1.994913in}}%
\pgfpathlineto{\pgfqpoint{0.822549in}{1.944549in}}%
\pgfpathlineto{\pgfqpoint{0.825068in}{1.960566in}}%
\pgfpathlineto{\pgfqpoint{0.827587in}{1.956803in}}%
\pgfpathlineto{\pgfqpoint{0.830106in}{1.955221in}}%
\pgfpathlineto{\pgfqpoint{0.832624in}{1.964491in}}%
\pgfpathlineto{\pgfqpoint{0.835143in}{2.024232in}}%
\pgfpathlineto{\pgfqpoint{0.837662in}{2.034215in}}%
\pgfpathlineto{\pgfqpoint{0.840181in}{2.039433in}}%
\pgfpathlineto{\pgfqpoint{0.842699in}{2.027365in}}%
\pgfpathlineto{\pgfqpoint{0.845218in}{2.032149in}}%
\pgfpathlineto{\pgfqpoint{0.847737in}{2.030803in}}%
\pgfpathlineto{\pgfqpoint{0.850256in}{2.051413in}}%
\pgfpathlineto{\pgfqpoint{0.852774in}{2.033088in}}%
\pgfpathlineto{\pgfqpoint{0.855293in}{2.019016in}}%
\pgfpathlineto{\pgfqpoint{0.857812in}{2.019462in}}%
\pgfpathlineto{\pgfqpoint{0.860331in}{1.988729in}}%
\pgfpathlineto{\pgfqpoint{0.862849in}{1.937145in}}%
\pgfpathlineto{\pgfqpoint{0.865368in}{1.944819in}}%
\pgfpathlineto{\pgfqpoint{0.867887in}{1.934916in}}%
\pgfpathlineto{\pgfqpoint{0.870406in}{1.967928in}}%
\pgfpathlineto{\pgfqpoint{0.872924in}{1.986591in}}%
\pgfpathlineto{\pgfqpoint{0.875443in}{2.045509in}}%
\pgfpathlineto{\pgfqpoint{0.877962in}{1.995627in}}%
\pgfpathlineto{\pgfqpoint{0.880481in}{1.996162in}}%
\pgfpathlineto{\pgfqpoint{0.882999in}{2.015002in}}%
\pgfpathlineto{\pgfqpoint{0.885518in}{1.999093in}}%
\pgfpathlineto{\pgfqpoint{0.888037in}{1.982079in}}%
\pgfpathlineto{\pgfqpoint{0.890556in}{1.942973in}}%
\pgfpathlineto{\pgfqpoint{0.893074in}{1.939981in}}%
\pgfpathlineto{\pgfqpoint{0.895593in}{1.914013in}}%
\pgfpathlineto{\pgfqpoint{0.898112in}{1.951403in}}%
\pgfpathlineto{\pgfqpoint{0.900631in}{1.943133in}}%
\pgfpathlineto{\pgfqpoint{0.903149in}{1.942053in}}%
\pgfpathlineto{\pgfqpoint{0.905668in}{1.899300in}}%
\pgfpathlineto{\pgfqpoint{0.908187in}{1.871516in}}%
\pgfpathlineto{\pgfqpoint{0.910706in}{1.894569in}}%
\pgfpathlineto{\pgfqpoint{0.913224in}{1.902775in}}%
\pgfpathlineto{\pgfqpoint{0.915743in}{1.878016in}}%
\pgfpathlineto{\pgfqpoint{0.918262in}{1.837537in}}%
\pgfpathlineto{\pgfqpoint{0.920781in}{1.794606in}}%
\pgfpathlineto{\pgfqpoint{0.923299in}{1.762857in}}%
\pgfpathlineto{\pgfqpoint{0.925818in}{1.783855in}}%
\pgfpathlineto{\pgfqpoint{0.928337in}{1.788060in}}%
\pgfpathlineto{\pgfqpoint{0.930856in}{1.791754in}}%
\pgfpathlineto{\pgfqpoint{0.933374in}{1.750664in}}%
\pgfpathlineto{\pgfqpoint{0.935893in}{1.717667in}}%
\pgfpathlineto{\pgfqpoint{0.940931in}{1.750255in}}%
\pgfpathlineto{\pgfqpoint{0.943449in}{1.752637in}}%
\pgfpathlineto{\pgfqpoint{0.945968in}{1.802967in}}%
\pgfpathlineto{\pgfqpoint{0.948487in}{1.818843in}}%
\pgfpathlineto{\pgfqpoint{0.951006in}{1.822323in}}%
\pgfpathlineto{\pgfqpoint{0.953524in}{1.839405in}}%
\pgfpathlineto{\pgfqpoint{0.956043in}{1.818865in}}%
\pgfpathlineto{\pgfqpoint{0.958562in}{1.837008in}}%
\pgfpathlineto{\pgfqpoint{0.961081in}{1.877249in}}%
\pgfpathlineto{\pgfqpoint{0.963599in}{1.835010in}}%
\pgfpathlineto{\pgfqpoint{0.966118in}{1.840217in}}%
\pgfpathlineto{\pgfqpoint{0.968637in}{1.812743in}}%
\pgfpathlineto{\pgfqpoint{0.971156in}{1.836036in}}%
\pgfpathlineto{\pgfqpoint{0.973674in}{1.822285in}}%
\pgfpathlineto{\pgfqpoint{0.976193in}{1.831327in}}%
\pgfpathlineto{\pgfqpoint{0.978712in}{1.821810in}}%
\pgfpathlineto{\pgfqpoint{0.981231in}{1.825253in}}%
\pgfpathlineto{\pgfqpoint{0.983749in}{1.842382in}}%
\pgfpathlineto{\pgfqpoint{0.986268in}{1.849396in}}%
\pgfpathlineto{\pgfqpoint{0.988787in}{1.840004in}}%
\pgfpathlineto{\pgfqpoint{0.991306in}{1.902813in}}%
\pgfpathlineto{\pgfqpoint{0.993824in}{1.871491in}}%
\pgfpathlineto{\pgfqpoint{0.996343in}{1.880734in}}%
\pgfpathlineto{\pgfqpoint{0.998862in}{1.905782in}}%
\pgfpathlineto{\pgfqpoint{1.001381in}{1.881622in}}%
\pgfpathlineto{\pgfqpoint{1.003899in}{1.882673in}}%
\pgfpathlineto{\pgfqpoint{1.006418in}{1.873149in}}%
\pgfpathlineto{\pgfqpoint{1.008937in}{1.855763in}}%
\pgfpathlineto{\pgfqpoint{1.011456in}{1.859535in}}%
\pgfpathlineto{\pgfqpoint{1.013974in}{1.840421in}}%
\pgfpathlineto{\pgfqpoint{1.016493in}{1.836363in}}%
\pgfpathlineto{\pgfqpoint{1.019012in}{1.763046in}}%
\pgfpathlineto{\pgfqpoint{1.021531in}{1.758003in}}%
\pgfpathlineto{\pgfqpoint{1.024049in}{1.739078in}}%
\pgfpathlineto{\pgfqpoint{1.026568in}{1.736572in}}%
\pgfpathlineto{\pgfqpoint{1.029087in}{1.725123in}}%
\pgfpathlineto{\pgfqpoint{1.031606in}{1.711098in}}%
\pgfpathlineto{\pgfqpoint{1.034124in}{1.708088in}}%
\pgfpathlineto{\pgfqpoint{1.036643in}{1.710364in}}%
\pgfpathlineto{\pgfqpoint{1.039162in}{1.741397in}}%
\pgfpathlineto{\pgfqpoint{1.041681in}{1.737652in}}%
\pgfpathlineto{\pgfqpoint{1.044199in}{1.718255in}}%
\pgfpathlineto{\pgfqpoint{1.046718in}{1.740868in}}%
\pgfpathlineto{\pgfqpoint{1.049237in}{1.704364in}}%
\pgfpathlineto{\pgfqpoint{1.051756in}{1.722635in}}%
\pgfpathlineto{\pgfqpoint{1.054274in}{1.689963in}}%
\pgfpathlineto{\pgfqpoint{1.056793in}{1.716690in}}%
\pgfpathlineto{\pgfqpoint{1.059312in}{1.718419in}}%
\pgfpathlineto{\pgfqpoint{1.061831in}{1.695060in}}%
\pgfpathlineto{\pgfqpoint{1.064349in}{1.682930in}}%
\pgfpathlineto{\pgfqpoint{1.066868in}{1.781400in}}%
\pgfpathlineto{\pgfqpoint{1.069387in}{1.791899in}}%
\pgfpathlineto{\pgfqpoint{1.071906in}{1.807398in}}%
\pgfpathlineto{\pgfqpoint{1.074424in}{1.862683in}}%
\pgfpathlineto{\pgfqpoint{1.076943in}{1.788511in}}%
\pgfpathlineto{\pgfqpoint{1.079462in}{1.812670in}}%
\pgfpathlineto{\pgfqpoint{1.081981in}{1.854634in}}%
\pgfpathlineto{\pgfqpoint{1.084499in}{1.852622in}}%
\pgfpathlineto{\pgfqpoint{1.087018in}{1.802112in}}%
\pgfpathlineto{\pgfqpoint{1.089537in}{1.815390in}}%
\pgfpathlineto{\pgfqpoint{1.092056in}{1.741522in}}%
\pgfpathlineto{\pgfqpoint{1.094574in}{1.728432in}}%
\pgfpathlineto{\pgfqpoint{1.097093in}{1.688306in}}%
\pgfpathlineto{\pgfqpoint{1.099612in}{1.713155in}}%
\pgfpathlineto{\pgfqpoint{1.102131in}{1.748715in}}%
\pgfpathlineto{\pgfqpoint{1.104649in}{1.751772in}}%
\pgfpathlineto{\pgfqpoint{1.107168in}{1.726360in}}%
\pgfpathlineto{\pgfqpoint{1.109687in}{1.751809in}}%
\pgfpathlineto{\pgfqpoint{1.112206in}{1.789085in}}%
\pgfpathlineto{\pgfqpoint{1.114724in}{1.836626in}}%
\pgfpathlineto{\pgfqpoint{1.117243in}{1.836803in}}%
\pgfpathlineto{\pgfqpoint{1.119762in}{1.828579in}}%
\pgfpathlineto{\pgfqpoint{1.122281in}{1.892734in}}%
\pgfpathlineto{\pgfqpoint{1.124799in}{1.887535in}}%
\pgfpathlineto{\pgfqpoint{1.127318in}{1.894481in}}%
\pgfpathlineto{\pgfqpoint{1.129837in}{1.929447in}}%
\pgfpathlineto{\pgfqpoint{1.132356in}{1.975583in}}%
\pgfpathlineto{\pgfqpoint{1.134874in}{1.995835in}}%
\pgfpathlineto{\pgfqpoint{1.137393in}{1.979285in}}%
\pgfpathlineto{\pgfqpoint{1.139912in}{1.961918in}}%
\pgfpathlineto{\pgfqpoint{1.142431in}{1.986517in}}%
\pgfpathlineto{\pgfqpoint{1.144949in}{1.943162in}}%
\pgfpathlineto{\pgfqpoint{1.147468in}{1.992056in}}%
\pgfpathlineto{\pgfqpoint{1.149987in}{1.965019in}}%
\pgfpathlineto{\pgfqpoint{1.152506in}{1.931779in}}%
\pgfpathlineto{\pgfqpoint{1.155024in}{1.917035in}}%
\pgfpathlineto{\pgfqpoint{1.157543in}{1.984278in}}%
\pgfpathlineto{\pgfqpoint{1.160062in}{1.956031in}}%
\pgfpathlineto{\pgfqpoint{1.162581in}{1.931328in}}%
\pgfpathlineto{\pgfqpoint{1.165099in}{1.916244in}}%
\pgfpathlineto{\pgfqpoint{1.167618in}{1.911968in}}%
\pgfpathlineto{\pgfqpoint{1.170137in}{1.917779in}}%
\pgfpathlineto{\pgfqpoint{1.172656in}{1.872329in}}%
\pgfpathlineto{\pgfqpoint{1.175174in}{1.852160in}}%
\pgfpathlineto{\pgfqpoint{1.177693in}{1.887298in}}%
\pgfpathlineto{\pgfqpoint{1.180212in}{1.853869in}}%
\pgfpathlineto{\pgfqpoint{1.182731in}{1.883886in}}%
\pgfpathlineto{\pgfqpoint{1.185249in}{1.877619in}}%
\pgfpathlineto{\pgfqpoint{1.187768in}{1.894482in}}%
\pgfpathlineto{\pgfqpoint{1.190287in}{1.861501in}}%
\pgfpathlineto{\pgfqpoint{1.192806in}{1.850490in}}%
\pgfpathlineto{\pgfqpoint{1.195324in}{1.866497in}}%
\pgfpathlineto{\pgfqpoint{1.197843in}{1.884089in}}%
\pgfpathlineto{\pgfqpoint{1.200362in}{1.892807in}}%
\pgfpathlineto{\pgfqpoint{1.202881in}{1.869636in}}%
\pgfpathlineto{\pgfqpoint{1.205399in}{1.861318in}}%
\pgfpathlineto{\pgfqpoint{1.207918in}{1.844194in}}%
\pgfpathlineto{\pgfqpoint{1.210437in}{1.837067in}}%
\pgfpathlineto{\pgfqpoint{1.212956in}{1.834762in}}%
\pgfpathlineto{\pgfqpoint{1.215474in}{1.819049in}}%
\pgfpathlineto{\pgfqpoint{1.217993in}{1.801566in}}%
\pgfpathlineto{\pgfqpoint{1.220512in}{1.813360in}}%
\pgfpathlineto{\pgfqpoint{1.223031in}{1.774974in}}%
\pgfpathlineto{\pgfqpoint{1.225549in}{1.752847in}}%
\pgfpathlineto{\pgfqpoint{1.228068in}{1.745741in}}%
\pgfpathlineto{\pgfqpoint{1.230587in}{1.775028in}}%
\pgfpathlineto{\pgfqpoint{1.233106in}{1.753197in}}%
\pgfpathlineto{\pgfqpoint{1.235624in}{1.787243in}}%
\pgfpathlineto{\pgfqpoint{1.238143in}{1.780072in}}%
\pgfpathlineto{\pgfqpoint{1.240662in}{1.841167in}}%
\pgfpathlineto{\pgfqpoint{1.243181in}{1.880720in}}%
\pgfpathlineto{\pgfqpoint{1.245699in}{1.891058in}}%
\pgfpathlineto{\pgfqpoint{1.248218in}{1.883145in}}%
\pgfpathlineto{\pgfqpoint{1.250737in}{1.856745in}}%
\pgfpathlineto{\pgfqpoint{1.253256in}{1.882759in}}%
\pgfpathlineto{\pgfqpoint{1.255774in}{1.877683in}}%
\pgfpathlineto{\pgfqpoint{1.258293in}{1.893661in}}%
\pgfpathlineto{\pgfqpoint{1.260812in}{1.928577in}}%
\pgfpathlineto{\pgfqpoint{1.263331in}{1.945322in}}%
\pgfpathlineto{\pgfqpoint{1.265849in}{1.881651in}}%
\pgfpathlineto{\pgfqpoint{1.268368in}{1.888689in}}%
\pgfpathlineto{\pgfqpoint{1.270887in}{1.833620in}}%
\pgfpathlineto{\pgfqpoint{1.273406in}{1.875113in}}%
\pgfpathlineto{\pgfqpoint{1.275924in}{1.865744in}}%
\pgfpathlineto{\pgfqpoint{1.278443in}{1.840127in}}%
\pgfpathlineto{\pgfqpoint{1.280962in}{1.880383in}}%
\pgfpathlineto{\pgfqpoint{1.283481in}{1.847248in}}%
\pgfpathlineto{\pgfqpoint{1.285999in}{1.779724in}}%
\pgfpathlineto{\pgfqpoint{1.288518in}{1.763859in}}%
\pgfpathlineto{\pgfqpoint{1.291037in}{1.748830in}}%
\pgfpathlineto{\pgfqpoint{1.293556in}{1.735412in}}%
\pgfpathlineto{\pgfqpoint{1.296074in}{1.725584in}}%
\pgfpathlineto{\pgfqpoint{1.298593in}{1.772720in}}%
\pgfpathlineto{\pgfqpoint{1.301112in}{1.747155in}}%
\pgfpathlineto{\pgfqpoint{1.303631in}{1.736182in}}%
\pgfpathlineto{\pgfqpoint{1.306149in}{1.714470in}}%
\pgfpathlineto{\pgfqpoint{1.308668in}{1.742413in}}%
\pgfpathlineto{\pgfqpoint{1.311187in}{1.733599in}}%
\pgfpathlineto{\pgfqpoint{1.313706in}{1.719663in}}%
\pgfpathlineto{\pgfqpoint{1.316224in}{1.728095in}}%
\pgfpathlineto{\pgfqpoint{1.318743in}{1.709778in}}%
\pgfpathlineto{\pgfqpoint{1.321262in}{1.695735in}}%
\pgfpathlineto{\pgfqpoint{1.323781in}{1.735546in}}%
\pgfpathlineto{\pgfqpoint{1.326299in}{1.710828in}}%
\pgfpathlineto{\pgfqpoint{1.328818in}{1.699038in}}%
\pgfpathlineto{\pgfqpoint{1.331337in}{1.681351in}}%
\pgfpathlineto{\pgfqpoint{1.333856in}{1.657001in}}%
\pgfpathlineto{\pgfqpoint{1.336374in}{1.662358in}}%
\pgfpathlineto{\pgfqpoint{1.338893in}{1.619645in}}%
\pgfpathlineto{\pgfqpoint{1.341412in}{1.656999in}}%
\pgfpathlineto{\pgfqpoint{1.343931in}{1.697956in}}%
\pgfpathlineto{\pgfqpoint{1.346449in}{1.692859in}}%
\pgfpathlineto{\pgfqpoint{1.348968in}{1.673194in}}%
\pgfpathlineto{\pgfqpoint{1.351487in}{1.740587in}}%
\pgfpathlineto{\pgfqpoint{1.354006in}{1.751098in}}%
\pgfpathlineto{\pgfqpoint{1.356524in}{1.736585in}}%
\pgfpathlineto{\pgfqpoint{1.359043in}{1.776190in}}%
\pgfpathlineto{\pgfqpoint{1.361562in}{1.750423in}}%
\pgfpathlineto{\pgfqpoint{1.364081in}{1.728368in}}%
\pgfpathlineto{\pgfqpoint{1.366599in}{1.736423in}}%
\pgfpathlineto{\pgfqpoint{1.369118in}{1.728993in}}%
\pgfpathlineto{\pgfqpoint{1.371637in}{1.723811in}}%
\pgfpathlineto{\pgfqpoint{1.374156in}{1.713402in}}%
\pgfpathlineto{\pgfqpoint{1.376674in}{1.718940in}}%
\pgfpathlineto{\pgfqpoint{1.379193in}{1.732573in}}%
\pgfpathlineto{\pgfqpoint{1.384231in}{1.646187in}}%
\pgfpathlineto{\pgfqpoint{1.386749in}{1.633417in}}%
\pgfpathlineto{\pgfqpoint{1.389268in}{1.645657in}}%
\pgfpathlineto{\pgfqpoint{1.391787in}{1.649221in}}%
\pgfpathlineto{\pgfqpoint{1.394306in}{1.720166in}}%
\pgfpathlineto{\pgfqpoint{1.396824in}{1.722605in}}%
\pgfpathlineto{\pgfqpoint{1.399343in}{1.724558in}}%
\pgfpathlineto{\pgfqpoint{1.401862in}{1.728019in}}%
\pgfpathlineto{\pgfqpoint{1.404381in}{1.710118in}}%
\pgfpathlineto{\pgfqpoint{1.406899in}{1.735179in}}%
\pgfpathlineto{\pgfqpoint{1.409418in}{1.691135in}}%
\pgfpathlineto{\pgfqpoint{1.411937in}{1.719304in}}%
\pgfpathlineto{\pgfqpoint{1.414456in}{1.775916in}}%
\pgfpathlineto{\pgfqpoint{1.416974in}{1.742021in}}%
\pgfpathlineto{\pgfqpoint{1.419493in}{1.748100in}}%
\pgfpathlineto{\pgfqpoint{1.422012in}{1.726031in}}%
\pgfpathlineto{\pgfqpoint{1.424531in}{1.712793in}}%
\pgfpathlineto{\pgfqpoint{1.427049in}{1.714693in}}%
\pgfpathlineto{\pgfqpoint{1.429568in}{1.716989in}}%
\pgfpathlineto{\pgfqpoint{1.432087in}{1.726308in}}%
\pgfpathlineto{\pgfqpoint{1.434606in}{1.755699in}}%
\pgfpathlineto{\pgfqpoint{1.437124in}{1.775079in}}%
\pgfpathlineto{\pgfqpoint{1.439643in}{1.753901in}}%
\pgfpathlineto{\pgfqpoint{1.442162in}{1.751746in}}%
\pgfpathlineto{\pgfqpoint{1.444681in}{1.732341in}}%
\pgfpathlineto{\pgfqpoint{1.447199in}{1.707363in}}%
\pgfpathlineto{\pgfqpoint{1.449718in}{1.713539in}}%
\pgfpathlineto{\pgfqpoint{1.452237in}{1.728955in}}%
\pgfpathlineto{\pgfqpoint{1.454756in}{1.753033in}}%
\pgfpathlineto{\pgfqpoint{1.457274in}{1.796872in}}%
\pgfpathlineto{\pgfqpoint{1.459793in}{1.837079in}}%
\pgfpathlineto{\pgfqpoint{1.462312in}{1.873086in}}%
\pgfpathlineto{\pgfqpoint{1.464831in}{1.863542in}}%
\pgfpathlineto{\pgfqpoint{1.467349in}{1.800906in}}%
\pgfpathlineto{\pgfqpoint{1.469868in}{1.816928in}}%
\pgfpathlineto{\pgfqpoint{1.472387in}{1.839623in}}%
\pgfpathlineto{\pgfqpoint{1.474906in}{1.851710in}}%
\pgfpathlineto{\pgfqpoint{1.477424in}{1.853528in}}%
\pgfpathlineto{\pgfqpoint{1.479943in}{1.905380in}}%
\pgfpathlineto{\pgfqpoint{1.482462in}{1.906875in}}%
\pgfpathlineto{\pgfqpoint{1.484981in}{1.896622in}}%
\pgfpathlineto{\pgfqpoint{1.487499in}{1.868497in}}%
\pgfpathlineto{\pgfqpoint{1.490018in}{1.881610in}}%
\pgfpathlineto{\pgfqpoint{1.492537in}{1.862979in}}%
\pgfpathlineto{\pgfqpoint{1.495056in}{1.876846in}}%
\pgfpathlineto{\pgfqpoint{1.497574in}{1.877372in}}%
\pgfpathlineto{\pgfqpoint{1.500093in}{1.875376in}}%
\pgfpathlineto{\pgfqpoint{1.502612in}{1.866881in}}%
\pgfpathlineto{\pgfqpoint{1.505131in}{1.885937in}}%
\pgfpathlineto{\pgfqpoint{1.507649in}{1.887932in}}%
\pgfpathlineto{\pgfqpoint{1.510168in}{1.931537in}}%
\pgfpathlineto{\pgfqpoint{1.512687in}{1.940665in}}%
\pgfpathlineto{\pgfqpoint{1.515206in}{1.948076in}}%
\pgfpathlineto{\pgfqpoint{1.517724in}{1.876806in}}%
\pgfpathlineto{\pgfqpoint{1.520243in}{1.878061in}}%
\pgfpathlineto{\pgfqpoint{1.522762in}{1.844311in}}%
\pgfpathlineto{\pgfqpoint{1.525281in}{1.832840in}}%
\pgfpathlineto{\pgfqpoint{1.527799in}{1.858594in}}%
\pgfpathlineto{\pgfqpoint{1.530318in}{1.868757in}}%
\pgfpathlineto{\pgfqpoint{1.532837in}{1.835498in}}%
\pgfpathlineto{\pgfqpoint{1.535356in}{1.797569in}}%
\pgfpathlineto{\pgfqpoint{1.537874in}{1.808967in}}%
\pgfpathlineto{\pgfqpoint{1.540393in}{1.830906in}}%
\pgfpathlineto{\pgfqpoint{1.542912in}{1.810586in}}%
\pgfpathlineto{\pgfqpoint{1.545431in}{1.819465in}}%
\pgfpathlineto{\pgfqpoint{1.547949in}{1.802719in}}%
\pgfpathlineto{\pgfqpoint{1.550468in}{1.784438in}}%
\pgfpathlineto{\pgfqpoint{1.552987in}{1.744690in}}%
\pgfpathlineto{\pgfqpoint{1.555506in}{1.723760in}}%
\pgfpathlineto{\pgfqpoint{1.558024in}{1.755438in}}%
\pgfpathlineto{\pgfqpoint{1.560543in}{1.764121in}}%
\pgfpathlineto{\pgfqpoint{1.563062in}{1.742234in}}%
\pgfpathlineto{\pgfqpoint{1.565581in}{1.726668in}}%
\pgfpathlineto{\pgfqpoint{1.568099in}{1.702157in}}%
\pgfpathlineto{\pgfqpoint{1.570618in}{1.720850in}}%
\pgfpathlineto{\pgfqpoint{1.573137in}{1.688091in}}%
\pgfpathlineto{\pgfqpoint{1.575656in}{1.661669in}}%
\pgfpathlineto{\pgfqpoint{1.578174in}{1.627354in}}%
\pgfpathlineto{\pgfqpoint{1.580693in}{1.635664in}}%
\pgfpathlineto{\pgfqpoint{1.583212in}{1.668206in}}%
\pgfpathlineto{\pgfqpoint{1.585731in}{1.649057in}}%
\pgfpathlineto{\pgfqpoint{1.588249in}{1.673611in}}%
\pgfpathlineto{\pgfqpoint{1.590768in}{1.718096in}}%
\pgfpathlineto{\pgfqpoint{1.593287in}{1.730127in}}%
\pgfpathlineto{\pgfqpoint{1.595806in}{1.736004in}}%
\pgfpathlineto{\pgfqpoint{1.598324in}{1.743105in}}%
\pgfpathlineto{\pgfqpoint{1.600843in}{1.739437in}}%
\pgfpathlineto{\pgfqpoint{1.603362in}{1.756192in}}%
\pgfpathlineto{\pgfqpoint{1.605881in}{1.743724in}}%
\pgfpathlineto{\pgfqpoint{1.608399in}{1.652955in}}%
\pgfpathlineto{\pgfqpoint{1.615956in}{1.597892in}}%
\pgfpathlineto{\pgfqpoint{1.618474in}{1.591655in}}%
\pgfpathlineto{\pgfqpoint{1.620993in}{1.530507in}}%
\pgfpathlineto{\pgfqpoint{1.623512in}{1.498491in}}%
\pgfpathlineto{\pgfqpoint{1.626031in}{1.496886in}}%
\pgfpathlineto{\pgfqpoint{1.628549in}{1.500663in}}%
\pgfpathlineto{\pgfqpoint{1.631068in}{1.499213in}}%
\pgfpathlineto{\pgfqpoint{1.633587in}{1.506863in}}%
\pgfpathlineto{\pgfqpoint{1.636106in}{1.531438in}}%
\pgfpathlineto{\pgfqpoint{1.638624in}{1.544616in}}%
\pgfpathlineto{\pgfqpoint{1.641143in}{1.572978in}}%
\pgfpathlineto{\pgfqpoint{1.643662in}{1.581733in}}%
\pgfpathlineto{\pgfqpoint{1.646181in}{1.608142in}}%
\pgfpathlineto{\pgfqpoint{1.648699in}{1.624597in}}%
\pgfpathlineto{\pgfqpoint{1.651218in}{1.647142in}}%
\pgfpathlineto{\pgfqpoint{1.653737in}{1.621488in}}%
\pgfpathlineto{\pgfqpoint{1.656256in}{1.630148in}}%
\pgfpathlineto{\pgfqpoint{1.658774in}{1.582456in}}%
\pgfpathlineto{\pgfqpoint{1.661293in}{1.569550in}}%
\pgfpathlineto{\pgfqpoint{1.663812in}{1.543373in}}%
\pgfpathlineto{\pgfqpoint{1.666331in}{1.538720in}}%
\pgfpathlineto{\pgfqpoint{1.668849in}{1.545754in}}%
\pgfpathlineto{\pgfqpoint{1.671368in}{1.543120in}}%
\pgfpathlineto{\pgfqpoint{1.673887in}{1.565227in}}%
\pgfpathlineto{\pgfqpoint{1.676406in}{1.584791in}}%
\pgfpathlineto{\pgfqpoint{1.678924in}{1.606989in}}%
\pgfpathlineto{\pgfqpoint{1.681443in}{1.606580in}}%
\pgfpathlineto{\pgfqpoint{1.683962in}{1.574632in}}%
\pgfpathlineto{\pgfqpoint{1.686481in}{1.557820in}}%
\pgfpathlineto{\pgfqpoint{1.688999in}{1.531738in}}%
\pgfpathlineto{\pgfqpoint{1.691518in}{1.513194in}}%
\pgfpathlineto{\pgfqpoint{1.694037in}{1.498233in}}%
\pgfpathlineto{\pgfqpoint{1.696556in}{1.480134in}}%
\pgfpathlineto{\pgfqpoint{1.699074in}{1.469593in}}%
\pgfpathlineto{\pgfqpoint{1.701593in}{1.509030in}}%
\pgfpathlineto{\pgfqpoint{1.704112in}{1.441019in}}%
\pgfpathlineto{\pgfqpoint{1.706631in}{1.430936in}}%
\pgfpathlineto{\pgfqpoint{1.709149in}{1.445679in}}%
\pgfpathlineto{\pgfqpoint{1.711668in}{1.400163in}}%
\pgfpathlineto{\pgfqpoint{1.714187in}{1.361731in}}%
\pgfpathlineto{\pgfqpoint{1.716706in}{1.332333in}}%
\pgfpathlineto{\pgfqpoint{1.719224in}{1.311626in}}%
\pgfpathlineto{\pgfqpoint{1.721743in}{1.274926in}}%
\pgfpathlineto{\pgfqpoint{1.724262in}{1.258770in}}%
\pgfpathlineto{\pgfqpoint{1.726781in}{1.248285in}}%
\pgfpathlineto{\pgfqpoint{1.729299in}{1.279870in}}%
\pgfpathlineto{\pgfqpoint{1.731818in}{1.271129in}}%
\pgfpathlineto{\pgfqpoint{1.734337in}{1.286076in}}%
\pgfpathlineto{\pgfqpoint{1.736856in}{1.257849in}}%
\pgfpathlineto{\pgfqpoint{1.739374in}{1.276137in}}%
\pgfpathlineto{\pgfqpoint{1.741893in}{1.276188in}}%
\pgfpathlineto{\pgfqpoint{1.744412in}{1.250771in}}%
\pgfpathlineto{\pgfqpoint{1.746931in}{1.228820in}}%
\pgfpathlineto{\pgfqpoint{1.749449in}{1.223285in}}%
\pgfpathlineto{\pgfqpoint{1.751968in}{1.228934in}}%
\pgfpathlineto{\pgfqpoint{1.754487in}{1.254113in}}%
\pgfpathlineto{\pgfqpoint{1.757006in}{1.277198in}}%
\pgfpathlineto{\pgfqpoint{1.759524in}{1.307971in}}%
\pgfpathlineto{\pgfqpoint{1.762043in}{1.283971in}}%
\pgfpathlineto{\pgfqpoint{1.764562in}{1.293600in}}%
\pgfpathlineto{\pgfqpoint{1.767081in}{1.278916in}}%
\pgfpathlineto{\pgfqpoint{1.769599in}{1.267515in}}%
\pgfpathlineto{\pgfqpoint{1.772118in}{1.324068in}}%
\pgfpathlineto{\pgfqpoint{1.774637in}{1.327665in}}%
\pgfpathlineto{\pgfqpoint{1.777156in}{1.338506in}}%
\pgfpathlineto{\pgfqpoint{1.779674in}{1.327853in}}%
\pgfpathlineto{\pgfqpoint{1.782193in}{1.363993in}}%
\pgfpathlineto{\pgfqpoint{1.784712in}{1.358518in}}%
\pgfpathlineto{\pgfqpoint{1.787231in}{1.298941in}}%
\pgfpathlineto{\pgfqpoint{1.789749in}{1.303767in}}%
\pgfpathlineto{\pgfqpoint{1.792268in}{1.318026in}}%
\pgfpathlineto{\pgfqpoint{1.794787in}{1.288032in}}%
\pgfpathlineto{\pgfqpoint{1.797306in}{1.255277in}}%
\pgfpathlineto{\pgfqpoint{1.799824in}{1.244972in}}%
\pgfpathlineto{\pgfqpoint{1.802343in}{1.233409in}}%
\pgfpathlineto{\pgfqpoint{1.804862in}{1.206800in}}%
\pgfpathlineto{\pgfqpoint{1.807381in}{1.216185in}}%
\pgfpathlineto{\pgfqpoint{1.809899in}{1.194273in}}%
\pgfpathlineto{\pgfqpoint{1.812418in}{1.161353in}}%
\pgfpathlineto{\pgfqpoint{1.814937in}{1.197544in}}%
\pgfpathlineto{\pgfqpoint{1.817456in}{1.226169in}}%
\pgfpathlineto{\pgfqpoint{1.819974in}{1.233607in}}%
\pgfpathlineto{\pgfqpoint{1.822493in}{1.210907in}}%
\pgfpathlineto{\pgfqpoint{1.825012in}{1.200437in}}%
\pgfpathlineto{\pgfqpoint{1.827531in}{1.178470in}}%
\pgfpathlineto{\pgfqpoint{1.830049in}{1.200534in}}%
\pgfpathlineto{\pgfqpoint{1.832568in}{1.239221in}}%
\pgfpathlineto{\pgfqpoint{1.835087in}{1.244780in}}%
\pgfpathlineto{\pgfqpoint{1.837606in}{1.228585in}}%
\pgfpathlineto{\pgfqpoint{1.840124in}{1.211424in}}%
\pgfpathlineto{\pgfqpoint{1.842643in}{1.259178in}}%
\pgfpathlineto{\pgfqpoint{1.845162in}{1.252142in}}%
\pgfpathlineto{\pgfqpoint{1.847681in}{1.245956in}}%
\pgfpathlineto{\pgfqpoint{1.850199in}{1.277871in}}%
\pgfpathlineto{\pgfqpoint{1.852718in}{1.274448in}}%
\pgfpathlineto{\pgfqpoint{1.855237in}{1.257849in}}%
\pgfpathlineto{\pgfqpoint{1.857756in}{1.288373in}}%
\pgfpathlineto{\pgfqpoint{1.860274in}{1.259627in}}%
\pgfpathlineto{\pgfqpoint{1.862793in}{1.268881in}}%
\pgfpathlineto{\pgfqpoint{1.865312in}{1.314371in}}%
\pgfpathlineto{\pgfqpoint{1.867831in}{1.287846in}}%
\pgfpathlineto{\pgfqpoint{1.870349in}{1.315250in}}%
\pgfpathlineto{\pgfqpoint{1.872868in}{1.338707in}}%
\pgfpathlineto{\pgfqpoint{1.877906in}{1.291662in}}%
\pgfpathlineto{\pgfqpoint{1.880424in}{1.327391in}}%
\pgfpathlineto{\pgfqpoint{1.882943in}{1.334180in}}%
\pgfpathlineto{\pgfqpoint{1.885462in}{1.295039in}}%
\pgfpathlineto{\pgfqpoint{1.887981in}{1.304025in}}%
\pgfpathlineto{\pgfqpoint{1.890499in}{1.300065in}}%
\pgfpathlineto{\pgfqpoint{1.893018in}{1.324597in}}%
\pgfpathlineto{\pgfqpoint{1.895537in}{1.290592in}}%
\pgfpathlineto{\pgfqpoint{1.898056in}{1.273413in}}%
\pgfpathlineto{\pgfqpoint{1.900574in}{1.201764in}}%
\pgfpathlineto{\pgfqpoint{1.903093in}{1.197163in}}%
\pgfpathlineto{\pgfqpoint{1.905612in}{1.196239in}}%
\pgfpathlineto{\pgfqpoint{1.908131in}{1.179390in}}%
\pgfpathlineto{\pgfqpoint{1.910649in}{1.177518in}}%
\pgfpathlineto{\pgfqpoint{1.913168in}{1.179827in}}%
\pgfpathlineto{\pgfqpoint{1.915687in}{1.192146in}}%
\pgfpathlineto{\pgfqpoint{1.918206in}{1.223561in}}%
\pgfpathlineto{\pgfqpoint{1.920724in}{1.197814in}}%
\pgfpathlineto{\pgfqpoint{1.923243in}{1.196240in}}%
\pgfpathlineto{\pgfqpoint{1.925762in}{1.130636in}}%
\pgfpathlineto{\pgfqpoint{1.928281in}{1.118172in}}%
\pgfpathlineto{\pgfqpoint{1.930799in}{1.048450in}}%
\pgfpathlineto{\pgfqpoint{1.933318in}{0.973641in}}%
\pgfpathlineto{\pgfqpoint{1.935837in}{0.989653in}}%
\pgfpathlineto{\pgfqpoint{1.938356in}{0.972507in}}%
\pgfpathlineto{\pgfqpoint{1.940874in}{0.971682in}}%
\pgfpathlineto{\pgfqpoint{1.943393in}{0.980969in}}%
\pgfpathlineto{\pgfqpoint{1.948431in}{0.940558in}}%
\pgfpathlineto{\pgfqpoint{1.950949in}{0.901855in}}%
\pgfpathlineto{\pgfqpoint{1.953468in}{0.914390in}}%
\pgfpathlineto{\pgfqpoint{1.955987in}{0.918480in}}%
\pgfpathlineto{\pgfqpoint{1.958506in}{0.919027in}}%
\pgfpathlineto{\pgfqpoint{1.961024in}{0.912348in}}%
\pgfpathlineto{\pgfqpoint{1.963543in}{0.890861in}}%
\pgfpathlineto{\pgfqpoint{1.966062in}{0.942956in}}%
\pgfpathlineto{\pgfqpoint{1.968581in}{0.912072in}}%
\pgfpathlineto{\pgfqpoint{1.971099in}{0.913151in}}%
\pgfpathlineto{\pgfqpoint{1.973618in}{0.955965in}}%
\pgfpathlineto{\pgfqpoint{1.976137in}{0.939546in}}%
\pgfpathlineto{\pgfqpoint{1.978656in}{0.920192in}}%
\pgfpathlineto{\pgfqpoint{1.983693in}{0.940536in}}%
\pgfpathlineto{\pgfqpoint{1.986212in}{0.953738in}}%
\pgfpathlineto{\pgfqpoint{1.988731in}{0.956835in}}%
\pgfpathlineto{\pgfqpoint{1.991249in}{0.968768in}}%
\pgfpathlineto{\pgfqpoint{1.993768in}{0.952753in}}%
\pgfpathlineto{\pgfqpoint{1.996287in}{0.965290in}}%
\pgfpathlineto{\pgfqpoint{1.998806in}{0.960326in}}%
\pgfpathlineto{\pgfqpoint{2.003843in}{0.905320in}}%
\pgfpathlineto{\pgfqpoint{2.006362in}{0.892253in}}%
\pgfpathlineto{\pgfqpoint{2.008881in}{0.854030in}}%
\pgfpathlineto{\pgfqpoint{2.011399in}{0.817933in}}%
\pgfpathlineto{\pgfqpoint{2.013918in}{0.831836in}}%
\pgfpathlineto{\pgfqpoint{2.016437in}{0.827928in}}%
\pgfpathlineto{\pgfqpoint{2.018956in}{0.847364in}}%
\pgfpathlineto{\pgfqpoint{2.021474in}{0.875989in}}%
\pgfpathlineto{\pgfqpoint{2.023993in}{0.830770in}}%
\pgfpathlineto{\pgfqpoint{2.026512in}{0.825807in}}%
\pgfpathlineto{\pgfqpoint{2.029031in}{0.858631in}}%
\pgfpathlineto{\pgfqpoint{2.031549in}{0.864921in}}%
\pgfpathlineto{\pgfqpoint{2.034068in}{0.844377in}}%
\pgfpathlineto{\pgfqpoint{2.036587in}{0.812966in}}%
\pgfpathlineto{\pgfqpoint{2.039106in}{0.818789in}}%
\pgfpathlineto{\pgfqpoint{2.041624in}{0.846848in}}%
\pgfpathlineto{\pgfqpoint{2.044143in}{0.824805in}}%
\pgfpathlineto{\pgfqpoint{2.046662in}{0.791502in}}%
\pgfpathlineto{\pgfqpoint{2.049181in}{0.767292in}}%
\pgfpathlineto{\pgfqpoint{2.051699in}{0.750578in}}%
\pgfpathlineto{\pgfqpoint{2.054218in}{0.756167in}}%
\pgfpathlineto{\pgfqpoint{2.056737in}{0.752390in}}%
\pgfpathlineto{\pgfqpoint{2.059256in}{0.708203in}}%
\pgfpathlineto{\pgfqpoint{2.061774in}{0.705575in}}%
\pgfpathlineto{\pgfqpoint{2.064293in}{0.696041in}}%
\pgfpathlineto{\pgfqpoint{2.066812in}{0.677701in}}%
\pgfpathlineto{\pgfqpoint{2.069331in}{0.682037in}}%
\pgfpathlineto{\pgfqpoint{2.071849in}{0.680280in}}%
\pgfpathlineto{\pgfqpoint{2.074368in}{0.668693in}}%
\pgfpathlineto{\pgfqpoint{2.079406in}{0.609673in}}%
\pgfpathlineto{\pgfqpoint{2.081924in}{0.583080in}}%
\pgfpathlineto{\pgfqpoint{2.084443in}{0.599568in}}%
\pgfpathlineto{\pgfqpoint{2.086962in}{0.589752in}}%
\pgfpathlineto{\pgfqpoint{2.089481in}{0.631278in}}%
\pgfpathlineto{\pgfqpoint{2.091999in}{0.648024in}}%
\pgfpathlineto{\pgfqpoint{2.094518in}{0.651855in}}%
\pgfpathlineto{\pgfqpoint{2.097037in}{0.684729in}}%
\pgfpathlineto{\pgfqpoint{2.099556in}{0.720039in}}%
\pgfpathlineto{\pgfqpoint{2.102074in}{0.706993in}}%
\pgfpathlineto{\pgfqpoint{2.104593in}{0.708011in}}%
\pgfpathlineto{\pgfqpoint{2.107112in}{0.716059in}}%
\pgfpathlineto{\pgfqpoint{2.109631in}{0.718785in}}%
\pgfpathlineto{\pgfqpoint{2.112149in}{0.767487in}}%
\pgfpathlineto{\pgfqpoint{2.114668in}{0.804250in}}%
\pgfpathlineto{\pgfqpoint{2.117187in}{0.833989in}}%
\pgfpathlineto{\pgfqpoint{2.119706in}{0.892647in}}%
\pgfpathlineto{\pgfqpoint{2.122224in}{0.924135in}}%
\pgfpathlineto{\pgfqpoint{2.124743in}{0.942533in}}%
\pgfpathlineto{\pgfqpoint{2.127262in}{0.941542in}}%
\pgfpathlineto{\pgfqpoint{2.129781in}{0.895662in}}%
\pgfpathlineto{\pgfqpoint{2.132299in}{0.877385in}}%
\pgfpathlineto{\pgfqpoint{2.134818in}{0.867080in}}%
\pgfpathlineto{\pgfqpoint{2.137337in}{0.834870in}}%
\pgfpathlineto{\pgfqpoint{2.139856in}{0.849907in}}%
\pgfpathlineto{\pgfqpoint{2.142374in}{0.823434in}}%
\pgfpathlineto{\pgfqpoint{2.144893in}{0.860955in}}%
\pgfpathlineto{\pgfqpoint{2.147412in}{0.878546in}}%
\pgfpathlineto{\pgfqpoint{2.149931in}{0.917349in}}%
\pgfpathlineto{\pgfqpoint{2.152449in}{0.891619in}}%
\pgfpathlineto{\pgfqpoint{2.154968in}{0.903081in}}%
\pgfpathlineto{\pgfqpoint{2.157487in}{0.887872in}}%
\pgfpathlineto{\pgfqpoint{2.160006in}{0.847354in}}%
\pgfpathlineto{\pgfqpoint{2.162524in}{0.886370in}}%
\pgfpathlineto{\pgfqpoint{2.167562in}{0.917107in}}%
\pgfpathlineto{\pgfqpoint{2.170081in}{0.918816in}}%
\pgfpathlineto{\pgfqpoint{2.172599in}{0.945430in}}%
\pgfpathlineto{\pgfqpoint{2.175118in}{0.968035in}}%
\pgfpathlineto{\pgfqpoint{2.177637in}{0.974196in}}%
\pgfpathlineto{\pgfqpoint{2.180156in}{1.039750in}}%
\pgfpathlineto{\pgfqpoint{2.182674in}{1.027798in}}%
\pgfpathlineto{\pgfqpoint{2.185193in}{1.057525in}}%
\pgfpathlineto{\pgfqpoint{2.187712in}{1.036759in}}%
\pgfpathlineto{\pgfqpoint{2.190231in}{1.059088in}}%
\pgfpathlineto{\pgfqpoint{2.192749in}{1.039572in}}%
\pgfpathlineto{\pgfqpoint{2.195268in}{1.063759in}}%
\pgfpathlineto{\pgfqpoint{2.197787in}{1.043877in}}%
\pgfpathlineto{\pgfqpoint{2.200306in}{0.999804in}}%
\pgfpathlineto{\pgfqpoint{2.202824in}{1.011255in}}%
\pgfpathlineto{\pgfqpoint{2.205343in}{0.992616in}}%
\pgfpathlineto{\pgfqpoint{2.207862in}{0.989995in}}%
\pgfpathlineto{\pgfqpoint{2.210381in}{0.972115in}}%
\pgfpathlineto{\pgfqpoint{2.212899in}{0.931122in}}%
\pgfpathlineto{\pgfqpoint{2.215418in}{0.962744in}}%
\pgfpathlineto{\pgfqpoint{2.217937in}{0.999701in}}%
\pgfpathlineto{\pgfqpoint{2.220456in}{0.990073in}}%
\pgfpathlineto{\pgfqpoint{2.222974in}{0.999815in}}%
\pgfpathlineto{\pgfqpoint{2.225493in}{1.035562in}}%
\pgfpathlineto{\pgfqpoint{2.228012in}{1.056179in}}%
\pgfpathlineto{\pgfqpoint{2.230531in}{1.019466in}}%
\pgfpathlineto{\pgfqpoint{2.233049in}{1.002173in}}%
\pgfpathlineto{\pgfqpoint{2.235568in}{0.974285in}}%
\pgfpathlineto{\pgfqpoint{2.238087in}{0.956275in}}%
\pgfpathlineto{\pgfqpoint{2.240606in}{0.973687in}}%
\pgfpathlineto{\pgfqpoint{2.243124in}{1.010165in}}%
\pgfpathlineto{\pgfqpoint{2.245643in}{0.983004in}}%
\pgfpathlineto{\pgfqpoint{2.248162in}{0.958102in}}%
\pgfpathlineto{\pgfqpoint{2.250681in}{0.966589in}}%
\pgfpathlineto{\pgfqpoint{2.253199in}{0.940445in}}%
\pgfpathlineto{\pgfqpoint{2.255718in}{0.944219in}}%
\pgfpathlineto{\pgfqpoint{2.258237in}{0.954254in}}%
\pgfpathlineto{\pgfqpoint{2.260756in}{0.910859in}}%
\pgfpathlineto{\pgfqpoint{2.263274in}{0.945217in}}%
\pgfpathlineto{\pgfqpoint{2.265793in}{0.964059in}}%
\pgfpathlineto{\pgfqpoint{2.268312in}{0.949404in}}%
\pgfpathlineto{\pgfqpoint{2.270831in}{0.948096in}}%
\pgfpathlineto{\pgfqpoint{2.273349in}{0.967913in}}%
\pgfpathlineto{\pgfqpoint{2.275868in}{0.934769in}}%
\pgfpathlineto{\pgfqpoint{2.278387in}{0.945082in}}%
\pgfpathlineto{\pgfqpoint{2.280906in}{0.881384in}}%
\pgfpathlineto{\pgfqpoint{2.283424in}{0.852691in}}%
\pgfpathlineto{\pgfqpoint{2.285943in}{0.800924in}}%
\pgfpathlineto{\pgfqpoint{2.288462in}{0.758379in}}%
\pgfpathlineto{\pgfqpoint{2.290981in}{0.773641in}}%
\pgfpathlineto{\pgfqpoint{2.293499in}{0.806390in}}%
\pgfpathlineto{\pgfqpoint{2.296018in}{0.822597in}}%
\pgfpathlineto{\pgfqpoint{2.298537in}{0.868846in}}%
\pgfpathlineto{\pgfqpoint{2.301056in}{0.891546in}}%
\pgfpathlineto{\pgfqpoint{2.303574in}{0.852710in}}%
\pgfpathlineto{\pgfqpoint{2.306093in}{0.845265in}}%
\pgfpathlineto{\pgfqpoint{2.308612in}{0.846433in}}%
\pgfpathlineto{\pgfqpoint{2.311131in}{0.854886in}}%
\pgfpathlineto{\pgfqpoint{2.313649in}{0.877922in}}%
\pgfpathlineto{\pgfqpoint{2.316168in}{0.875364in}}%
\pgfpathlineto{\pgfqpoint{2.318687in}{0.869484in}}%
\pgfpathlineto{\pgfqpoint{2.321206in}{0.912008in}}%
\pgfpathlineto{\pgfqpoint{2.323724in}{0.933293in}}%
\pgfpathlineto{\pgfqpoint{2.326243in}{0.893352in}}%
\pgfpathlineto{\pgfqpoint{2.328762in}{0.907983in}}%
\pgfpathlineto{\pgfqpoint{2.331281in}{0.918312in}}%
\pgfpathlineto{\pgfqpoint{2.333799in}{0.894277in}}%
\pgfpathlineto{\pgfqpoint{2.336318in}{0.914143in}}%
\pgfpathlineto{\pgfqpoint{2.338837in}{0.958128in}}%
\pgfpathlineto{\pgfqpoint{2.341356in}{0.990118in}}%
\pgfpathlineto{\pgfqpoint{2.343874in}{0.985512in}}%
\pgfpathlineto{\pgfqpoint{2.346393in}{0.996031in}}%
\pgfpathlineto{\pgfqpoint{2.348912in}{1.008656in}}%
\pgfpathlineto{\pgfqpoint{2.351431in}{1.030463in}}%
\pgfpathlineto{\pgfqpoint{2.353949in}{1.025441in}}%
\pgfpathlineto{\pgfqpoint{2.356468in}{1.055193in}}%
\pgfpathlineto{\pgfqpoint{2.358987in}{1.045885in}}%
\pgfpathlineto{\pgfqpoint{2.361506in}{1.071883in}}%
\pgfpathlineto{\pgfqpoint{2.364024in}{1.125386in}}%
\pgfpathlineto{\pgfqpoint{2.366543in}{1.104692in}}%
\pgfpathlineto{\pgfqpoint{2.369062in}{1.094777in}}%
\pgfpathlineto{\pgfqpoint{2.371581in}{1.132989in}}%
\pgfpathlineto{\pgfqpoint{2.374099in}{1.106039in}}%
\pgfpathlineto{\pgfqpoint{2.376618in}{1.093618in}}%
\pgfpathlineto{\pgfqpoint{2.379137in}{1.061471in}}%
\pgfpathlineto{\pgfqpoint{2.381656in}{1.098981in}}%
\pgfpathlineto{\pgfqpoint{2.384174in}{1.138461in}}%
\pgfpathlineto{\pgfqpoint{2.386693in}{1.164586in}}%
\pgfpathlineto{\pgfqpoint{2.389212in}{1.200297in}}%
\pgfpathlineto{\pgfqpoint{2.391731in}{1.176746in}}%
\pgfpathlineto{\pgfqpoint{2.394249in}{1.154613in}}%
\pgfpathlineto{\pgfqpoint{2.396768in}{1.143903in}}%
\pgfpathlineto{\pgfqpoint{2.399287in}{1.099245in}}%
\pgfpathlineto{\pgfqpoint{2.401806in}{1.095511in}}%
\pgfpathlineto{\pgfqpoint{2.404324in}{1.106350in}}%
\pgfpathlineto{\pgfqpoint{2.406843in}{1.089478in}}%
\pgfpathlineto{\pgfqpoint{2.409362in}{1.087769in}}%
\pgfpathlineto{\pgfqpoint{2.411881in}{1.080160in}}%
\pgfpathlineto{\pgfqpoint{2.414399in}{1.086579in}}%
\pgfpathlineto{\pgfqpoint{2.416918in}{1.095595in}}%
\pgfpathlineto{\pgfqpoint{2.419437in}{1.020603in}}%
\pgfpathlineto{\pgfqpoint{2.421956in}{0.995658in}}%
\pgfpathlineto{\pgfqpoint{2.424474in}{1.012895in}}%
\pgfpathlineto{\pgfqpoint{2.426993in}{1.013819in}}%
\pgfpathlineto{\pgfqpoint{2.429512in}{1.014221in}}%
\pgfpathlineto{\pgfqpoint{2.432031in}{1.007243in}}%
\pgfpathlineto{\pgfqpoint{2.434549in}{1.030229in}}%
\pgfpathlineto{\pgfqpoint{2.437068in}{1.056556in}}%
\pgfpathlineto{\pgfqpoint{2.439587in}{1.097076in}}%
\pgfpathlineto{\pgfqpoint{2.442106in}{1.087275in}}%
\pgfpathlineto{\pgfqpoint{2.444624in}{1.112271in}}%
\pgfpathlineto{\pgfqpoint{2.447143in}{1.130185in}}%
\pgfpathlineto{\pgfqpoint{2.449662in}{1.125156in}}%
\pgfpathlineto{\pgfqpoint{2.452181in}{1.093693in}}%
\pgfpathlineto{\pgfqpoint{2.454699in}{1.104616in}}%
\pgfpathlineto{\pgfqpoint{2.457218in}{1.084783in}}%
\pgfpathlineto{\pgfqpoint{2.459737in}{1.108018in}}%
\pgfpathlineto{\pgfqpoint{2.462256in}{1.073733in}}%
\pgfpathlineto{\pgfqpoint{2.464774in}{1.115599in}}%
\pgfpathlineto{\pgfqpoint{2.467293in}{1.140013in}}%
\pgfpathlineto{\pgfqpoint{2.469812in}{1.160900in}}%
\pgfpathlineto{\pgfqpoint{2.472331in}{1.173291in}}%
\pgfpathlineto{\pgfqpoint{2.474849in}{1.188191in}}%
\pgfpathlineto{\pgfqpoint{2.477368in}{1.184102in}}%
\pgfpathlineto{\pgfqpoint{2.479887in}{1.163391in}}%
\pgfpathlineto{\pgfqpoint{2.482406in}{1.178170in}}%
\pgfpathlineto{\pgfqpoint{2.484924in}{1.178793in}}%
\pgfpathlineto{\pgfqpoint{2.487443in}{1.194851in}}%
\pgfpathlineto{\pgfqpoint{2.489962in}{1.227821in}}%
\pgfpathlineto{\pgfqpoint{2.492481in}{1.203916in}}%
\pgfpathlineto{\pgfqpoint{2.494999in}{1.187254in}}%
\pgfpathlineto{\pgfqpoint{2.497518in}{1.224918in}}%
\pgfpathlineto{\pgfqpoint{2.500037in}{1.239951in}}%
\pgfpathlineto{\pgfqpoint{2.502556in}{1.243488in}}%
\pgfpathlineto{\pgfqpoint{2.505074in}{1.241295in}}%
\pgfpathlineto{\pgfqpoint{2.507593in}{1.227782in}}%
\pgfpathlineto{\pgfqpoint{2.510112in}{1.255626in}}%
\pgfpathlineto{\pgfqpoint{2.512631in}{1.264811in}}%
\pgfpathlineto{\pgfqpoint{2.515149in}{1.267179in}}%
\pgfpathlineto{\pgfqpoint{2.517668in}{1.265318in}}%
\pgfpathlineto{\pgfqpoint{2.520187in}{1.284997in}}%
\pgfpathlineto{\pgfqpoint{2.522706in}{1.266064in}}%
\pgfpathlineto{\pgfqpoint{2.525224in}{1.307890in}}%
\pgfpathlineto{\pgfqpoint{2.527743in}{1.283383in}}%
\pgfpathlineto{\pgfqpoint{2.530262in}{1.312792in}}%
\pgfpathlineto{\pgfqpoint{2.532781in}{1.343760in}}%
\pgfpathlineto{\pgfqpoint{2.535299in}{1.346365in}}%
\pgfpathlineto{\pgfqpoint{2.537818in}{1.336411in}}%
\pgfpathlineto{\pgfqpoint{2.540337in}{1.399028in}}%
\pgfpathlineto{\pgfqpoint{2.542856in}{1.398845in}}%
\pgfpathlineto{\pgfqpoint{2.545374in}{1.368049in}}%
\pgfpathlineto{\pgfqpoint{2.547893in}{1.402865in}}%
\pgfpathlineto{\pgfqpoint{2.550412in}{1.395002in}}%
\pgfpathlineto{\pgfqpoint{2.552931in}{1.351850in}}%
\pgfpathlineto{\pgfqpoint{2.555449in}{1.354576in}}%
\pgfpathlineto{\pgfqpoint{2.557968in}{1.352749in}}%
\pgfpathlineto{\pgfqpoint{2.560487in}{1.367400in}}%
\pgfpathlineto{\pgfqpoint{2.563006in}{1.384146in}}%
\pgfpathlineto{\pgfqpoint{2.565524in}{1.373657in}}%
\pgfpathlineto{\pgfqpoint{2.568043in}{1.364802in}}%
\pgfpathlineto{\pgfqpoint{2.570562in}{1.396856in}}%
\pgfpathlineto{\pgfqpoint{2.573081in}{1.395466in}}%
\pgfpathlineto{\pgfqpoint{2.575599in}{1.470793in}}%
\pgfpathlineto{\pgfqpoint{2.578118in}{1.492297in}}%
\pgfpathlineto{\pgfqpoint{2.580637in}{1.450270in}}%
\pgfpathlineto{\pgfqpoint{2.583156in}{1.486184in}}%
\pgfpathlineto{\pgfqpoint{2.585674in}{1.488141in}}%
\pgfpathlineto{\pgfqpoint{2.588193in}{1.471465in}}%
\pgfpathlineto{\pgfqpoint{2.590712in}{1.438274in}}%
\pgfpathlineto{\pgfqpoint{2.593231in}{1.440420in}}%
\pgfpathlineto{\pgfqpoint{2.595749in}{1.469960in}}%
\pgfpathlineto{\pgfqpoint{2.598268in}{1.461804in}}%
\pgfpathlineto{\pgfqpoint{2.600787in}{1.479150in}}%
\pgfpathlineto{\pgfqpoint{2.603306in}{1.433798in}}%
\pgfpathlineto{\pgfqpoint{2.605824in}{1.419081in}}%
\pgfpathlineto{\pgfqpoint{2.608343in}{1.455511in}}%
\pgfpathlineto{\pgfqpoint{2.610862in}{1.417508in}}%
\pgfpathlineto{\pgfqpoint{2.613381in}{1.474939in}}%
\pgfpathlineto{\pgfqpoint{2.615899in}{1.464628in}}%
\pgfpathlineto{\pgfqpoint{2.618418in}{1.444993in}}%
\pgfpathlineto{\pgfqpoint{2.620937in}{1.456267in}}%
\pgfpathlineto{\pgfqpoint{2.623456in}{1.463984in}}%
\pgfpathlineto{\pgfqpoint{2.625974in}{1.455223in}}%
\pgfpathlineto{\pgfqpoint{2.628493in}{1.448029in}}%
\pgfpathlineto{\pgfqpoint{2.631012in}{1.425846in}}%
\pgfpathlineto{\pgfqpoint{2.633531in}{1.472403in}}%
\pgfpathlineto{\pgfqpoint{2.636049in}{1.398517in}}%
\pgfpathlineto{\pgfqpoint{2.638568in}{1.385154in}}%
\pgfpathlineto{\pgfqpoint{2.641087in}{1.321892in}}%
\pgfpathlineto{\pgfqpoint{2.643606in}{1.313913in}}%
\pgfpathlineto{\pgfqpoint{2.646124in}{1.301261in}}%
\pgfpathlineto{\pgfqpoint{2.648643in}{1.301617in}}%
\pgfpathlineto{\pgfqpoint{2.651162in}{1.321688in}}%
\pgfpathlineto{\pgfqpoint{2.653681in}{1.293705in}}%
\pgfpathlineto{\pgfqpoint{2.656199in}{1.272273in}}%
\pgfpathlineto{\pgfqpoint{2.658718in}{1.268734in}}%
\pgfpathlineto{\pgfqpoint{2.661237in}{1.227200in}}%
\pgfpathlineto{\pgfqpoint{2.663756in}{1.226775in}}%
\pgfpathlineto{\pgfqpoint{2.666274in}{1.263542in}}%
\pgfpathlineto{\pgfqpoint{2.668793in}{1.209613in}}%
\pgfpathlineto{\pgfqpoint{2.671312in}{1.235654in}}%
\pgfpathlineto{\pgfqpoint{2.673831in}{1.263029in}}%
\pgfpathlineto{\pgfqpoint{2.676349in}{1.251475in}}%
\pgfpathlineto{\pgfqpoint{2.678868in}{1.272147in}}%
\pgfpathlineto{\pgfqpoint{2.681387in}{1.231523in}}%
\pgfpathlineto{\pgfqpoint{2.683906in}{1.207974in}}%
\pgfpathlineto{\pgfqpoint{2.686424in}{1.216470in}}%
\pgfpathlineto{\pgfqpoint{2.688943in}{1.185660in}}%
\pgfpathlineto{\pgfqpoint{2.691462in}{1.164114in}}%
\pgfpathlineto{\pgfqpoint{2.693981in}{1.190123in}}%
\pgfpathlineto{\pgfqpoint{2.696499in}{1.211339in}}%
\pgfpathlineto{\pgfqpoint{2.699018in}{1.160415in}}%
\pgfpathlineto{\pgfqpoint{2.701537in}{1.191776in}}%
\pgfpathlineto{\pgfqpoint{2.704056in}{1.201194in}}%
\pgfpathlineto{\pgfqpoint{2.706574in}{1.200976in}}%
\pgfpathlineto{\pgfqpoint{2.709093in}{1.240847in}}%
\pgfpathlineto{\pgfqpoint{2.711612in}{1.242193in}}%
\pgfpathlineto{\pgfqpoint{2.714131in}{1.207421in}}%
\pgfpathlineto{\pgfqpoint{2.716649in}{1.193784in}}%
\pgfpathlineto{\pgfqpoint{2.719168in}{1.192853in}}%
\pgfpathlineto{\pgfqpoint{2.721687in}{1.201842in}}%
\pgfpathlineto{\pgfqpoint{2.724206in}{1.197227in}}%
\pgfpathlineto{\pgfqpoint{2.726724in}{1.185793in}}%
\pgfpathlineto{\pgfqpoint{2.729243in}{1.190889in}}%
\pgfpathlineto{\pgfqpoint{2.731762in}{1.170255in}}%
\pgfpathlineto{\pgfqpoint{2.734281in}{1.175855in}}%
\pgfpathlineto{\pgfqpoint{2.736799in}{1.197704in}}%
\pgfpathlineto{\pgfqpoint{2.739318in}{1.179446in}}%
\pgfpathlineto{\pgfqpoint{2.741837in}{1.158525in}}%
\pgfpathlineto{\pgfqpoint{2.744356in}{1.133316in}}%
\pgfpathlineto{\pgfqpoint{2.746874in}{1.149302in}}%
\pgfpathlineto{\pgfqpoint{2.749393in}{1.193748in}}%
\pgfpathlineto{\pgfqpoint{2.751912in}{1.137979in}}%
\pgfpathlineto{\pgfqpoint{2.754431in}{1.154269in}}%
\pgfpathlineto{\pgfqpoint{2.756949in}{1.126779in}}%
\pgfpathlineto{\pgfqpoint{2.759468in}{1.135661in}}%
\pgfpathlineto{\pgfqpoint{2.761987in}{1.112249in}}%
\pgfpathlineto{\pgfqpoint{2.764506in}{1.123504in}}%
\pgfpathlineto{\pgfqpoint{2.767024in}{1.132831in}}%
\pgfpathlineto{\pgfqpoint{2.769543in}{1.099604in}}%
\pgfpathlineto{\pgfqpoint{2.772062in}{1.053867in}}%
\pgfpathlineto{\pgfqpoint{2.774581in}{1.094368in}}%
\pgfpathlineto{\pgfqpoint{2.777099in}{1.111916in}}%
\pgfpathlineto{\pgfqpoint{2.779618in}{1.139401in}}%
\pgfpathlineto{\pgfqpoint{2.782137in}{1.134089in}}%
\pgfpathlineto{\pgfqpoint{2.784656in}{1.119409in}}%
\pgfpathlineto{\pgfqpoint{2.787174in}{1.077117in}}%
\pgfpathlineto{\pgfqpoint{2.789693in}{1.115261in}}%
\pgfpathlineto{\pgfqpoint{2.792212in}{1.184158in}}%
\pgfpathlineto{\pgfqpoint{2.794731in}{1.102423in}}%
\pgfpathlineto{\pgfqpoint{2.797249in}{1.059073in}}%
\pgfpathlineto{\pgfqpoint{2.799768in}{1.065368in}}%
\pgfpathlineto{\pgfqpoint{2.802287in}{1.077023in}}%
\pgfpathlineto{\pgfqpoint{2.804806in}{1.116490in}}%
\pgfpathlineto{\pgfqpoint{2.807324in}{1.084999in}}%
\pgfpathlineto{\pgfqpoint{2.809843in}{1.122475in}}%
\pgfpathlineto{\pgfqpoint{2.812362in}{1.127949in}}%
\pgfpathlineto{\pgfqpoint{2.814881in}{1.113331in}}%
\pgfpathlineto{\pgfqpoint{2.817399in}{1.071891in}}%
\pgfpathlineto{\pgfqpoint{2.819918in}{1.046572in}}%
\pgfpathlineto{\pgfqpoint{2.822437in}{1.075042in}}%
\pgfpathlineto{\pgfqpoint{2.824956in}{1.080738in}}%
\pgfpathlineto{\pgfqpoint{2.827474in}{1.029130in}}%
\pgfpathlineto{\pgfqpoint{2.829993in}{0.986453in}}%
\pgfpathlineto{\pgfqpoint{2.832512in}{0.962069in}}%
\pgfpathlineto{\pgfqpoint{2.835031in}{0.973629in}}%
\pgfpathlineto{\pgfqpoint{2.837549in}{0.970986in}}%
\pgfpathlineto{\pgfqpoint{2.840068in}{1.008178in}}%
\pgfpathlineto{\pgfqpoint{2.842587in}{1.018470in}}%
\pgfpathlineto{\pgfqpoint{2.845106in}{1.042075in}}%
\pgfpathlineto{\pgfqpoint{2.847624in}{1.073932in}}%
\pgfpathlineto{\pgfqpoint{2.850143in}{1.102699in}}%
\pgfpathlineto{\pgfqpoint{2.852662in}{1.086346in}}%
\pgfpathlineto{\pgfqpoint{2.855181in}{1.094198in}}%
\pgfpathlineto{\pgfqpoint{2.857699in}{1.091152in}}%
\pgfpathlineto{\pgfqpoint{2.860218in}{1.116521in}}%
\pgfpathlineto{\pgfqpoint{2.862737in}{1.107429in}}%
\pgfpathlineto{\pgfqpoint{2.865256in}{1.130072in}}%
\pgfpathlineto{\pgfqpoint{2.867774in}{1.100823in}}%
\pgfpathlineto{\pgfqpoint{2.870293in}{1.148955in}}%
\pgfpathlineto{\pgfqpoint{2.872812in}{1.140887in}}%
\pgfpathlineto{\pgfqpoint{2.875331in}{1.191671in}}%
\pgfpathlineto{\pgfqpoint{2.877849in}{1.152367in}}%
\pgfpathlineto{\pgfqpoint{2.880368in}{1.160042in}}%
\pgfpathlineto{\pgfqpoint{2.882887in}{1.171735in}}%
\pgfpathlineto{\pgfqpoint{2.885406in}{1.177899in}}%
\pgfpathlineto{\pgfqpoint{2.887924in}{1.151579in}}%
\pgfpathlineto{\pgfqpoint{2.890443in}{1.175649in}}%
\pgfpathlineto{\pgfqpoint{2.892962in}{1.176965in}}%
\pgfpathlineto{\pgfqpoint{2.895481in}{1.192576in}}%
\pgfpathlineto{\pgfqpoint{2.897999in}{1.184438in}}%
\pgfpathlineto{\pgfqpoint{2.900518in}{1.213673in}}%
\pgfpathlineto{\pgfqpoint{2.903037in}{1.173783in}}%
\pgfpathlineto{\pgfqpoint{2.905556in}{1.195814in}}%
\pgfpathlineto{\pgfqpoint{2.908074in}{1.224117in}}%
\pgfpathlineto{\pgfqpoint{2.910593in}{1.239137in}}%
\pgfpathlineto{\pgfqpoint{2.913112in}{1.203871in}}%
\pgfpathlineto{\pgfqpoint{2.915631in}{1.201932in}}%
\pgfpathlineto{\pgfqpoint{2.918149in}{1.206960in}}%
\pgfpathlineto{\pgfqpoint{2.920668in}{1.235670in}}%
\pgfpathlineto{\pgfqpoint{2.923187in}{1.217831in}}%
\pgfpathlineto{\pgfqpoint{2.925706in}{1.215210in}}%
\pgfpathlineto{\pgfqpoint{2.928224in}{1.216279in}}%
\pgfpathlineto{\pgfqpoint{2.930743in}{1.218131in}}%
\pgfpathlineto{\pgfqpoint{2.933262in}{1.236610in}}%
\pgfpathlineto{\pgfqpoint{2.935781in}{1.252871in}}%
\pgfpathlineto{\pgfqpoint{2.938299in}{1.235475in}}%
\pgfpathlineto{\pgfqpoint{2.940818in}{1.266941in}}%
\pgfpathlineto{\pgfqpoint{2.943337in}{1.240806in}}%
\pgfpathlineto{\pgfqpoint{2.945856in}{1.249211in}}%
\pgfpathlineto{\pgfqpoint{2.948374in}{1.269399in}}%
\pgfpathlineto{\pgfqpoint{2.950893in}{1.285479in}}%
\pgfpathlineto{\pgfqpoint{2.955931in}{1.371877in}}%
\pgfpathlineto{\pgfqpoint{2.958449in}{1.346820in}}%
\pgfpathlineto{\pgfqpoint{2.960968in}{1.343479in}}%
\pgfpathlineto{\pgfqpoint{2.963487in}{1.358493in}}%
\pgfpathlineto{\pgfqpoint{2.966006in}{1.354667in}}%
\pgfpathlineto{\pgfqpoint{2.968524in}{1.375580in}}%
\pgfpathlineto{\pgfqpoint{2.971043in}{1.391355in}}%
\pgfpathlineto{\pgfqpoint{2.973562in}{1.443180in}}%
\pgfpathlineto{\pgfqpoint{2.976081in}{1.434486in}}%
\pgfpathlineto{\pgfqpoint{2.978599in}{1.421899in}}%
\pgfpathlineto{\pgfqpoint{2.981118in}{1.430644in}}%
\pgfpathlineto{\pgfqpoint{2.983637in}{1.430965in}}%
\pgfpathlineto{\pgfqpoint{2.986156in}{1.465647in}}%
\pgfpathlineto{\pgfqpoint{2.988674in}{1.445270in}}%
\pgfpathlineto{\pgfqpoint{2.991193in}{1.448703in}}%
\pgfpathlineto{\pgfqpoint{2.993712in}{1.433839in}}%
\pgfpathlineto{\pgfqpoint{2.996231in}{1.409330in}}%
\pgfpathlineto{\pgfqpoint{2.998749in}{1.454704in}}%
\pgfpathlineto{\pgfqpoint{3.001268in}{1.393044in}}%
\pgfpathlineto{\pgfqpoint{3.003787in}{1.471257in}}%
\pgfpathlineto{\pgfqpoint{3.006306in}{1.461767in}}%
\pgfpathlineto{\pgfqpoint{3.008824in}{1.435556in}}%
\pgfpathlineto{\pgfqpoint{3.011343in}{1.413179in}}%
\pgfpathlineto{\pgfqpoint{3.013862in}{1.437906in}}%
\pgfpathlineto{\pgfqpoint{3.016381in}{1.402502in}}%
\pgfpathlineto{\pgfqpoint{3.018899in}{1.391310in}}%
\pgfpathlineto{\pgfqpoint{3.021418in}{1.346781in}}%
\pgfpathlineto{\pgfqpoint{3.023937in}{1.394949in}}%
\pgfpathlineto{\pgfqpoint{3.026456in}{1.398150in}}%
\pgfpathlineto{\pgfqpoint{3.028974in}{1.357871in}}%
\pgfpathlineto{\pgfqpoint{3.031493in}{1.369044in}}%
\pgfpathlineto{\pgfqpoint{3.034012in}{1.385143in}}%
\pgfpathlineto{\pgfqpoint{3.036531in}{1.417134in}}%
\pgfpathlineto{\pgfqpoint{3.039049in}{1.401641in}}%
\pgfpathlineto{\pgfqpoint{3.041568in}{1.404630in}}%
\pgfpathlineto{\pgfqpoint{3.044087in}{1.428311in}}%
\pgfpathlineto{\pgfqpoint{3.046606in}{1.414968in}}%
\pgfpathlineto{\pgfqpoint{3.049124in}{1.409519in}}%
\pgfpathlineto{\pgfqpoint{3.051643in}{1.349784in}}%
\pgfpathlineto{\pgfqpoint{3.054162in}{1.338354in}}%
\pgfpathlineto{\pgfqpoint{3.056681in}{1.300360in}}%
\pgfpathlineto{\pgfqpoint{3.059199in}{1.355204in}}%
\pgfpathlineto{\pgfqpoint{3.061718in}{1.365429in}}%
\pgfpathlineto{\pgfqpoint{3.064237in}{1.360352in}}%
\pgfpathlineto{\pgfqpoint{3.066756in}{1.328736in}}%
\pgfpathlineto{\pgfqpoint{3.069274in}{1.374096in}}%
\pgfpathlineto{\pgfqpoint{3.071793in}{1.377355in}}%
\pgfpathlineto{\pgfqpoint{3.074312in}{1.405303in}}%
\pgfpathlineto{\pgfqpoint{3.076831in}{1.374448in}}%
\pgfpathlineto{\pgfqpoint{3.079349in}{1.359814in}}%
\pgfpathlineto{\pgfqpoint{3.081868in}{1.357343in}}%
\pgfpathlineto{\pgfqpoint{3.084387in}{1.359111in}}%
\pgfpathlineto{\pgfqpoint{3.086906in}{1.365380in}}%
\pgfpathlineto{\pgfqpoint{3.089424in}{1.364303in}}%
\pgfpathlineto{\pgfqpoint{3.091943in}{1.307681in}}%
\pgfpathlineto{\pgfqpoint{3.094462in}{1.320089in}}%
\pgfpathlineto{\pgfqpoint{3.096981in}{1.304604in}}%
\pgfpathlineto{\pgfqpoint{3.099499in}{1.340603in}}%
\pgfpathlineto{\pgfqpoint{3.102018in}{1.359764in}}%
\pgfpathlineto{\pgfqpoint{3.104537in}{1.407339in}}%
\pgfpathlineto{\pgfqpoint{3.107056in}{1.429977in}}%
\pgfpathlineto{\pgfqpoint{3.109574in}{1.474204in}}%
\pgfpathlineto{\pgfqpoint{3.112093in}{1.473823in}}%
\pgfpathlineto{\pgfqpoint{3.114612in}{1.524096in}}%
\pgfpathlineto{\pgfqpoint{3.117131in}{1.484014in}}%
\pgfpathlineto{\pgfqpoint{3.119649in}{1.507544in}}%
\pgfpathlineto{\pgfqpoint{3.122168in}{1.504117in}}%
\pgfpathlineto{\pgfqpoint{3.124687in}{1.543986in}}%
\pgfpathlineto{\pgfqpoint{3.127206in}{1.597900in}}%
\pgfpathlineto{\pgfqpoint{3.129724in}{1.562032in}}%
\pgfpathlineto{\pgfqpoint{3.132243in}{1.520807in}}%
\pgfpathlineto{\pgfqpoint{3.134762in}{1.555307in}}%
\pgfpathlineto{\pgfqpoint{3.137281in}{1.519396in}}%
\pgfpathlineto{\pgfqpoint{3.139799in}{1.499625in}}%
\pgfpathlineto{\pgfqpoint{3.142318in}{1.493582in}}%
\pgfpathlineto{\pgfqpoint{3.147356in}{1.571130in}}%
\pgfpathlineto{\pgfqpoint{3.149874in}{1.554994in}}%
\pgfpathlineto{\pgfqpoint{3.152393in}{1.573957in}}%
\pgfpathlineto{\pgfqpoint{3.154912in}{1.573901in}}%
\pgfpathlineto{\pgfqpoint{3.157431in}{1.556932in}}%
\pgfpathlineto{\pgfqpoint{3.159949in}{1.610008in}}%
\pgfpathlineto{\pgfqpoint{3.162468in}{1.659264in}}%
\pgfpathlineto{\pgfqpoint{3.164987in}{1.689988in}}%
\pgfpathlineto{\pgfqpoint{3.167506in}{1.697947in}}%
\pgfpathlineto{\pgfqpoint{3.170024in}{1.687471in}}%
\pgfpathlineto{\pgfqpoint{3.172543in}{1.714871in}}%
\pgfpathlineto{\pgfqpoint{3.175062in}{1.725535in}}%
\pgfpathlineto{\pgfqpoint{3.177581in}{1.749939in}}%
\pgfpathlineto{\pgfqpoint{3.180099in}{1.766473in}}%
\pgfpathlineto{\pgfqpoint{3.182618in}{1.783647in}}%
\pgfpathlineto{\pgfqpoint{3.185137in}{1.823485in}}%
\pgfpathlineto{\pgfqpoint{3.187656in}{1.840935in}}%
\pgfpathlineto{\pgfqpoint{3.190174in}{1.831413in}}%
\pgfpathlineto{\pgfqpoint{3.192693in}{1.763601in}}%
\pgfpathlineto{\pgfqpoint{3.195212in}{1.780237in}}%
\pgfpathlineto{\pgfqpoint{3.197731in}{1.749617in}}%
\pgfpathlineto{\pgfqpoint{3.200249in}{1.684468in}}%
\pgfpathlineto{\pgfqpoint{3.202768in}{1.698207in}}%
\pgfpathlineto{\pgfqpoint{3.205287in}{1.662432in}}%
\pgfpathlineto{\pgfqpoint{3.207806in}{1.659557in}}%
\pgfpathlineto{\pgfqpoint{3.210324in}{1.638386in}}%
\pgfpathlineto{\pgfqpoint{3.212843in}{1.697409in}}%
\pgfpathlineto{\pgfqpoint{3.215362in}{1.674201in}}%
\pgfpathlineto{\pgfqpoint{3.217881in}{1.647671in}}%
\pgfpathlineto{\pgfqpoint{3.220399in}{1.648729in}}%
\pgfpathlineto{\pgfqpoint{3.222918in}{1.636487in}}%
\pgfpathlineto{\pgfqpoint{3.225437in}{1.661278in}}%
\pgfpathlineto{\pgfqpoint{3.227956in}{1.666181in}}%
\pgfpathlineto{\pgfqpoint{3.230474in}{1.620909in}}%
\pgfpathlineto{\pgfqpoint{3.232993in}{1.635763in}}%
\pgfpathlineto{\pgfqpoint{3.235512in}{1.588813in}}%
\pgfpathlineto{\pgfqpoint{3.238031in}{1.547561in}}%
\pgfpathlineto{\pgfqpoint{3.240549in}{1.554173in}}%
\pgfpathlineto{\pgfqpoint{3.243068in}{1.534845in}}%
\pgfpathlineto{\pgfqpoint{3.245587in}{1.535097in}}%
\pgfpathlineto{\pgfqpoint{3.248106in}{1.527107in}}%
\pgfpathlineto{\pgfqpoint{3.250624in}{1.532285in}}%
\pgfpathlineto{\pgfqpoint{3.253143in}{1.593536in}}%
\pgfpathlineto{\pgfqpoint{3.255662in}{1.587272in}}%
\pgfpathlineto{\pgfqpoint{3.258181in}{1.561431in}}%
\pgfpathlineto{\pgfqpoint{3.260699in}{1.575974in}}%
\pgfpathlineto{\pgfqpoint{3.263218in}{1.606922in}}%
\pgfpathlineto{\pgfqpoint{3.265737in}{1.622357in}}%
\pgfpathlineto{\pgfqpoint{3.268256in}{1.627159in}}%
\pgfpathlineto{\pgfqpoint{3.270774in}{1.591536in}}%
\pgfpathlineto{\pgfqpoint{3.273293in}{1.605998in}}%
\pgfpathlineto{\pgfqpoint{3.275812in}{1.616976in}}%
\pgfpathlineto{\pgfqpoint{3.278331in}{1.582429in}}%
\pgfpathlineto{\pgfqpoint{3.280849in}{1.590954in}}%
\pgfpathlineto{\pgfqpoint{3.283368in}{1.605352in}}%
\pgfpathlineto{\pgfqpoint{3.285887in}{1.623886in}}%
\pgfpathlineto{\pgfqpoint{3.288406in}{1.633577in}}%
\pgfpathlineto{\pgfqpoint{3.290924in}{1.638060in}}%
\pgfpathlineto{\pgfqpoint{3.293443in}{1.627360in}}%
\pgfpathlineto{\pgfqpoint{3.295962in}{1.629126in}}%
\pgfpathlineto{\pgfqpoint{3.298481in}{1.633693in}}%
\pgfpathlineto{\pgfqpoint{3.300999in}{1.610674in}}%
\pgfpathlineto{\pgfqpoint{3.303518in}{1.619926in}}%
\pgfpathlineto{\pgfqpoint{3.306037in}{1.588232in}}%
\pgfpathlineto{\pgfqpoint{3.308556in}{1.528690in}}%
\pgfpathlineto{\pgfqpoint{3.311074in}{1.510068in}}%
\pgfpathlineto{\pgfqpoint{3.313593in}{1.555714in}}%
\pgfpathlineto{\pgfqpoint{3.316112in}{1.567918in}}%
\pgfpathlineto{\pgfqpoint{3.318631in}{1.582885in}}%
\pgfpathlineto{\pgfqpoint{3.321149in}{1.549283in}}%
\pgfpathlineto{\pgfqpoint{3.323668in}{1.531319in}}%
\pgfpathlineto{\pgfqpoint{3.326187in}{1.478840in}}%
\pgfpathlineto{\pgfqpoint{3.328706in}{1.526514in}}%
\pgfpathlineto{\pgfqpoint{3.331224in}{1.542924in}}%
\pgfpathlineto{\pgfqpoint{3.333743in}{1.562357in}}%
\pgfpathlineto{\pgfqpoint{3.336262in}{1.517400in}}%
\pgfpathlineto{\pgfqpoint{3.338781in}{1.565535in}}%
\pgfpathlineto{\pgfqpoint{3.341299in}{1.567515in}}%
\pgfpathlineto{\pgfqpoint{3.343818in}{1.583811in}}%
\pgfpathlineto{\pgfqpoint{3.346337in}{1.597041in}}%
\pgfpathlineto{\pgfqpoint{3.348856in}{1.591355in}}%
\pgfpathlineto{\pgfqpoint{3.351374in}{1.598838in}}%
\pgfpathlineto{\pgfqpoint{3.353893in}{1.592543in}}%
\pgfpathlineto{\pgfqpoint{3.356412in}{1.572198in}}%
\pgfpathlineto{\pgfqpoint{3.358931in}{1.564292in}}%
\pgfpathlineto{\pgfqpoint{3.361449in}{1.518396in}}%
\pgfpathlineto{\pgfqpoint{3.363968in}{1.558525in}}%
\pgfpathlineto{\pgfqpoint{3.366487in}{1.552226in}}%
\pgfpathlineto{\pgfqpoint{3.369006in}{1.575527in}}%
\pgfpathlineto{\pgfqpoint{3.371524in}{1.511154in}}%
\pgfpathlineto{\pgfqpoint{3.374043in}{1.514365in}}%
\pgfpathlineto{\pgfqpoint{3.376562in}{1.491902in}}%
\pgfpathlineto{\pgfqpoint{3.379081in}{1.473217in}}%
\pgfpathlineto{\pgfqpoint{3.381599in}{1.521151in}}%
\pgfpathlineto{\pgfqpoint{3.384118in}{1.532608in}}%
\pgfpathlineto{\pgfqpoint{3.386637in}{1.536844in}}%
\pgfpathlineto{\pgfqpoint{3.389156in}{1.590478in}}%
\pgfpathlineto{\pgfqpoint{3.391674in}{1.582805in}}%
\pgfpathlineto{\pgfqpoint{3.394193in}{1.608359in}}%
\pgfpathlineto{\pgfqpoint{3.396712in}{1.617066in}}%
\pgfpathlineto{\pgfqpoint{3.399231in}{1.582390in}}%
\pgfpathlineto{\pgfqpoint{3.401749in}{1.621649in}}%
\pgfpathlineto{\pgfqpoint{3.404268in}{1.621428in}}%
\pgfpathlineto{\pgfqpoint{3.406787in}{1.674349in}}%
\pgfpathlineto{\pgfqpoint{3.409306in}{1.716192in}}%
\pgfpathlineto{\pgfqpoint{3.411824in}{1.681591in}}%
\pgfpathlineto{\pgfqpoint{3.414343in}{1.697872in}}%
\pgfpathlineto{\pgfqpoint{3.416862in}{1.715128in}}%
\pgfpathlineto{\pgfqpoint{3.419381in}{1.688048in}}%
\pgfpathlineto{\pgfqpoint{3.421899in}{1.669864in}}%
\pgfpathlineto{\pgfqpoint{3.424418in}{1.654944in}}%
\pgfpathlineto{\pgfqpoint{3.426937in}{1.663608in}}%
\pgfpathlineto{\pgfqpoint{3.429456in}{1.645111in}}%
\pgfpathlineto{\pgfqpoint{3.431974in}{1.625744in}}%
\pgfpathlineto{\pgfqpoint{3.434493in}{1.584307in}}%
\pgfpathlineto{\pgfqpoint{3.437012in}{1.571153in}}%
\pgfpathlineto{\pgfqpoint{3.439531in}{1.643379in}}%
\pgfpathlineto{\pgfqpoint{3.442049in}{1.616437in}}%
\pgfpathlineto{\pgfqpoint{3.444568in}{1.626806in}}%
\pgfpathlineto{\pgfqpoint{3.447087in}{1.619080in}}%
\pgfpathlineto{\pgfqpoint{3.449606in}{1.639664in}}%
\pgfpathlineto{\pgfqpoint{3.452124in}{1.624232in}}%
\pgfpathlineto{\pgfqpoint{3.454643in}{1.641951in}}%
\pgfpathlineto{\pgfqpoint{3.457162in}{1.661954in}}%
\pgfpathlineto{\pgfqpoint{3.459681in}{1.666772in}}%
\pgfpathlineto{\pgfqpoint{3.462199in}{1.645593in}}%
\pgfpathlineto{\pgfqpoint{3.464718in}{1.595790in}}%
\pgfpathlineto{\pgfqpoint{3.467237in}{1.603910in}}%
\pgfpathlineto{\pgfqpoint{3.469756in}{1.652649in}}%
\pgfpathlineto{\pgfqpoint{3.472274in}{1.624411in}}%
\pgfpathlineto{\pgfqpoint{3.474793in}{1.629331in}}%
\pgfpathlineto{\pgfqpoint{3.477312in}{1.641088in}}%
\pgfpathlineto{\pgfqpoint{3.479831in}{1.704198in}}%
\pgfpathlineto{\pgfqpoint{3.482349in}{1.671007in}}%
\pgfpathlineto{\pgfqpoint{3.484868in}{1.733844in}}%
\pgfpathlineto{\pgfqpoint{3.487387in}{1.740128in}}%
\pgfpathlineto{\pgfqpoint{3.489906in}{1.737087in}}%
\pgfpathlineto{\pgfqpoint{3.492424in}{1.717509in}}%
\pgfpathlineto{\pgfqpoint{3.494943in}{1.696534in}}%
\pgfpathlineto{\pgfqpoint{3.497462in}{1.728664in}}%
\pgfpathlineto{\pgfqpoint{3.499981in}{1.742359in}}%
\pgfpathlineto{\pgfqpoint{3.502499in}{1.771794in}}%
\pgfpathlineto{\pgfqpoint{3.505018in}{1.741291in}}%
\pgfpathlineto{\pgfqpoint{3.507537in}{1.742113in}}%
\pgfpathlineto{\pgfqpoint{3.510056in}{1.730109in}}%
\pgfpathlineto{\pgfqpoint{3.512574in}{1.775102in}}%
\pgfpathlineto{\pgfqpoint{3.517612in}{1.871822in}}%
\pgfpathlineto{\pgfqpoint{3.520131in}{1.884214in}}%
\pgfpathlineto{\pgfqpoint{3.522649in}{1.871846in}}%
\pgfpathlineto{\pgfqpoint{3.525168in}{1.841911in}}%
\pgfpathlineto{\pgfqpoint{3.527687in}{1.785930in}}%
\pgfpathlineto{\pgfqpoint{3.530206in}{1.765536in}}%
\pgfpathlineto{\pgfqpoint{3.532724in}{1.786351in}}%
\pgfpathlineto{\pgfqpoint{3.535243in}{1.831520in}}%
\pgfpathlineto{\pgfqpoint{3.537762in}{1.808884in}}%
\pgfpathlineto{\pgfqpoint{3.540281in}{1.790809in}}%
\pgfpathlineto{\pgfqpoint{3.542799in}{1.782519in}}%
\pgfpathlineto{\pgfqpoint{3.545318in}{1.850977in}}%
\pgfpathlineto{\pgfqpoint{3.547837in}{1.826556in}}%
\pgfpathlineto{\pgfqpoint{3.550356in}{1.774391in}}%
\pgfpathlineto{\pgfqpoint{3.552874in}{1.785918in}}%
\pgfpathlineto{\pgfqpoint{3.555393in}{1.764341in}}%
\pgfpathlineto{\pgfqpoint{3.557912in}{1.804490in}}%
\pgfpathlineto{\pgfqpoint{3.560431in}{1.786364in}}%
\pgfpathlineto{\pgfqpoint{3.562949in}{1.792606in}}%
\pgfpathlineto{\pgfqpoint{3.565468in}{1.795991in}}%
\pgfpathlineto{\pgfqpoint{3.567987in}{1.783771in}}%
\pgfpathlineto{\pgfqpoint{3.570506in}{1.758501in}}%
\pgfpathlineto{\pgfqpoint{3.573024in}{1.762392in}}%
\pgfpathlineto{\pgfqpoint{3.575543in}{1.778649in}}%
\pgfpathlineto{\pgfqpoint{3.578062in}{1.774075in}}%
\pgfpathlineto{\pgfqpoint{3.580581in}{1.815449in}}%
\pgfpathlineto{\pgfqpoint{3.583099in}{1.841181in}}%
\pgfpathlineto{\pgfqpoint{3.585618in}{1.809965in}}%
\pgfpathlineto{\pgfqpoint{3.588137in}{1.809112in}}%
\pgfpathlineto{\pgfqpoint{3.590656in}{1.832092in}}%
\pgfpathlineto{\pgfqpoint{3.593174in}{1.843690in}}%
\pgfpathlineto{\pgfqpoint{3.595693in}{1.897364in}}%
\pgfpathlineto{\pgfqpoint{3.598212in}{1.902549in}}%
\pgfpathlineto{\pgfqpoint{3.600731in}{1.899888in}}%
\pgfpathlineto{\pgfqpoint{3.603249in}{1.879387in}}%
\pgfpathlineto{\pgfqpoint{3.605768in}{1.885596in}}%
\pgfpathlineto{\pgfqpoint{3.608287in}{1.897794in}}%
\pgfpathlineto{\pgfqpoint{3.610806in}{1.928527in}}%
\pgfpathlineto{\pgfqpoint{3.613324in}{1.917976in}}%
\pgfpathlineto{\pgfqpoint{3.615843in}{1.925303in}}%
\pgfpathlineto{\pgfqpoint{3.618362in}{1.887977in}}%
\pgfpathlineto{\pgfqpoint{3.620881in}{1.865636in}}%
\pgfpathlineto{\pgfqpoint{3.623399in}{1.832071in}}%
\pgfpathlineto{\pgfqpoint{3.625918in}{1.788234in}}%
\pgfpathlineto{\pgfqpoint{3.628437in}{1.760550in}}%
\pgfpathlineto{\pgfqpoint{3.630956in}{1.748217in}}%
\pgfpathlineto{\pgfqpoint{3.633474in}{1.745869in}}%
\pgfpathlineto{\pgfqpoint{3.635993in}{1.768340in}}%
\pgfpathlineto{\pgfqpoint{3.638512in}{1.767338in}}%
\pgfpathlineto{\pgfqpoint{3.643549in}{1.649366in}}%
\pgfpathlineto{\pgfqpoint{3.646068in}{1.674545in}}%
\pgfpathlineto{\pgfqpoint{3.648587in}{1.636819in}}%
\pgfpathlineto{\pgfqpoint{3.651106in}{1.634973in}}%
\pgfpathlineto{\pgfqpoint{3.653624in}{1.626963in}}%
\pgfpathlineto{\pgfqpoint{3.656143in}{1.640482in}}%
\pgfpathlineto{\pgfqpoint{3.658662in}{1.635382in}}%
\pgfpathlineto{\pgfqpoint{3.661181in}{1.682970in}}%
\pgfpathlineto{\pgfqpoint{3.663699in}{1.724077in}}%
\pgfpathlineto{\pgfqpoint{3.666218in}{1.693727in}}%
\pgfpathlineto{\pgfqpoint{3.668737in}{1.713853in}}%
\pgfpathlineto{\pgfqpoint{3.671256in}{1.690082in}}%
\pgfpathlineto{\pgfqpoint{3.673774in}{1.715128in}}%
\pgfpathlineto{\pgfqpoint{3.676293in}{1.724800in}}%
\pgfpathlineto{\pgfqpoint{3.678812in}{1.769627in}}%
\pgfpathlineto{\pgfqpoint{3.681331in}{1.740065in}}%
\pgfpathlineto{\pgfqpoint{3.683849in}{1.741272in}}%
\pgfpathlineto{\pgfqpoint{3.686368in}{1.716804in}}%
\pgfpathlineto{\pgfqpoint{3.688887in}{1.715260in}}%
\pgfpathlineto{\pgfqpoint{3.691406in}{1.684462in}}%
\pgfpathlineto{\pgfqpoint{3.693924in}{1.695109in}}%
\pgfpathlineto{\pgfqpoint{3.696443in}{1.780826in}}%
\pgfpathlineto{\pgfqpoint{3.698962in}{1.778310in}}%
\pgfpathlineto{\pgfqpoint{3.701481in}{1.780185in}}%
\pgfpathlineto{\pgfqpoint{3.703999in}{1.827872in}}%
\pgfpathlineto{\pgfqpoint{3.706518in}{1.852775in}}%
\pgfpathlineto{\pgfqpoint{3.709037in}{1.825408in}}%
\pgfpathlineto{\pgfqpoint{3.711556in}{1.841486in}}%
\pgfpathlineto{\pgfqpoint{3.714074in}{1.846321in}}%
\pgfpathlineto{\pgfqpoint{3.716593in}{1.839602in}}%
\pgfpathlineto{\pgfqpoint{3.719112in}{1.855933in}}%
\pgfpathlineto{\pgfqpoint{3.721631in}{1.907486in}}%
\pgfpathlineto{\pgfqpoint{3.724149in}{1.896687in}}%
\pgfpathlineto{\pgfqpoint{3.726668in}{1.888446in}}%
\pgfpathlineto{\pgfqpoint{3.729187in}{1.883812in}}%
\pgfpathlineto{\pgfqpoint{3.731706in}{1.917933in}}%
\pgfpathlineto{\pgfqpoint{3.734224in}{1.874970in}}%
\pgfpathlineto{\pgfqpoint{3.736743in}{1.824836in}}%
\pgfpathlineto{\pgfqpoint{3.739262in}{1.834228in}}%
\pgfpathlineto{\pgfqpoint{3.741781in}{1.839603in}}%
\pgfpathlineto{\pgfqpoint{3.744299in}{1.887407in}}%
\pgfpathlineto{\pgfqpoint{3.746818in}{1.880227in}}%
\pgfpathlineto{\pgfqpoint{3.749337in}{1.849442in}}%
\pgfpathlineto{\pgfqpoint{3.751856in}{1.890931in}}%
\pgfpathlineto{\pgfqpoint{3.754374in}{1.937901in}}%
\pgfpathlineto{\pgfqpoint{3.756893in}{1.954838in}}%
\pgfpathlineto{\pgfqpoint{3.761931in}{1.985836in}}%
\pgfpathlineto{\pgfqpoint{3.764449in}{2.018322in}}%
\pgfpathlineto{\pgfqpoint{3.766968in}{1.985461in}}%
\pgfpathlineto{\pgfqpoint{3.769487in}{1.980266in}}%
\pgfpathlineto{\pgfqpoint{3.772006in}{1.977631in}}%
\pgfpathlineto{\pgfqpoint{3.774524in}{1.956536in}}%
\pgfpathlineto{\pgfqpoint{3.777043in}{1.944292in}}%
\pgfpathlineto{\pgfqpoint{3.779562in}{1.959544in}}%
\pgfpathlineto{\pgfqpoint{3.782081in}{1.923249in}}%
\pgfpathlineto{\pgfqpoint{3.784599in}{1.917934in}}%
\pgfpathlineto{\pgfqpoint{3.787118in}{1.879770in}}%
\pgfpathlineto{\pgfqpoint{3.789637in}{1.878331in}}%
\pgfpathlineto{\pgfqpoint{3.792156in}{1.830083in}}%
\pgfpathlineto{\pgfqpoint{3.794674in}{1.830235in}}%
\pgfpathlineto{\pgfqpoint{3.797193in}{1.859951in}}%
\pgfpathlineto{\pgfqpoint{3.799712in}{1.853683in}}%
\pgfpathlineto{\pgfqpoint{3.802231in}{1.839351in}}%
\pgfpathlineto{\pgfqpoint{3.804749in}{1.838656in}}%
\pgfpathlineto{\pgfqpoint{3.807268in}{1.805511in}}%
\pgfpathlineto{\pgfqpoint{3.809787in}{1.844538in}}%
\pgfpathlineto{\pgfqpoint{3.812306in}{1.819747in}}%
\pgfpathlineto{\pgfqpoint{3.814824in}{1.835860in}}%
\pgfpathlineto{\pgfqpoint{3.817343in}{1.824193in}}%
\pgfpathlineto{\pgfqpoint{3.819862in}{1.869854in}}%
\pgfpathlineto{\pgfqpoint{3.822381in}{1.875612in}}%
\pgfpathlineto{\pgfqpoint{3.824899in}{1.882319in}}%
\pgfpathlineto{\pgfqpoint{3.827418in}{1.875493in}}%
\pgfpathlineto{\pgfqpoint{3.829937in}{1.891402in}}%
\pgfpathlineto{\pgfqpoint{3.832456in}{1.919626in}}%
\pgfpathlineto{\pgfqpoint{3.834974in}{1.911370in}}%
\pgfpathlineto{\pgfqpoint{3.837493in}{1.936720in}}%
\pgfpathlineto{\pgfqpoint{3.840012in}{1.907945in}}%
\pgfpathlineto{\pgfqpoint{3.842531in}{1.867581in}}%
\pgfpathlineto{\pgfqpoint{3.845049in}{1.854869in}}%
\pgfpathlineto{\pgfqpoint{3.847568in}{1.855201in}}%
\pgfpathlineto{\pgfqpoint{3.850087in}{1.887487in}}%
\pgfpathlineto{\pgfqpoint{3.852606in}{1.909342in}}%
\pgfpathlineto{\pgfqpoint{3.857643in}{1.995465in}}%
\pgfpathlineto{\pgfqpoint{3.860162in}{1.976456in}}%
\pgfpathlineto{\pgfqpoint{3.862681in}{1.965208in}}%
\pgfpathlineto{\pgfqpoint{3.865199in}{1.969692in}}%
\pgfpathlineto{\pgfqpoint{3.867718in}{1.987803in}}%
\pgfpathlineto{\pgfqpoint{3.870237in}{1.993895in}}%
\pgfpathlineto{\pgfqpoint{3.872756in}{1.973665in}}%
\pgfpathlineto{\pgfqpoint{3.875274in}{1.952281in}}%
\pgfpathlineto{\pgfqpoint{3.877793in}{1.924588in}}%
\pgfpathlineto{\pgfqpoint{3.880312in}{1.901777in}}%
\pgfpathlineto{\pgfqpoint{3.882831in}{1.903502in}}%
\pgfpathlineto{\pgfqpoint{3.885349in}{1.914010in}}%
\pgfpathlineto{\pgfqpoint{3.887868in}{1.918917in}}%
\pgfpathlineto{\pgfqpoint{3.890387in}{1.921756in}}%
\pgfpathlineto{\pgfqpoint{3.892906in}{1.958175in}}%
\pgfpathlineto{\pgfqpoint{3.895424in}{1.975840in}}%
\pgfpathlineto{\pgfqpoint{3.897943in}{1.969462in}}%
\pgfpathlineto{\pgfqpoint{3.900462in}{1.959842in}}%
\pgfpathlineto{\pgfqpoint{3.902981in}{1.947082in}}%
\pgfpathlineto{\pgfqpoint{3.905499in}{1.909871in}}%
\pgfpathlineto{\pgfqpoint{3.908018in}{1.908944in}}%
\pgfpathlineto{\pgfqpoint{3.910537in}{1.905647in}}%
\pgfpathlineto{\pgfqpoint{3.913056in}{1.908304in}}%
\pgfpathlineto{\pgfqpoint{3.915574in}{1.861003in}}%
\pgfpathlineto{\pgfqpoint{3.918093in}{1.890191in}}%
\pgfpathlineto{\pgfqpoint{3.920612in}{1.869133in}}%
\pgfpathlineto{\pgfqpoint{3.923131in}{1.882894in}}%
\pgfpathlineto{\pgfqpoint{3.925649in}{1.897635in}}%
\pgfpathlineto{\pgfqpoint{3.928168in}{1.898778in}}%
\pgfpathlineto{\pgfqpoint{3.930687in}{1.946169in}}%
\pgfpathlineto{\pgfqpoint{3.933206in}{1.943228in}}%
\pgfpathlineto{\pgfqpoint{3.935724in}{1.956748in}}%
\pgfpathlineto{\pgfqpoint{3.938243in}{2.039619in}}%
\pgfpathlineto{\pgfqpoint{3.940762in}{2.045901in}}%
\pgfpathlineto{\pgfqpoint{3.943281in}{2.016673in}}%
\pgfpathlineto{\pgfqpoint{3.945799in}{1.973210in}}%
\pgfpathlineto{\pgfqpoint{3.948318in}{1.957536in}}%
\pgfpathlineto{\pgfqpoint{3.950837in}{1.921538in}}%
\pgfpathlineto{\pgfqpoint{3.953356in}{1.891724in}}%
\pgfpathlineto{\pgfqpoint{3.955874in}{1.870555in}}%
\pgfpathlineto{\pgfqpoint{3.958393in}{1.823937in}}%
\pgfpathlineto{\pgfqpoint{3.960912in}{1.814849in}}%
\pgfpathlineto{\pgfqpoint{3.963431in}{1.791639in}}%
\pgfpathlineto{\pgfqpoint{3.965949in}{1.805065in}}%
\pgfpathlineto{\pgfqpoint{3.968468in}{1.819916in}}%
\pgfpathlineto{\pgfqpoint{3.970987in}{1.801865in}}%
\pgfpathlineto{\pgfqpoint{3.973506in}{1.824621in}}%
\pgfpathlineto{\pgfqpoint{3.976024in}{1.800742in}}%
\pgfpathlineto{\pgfqpoint{3.978543in}{1.783246in}}%
\pgfpathlineto{\pgfqpoint{3.981062in}{1.783309in}}%
\pgfpathlineto{\pgfqpoint{3.983581in}{1.714871in}}%
\pgfpathlineto{\pgfqpoint{3.986099in}{1.714283in}}%
\pgfpathlineto{\pgfqpoint{3.988618in}{1.688974in}}%
\pgfpathlineto{\pgfqpoint{3.991137in}{1.691041in}}%
\pgfpathlineto{\pgfqpoint{3.993656in}{1.718078in}}%
\pgfpathlineto{\pgfqpoint{3.996174in}{1.736381in}}%
\pgfpathlineto{\pgfqpoint{3.998693in}{1.721547in}}%
\pgfpathlineto{\pgfqpoint{4.001212in}{1.675844in}}%
\pgfpathlineto{\pgfqpoint{4.003731in}{1.659195in}}%
\pgfpathlineto{\pgfqpoint{4.006249in}{1.701545in}}%
\pgfpathlineto{\pgfqpoint{4.008768in}{1.681602in}}%
\pgfpathlineto{\pgfqpoint{4.011287in}{1.698895in}}%
\pgfpathlineto{\pgfqpoint{4.013806in}{1.708982in}}%
\pgfpathlineto{\pgfqpoint{4.016324in}{1.668581in}}%
\pgfpathlineto{\pgfqpoint{4.018843in}{1.682020in}}%
\pgfpathlineto{\pgfqpoint{4.021362in}{1.681295in}}%
\pgfpathlineto{\pgfqpoint{4.023881in}{1.710376in}}%
\pgfpathlineto{\pgfqpoint{4.026399in}{1.705272in}}%
\pgfpathlineto{\pgfqpoint{4.028918in}{1.714884in}}%
\pgfpathlineto{\pgfqpoint{4.031437in}{1.718776in}}%
\pgfpathlineto{\pgfqpoint{4.033956in}{1.740888in}}%
\pgfpathlineto{\pgfqpoint{4.036474in}{1.734954in}}%
\pgfpathlineto{\pgfqpoint{4.038993in}{1.696584in}}%
\pgfpathlineto{\pgfqpoint{4.041512in}{1.624777in}}%
\pgfpathlineto{\pgfqpoint{4.044031in}{1.587948in}}%
\pgfpathlineto{\pgfqpoint{4.046549in}{1.575134in}}%
\pgfpathlineto{\pgfqpoint{4.049068in}{1.568481in}}%
\pgfpathlineto{\pgfqpoint{4.051587in}{1.575430in}}%
\pgfpathlineto{\pgfqpoint{4.054106in}{1.584937in}}%
\pgfpathlineto{\pgfqpoint{4.056624in}{1.639505in}}%
\pgfpathlineto{\pgfqpoint{4.059143in}{1.656214in}}%
\pgfpathlineto{\pgfqpoint{4.061662in}{1.712024in}}%
\pgfpathlineto{\pgfqpoint{4.064181in}{1.674885in}}%
\pgfpathlineto{\pgfqpoint{4.066699in}{1.666688in}}%
\pgfpathlineto{\pgfqpoint{4.069218in}{1.658100in}}%
\pgfpathlineto{\pgfqpoint{4.071737in}{1.657665in}}%
\pgfpathlineto{\pgfqpoint{4.074256in}{1.657561in}}%
\pgfpathlineto{\pgfqpoint{4.076774in}{1.658457in}}%
\pgfpathlineto{\pgfqpoint{4.079293in}{1.653900in}}%
\pgfpathlineto{\pgfqpoint{4.081812in}{1.657889in}}%
\pgfpathlineto{\pgfqpoint{4.084331in}{1.703833in}}%
\pgfpathlineto{\pgfqpoint{4.086849in}{1.721854in}}%
\pgfpathlineto{\pgfqpoint{4.089368in}{1.678224in}}%
\pgfpathlineto{\pgfqpoint{4.091887in}{1.637761in}}%
\pgfpathlineto{\pgfqpoint{4.094406in}{1.659125in}}%
\pgfpathlineto{\pgfqpoint{4.096924in}{1.699768in}}%
\pgfpathlineto{\pgfqpoint{4.099443in}{1.679827in}}%
\pgfpathlineto{\pgfqpoint{4.101962in}{1.677431in}}%
\pgfpathlineto{\pgfqpoint{4.104481in}{1.644657in}}%
\pgfpathlineto{\pgfqpoint{4.106999in}{1.625827in}}%
\pgfpathlineto{\pgfqpoint{4.112037in}{1.650088in}}%
\pgfpathlineto{\pgfqpoint{4.114556in}{1.647034in}}%
\pgfpathlineto{\pgfqpoint{4.117074in}{1.590455in}}%
\pgfpathlineto{\pgfqpoint{4.119593in}{1.606850in}}%
\pgfpathlineto{\pgfqpoint{4.122112in}{1.632774in}}%
\pgfpathlineto{\pgfqpoint{4.124631in}{1.678947in}}%
\pgfpathlineto{\pgfqpoint{4.127149in}{1.750824in}}%
\pgfpathlineto{\pgfqpoint{4.129668in}{1.759205in}}%
\pgfpathlineto{\pgfqpoint{4.132187in}{1.791218in}}%
\pgfpathlineto{\pgfqpoint{4.134706in}{1.861390in}}%
\pgfpathlineto{\pgfqpoint{4.137224in}{1.893517in}}%
\pgfpathlineto{\pgfqpoint{4.139743in}{1.928580in}}%
\pgfpathlineto{\pgfqpoint{4.142262in}{2.011464in}}%
\pgfpathlineto{\pgfqpoint{4.144781in}{2.012214in}}%
\pgfpathlineto{\pgfqpoint{4.147299in}{2.018476in}}%
\pgfpathlineto{\pgfqpoint{4.149818in}{1.989607in}}%
\pgfpathlineto{\pgfqpoint{4.152337in}{2.005134in}}%
\pgfpathlineto{\pgfqpoint{4.154856in}{2.040166in}}%
\pgfpathlineto{\pgfqpoint{4.157374in}{2.008861in}}%
\pgfpathlineto{\pgfqpoint{4.159893in}{2.019995in}}%
\pgfpathlineto{\pgfqpoint{4.162412in}{2.045137in}}%
\pgfpathlineto{\pgfqpoint{4.164931in}{2.021945in}}%
\pgfpathlineto{\pgfqpoint{4.167449in}{2.046286in}}%
\pgfpathlineto{\pgfqpoint{4.169968in}{2.048549in}}%
\pgfpathlineto{\pgfqpoint{4.172487in}{2.041474in}}%
\pgfpathlineto{\pgfqpoint{4.175006in}{2.033089in}}%
\pgfpathlineto{\pgfqpoint{4.177524in}{2.002922in}}%
\pgfpathlineto{\pgfqpoint{4.180043in}{2.013157in}}%
\pgfpathlineto{\pgfqpoint{4.182562in}{2.006878in}}%
\pgfpathlineto{\pgfqpoint{4.185081in}{2.028285in}}%
\pgfpathlineto{\pgfqpoint{4.187599in}{2.030312in}}%
\pgfpathlineto{\pgfqpoint{4.190118in}{1.965235in}}%
\pgfpathlineto{\pgfqpoint{4.192637in}{1.969678in}}%
\pgfpathlineto{\pgfqpoint{4.195156in}{1.965213in}}%
\pgfpathlineto{\pgfqpoint{4.197674in}{1.998446in}}%
\pgfpathlineto{\pgfqpoint{4.200193in}{2.025625in}}%
\pgfpathlineto{\pgfqpoint{4.202712in}{2.049055in}}%
\pgfpathlineto{\pgfqpoint{4.205231in}{1.988848in}}%
\pgfpathlineto{\pgfqpoint{4.207749in}{1.973684in}}%
\pgfpathlineto{\pgfqpoint{4.210268in}{1.965971in}}%
\pgfpathlineto{\pgfqpoint{4.212787in}{1.956756in}}%
\pgfpathlineto{\pgfqpoint{4.215306in}{1.977411in}}%
\pgfpathlineto{\pgfqpoint{4.217824in}{2.030893in}}%
\pgfpathlineto{\pgfqpoint{4.220343in}{2.034083in}}%
\pgfpathlineto{\pgfqpoint{4.222862in}{2.025435in}}%
\pgfpathlineto{\pgfqpoint{4.225381in}{2.047044in}}%
\pgfpathlineto{\pgfqpoint{4.227899in}{1.996758in}}%
\pgfpathlineto{\pgfqpoint{4.230418in}{2.024608in}}%
\pgfpathlineto{\pgfqpoint{4.232937in}{2.018486in}}%
\pgfpathlineto{\pgfqpoint{4.235456in}{2.027985in}}%
\pgfpathlineto{\pgfqpoint{4.237974in}{2.079677in}}%
\pgfpathlineto{\pgfqpoint{4.240493in}{2.108698in}}%
\pgfpathlineto{\pgfqpoint{4.243012in}{2.154530in}}%
\pgfpathlineto{\pgfqpoint{4.245531in}{2.121489in}}%
\pgfpathlineto{\pgfqpoint{4.248049in}{2.101111in}}%
\pgfpathlineto{\pgfqpoint{4.250568in}{2.048882in}}%
\pgfpathlineto{\pgfqpoint{4.253087in}{2.100825in}}%
\pgfpathlineto{\pgfqpoint{4.255606in}{2.090254in}}%
\pgfpathlineto{\pgfqpoint{4.258124in}{2.089311in}}%
\pgfpathlineto{\pgfqpoint{4.260643in}{2.107732in}}%
\pgfpathlineto{\pgfqpoint{4.263162in}{2.120479in}}%
\pgfpathlineto{\pgfqpoint{4.265681in}{2.169853in}}%
\pgfpathlineto{\pgfqpoint{4.268199in}{2.161496in}}%
\pgfpathlineto{\pgfqpoint{4.270718in}{2.187096in}}%
\pgfpathlineto{\pgfqpoint{4.273237in}{2.199326in}}%
\pgfpathlineto{\pgfqpoint{4.275756in}{2.113634in}}%
\pgfpathlineto{\pgfqpoint{4.278274in}{2.098000in}}%
\pgfpathlineto{\pgfqpoint{4.280793in}{2.088637in}}%
\pgfpathlineto{\pgfqpoint{4.283312in}{2.049757in}}%
\pgfpathlineto{\pgfqpoint{4.285831in}{2.090068in}}%
\pgfpathlineto{\pgfqpoint{4.288349in}{2.071934in}}%
\pgfpathlineto{\pgfqpoint{4.290868in}{2.030158in}}%
\pgfpathlineto{\pgfqpoint{4.293387in}{2.072279in}}%
\pgfpathlineto{\pgfqpoint{4.295906in}{2.097051in}}%
\pgfpathlineto{\pgfqpoint{4.298424in}{2.102971in}}%
\pgfpathlineto{\pgfqpoint{4.300943in}{2.078066in}}%
\pgfpathlineto{\pgfqpoint{4.303462in}{2.057556in}}%
\pgfpathlineto{\pgfqpoint{4.305981in}{2.063882in}}%
\pgfpathlineto{\pgfqpoint{4.308499in}{2.060452in}}%
\pgfpathlineto{\pgfqpoint{4.311018in}{2.039520in}}%
\pgfpathlineto{\pgfqpoint{4.316056in}{1.917394in}}%
\pgfpathlineto{\pgfqpoint{4.318574in}{1.867628in}}%
\pgfpathlineto{\pgfqpoint{4.321093in}{1.861133in}}%
\pgfpathlineto{\pgfqpoint{4.323612in}{1.820946in}}%
\pgfpathlineto{\pgfqpoint{4.326131in}{1.798468in}}%
\pgfpathlineto{\pgfqpoint{4.328649in}{1.782873in}}%
\pgfpathlineto{\pgfqpoint{4.331168in}{1.813325in}}%
\pgfpathlineto{\pgfqpoint{4.333687in}{1.741711in}}%
\pgfpathlineto{\pgfqpoint{4.336206in}{1.758607in}}%
\pgfpathlineto{\pgfqpoint{4.338724in}{1.747998in}}%
\pgfpathlineto{\pgfqpoint{4.341243in}{1.720492in}}%
\pgfpathlineto{\pgfqpoint{4.343762in}{1.729109in}}%
\pgfpathlineto{\pgfqpoint{4.346281in}{1.787977in}}%
\pgfpathlineto{\pgfqpoint{4.348799in}{1.753534in}}%
\pgfpathlineto{\pgfqpoint{4.351318in}{1.763264in}}%
\pgfpathlineto{\pgfqpoint{4.353837in}{1.769499in}}%
\pgfpathlineto{\pgfqpoint{4.356356in}{1.769736in}}%
\pgfpathlineto{\pgfqpoint{4.358874in}{1.735300in}}%
\pgfpathlineto{\pgfqpoint{4.361393in}{1.711016in}}%
\pgfpathlineto{\pgfqpoint{4.363912in}{1.720366in}}%
\pgfpathlineto{\pgfqpoint{4.366431in}{1.677743in}}%
\pgfpathlineto{\pgfqpoint{4.368949in}{1.689934in}}%
\pgfpathlineto{\pgfqpoint{4.371468in}{1.647071in}}%
\pgfpathlineto{\pgfqpoint{4.373987in}{1.642750in}}%
\pgfpathlineto{\pgfqpoint{4.376506in}{1.640913in}}%
\pgfpathlineto{\pgfqpoint{4.379024in}{1.655665in}}%
\pgfpathlineto{\pgfqpoint{4.381543in}{1.631016in}}%
\pgfpathlineto{\pgfqpoint{4.384062in}{1.598244in}}%
\pgfpathlineto{\pgfqpoint{4.386581in}{1.612165in}}%
\pgfpathlineto{\pgfqpoint{4.391618in}{1.668099in}}%
\pgfpathlineto{\pgfqpoint{4.394137in}{1.684509in}}%
\pgfpathlineto{\pgfqpoint{4.396656in}{1.686645in}}%
\pgfpathlineto{\pgfqpoint{4.399174in}{1.707686in}}%
\pgfpathlineto{\pgfqpoint{4.401693in}{1.717726in}}%
\pgfpathlineto{\pgfqpoint{4.404212in}{1.686915in}}%
\pgfpathlineto{\pgfqpoint{4.406731in}{1.687123in}}%
\pgfpathlineto{\pgfqpoint{4.409249in}{1.691857in}}%
\pgfpathlineto{\pgfqpoint{4.411768in}{1.690483in}}%
\pgfpathlineto{\pgfqpoint{4.414287in}{1.676354in}}%
\pgfpathlineto{\pgfqpoint{4.416806in}{1.709033in}}%
\pgfpathlineto{\pgfqpoint{4.419324in}{1.683992in}}%
\pgfpathlineto{\pgfqpoint{4.421843in}{1.629701in}}%
\pgfpathlineto{\pgfqpoint{4.424362in}{1.614548in}}%
\pgfpathlineto{\pgfqpoint{4.426881in}{1.619326in}}%
\pgfpathlineto{\pgfqpoint{4.429399in}{1.632744in}}%
\pgfpathlineto{\pgfqpoint{4.431918in}{1.639193in}}%
\pgfpathlineto{\pgfqpoint{4.434437in}{1.632818in}}%
\pgfpathlineto{\pgfqpoint{4.436956in}{1.636760in}}%
\pgfpathlineto{\pgfqpoint{4.439474in}{1.673668in}}%
\pgfpathlineto{\pgfqpoint{4.444512in}{1.640324in}}%
\pgfpathlineto{\pgfqpoint{4.447031in}{1.633242in}}%
\pgfpathlineto{\pgfqpoint{4.449549in}{1.664964in}}%
\pgfpathlineto{\pgfqpoint{4.452068in}{1.644914in}}%
\pgfpathlineto{\pgfqpoint{4.454587in}{1.642839in}}%
\pgfpathlineto{\pgfqpoint{4.457106in}{1.659382in}}%
\pgfpathlineto{\pgfqpoint{4.459624in}{1.652675in}}%
\pgfpathlineto{\pgfqpoint{4.462143in}{1.636512in}}%
\pgfpathlineto{\pgfqpoint{4.464662in}{1.663575in}}%
\pgfpathlineto{\pgfqpoint{4.467181in}{1.703837in}}%
\pgfpathlineto{\pgfqpoint{4.469699in}{1.682290in}}%
\pgfpathlineto{\pgfqpoint{4.472218in}{1.666610in}}%
\pgfpathlineto{\pgfqpoint{4.474737in}{1.669564in}}%
\pgfpathlineto{\pgfqpoint{4.477256in}{1.654785in}}%
\pgfpathlineto{\pgfqpoint{4.479774in}{1.632457in}}%
\pgfpathlineto{\pgfqpoint{4.482293in}{1.600259in}}%
\pgfpathlineto{\pgfqpoint{4.484812in}{1.651609in}}%
\pgfpathlineto{\pgfqpoint{4.487331in}{1.649117in}}%
\pgfpathlineto{\pgfqpoint{4.489849in}{1.631893in}}%
\pgfpathlineto{\pgfqpoint{4.492368in}{1.648892in}}%
\pgfpathlineto{\pgfqpoint{4.494887in}{1.623360in}}%
\pgfpathlineto{\pgfqpoint{4.497406in}{1.703826in}}%
\pgfpathlineto{\pgfqpoint{4.499924in}{1.713501in}}%
\pgfpathlineto{\pgfqpoint{4.502443in}{1.747411in}}%
\pgfpathlineto{\pgfqpoint{4.504962in}{1.720806in}}%
\pgfpathlineto{\pgfqpoint{4.507481in}{1.731317in}}%
\pgfpathlineto{\pgfqpoint{4.512518in}{1.647156in}}%
\pgfpathlineto{\pgfqpoint{4.515037in}{1.659239in}}%
\pgfpathlineto{\pgfqpoint{4.517556in}{1.595567in}}%
\pgfpathlineto{\pgfqpoint{4.520074in}{1.558577in}}%
\pgfpathlineto{\pgfqpoint{4.522593in}{1.523306in}}%
\pgfpathlineto{\pgfqpoint{4.525112in}{1.533268in}}%
\pgfpathlineto{\pgfqpoint{4.527631in}{1.544847in}}%
\pgfpathlineto{\pgfqpoint{4.530149in}{1.596238in}}%
\pgfpathlineto{\pgfqpoint{4.532668in}{1.598878in}}%
\pgfpathlineto{\pgfqpoint{4.535187in}{1.588973in}}%
\pgfpathlineto{\pgfqpoint{4.537706in}{1.614769in}}%
\pgfpathlineto{\pgfqpoint{4.540224in}{1.567345in}}%
\pgfpathlineto{\pgfqpoint{4.542743in}{1.601204in}}%
\pgfpathlineto{\pgfqpoint{4.545262in}{1.600100in}}%
\pgfpathlineto{\pgfqpoint{4.547781in}{1.613218in}}%
\pgfpathlineto{\pgfqpoint{4.550299in}{1.591740in}}%
\pgfpathlineto{\pgfqpoint{4.552818in}{1.534822in}}%
\pgfpathlineto{\pgfqpoint{4.555337in}{1.571785in}}%
\pgfpathlineto{\pgfqpoint{4.557856in}{1.599298in}}%
\pgfpathlineto{\pgfqpoint{4.560374in}{1.583906in}}%
\pgfpathlineto{\pgfqpoint{4.562893in}{1.617553in}}%
\pgfpathlineto{\pgfqpoint{4.565412in}{1.598806in}}%
\pgfpathlineto{\pgfqpoint{4.567931in}{1.605903in}}%
\pgfpathlineto{\pgfqpoint{4.570449in}{1.591508in}}%
\pgfpathlineto{\pgfqpoint{4.572968in}{1.578106in}}%
\pgfpathlineto{\pgfqpoint{4.575487in}{1.612167in}}%
\pgfpathlineto{\pgfqpoint{4.578006in}{1.605387in}}%
\pgfpathlineto{\pgfqpoint{4.580524in}{1.607339in}}%
\pgfpathlineto{\pgfqpoint{4.583043in}{1.542598in}}%
\pgfpathlineto{\pgfqpoint{4.585562in}{1.594810in}}%
\pgfpathlineto{\pgfqpoint{4.590599in}{1.531544in}}%
\pgfpathlineto{\pgfqpoint{4.593118in}{1.523338in}}%
\pgfpathlineto{\pgfqpoint{4.595637in}{1.615753in}}%
\pgfpathlineto{\pgfqpoint{4.598156in}{1.625959in}}%
\pgfpathlineto{\pgfqpoint{4.600674in}{1.623382in}}%
\pgfpathlineto{\pgfqpoint{4.603193in}{1.595742in}}%
\pgfpathlineto{\pgfqpoint{4.605712in}{1.603846in}}%
\pgfpathlineto{\pgfqpoint{4.608231in}{1.591085in}}%
\pgfpathlineto{\pgfqpoint{4.613268in}{1.611314in}}%
\pgfpathlineto{\pgfqpoint{4.615787in}{1.624288in}}%
\pgfpathlineto{\pgfqpoint{4.618306in}{1.648112in}}%
\pgfpathlineto{\pgfqpoint{4.620824in}{1.706753in}}%
\pgfpathlineto{\pgfqpoint{4.623343in}{1.718730in}}%
\pgfpathlineto{\pgfqpoint{4.625862in}{1.725321in}}%
\pgfpathlineto{\pgfqpoint{4.628381in}{1.731370in}}%
\pgfpathlineto{\pgfqpoint{4.630899in}{1.702730in}}%
\pgfpathlineto{\pgfqpoint{4.633418in}{1.679738in}}%
\pgfpathlineto{\pgfqpoint{4.635937in}{1.673911in}}%
\pgfpathlineto{\pgfqpoint{4.638456in}{1.636810in}}%
\pgfpathlineto{\pgfqpoint{4.640974in}{1.704030in}}%
\pgfpathlineto{\pgfqpoint{4.643493in}{1.730383in}}%
\pgfpathlineto{\pgfqpoint{4.646012in}{1.754090in}}%
\pgfpathlineto{\pgfqpoint{4.648531in}{1.779670in}}%
\pgfpathlineto{\pgfqpoint{4.651049in}{1.790603in}}%
\pgfpathlineto{\pgfqpoint{4.653568in}{1.762553in}}%
\pgfpathlineto{\pgfqpoint{4.656087in}{1.749381in}}%
\pgfpathlineto{\pgfqpoint{4.658606in}{1.708058in}}%
\pgfpathlineto{\pgfqpoint{4.661124in}{1.739756in}}%
\pgfpathlineto{\pgfqpoint{4.663643in}{1.719125in}}%
\pgfpathlineto{\pgfqpoint{4.666162in}{1.691401in}}%
\pgfpathlineto{\pgfqpoint{4.668681in}{1.684073in}}%
\pgfpathlineto{\pgfqpoint{4.671199in}{1.725678in}}%
\pgfpathlineto{\pgfqpoint{4.673718in}{1.710062in}}%
\pgfpathlineto{\pgfqpoint{4.676237in}{1.723176in}}%
\pgfpathlineto{\pgfqpoint{4.678756in}{1.723838in}}%
\pgfpathlineto{\pgfqpoint{4.681274in}{1.711824in}}%
\pgfpathlineto{\pgfqpoint{4.683793in}{1.661094in}}%
\pgfpathlineto{\pgfqpoint{4.686312in}{1.646452in}}%
\pgfpathlineto{\pgfqpoint{4.688831in}{1.664720in}}%
\pgfpathlineto{\pgfqpoint{4.691349in}{1.705250in}}%
\pgfpathlineto{\pgfqpoint{4.693868in}{1.693366in}}%
\pgfpathlineto{\pgfqpoint{4.696387in}{1.685294in}}%
\pgfpathlineto{\pgfqpoint{4.698906in}{1.673553in}}%
\pgfpathlineto{\pgfqpoint{4.701424in}{1.636778in}}%
\pgfpathlineto{\pgfqpoint{4.703943in}{1.619258in}}%
\pgfpathlineto{\pgfqpoint{4.706462in}{1.631975in}}%
\pgfpathlineto{\pgfqpoint{4.708981in}{1.643775in}}%
\pgfpathlineto{\pgfqpoint{4.711499in}{1.671390in}}%
\pgfpathlineto{\pgfqpoint{4.714018in}{1.635557in}}%
\pgfpathlineto{\pgfqpoint{4.716537in}{1.589935in}}%
\pgfpathlineto{\pgfqpoint{4.719056in}{1.584122in}}%
\pgfpathlineto{\pgfqpoint{4.721574in}{1.595625in}}%
\pgfpathlineto{\pgfqpoint{4.724093in}{1.629944in}}%
\pgfpathlineto{\pgfqpoint{4.726612in}{1.660824in}}%
\pgfpathlineto{\pgfqpoint{4.729131in}{1.694974in}}%
\pgfpathlineto{\pgfqpoint{4.731649in}{1.636735in}}%
\pgfpathlineto{\pgfqpoint{4.734168in}{1.591442in}}%
\pgfpathlineto{\pgfqpoint{4.736687in}{1.596621in}}%
\pgfpathlineto{\pgfqpoint{4.739206in}{1.583092in}}%
\pgfpathlineto{\pgfqpoint{4.744243in}{1.629610in}}%
\pgfpathlineto{\pgfqpoint{4.746762in}{1.650900in}}%
\pgfpathlineto{\pgfqpoint{4.749281in}{1.689316in}}%
\pgfpathlineto{\pgfqpoint{4.751799in}{1.685414in}}%
\pgfpathlineto{\pgfqpoint{4.754318in}{1.736556in}}%
\pgfpathlineto{\pgfqpoint{4.756837in}{1.722466in}}%
\pgfpathlineto{\pgfqpoint{4.759356in}{1.728184in}}%
\pgfpathlineto{\pgfqpoint{4.761874in}{1.688525in}}%
\pgfpathlineto{\pgfqpoint{4.764393in}{1.665464in}}%
\pgfpathlineto{\pgfqpoint{4.766912in}{1.695568in}}%
\pgfpathlineto{\pgfqpoint{4.769431in}{1.689038in}}%
\pgfpathlineto{\pgfqpoint{4.774468in}{1.623537in}}%
\pgfpathlineto{\pgfqpoint{4.776987in}{1.603398in}}%
\pgfpathlineto{\pgfqpoint{4.779506in}{1.634650in}}%
\pgfpathlineto{\pgfqpoint{4.782024in}{1.621104in}}%
\pgfpathlineto{\pgfqpoint{4.784543in}{1.651803in}}%
\pgfpathlineto{\pgfqpoint{4.787062in}{1.656677in}}%
\pgfpathlineto{\pgfqpoint{4.789581in}{1.658056in}}%
\pgfpathlineto{\pgfqpoint{4.792099in}{1.671433in}}%
\pgfpathlineto{\pgfqpoint{4.794618in}{1.670416in}}%
\pgfpathlineto{\pgfqpoint{4.797137in}{1.633861in}}%
\pgfpathlineto{\pgfqpoint{4.799656in}{1.634627in}}%
\pgfpathlineto{\pgfqpoint{4.802174in}{1.665840in}}%
\pgfpathlineto{\pgfqpoint{4.804693in}{1.641689in}}%
\pgfpathlineto{\pgfqpoint{4.807212in}{1.638563in}}%
\pgfpathlineto{\pgfqpoint{4.809731in}{1.642327in}}%
\pgfpathlineto{\pgfqpoint{4.812249in}{1.629115in}}%
\pgfpathlineto{\pgfqpoint{4.814768in}{1.600615in}}%
\pgfpathlineto{\pgfqpoint{4.817287in}{1.589211in}}%
\pgfpathlineto{\pgfqpoint{4.819806in}{1.594217in}}%
\pgfpathlineto{\pgfqpoint{4.822324in}{1.640061in}}%
\pgfpathlineto{\pgfqpoint{4.824843in}{1.617710in}}%
\pgfpathlineto{\pgfqpoint{4.827362in}{1.646154in}}%
\pgfpathlineto{\pgfqpoint{4.829881in}{1.672571in}}%
\pgfpathlineto{\pgfqpoint{4.832399in}{1.731888in}}%
\pgfpathlineto{\pgfqpoint{4.834918in}{1.729964in}}%
\pgfpathlineto{\pgfqpoint{4.837437in}{1.759777in}}%
\pgfpathlineto{\pgfqpoint{4.839956in}{1.798897in}}%
\pgfpathlineto{\pgfqpoint{4.842474in}{1.840705in}}%
\pgfpathlineto{\pgfqpoint{4.844993in}{1.763502in}}%
\pgfpathlineto{\pgfqpoint{4.847512in}{1.786528in}}%
\pgfpathlineto{\pgfqpoint{4.850031in}{1.778334in}}%
\pgfpathlineto{\pgfqpoint{4.852549in}{1.807219in}}%
\pgfpathlineto{\pgfqpoint{4.855068in}{1.814003in}}%
\pgfpathlineto{\pgfqpoint{4.857587in}{1.838909in}}%
\pgfpathlineto{\pgfqpoint{4.860106in}{1.826007in}}%
\pgfpathlineto{\pgfqpoint{4.862624in}{1.845270in}}%
\pgfpathlineto{\pgfqpoint{4.865143in}{1.872209in}}%
\pgfpathlineto{\pgfqpoint{4.867662in}{1.871996in}}%
\pgfpathlineto{\pgfqpoint{4.870181in}{1.868442in}}%
\pgfpathlineto{\pgfqpoint{4.872699in}{1.863719in}}%
\pgfpathlineto{\pgfqpoint{4.875218in}{1.899718in}}%
\pgfpathlineto{\pgfqpoint{4.877737in}{1.850360in}}%
\pgfpathlineto{\pgfqpoint{4.880256in}{1.803868in}}%
\pgfpathlineto{\pgfqpoint{4.882774in}{1.826719in}}%
\pgfpathlineto{\pgfqpoint{4.885293in}{1.858416in}}%
\pgfpathlineto{\pgfqpoint{4.887812in}{1.846137in}}%
\pgfpathlineto{\pgfqpoint{4.890331in}{1.843088in}}%
\pgfpathlineto{\pgfqpoint{4.892849in}{1.836738in}}%
\pgfpathlineto{\pgfqpoint{4.895368in}{1.861913in}}%
\pgfpathlineto{\pgfqpoint{4.897887in}{1.857121in}}%
\pgfpathlineto{\pgfqpoint{4.900406in}{1.823256in}}%
\pgfpathlineto{\pgfqpoint{4.902924in}{1.784818in}}%
\pgfpathlineto{\pgfqpoint{4.905443in}{1.801458in}}%
\pgfpathlineto{\pgfqpoint{4.907962in}{1.831095in}}%
\pgfpathlineto{\pgfqpoint{4.910481in}{1.835953in}}%
\pgfpathlineto{\pgfqpoint{4.912999in}{1.840320in}}%
\pgfpathlineto{\pgfqpoint{4.915518in}{1.798779in}}%
\pgfpathlineto{\pgfqpoint{4.918037in}{1.794591in}}%
\pgfpathlineto{\pgfqpoint{4.920556in}{1.759735in}}%
\pgfpathlineto{\pgfqpoint{4.923074in}{1.739050in}}%
\pgfpathlineto{\pgfqpoint{4.925593in}{1.727929in}}%
\pgfpathlineto{\pgfqpoint{4.928112in}{1.783156in}}%
\pgfpathlineto{\pgfqpoint{4.930631in}{1.750578in}}%
\pgfpathlineto{\pgfqpoint{4.933149in}{1.796161in}}%
\pgfpathlineto{\pgfqpoint{4.935668in}{1.834252in}}%
\pgfpathlineto{\pgfqpoint{4.938187in}{1.815795in}}%
\pgfpathlineto{\pgfqpoint{4.940706in}{1.806885in}}%
\pgfpathlineto{\pgfqpoint{4.943224in}{1.810170in}}%
\pgfpathlineto{\pgfqpoint{4.945743in}{1.818297in}}%
\pgfpathlineto{\pgfqpoint{4.948262in}{1.823688in}}%
\pgfpathlineto{\pgfqpoint{4.950781in}{1.789674in}}%
\pgfpathlineto{\pgfqpoint{4.953299in}{1.830231in}}%
\pgfpathlineto{\pgfqpoint{4.955818in}{1.818598in}}%
\pgfpathlineto{\pgfqpoint{4.958337in}{1.804654in}}%
\pgfpathlineto{\pgfqpoint{4.960856in}{1.778980in}}%
\pgfpathlineto{\pgfqpoint{4.963374in}{1.785630in}}%
\pgfpathlineto{\pgfqpoint{4.965893in}{1.804142in}}%
\pgfpathlineto{\pgfqpoint{4.968412in}{1.819354in}}%
\pgfpathlineto{\pgfqpoint{4.970931in}{1.801286in}}%
\pgfpathlineto{\pgfqpoint{4.973449in}{1.814281in}}%
\pgfpathlineto{\pgfqpoint{4.975968in}{1.791258in}}%
\pgfpathlineto{\pgfqpoint{4.978487in}{1.782984in}}%
\pgfpathlineto{\pgfqpoint{4.981006in}{1.777755in}}%
\pgfpathlineto{\pgfqpoint{4.983524in}{1.785899in}}%
\pgfpathlineto{\pgfqpoint{4.986043in}{1.803757in}}%
\pgfpathlineto{\pgfqpoint{4.988562in}{1.750149in}}%
\pgfpathlineto{\pgfqpoint{4.991081in}{1.744338in}}%
\pgfpathlineto{\pgfqpoint{4.993599in}{1.746795in}}%
\pgfpathlineto{\pgfqpoint{4.996118in}{1.764708in}}%
\pgfpathlineto{\pgfqpoint{4.998637in}{1.723062in}}%
\pgfpathlineto{\pgfqpoint{5.001156in}{1.712108in}}%
\pgfpathlineto{\pgfqpoint{5.003674in}{1.692410in}}%
\pgfpathlineto{\pgfqpoint{5.006193in}{1.656650in}}%
\pgfpathlineto{\pgfqpoint{5.008712in}{1.652549in}}%
\pgfpathlineto{\pgfqpoint{5.011231in}{1.635219in}}%
\pgfpathlineto{\pgfqpoint{5.013749in}{1.642413in}}%
\pgfpathlineto{\pgfqpoint{5.016268in}{1.637314in}}%
\pgfpathlineto{\pgfqpoint{5.018787in}{1.639822in}}%
\pgfpathlineto{\pgfqpoint{5.021306in}{1.623903in}}%
\pgfpathlineto{\pgfqpoint{5.023824in}{1.656226in}}%
\pgfpathlineto{\pgfqpoint{5.026343in}{1.636150in}}%
\pgfpathlineto{\pgfqpoint{5.028862in}{1.621628in}}%
\pgfpathlineto{\pgfqpoint{5.031381in}{1.666089in}}%
\pgfpathlineto{\pgfqpoint{5.033899in}{1.649018in}}%
\pgfpathlineto{\pgfqpoint{5.036418in}{1.613157in}}%
\pgfpathlineto{\pgfqpoint{5.038937in}{1.660776in}}%
\pgfpathlineto{\pgfqpoint{5.041456in}{1.705261in}}%
\pgfpathlineto{\pgfqpoint{5.043974in}{1.700409in}}%
\pgfpathlineto{\pgfqpoint{5.046493in}{1.686252in}}%
\pgfpathlineto{\pgfqpoint{5.049012in}{1.631489in}}%
\pgfpathlineto{\pgfqpoint{5.051531in}{1.634496in}}%
\pgfpathlineto{\pgfqpoint{5.054049in}{1.662018in}}%
\pgfpathlineto{\pgfqpoint{5.056568in}{1.667404in}}%
\pgfpathlineto{\pgfqpoint{5.059087in}{1.660266in}}%
\pgfpathlineto{\pgfqpoint{5.061606in}{1.692117in}}%
\pgfpathlineto{\pgfqpoint{5.064124in}{1.689273in}}%
\pgfpathlineto{\pgfqpoint{5.066643in}{1.696645in}}%
\pgfpathlineto{\pgfqpoint{5.069162in}{1.702083in}}%
\pgfpathlineto{\pgfqpoint{5.071681in}{1.682765in}}%
\pgfpathlineto{\pgfqpoint{5.074199in}{1.684322in}}%
\pgfpathlineto{\pgfqpoint{5.076718in}{1.698980in}}%
\pgfpathlineto{\pgfqpoint{5.079237in}{1.676730in}}%
\pgfpathlineto{\pgfqpoint{5.081756in}{1.638540in}}%
\pgfpathlineto{\pgfqpoint{5.084274in}{1.614785in}}%
\pgfpathlineto{\pgfqpoint{5.086793in}{1.612704in}}%
\pgfpathlineto{\pgfqpoint{5.089312in}{1.616081in}}%
\pgfpathlineto{\pgfqpoint{5.091831in}{1.575036in}}%
\pgfpathlineto{\pgfqpoint{5.094349in}{1.592569in}}%
\pgfpathlineto{\pgfqpoint{5.096868in}{1.570295in}}%
\pgfpathlineto{\pgfqpoint{5.099387in}{1.618721in}}%
\pgfpathlineto{\pgfqpoint{5.101906in}{1.584386in}}%
\pgfpathlineto{\pgfqpoint{5.104424in}{1.567289in}}%
\pgfpathlineto{\pgfqpoint{5.106943in}{1.535028in}}%
\pgfpathlineto{\pgfqpoint{5.109462in}{1.512078in}}%
\pgfpathlineto{\pgfqpoint{5.111981in}{1.553043in}}%
\pgfpathlineto{\pgfqpoint{5.114499in}{1.542309in}}%
\pgfpathlineto{\pgfqpoint{5.117018in}{1.506600in}}%
\pgfpathlineto{\pgfqpoint{5.122056in}{1.589532in}}%
\pgfpathlineto{\pgfqpoint{5.124574in}{1.590210in}}%
\pgfpathlineto{\pgfqpoint{5.127093in}{1.592361in}}%
\pgfpathlineto{\pgfqpoint{5.129612in}{1.552140in}}%
\pgfpathlineto{\pgfqpoint{5.132131in}{1.526148in}}%
\pgfpathlineto{\pgfqpoint{5.134649in}{1.532968in}}%
\pgfpathlineto{\pgfqpoint{5.137168in}{1.496867in}}%
\pgfpathlineto{\pgfqpoint{5.139687in}{1.511661in}}%
\pgfpathlineto{\pgfqpoint{5.142206in}{1.490459in}}%
\pgfpathlineto{\pgfqpoint{5.144724in}{1.525780in}}%
\pgfpathlineto{\pgfqpoint{5.147243in}{1.494695in}}%
\pgfpathlineto{\pgfqpoint{5.152281in}{1.517660in}}%
\pgfpathlineto{\pgfqpoint{5.154799in}{1.498037in}}%
\pgfpathlineto{\pgfqpoint{5.157318in}{1.488604in}}%
\pgfpathlineto{\pgfqpoint{5.159837in}{1.500110in}}%
\pgfpathlineto{\pgfqpoint{5.162356in}{1.510849in}}%
\pgfpathlineto{\pgfqpoint{5.164874in}{1.497005in}}%
\pgfpathlineto{\pgfqpoint{5.167393in}{1.493084in}}%
\pgfpathlineto{\pgfqpoint{5.169912in}{1.544532in}}%
\pgfpathlineto{\pgfqpoint{5.172431in}{1.564597in}}%
\pgfpathlineto{\pgfqpoint{5.174949in}{1.602961in}}%
\pgfpathlineto{\pgfqpoint{5.177468in}{1.580746in}}%
\pgfpathlineto{\pgfqpoint{5.179987in}{1.609124in}}%
\pgfpathlineto{\pgfqpoint{5.182506in}{1.607014in}}%
\pgfpathlineto{\pgfqpoint{5.185024in}{1.610665in}}%
\pgfpathlineto{\pgfqpoint{5.187543in}{1.638086in}}%
\pgfpathlineto{\pgfqpoint{5.190062in}{1.604691in}}%
\pgfpathlineto{\pgfqpoint{5.192581in}{1.587107in}}%
\pgfpathlineto{\pgfqpoint{5.195099in}{1.579983in}}%
\pgfpathlineto{\pgfqpoint{5.197618in}{1.561159in}}%
\pgfpathlineto{\pgfqpoint{5.200137in}{1.608409in}}%
\pgfpathlineto{\pgfqpoint{5.202656in}{1.601990in}}%
\pgfpathlineto{\pgfqpoint{5.205174in}{1.567800in}}%
\pgfpathlineto{\pgfqpoint{5.207693in}{1.562782in}}%
\pgfpathlineto{\pgfqpoint{5.210212in}{1.610373in}}%
\pgfpathlineto{\pgfqpoint{5.212731in}{1.620192in}}%
\pgfpathlineto{\pgfqpoint{5.215249in}{1.634706in}}%
\pgfpathlineto{\pgfqpoint{5.217768in}{1.668359in}}%
\pgfpathlineto{\pgfqpoint{5.220287in}{1.690874in}}%
\pgfpathlineto{\pgfqpoint{5.222806in}{1.699128in}}%
\pgfpathlineto{\pgfqpoint{5.225324in}{1.699483in}}%
\pgfpathlineto{\pgfqpoint{5.227843in}{1.693023in}}%
\pgfpathlineto{\pgfqpoint{5.230362in}{1.718849in}}%
\pgfpathlineto{\pgfqpoint{5.232881in}{1.705174in}}%
\pgfpathlineto{\pgfqpoint{5.235399in}{1.703904in}}%
\pgfpathlineto{\pgfqpoint{5.237918in}{1.703954in}}%
\pgfpathlineto{\pgfqpoint{5.240437in}{1.695286in}}%
\pgfpathlineto{\pgfqpoint{5.242956in}{1.714893in}}%
\pgfpathlineto{\pgfqpoint{5.245474in}{1.693793in}}%
\pgfpathlineto{\pgfqpoint{5.247993in}{1.679085in}}%
\pgfpathlineto{\pgfqpoint{5.250512in}{1.710578in}}%
\pgfpathlineto{\pgfqpoint{5.253031in}{1.663986in}}%
\pgfpathlineto{\pgfqpoint{5.255549in}{1.679277in}}%
\pgfpathlineto{\pgfqpoint{5.258068in}{1.704368in}}%
\pgfpathlineto{\pgfqpoint{5.260587in}{1.649679in}}%
\pgfpathlineto{\pgfqpoint{5.263106in}{1.602899in}}%
\pgfpathlineto{\pgfqpoint{5.265624in}{1.623562in}}%
\pgfpathlineto{\pgfqpoint{5.268143in}{1.604588in}}%
\pgfpathlineto{\pgfqpoint{5.270662in}{1.609512in}}%
\pgfpathlineto{\pgfqpoint{5.275699in}{1.628423in}}%
\pgfpathlineto{\pgfqpoint{5.278218in}{1.628054in}}%
\pgfpathlineto{\pgfqpoint{5.280737in}{1.676317in}}%
\pgfpathlineto{\pgfqpoint{5.283256in}{1.693372in}}%
\pgfpathlineto{\pgfqpoint{5.285774in}{1.696142in}}%
\pgfpathlineto{\pgfqpoint{5.288293in}{1.693663in}}%
\pgfpathlineto{\pgfqpoint{5.290812in}{1.661705in}}%
\pgfpathlineto{\pgfqpoint{5.293331in}{1.687923in}}%
\pgfpathlineto{\pgfqpoint{5.295849in}{1.687514in}}%
\pgfpathlineto{\pgfqpoint{5.298368in}{1.678430in}}%
\pgfpathlineto{\pgfqpoint{5.300887in}{1.667772in}}%
\pgfpathlineto{\pgfqpoint{5.303406in}{1.680868in}}%
\pgfpathlineto{\pgfqpoint{5.305924in}{1.731014in}}%
\pgfpathlineto{\pgfqpoint{5.308443in}{1.755299in}}%
\pgfpathlineto{\pgfqpoint{5.310962in}{1.800539in}}%
\pgfpathlineto{\pgfqpoint{5.313481in}{1.762871in}}%
\pgfpathlineto{\pgfqpoint{5.315999in}{1.733608in}}%
\pgfpathlineto{\pgfqpoint{5.318518in}{1.724928in}}%
\pgfpathlineto{\pgfqpoint{5.321037in}{1.727805in}}%
\pgfpathlineto{\pgfqpoint{5.323556in}{1.761336in}}%
\pgfpathlineto{\pgfqpoint{5.328593in}{1.778816in}}%
\pgfpathlineto{\pgfqpoint{5.331112in}{1.718142in}}%
\pgfpathlineto{\pgfqpoint{5.333631in}{1.697564in}}%
\pgfpathlineto{\pgfqpoint{5.336149in}{1.694674in}}%
\pgfpathlineto{\pgfqpoint{5.338668in}{1.694079in}}%
\pgfpathlineto{\pgfqpoint{5.341187in}{1.727672in}}%
\pgfpathlineto{\pgfqpoint{5.343706in}{1.748188in}}%
\pgfpathlineto{\pgfqpoint{5.346224in}{1.735433in}}%
\pgfpathlineto{\pgfqpoint{5.348743in}{1.730815in}}%
\pgfpathlineto{\pgfqpoint{5.351262in}{1.743139in}}%
\pgfpathlineto{\pgfqpoint{5.353781in}{1.732135in}}%
\pgfpathlineto{\pgfqpoint{5.356299in}{1.738591in}}%
\pgfpathlineto{\pgfqpoint{5.358818in}{1.787988in}}%
\pgfpathlineto{\pgfqpoint{5.361337in}{1.784046in}}%
\pgfpathlineto{\pgfqpoint{5.363856in}{1.781094in}}%
\pgfpathlineto{\pgfqpoint{5.366374in}{1.796383in}}%
\pgfpathlineto{\pgfqpoint{5.368893in}{1.805157in}}%
\pgfpathlineto{\pgfqpoint{5.371412in}{1.838471in}}%
\pgfpathlineto{\pgfqpoint{5.373931in}{1.786741in}}%
\pgfpathlineto{\pgfqpoint{5.376449in}{1.745239in}}%
\pgfpathlineto{\pgfqpoint{5.378968in}{1.760776in}}%
\pgfpathlineto{\pgfqpoint{5.381487in}{1.778627in}}%
\pgfpathlineto{\pgfqpoint{5.384006in}{1.767158in}}%
\pgfpathlineto{\pgfqpoint{5.386524in}{1.787174in}}%
\pgfpathlineto{\pgfqpoint{5.389043in}{1.769740in}}%
\pgfpathlineto{\pgfqpoint{5.391562in}{1.777180in}}%
\pgfpathlineto{\pgfqpoint{5.394081in}{1.822255in}}%
\pgfpathlineto{\pgfqpoint{5.396599in}{1.825411in}}%
\pgfpathlineto{\pgfqpoint{5.399118in}{1.781960in}}%
\pgfpathlineto{\pgfqpoint{5.401637in}{1.784877in}}%
\pgfpathlineto{\pgfqpoint{5.404156in}{1.805288in}}%
\pgfpathlineto{\pgfqpoint{5.406674in}{1.793923in}}%
\pgfpathlineto{\pgfqpoint{5.409193in}{1.809210in}}%
\pgfpathlineto{\pgfqpoint{5.411712in}{1.821546in}}%
\pgfpathlineto{\pgfqpoint{5.414231in}{1.804127in}}%
\pgfpathlineto{\pgfqpoint{5.416749in}{1.781515in}}%
\pgfpathlineto{\pgfqpoint{5.419268in}{1.761318in}}%
\pgfpathlineto{\pgfqpoint{5.421787in}{1.722637in}}%
\pgfpathlineto{\pgfqpoint{5.424306in}{1.757573in}}%
\pgfpathlineto{\pgfqpoint{5.426824in}{1.742567in}}%
\pgfpathlineto{\pgfqpoint{5.429343in}{1.722527in}}%
\pgfpathlineto{\pgfqpoint{5.431862in}{1.691406in}}%
\pgfpathlineto{\pgfqpoint{5.434381in}{1.701437in}}%
\pgfpathlineto{\pgfqpoint{5.436899in}{1.634312in}}%
\pgfpathlineto{\pgfqpoint{5.439418in}{1.650197in}}%
\pgfpathlineto{\pgfqpoint{5.441937in}{1.734069in}}%
\pgfpathlineto{\pgfqpoint{5.444456in}{1.751575in}}%
\pgfpathlineto{\pgfqpoint{5.446974in}{1.702943in}}%
\pgfpathlineto{\pgfqpoint{5.449493in}{1.710344in}}%
\pgfpathlineto{\pgfqpoint{5.452012in}{1.651566in}}%
\pgfpathlineto{\pgfqpoint{5.457049in}{1.721455in}}%
\pgfpathlineto{\pgfqpoint{5.459568in}{1.717553in}}%
\pgfpathlineto{\pgfqpoint{5.462087in}{1.716742in}}%
\pgfpathlineto{\pgfqpoint{5.464606in}{1.727520in}}%
\pgfpathlineto{\pgfqpoint{5.467124in}{1.736599in}}%
\pgfpathlineto{\pgfqpoint{5.469643in}{1.706783in}}%
\pgfpathlineto{\pgfqpoint{5.472162in}{1.710950in}}%
\pgfpathlineto{\pgfqpoint{5.474681in}{1.730308in}}%
\pgfpathlineto{\pgfqpoint{5.477199in}{1.726152in}}%
\pgfpathlineto{\pgfqpoint{5.479718in}{1.737909in}}%
\pgfpathlineto{\pgfqpoint{5.482237in}{1.751950in}}%
\pgfpathlineto{\pgfqpoint{5.484756in}{1.704257in}}%
\pgfpathlineto{\pgfqpoint{5.487274in}{1.694835in}}%
\pgfpathlineto{\pgfqpoint{5.489793in}{1.660154in}}%
\pgfpathlineto{\pgfqpoint{5.492312in}{1.695754in}}%
\pgfpathlineto{\pgfqpoint{5.494831in}{1.663035in}}%
\pgfpathlineto{\pgfqpoint{5.497349in}{1.622038in}}%
\pgfpathlineto{\pgfqpoint{5.499868in}{1.615469in}}%
\pgfpathlineto{\pgfqpoint{5.502387in}{1.594249in}}%
\pgfpathlineto{\pgfqpoint{5.504906in}{1.606973in}}%
\pgfpathlineto{\pgfqpoint{5.507424in}{1.632628in}}%
\pgfpathlineto{\pgfqpoint{5.509943in}{1.619517in}}%
\pgfpathlineto{\pgfqpoint{5.512462in}{1.584953in}}%
\pgfpathlineto{\pgfqpoint{5.514981in}{1.590107in}}%
\pgfpathlineto{\pgfqpoint{5.517499in}{1.571057in}}%
\pgfpathlineto{\pgfqpoint{5.520018in}{1.519332in}}%
\pgfpathlineto{\pgfqpoint{5.522537in}{1.546374in}}%
\pgfpathlineto{\pgfqpoint{5.525056in}{1.581646in}}%
\pgfpathlineto{\pgfqpoint{5.527574in}{1.591636in}}%
\pgfpathlineto{\pgfqpoint{5.530093in}{1.553414in}}%
\pgfpathlineto{\pgfqpoint{5.532612in}{1.580886in}}%
\pgfpathlineto{\pgfqpoint{5.535131in}{1.565292in}}%
\pgfpathlineto{\pgfqpoint{5.537649in}{1.599305in}}%
\pgfpathlineto{\pgfqpoint{5.540168in}{1.553763in}}%
\pgfpathlineto{\pgfqpoint{5.542687in}{1.554401in}}%
\pgfpathlineto{\pgfqpoint{5.547724in}{1.540290in}}%
\pgfpathlineto{\pgfqpoint{5.550243in}{1.549160in}}%
\pgfpathlineto{\pgfqpoint{5.552762in}{1.563604in}}%
\pgfpathlineto{\pgfqpoint{5.555281in}{1.569619in}}%
\pgfpathlineto{\pgfqpoint{5.557799in}{1.568371in}}%
\pgfpathlineto{\pgfqpoint{5.560318in}{1.597620in}}%
\pgfpathlineto{\pgfqpoint{5.562837in}{1.591298in}}%
\pgfpathlineto{\pgfqpoint{5.565356in}{1.632656in}}%
\pgfpathlineto{\pgfqpoint{5.570393in}{1.583072in}}%
\pgfpathlineto{\pgfqpoint{5.572912in}{1.520911in}}%
\pgfpathlineto{\pgfqpoint{5.575431in}{1.533688in}}%
\pgfpathlineto{\pgfqpoint{5.577949in}{1.523396in}}%
\pgfpathlineto{\pgfqpoint{5.580468in}{1.552858in}}%
\pgfpathlineto{\pgfqpoint{5.582987in}{1.577158in}}%
\pgfpathlineto{\pgfqpoint{5.585506in}{1.607033in}}%
\pgfpathlineto{\pgfqpoint{5.588024in}{1.642615in}}%
\pgfpathlineto{\pgfqpoint{5.590543in}{1.628097in}}%
\pgfpathlineto{\pgfqpoint{5.593062in}{1.603735in}}%
\pgfpathlineto{\pgfqpoint{5.595581in}{1.585335in}}%
\pgfpathlineto{\pgfqpoint{5.598099in}{1.565367in}}%
\pgfpathlineto{\pgfqpoint{5.600618in}{1.541394in}}%
\pgfpathlineto{\pgfqpoint{5.600618in}{1.541394in}}%
\pgfusepath{stroke}%
\end{pgfscope}%
\begin{pgfscope}%
\pgfpathrectangle{\pgfqpoint{0.563118in}{0.401080in}}{\pgfqpoint{5.037500in}{4.004000in}}%
\pgfusepath{clip}%
\pgfsetbuttcap%
\pgfsetroundjoin%
\pgfsetlinewidth{0.803000pt}%
\definecolor{currentstroke}{rgb}{0.690196,0.188235,0.376471}%
\pgfsetstrokecolor{currentstroke}%
\pgfsetdash{{2.960000pt}{1.280000pt}}{0.000000pt}%
\pgfpathmoveto{\pgfqpoint{0.563118in}{1.988980in}}%
\pgfpathlineto{\pgfqpoint{0.565637in}{1.983653in}}%
\pgfpathlineto{\pgfqpoint{0.568156in}{2.010483in}}%
\pgfpathlineto{\pgfqpoint{0.570674in}{2.025389in}}%
\pgfpathlineto{\pgfqpoint{0.573193in}{2.014425in}}%
\pgfpathlineto{\pgfqpoint{0.575712in}{2.010472in}}%
\pgfpathlineto{\pgfqpoint{0.578231in}{2.022732in}}%
\pgfpathlineto{\pgfqpoint{0.580749in}{2.026260in}}%
\pgfpathlineto{\pgfqpoint{0.583268in}{2.070134in}}%
\pgfpathlineto{\pgfqpoint{0.585787in}{2.076010in}}%
\pgfpathlineto{\pgfqpoint{0.588306in}{2.106502in}}%
\pgfpathlineto{\pgfqpoint{0.590824in}{2.115571in}}%
\pgfpathlineto{\pgfqpoint{0.593343in}{2.092978in}}%
\pgfpathlineto{\pgfqpoint{0.595862in}{2.052048in}}%
\pgfpathlineto{\pgfqpoint{0.598381in}{2.097724in}}%
\pgfpathlineto{\pgfqpoint{0.600899in}{2.123760in}}%
\pgfpathlineto{\pgfqpoint{0.603418in}{2.125374in}}%
\pgfpathlineto{\pgfqpoint{0.605937in}{2.125289in}}%
\pgfpathlineto{\pgfqpoint{0.608456in}{2.159633in}}%
\pgfpathlineto{\pgfqpoint{0.610974in}{2.197571in}}%
\pgfpathlineto{\pgfqpoint{0.613493in}{2.166988in}}%
\pgfpathlineto{\pgfqpoint{0.616012in}{2.122725in}}%
\pgfpathlineto{\pgfqpoint{0.618531in}{2.141602in}}%
\pgfpathlineto{\pgfqpoint{0.621049in}{2.112312in}}%
\pgfpathlineto{\pgfqpoint{0.623568in}{2.087163in}}%
\pgfpathlineto{\pgfqpoint{0.626087in}{2.060707in}}%
\pgfpathlineto{\pgfqpoint{0.628606in}{2.060082in}}%
\pgfpathlineto{\pgfqpoint{0.631124in}{2.065265in}}%
\pgfpathlineto{\pgfqpoint{0.633643in}{2.083029in}}%
\pgfpathlineto{\pgfqpoint{0.636162in}{2.092722in}}%
\pgfpathlineto{\pgfqpoint{0.638681in}{2.154576in}}%
\pgfpathlineto{\pgfqpoint{0.641199in}{2.165123in}}%
\pgfpathlineto{\pgfqpoint{0.643718in}{2.187273in}}%
\pgfpathlineto{\pgfqpoint{0.646237in}{2.191746in}}%
\pgfpathlineto{\pgfqpoint{0.648756in}{2.162284in}}%
\pgfpathlineto{\pgfqpoint{0.651274in}{2.199912in}}%
\pgfpathlineto{\pgfqpoint{0.653793in}{2.213838in}}%
\pgfpathlineto{\pgfqpoint{0.656312in}{2.220086in}}%
\pgfpathlineto{\pgfqpoint{0.658831in}{2.147448in}}%
\pgfpathlineto{\pgfqpoint{0.661349in}{2.134851in}}%
\pgfpathlineto{\pgfqpoint{0.663868in}{2.130112in}}%
\pgfpathlineto{\pgfqpoint{0.666387in}{2.143754in}}%
\pgfpathlineto{\pgfqpoint{0.668906in}{2.154422in}}%
\pgfpathlineto{\pgfqpoint{0.671424in}{2.136336in}}%
\pgfpathlineto{\pgfqpoint{0.673943in}{2.075608in}}%
\pgfpathlineto{\pgfqpoint{0.676462in}{2.056422in}}%
\pgfpathlineto{\pgfqpoint{0.678981in}{2.057195in}}%
\pgfpathlineto{\pgfqpoint{0.681499in}{2.026949in}}%
\pgfpathlineto{\pgfqpoint{0.684018in}{2.005256in}}%
\pgfpathlineto{\pgfqpoint{0.686537in}{1.964089in}}%
\pgfpathlineto{\pgfqpoint{0.689056in}{2.003706in}}%
\pgfpathlineto{\pgfqpoint{0.691574in}{1.994688in}}%
\pgfpathlineto{\pgfqpoint{0.694093in}{1.994737in}}%
\pgfpathlineto{\pgfqpoint{0.696612in}{2.021430in}}%
\pgfpathlineto{\pgfqpoint{0.699131in}{1.990996in}}%
\pgfpathlineto{\pgfqpoint{0.701649in}{1.983438in}}%
\pgfpathlineto{\pgfqpoint{0.704168in}{1.953461in}}%
\pgfpathlineto{\pgfqpoint{0.706687in}{1.952170in}}%
\pgfpathlineto{\pgfqpoint{0.709206in}{1.941883in}}%
\pgfpathlineto{\pgfqpoint{0.711724in}{1.949214in}}%
\pgfpathlineto{\pgfqpoint{0.714243in}{1.933254in}}%
\pgfpathlineto{\pgfqpoint{0.716762in}{1.954150in}}%
\pgfpathlineto{\pgfqpoint{0.719281in}{1.949732in}}%
\pgfpathlineto{\pgfqpoint{0.721799in}{1.967051in}}%
\pgfpathlineto{\pgfqpoint{0.724318in}{1.901958in}}%
\pgfpathlineto{\pgfqpoint{0.726837in}{1.858684in}}%
\pgfpathlineto{\pgfqpoint{0.729356in}{1.855481in}}%
\pgfpathlineto{\pgfqpoint{0.731874in}{1.866894in}}%
\pgfpathlineto{\pgfqpoint{0.734393in}{1.856647in}}%
\pgfpathlineto{\pgfqpoint{0.736912in}{1.842420in}}%
\pgfpathlineto{\pgfqpoint{0.739431in}{1.832267in}}%
\pgfpathlineto{\pgfqpoint{0.741949in}{1.842561in}}%
\pgfpathlineto{\pgfqpoint{0.744468in}{1.830140in}}%
\pgfpathlineto{\pgfqpoint{0.746987in}{1.809100in}}%
\pgfpathlineto{\pgfqpoint{0.749506in}{1.783809in}}%
\pgfpathlineto{\pgfqpoint{0.752024in}{1.784200in}}%
\pgfpathlineto{\pgfqpoint{0.754543in}{1.736483in}}%
\pgfpathlineto{\pgfqpoint{0.757062in}{1.720613in}}%
\pgfpathlineto{\pgfqpoint{0.759581in}{1.698012in}}%
\pgfpathlineto{\pgfqpoint{0.762099in}{1.696460in}}%
\pgfpathlineto{\pgfqpoint{0.764618in}{1.669080in}}%
\pgfpathlineto{\pgfqpoint{0.767137in}{1.630685in}}%
\pgfpathlineto{\pgfqpoint{0.769656in}{1.643384in}}%
\pgfpathlineto{\pgfqpoint{0.772174in}{1.666742in}}%
\pgfpathlineto{\pgfqpoint{0.774693in}{1.689194in}}%
\pgfpathlineto{\pgfqpoint{0.777212in}{1.630844in}}%
\pgfpathlineto{\pgfqpoint{0.779731in}{1.599293in}}%
\pgfpathlineto{\pgfqpoint{0.782249in}{1.644580in}}%
\pgfpathlineto{\pgfqpoint{0.784768in}{1.632426in}}%
\pgfpathlineto{\pgfqpoint{0.787287in}{1.567174in}}%
\pgfpathlineto{\pgfqpoint{0.789806in}{1.565697in}}%
\pgfpathlineto{\pgfqpoint{0.792324in}{1.577677in}}%
\pgfpathlineto{\pgfqpoint{0.794843in}{1.579403in}}%
\pgfpathlineto{\pgfqpoint{0.797362in}{1.553097in}}%
\pgfpathlineto{\pgfqpoint{0.799881in}{1.572622in}}%
\pgfpathlineto{\pgfqpoint{0.802399in}{1.548232in}}%
\pgfpathlineto{\pgfqpoint{0.804918in}{1.578607in}}%
\pgfpathlineto{\pgfqpoint{0.807437in}{1.597159in}}%
\pgfpathlineto{\pgfqpoint{0.809956in}{1.591205in}}%
\pgfpathlineto{\pgfqpoint{0.812474in}{1.568849in}}%
\pgfpathlineto{\pgfqpoint{0.814993in}{1.626202in}}%
\pgfpathlineto{\pgfqpoint{0.817512in}{1.622199in}}%
\pgfpathlineto{\pgfqpoint{0.820031in}{1.623902in}}%
\pgfpathlineto{\pgfqpoint{0.822549in}{1.624113in}}%
\pgfpathlineto{\pgfqpoint{0.825068in}{1.631800in}}%
\pgfpathlineto{\pgfqpoint{0.827587in}{1.614845in}}%
\pgfpathlineto{\pgfqpoint{0.830106in}{1.610316in}}%
\pgfpathlineto{\pgfqpoint{0.832624in}{1.620810in}}%
\pgfpathlineto{\pgfqpoint{0.835143in}{1.601564in}}%
\pgfpathlineto{\pgfqpoint{0.837662in}{1.691248in}}%
\pgfpathlineto{\pgfqpoint{0.840181in}{1.686067in}}%
\pgfpathlineto{\pgfqpoint{0.842699in}{1.742192in}}%
\pgfpathlineto{\pgfqpoint{0.845218in}{1.728353in}}%
\pgfpathlineto{\pgfqpoint{0.847737in}{1.759500in}}%
\pgfpathlineto{\pgfqpoint{0.850256in}{1.770499in}}%
\pgfpathlineto{\pgfqpoint{0.852774in}{1.795895in}}%
\pgfpathlineto{\pgfqpoint{0.855293in}{1.826003in}}%
\pgfpathlineto{\pgfqpoint{0.857812in}{1.799415in}}%
\pgfpathlineto{\pgfqpoint{0.860331in}{1.805728in}}%
\pgfpathlineto{\pgfqpoint{0.862849in}{1.793203in}}%
\pgfpathlineto{\pgfqpoint{0.865368in}{1.742866in}}%
\pgfpathlineto{\pgfqpoint{0.867887in}{1.785661in}}%
\pgfpathlineto{\pgfqpoint{0.870406in}{1.794771in}}%
\pgfpathlineto{\pgfqpoint{0.872924in}{1.831074in}}%
\pgfpathlineto{\pgfqpoint{0.875443in}{1.853840in}}%
\pgfpathlineto{\pgfqpoint{0.877962in}{1.846805in}}%
\pgfpathlineto{\pgfqpoint{0.880481in}{1.832960in}}%
\pgfpathlineto{\pgfqpoint{0.882999in}{1.828777in}}%
\pgfpathlineto{\pgfqpoint{0.885518in}{1.852353in}}%
\pgfpathlineto{\pgfqpoint{0.888037in}{1.877155in}}%
\pgfpathlineto{\pgfqpoint{0.890556in}{1.849752in}}%
\pgfpathlineto{\pgfqpoint{0.893074in}{1.817802in}}%
\pgfpathlineto{\pgfqpoint{0.895593in}{1.840593in}}%
\pgfpathlineto{\pgfqpoint{0.898112in}{1.878457in}}%
\pgfpathlineto{\pgfqpoint{0.900631in}{1.864003in}}%
\pgfpathlineto{\pgfqpoint{0.903149in}{1.873483in}}%
\pgfpathlineto{\pgfqpoint{0.905668in}{1.880575in}}%
\pgfpathlineto{\pgfqpoint{0.908187in}{1.845230in}}%
\pgfpathlineto{\pgfqpoint{0.910706in}{1.835302in}}%
\pgfpathlineto{\pgfqpoint{0.913224in}{1.816932in}}%
\pgfpathlineto{\pgfqpoint{0.915743in}{1.811104in}}%
\pgfpathlineto{\pgfqpoint{0.918262in}{1.854951in}}%
\pgfpathlineto{\pgfqpoint{0.920781in}{1.878176in}}%
\pgfpathlineto{\pgfqpoint{0.923299in}{1.857102in}}%
\pgfpathlineto{\pgfqpoint{0.925818in}{1.888147in}}%
\pgfpathlineto{\pgfqpoint{0.928337in}{1.858564in}}%
\pgfpathlineto{\pgfqpoint{0.930856in}{1.845272in}}%
\pgfpathlineto{\pgfqpoint{0.933374in}{1.844816in}}%
\pgfpathlineto{\pgfqpoint{0.935893in}{1.875634in}}%
\pgfpathlineto{\pgfqpoint{0.938412in}{1.903930in}}%
\pgfpathlineto{\pgfqpoint{0.940931in}{1.943022in}}%
\pgfpathlineto{\pgfqpoint{0.943449in}{1.948482in}}%
\pgfpathlineto{\pgfqpoint{0.945968in}{1.965028in}}%
\pgfpathlineto{\pgfqpoint{0.948487in}{1.993836in}}%
\pgfpathlineto{\pgfqpoint{0.951006in}{1.960358in}}%
\pgfpathlineto{\pgfqpoint{0.953524in}{1.972588in}}%
\pgfpathlineto{\pgfqpoint{0.956043in}{1.986800in}}%
\pgfpathlineto{\pgfqpoint{0.958562in}{1.971867in}}%
\pgfpathlineto{\pgfqpoint{0.961081in}{1.962545in}}%
\pgfpathlineto{\pgfqpoint{0.963599in}{1.972130in}}%
\pgfpathlineto{\pgfqpoint{0.966118in}{1.942995in}}%
\pgfpathlineto{\pgfqpoint{0.968637in}{1.966883in}}%
\pgfpathlineto{\pgfqpoint{0.971156in}{1.964356in}}%
\pgfpathlineto{\pgfqpoint{0.973674in}{1.961179in}}%
\pgfpathlineto{\pgfqpoint{0.976193in}{1.967099in}}%
\pgfpathlineto{\pgfqpoint{0.978712in}{1.980823in}}%
\pgfpathlineto{\pgfqpoint{0.981231in}{1.918046in}}%
\pgfpathlineto{\pgfqpoint{0.983749in}{1.937860in}}%
\pgfpathlineto{\pgfqpoint{0.986268in}{1.954380in}}%
\pgfpathlineto{\pgfqpoint{0.988787in}{1.956604in}}%
\pgfpathlineto{\pgfqpoint{0.991306in}{2.027920in}}%
\pgfpathlineto{\pgfqpoint{0.993824in}{2.031687in}}%
\pgfpathlineto{\pgfqpoint{0.996343in}{2.032239in}}%
\pgfpathlineto{\pgfqpoint{0.998862in}{1.988436in}}%
\pgfpathlineto{\pgfqpoint{1.001381in}{1.950749in}}%
\pgfpathlineto{\pgfqpoint{1.003899in}{1.962881in}}%
\pgfpathlineto{\pgfqpoint{1.006418in}{1.949090in}}%
\pgfpathlineto{\pgfqpoint{1.008937in}{1.997322in}}%
\pgfpathlineto{\pgfqpoint{1.011456in}{2.051783in}}%
\pgfpathlineto{\pgfqpoint{1.013974in}{2.007593in}}%
\pgfpathlineto{\pgfqpoint{1.016493in}{2.008924in}}%
\pgfpathlineto{\pgfqpoint{1.019012in}{2.033975in}}%
\pgfpathlineto{\pgfqpoint{1.021531in}{2.069655in}}%
\pgfpathlineto{\pgfqpoint{1.024049in}{2.075038in}}%
\pgfpathlineto{\pgfqpoint{1.026568in}{2.079324in}}%
\pgfpathlineto{\pgfqpoint{1.029087in}{2.119338in}}%
\pgfpathlineto{\pgfqpoint{1.031606in}{2.162249in}}%
\pgfpathlineto{\pgfqpoint{1.034124in}{2.158227in}}%
\pgfpathlineto{\pgfqpoint{1.036643in}{2.143655in}}%
\pgfpathlineto{\pgfqpoint{1.039162in}{2.195649in}}%
\pgfpathlineto{\pgfqpoint{1.041681in}{2.172350in}}%
\pgfpathlineto{\pgfqpoint{1.044199in}{2.133192in}}%
\pgfpathlineto{\pgfqpoint{1.046718in}{2.140405in}}%
\pgfpathlineto{\pgfqpoint{1.049237in}{2.127330in}}%
\pgfpathlineto{\pgfqpoint{1.051756in}{2.168490in}}%
\pgfpathlineto{\pgfqpoint{1.054274in}{2.149767in}}%
\pgfpathlineto{\pgfqpoint{1.056793in}{2.157033in}}%
\pgfpathlineto{\pgfqpoint{1.059312in}{2.133731in}}%
\pgfpathlineto{\pgfqpoint{1.061831in}{2.132029in}}%
\pgfpathlineto{\pgfqpoint{1.066868in}{2.102305in}}%
\pgfpathlineto{\pgfqpoint{1.069387in}{2.126906in}}%
\pgfpathlineto{\pgfqpoint{1.071906in}{2.152462in}}%
\pgfpathlineto{\pgfqpoint{1.074424in}{2.162276in}}%
\pgfpathlineto{\pgfqpoint{1.076943in}{2.148478in}}%
\pgfpathlineto{\pgfqpoint{1.079462in}{2.165521in}}%
\pgfpathlineto{\pgfqpoint{1.081981in}{2.169764in}}%
\pgfpathlineto{\pgfqpoint{1.084499in}{2.132247in}}%
\pgfpathlineto{\pgfqpoint{1.087018in}{2.127513in}}%
\pgfpathlineto{\pgfqpoint{1.089537in}{2.150452in}}%
\pgfpathlineto{\pgfqpoint{1.092056in}{2.165042in}}%
\pgfpathlineto{\pgfqpoint{1.094574in}{2.192897in}}%
\pgfpathlineto{\pgfqpoint{1.097093in}{2.210453in}}%
\pgfpathlineto{\pgfqpoint{1.099612in}{2.210131in}}%
\pgfpathlineto{\pgfqpoint{1.102131in}{2.174634in}}%
\pgfpathlineto{\pgfqpoint{1.104649in}{2.235978in}}%
\pgfpathlineto{\pgfqpoint{1.107168in}{2.212202in}}%
\pgfpathlineto{\pgfqpoint{1.109687in}{2.169533in}}%
\pgfpathlineto{\pgfqpoint{1.112206in}{2.150782in}}%
\pgfpathlineto{\pgfqpoint{1.114724in}{2.232199in}}%
\pgfpathlineto{\pgfqpoint{1.117243in}{2.273230in}}%
\pgfpathlineto{\pgfqpoint{1.119762in}{2.203140in}}%
\pgfpathlineto{\pgfqpoint{1.122281in}{2.159871in}}%
\pgfpathlineto{\pgfqpoint{1.124799in}{2.217534in}}%
\pgfpathlineto{\pgfqpoint{1.127318in}{2.192574in}}%
\pgfpathlineto{\pgfqpoint{1.129837in}{2.186255in}}%
\pgfpathlineto{\pgfqpoint{1.132356in}{2.253941in}}%
\pgfpathlineto{\pgfqpoint{1.134874in}{2.273223in}}%
\pgfpathlineto{\pgfqpoint{1.137393in}{2.224306in}}%
\pgfpathlineto{\pgfqpoint{1.139912in}{2.235275in}}%
\pgfpathlineto{\pgfqpoint{1.142431in}{2.270169in}}%
\pgfpathlineto{\pgfqpoint{1.144949in}{2.292781in}}%
\pgfpathlineto{\pgfqpoint{1.147468in}{2.277439in}}%
\pgfpathlineto{\pgfqpoint{1.149987in}{2.294791in}}%
\pgfpathlineto{\pgfqpoint{1.152506in}{2.285726in}}%
\pgfpathlineto{\pgfqpoint{1.155024in}{2.348835in}}%
\pgfpathlineto{\pgfqpoint{1.157543in}{2.287584in}}%
\pgfpathlineto{\pgfqpoint{1.160062in}{2.258065in}}%
\pgfpathlineto{\pgfqpoint{1.162581in}{2.216714in}}%
\pgfpathlineto{\pgfqpoint{1.165099in}{2.216281in}}%
\pgfpathlineto{\pgfqpoint{1.167618in}{2.228536in}}%
\pgfpathlineto{\pgfqpoint{1.170137in}{2.233233in}}%
\pgfpathlineto{\pgfqpoint{1.172656in}{2.194639in}}%
\pgfpathlineto{\pgfqpoint{1.175174in}{2.182364in}}%
\pgfpathlineto{\pgfqpoint{1.177693in}{2.212115in}}%
\pgfpathlineto{\pgfqpoint{1.180212in}{2.196708in}}%
\pgfpathlineto{\pgfqpoint{1.182731in}{2.153614in}}%
\pgfpathlineto{\pgfqpoint{1.185249in}{2.158966in}}%
\pgfpathlineto{\pgfqpoint{1.187768in}{2.174384in}}%
\pgfpathlineto{\pgfqpoint{1.190287in}{2.221960in}}%
\pgfpathlineto{\pgfqpoint{1.192806in}{2.194364in}}%
\pgfpathlineto{\pgfqpoint{1.195324in}{2.226921in}}%
\pgfpathlineto{\pgfqpoint{1.197843in}{2.271532in}}%
\pgfpathlineto{\pgfqpoint{1.200362in}{2.272387in}}%
\pgfpathlineto{\pgfqpoint{1.202881in}{2.342897in}}%
\pgfpathlineto{\pgfqpoint{1.205399in}{2.332313in}}%
\pgfpathlineto{\pgfqpoint{1.207918in}{2.294604in}}%
\pgfpathlineto{\pgfqpoint{1.210437in}{2.366782in}}%
\pgfpathlineto{\pgfqpoint{1.212956in}{2.370244in}}%
\pgfpathlineto{\pgfqpoint{1.215474in}{2.337703in}}%
\pgfpathlineto{\pgfqpoint{1.217993in}{2.377641in}}%
\pgfpathlineto{\pgfqpoint{1.220512in}{2.365552in}}%
\pgfpathlineto{\pgfqpoint{1.223031in}{2.347351in}}%
\pgfpathlineto{\pgfqpoint{1.225549in}{2.339127in}}%
\pgfpathlineto{\pgfqpoint{1.228068in}{2.365743in}}%
\pgfpathlineto{\pgfqpoint{1.230587in}{2.350715in}}%
\pgfpathlineto{\pgfqpoint{1.233106in}{2.387283in}}%
\pgfpathlineto{\pgfqpoint{1.235624in}{2.387764in}}%
\pgfpathlineto{\pgfqpoint{1.238143in}{2.406788in}}%
\pgfpathlineto{\pgfqpoint{1.240662in}{2.429517in}}%
\pgfpathlineto{\pgfqpoint{1.243181in}{2.416739in}}%
\pgfpathlineto{\pgfqpoint{1.245699in}{2.434658in}}%
\pgfpathlineto{\pgfqpoint{1.248218in}{2.428703in}}%
\pgfpathlineto{\pgfqpoint{1.250737in}{2.401387in}}%
\pgfpathlineto{\pgfqpoint{1.253256in}{2.379799in}}%
\pgfpathlineto{\pgfqpoint{1.255774in}{2.368808in}}%
\pgfpathlineto{\pgfqpoint{1.258293in}{2.366873in}}%
\pgfpathlineto{\pgfqpoint{1.260812in}{2.352913in}}%
\pgfpathlineto{\pgfqpoint{1.263331in}{2.330330in}}%
\pgfpathlineto{\pgfqpoint{1.265849in}{2.304473in}}%
\pgfpathlineto{\pgfqpoint{1.268368in}{2.294265in}}%
\pgfpathlineto{\pgfqpoint{1.270887in}{2.325493in}}%
\pgfpathlineto{\pgfqpoint{1.273406in}{2.326460in}}%
\pgfpathlineto{\pgfqpoint{1.275924in}{2.331803in}}%
\pgfpathlineto{\pgfqpoint{1.278443in}{2.316559in}}%
\pgfpathlineto{\pgfqpoint{1.280962in}{2.334010in}}%
\pgfpathlineto{\pgfqpoint{1.283481in}{2.286514in}}%
\pgfpathlineto{\pgfqpoint{1.285999in}{2.296862in}}%
\pgfpathlineto{\pgfqpoint{1.288518in}{2.319340in}}%
\pgfpathlineto{\pgfqpoint{1.291037in}{2.314107in}}%
\pgfpathlineto{\pgfqpoint{1.293556in}{2.319886in}}%
\pgfpathlineto{\pgfqpoint{1.296074in}{2.251020in}}%
\pgfpathlineto{\pgfqpoint{1.298593in}{2.273447in}}%
\pgfpathlineto{\pgfqpoint{1.301112in}{2.297828in}}%
\pgfpathlineto{\pgfqpoint{1.303631in}{2.345807in}}%
\pgfpathlineto{\pgfqpoint{1.306149in}{2.376567in}}%
\pgfpathlineto{\pgfqpoint{1.308668in}{2.398239in}}%
\pgfpathlineto{\pgfqpoint{1.311187in}{2.380150in}}%
\pgfpathlineto{\pgfqpoint{1.313706in}{2.379858in}}%
\pgfpathlineto{\pgfqpoint{1.316224in}{2.384556in}}%
\pgfpathlineto{\pgfqpoint{1.318743in}{2.354472in}}%
\pgfpathlineto{\pgfqpoint{1.321262in}{2.375855in}}%
\pgfpathlineto{\pgfqpoint{1.323781in}{2.464020in}}%
\pgfpathlineto{\pgfqpoint{1.326299in}{2.480322in}}%
\pgfpathlineto{\pgfqpoint{1.328818in}{2.485436in}}%
\pgfpathlineto{\pgfqpoint{1.331337in}{2.503987in}}%
\pgfpathlineto{\pgfqpoint{1.333856in}{2.513426in}}%
\pgfpathlineto{\pgfqpoint{1.336374in}{2.491437in}}%
\pgfpathlineto{\pgfqpoint{1.338893in}{2.520474in}}%
\pgfpathlineto{\pgfqpoint{1.341412in}{2.462732in}}%
\pgfpathlineto{\pgfqpoint{1.343931in}{2.440515in}}%
\pgfpathlineto{\pgfqpoint{1.346449in}{2.458788in}}%
\pgfpathlineto{\pgfqpoint{1.348968in}{2.434699in}}%
\pgfpathlineto{\pgfqpoint{1.351487in}{2.390214in}}%
\pgfpathlineto{\pgfqpoint{1.354006in}{2.367919in}}%
\pgfpathlineto{\pgfqpoint{1.356524in}{2.321936in}}%
\pgfpathlineto{\pgfqpoint{1.359043in}{2.311110in}}%
\pgfpathlineto{\pgfqpoint{1.361562in}{2.302614in}}%
\pgfpathlineto{\pgfqpoint{1.364081in}{2.265300in}}%
\pgfpathlineto{\pgfqpoint{1.366599in}{2.229713in}}%
\pgfpathlineto{\pgfqpoint{1.369118in}{2.225216in}}%
\pgfpathlineto{\pgfqpoint{1.371637in}{2.230303in}}%
\pgfpathlineto{\pgfqpoint{1.374156in}{2.207628in}}%
\pgfpathlineto{\pgfqpoint{1.376674in}{2.209035in}}%
\pgfpathlineto{\pgfqpoint{1.379193in}{2.211576in}}%
\pgfpathlineto{\pgfqpoint{1.381712in}{2.197471in}}%
\pgfpathlineto{\pgfqpoint{1.384231in}{2.176539in}}%
\pgfpathlineto{\pgfqpoint{1.386749in}{2.207345in}}%
\pgfpathlineto{\pgfqpoint{1.389268in}{2.202773in}}%
\pgfpathlineto{\pgfqpoint{1.391787in}{2.191280in}}%
\pgfpathlineto{\pgfqpoint{1.394306in}{2.223901in}}%
\pgfpathlineto{\pgfqpoint{1.396824in}{2.218831in}}%
\pgfpathlineto{\pgfqpoint{1.399343in}{2.261417in}}%
\pgfpathlineto{\pgfqpoint{1.401862in}{2.226290in}}%
\pgfpathlineto{\pgfqpoint{1.404381in}{2.242868in}}%
\pgfpathlineto{\pgfqpoint{1.406899in}{2.225594in}}%
\pgfpathlineto{\pgfqpoint{1.409418in}{2.239633in}}%
\pgfpathlineto{\pgfqpoint{1.411937in}{2.278032in}}%
\pgfpathlineto{\pgfqpoint{1.414456in}{2.246881in}}%
\pgfpathlineto{\pgfqpoint{1.416974in}{2.245498in}}%
\pgfpathlineto{\pgfqpoint{1.419493in}{2.273197in}}%
\pgfpathlineto{\pgfqpoint{1.422012in}{2.253997in}}%
\pgfpathlineto{\pgfqpoint{1.424531in}{2.280440in}}%
\pgfpathlineto{\pgfqpoint{1.427049in}{2.229410in}}%
\pgfpathlineto{\pgfqpoint{1.429568in}{2.219250in}}%
\pgfpathlineto{\pgfqpoint{1.432087in}{2.230560in}}%
\pgfpathlineto{\pgfqpoint{1.434606in}{2.192781in}}%
\pgfpathlineto{\pgfqpoint{1.437124in}{2.169292in}}%
\pgfpathlineto{\pgfqpoint{1.439643in}{2.139741in}}%
\pgfpathlineto{\pgfqpoint{1.442162in}{2.174681in}}%
\pgfpathlineto{\pgfqpoint{1.444681in}{2.177240in}}%
\pgfpathlineto{\pgfqpoint{1.447199in}{2.142558in}}%
\pgfpathlineto{\pgfqpoint{1.449718in}{2.160866in}}%
\pgfpathlineto{\pgfqpoint{1.452237in}{2.182581in}}%
\pgfpathlineto{\pgfqpoint{1.454756in}{2.168917in}}%
\pgfpathlineto{\pgfqpoint{1.457274in}{2.204983in}}%
\pgfpathlineto{\pgfqpoint{1.459793in}{2.193680in}}%
\pgfpathlineto{\pgfqpoint{1.462312in}{2.167531in}}%
\pgfpathlineto{\pgfqpoint{1.464831in}{2.201614in}}%
\pgfpathlineto{\pgfqpoint{1.467349in}{2.192885in}}%
\pgfpathlineto{\pgfqpoint{1.469868in}{2.246731in}}%
\pgfpathlineto{\pgfqpoint{1.472387in}{2.296103in}}%
\pgfpathlineto{\pgfqpoint{1.474906in}{2.336814in}}%
\pgfpathlineto{\pgfqpoint{1.477424in}{2.369818in}}%
\pgfpathlineto{\pgfqpoint{1.479943in}{2.375959in}}%
\pgfpathlineto{\pgfqpoint{1.482462in}{2.385468in}}%
\pgfpathlineto{\pgfqpoint{1.484981in}{2.380924in}}%
\pgfpathlineto{\pgfqpoint{1.487499in}{2.436813in}}%
\pgfpathlineto{\pgfqpoint{1.490018in}{2.477228in}}%
\pgfpathlineto{\pgfqpoint{1.492537in}{2.466031in}}%
\pgfpathlineto{\pgfqpoint{1.495056in}{2.506895in}}%
\pgfpathlineto{\pgfqpoint{1.497574in}{2.511047in}}%
\pgfpathlineto{\pgfqpoint{1.500093in}{2.496457in}}%
\pgfpathlineto{\pgfqpoint{1.502612in}{2.480737in}}%
\pgfpathlineto{\pgfqpoint{1.505131in}{2.477845in}}%
\pgfpathlineto{\pgfqpoint{1.507649in}{2.548343in}}%
\pgfpathlineto{\pgfqpoint{1.510168in}{2.490218in}}%
\pgfpathlineto{\pgfqpoint{1.512687in}{2.539056in}}%
\pgfpathlineto{\pgfqpoint{1.515206in}{2.530207in}}%
\pgfpathlineto{\pgfqpoint{1.517724in}{2.535508in}}%
\pgfpathlineto{\pgfqpoint{1.520243in}{2.527635in}}%
\pgfpathlineto{\pgfqpoint{1.522762in}{2.540905in}}%
\pgfpathlineto{\pgfqpoint{1.525281in}{2.481060in}}%
\pgfpathlineto{\pgfqpoint{1.527799in}{2.482612in}}%
\pgfpathlineto{\pgfqpoint{1.530318in}{2.433908in}}%
\pgfpathlineto{\pgfqpoint{1.532837in}{2.452244in}}%
\pgfpathlineto{\pgfqpoint{1.535356in}{2.437327in}}%
\pgfpathlineto{\pgfqpoint{1.537874in}{2.442460in}}%
\pgfpathlineto{\pgfqpoint{1.540393in}{2.392915in}}%
\pgfpathlineto{\pgfqpoint{1.542912in}{2.379496in}}%
\pgfpathlineto{\pgfqpoint{1.547949in}{2.391325in}}%
\pgfpathlineto{\pgfqpoint{1.550468in}{2.401211in}}%
\pgfpathlineto{\pgfqpoint{1.552987in}{2.451945in}}%
\pgfpathlineto{\pgfqpoint{1.555506in}{2.438044in}}%
\pgfpathlineto{\pgfqpoint{1.558024in}{2.465980in}}%
\pgfpathlineto{\pgfqpoint{1.560543in}{2.452831in}}%
\pgfpathlineto{\pgfqpoint{1.563062in}{2.471764in}}%
\pgfpathlineto{\pgfqpoint{1.565581in}{2.441731in}}%
\pgfpathlineto{\pgfqpoint{1.568099in}{2.418565in}}%
\pgfpathlineto{\pgfqpoint{1.570618in}{2.488555in}}%
\pgfpathlineto{\pgfqpoint{1.573137in}{2.492519in}}%
\pgfpathlineto{\pgfqpoint{1.575656in}{2.467334in}}%
\pgfpathlineto{\pgfqpoint{1.578174in}{2.409860in}}%
\pgfpathlineto{\pgfqpoint{1.580693in}{2.381868in}}%
\pgfpathlineto{\pgfqpoint{1.583212in}{2.355155in}}%
\pgfpathlineto{\pgfqpoint{1.585731in}{2.329479in}}%
\pgfpathlineto{\pgfqpoint{1.588249in}{2.326123in}}%
\pgfpathlineto{\pgfqpoint{1.590768in}{2.316916in}}%
\pgfpathlineto{\pgfqpoint{1.593287in}{2.349341in}}%
\pgfpathlineto{\pgfqpoint{1.595806in}{2.362328in}}%
\pgfpathlineto{\pgfqpoint{1.598324in}{2.414362in}}%
\pgfpathlineto{\pgfqpoint{1.600843in}{2.376397in}}%
\pgfpathlineto{\pgfqpoint{1.603362in}{2.378379in}}%
\pgfpathlineto{\pgfqpoint{1.605881in}{2.407882in}}%
\pgfpathlineto{\pgfqpoint{1.608399in}{2.454579in}}%
\pgfpathlineto{\pgfqpoint{1.610918in}{2.411715in}}%
\pgfpathlineto{\pgfqpoint{1.613437in}{2.416300in}}%
\pgfpathlineto{\pgfqpoint{1.615956in}{2.410328in}}%
\pgfpathlineto{\pgfqpoint{1.618474in}{2.408585in}}%
\pgfpathlineto{\pgfqpoint{1.620993in}{2.389113in}}%
\pgfpathlineto{\pgfqpoint{1.623512in}{2.384158in}}%
\pgfpathlineto{\pgfqpoint{1.626031in}{2.332898in}}%
\pgfpathlineto{\pgfqpoint{1.628549in}{2.334419in}}%
\pgfpathlineto{\pgfqpoint{1.631068in}{2.280599in}}%
\pgfpathlineto{\pgfqpoint{1.633587in}{2.308843in}}%
\pgfpathlineto{\pgfqpoint{1.636106in}{2.310481in}}%
\pgfpathlineto{\pgfqpoint{1.638624in}{2.327191in}}%
\pgfpathlineto{\pgfqpoint{1.641143in}{2.392747in}}%
\pgfpathlineto{\pgfqpoint{1.643662in}{2.396478in}}%
\pgfpathlineto{\pgfqpoint{1.646181in}{2.476843in}}%
\pgfpathlineto{\pgfqpoint{1.648699in}{2.480983in}}%
\pgfpathlineto{\pgfqpoint{1.651218in}{2.473401in}}%
\pgfpathlineto{\pgfqpoint{1.653737in}{2.492092in}}%
\pgfpathlineto{\pgfqpoint{1.656256in}{2.516799in}}%
\pgfpathlineto{\pgfqpoint{1.658774in}{2.547475in}}%
\pgfpathlineto{\pgfqpoint{1.661293in}{2.549481in}}%
\pgfpathlineto{\pgfqpoint{1.663812in}{2.514748in}}%
\pgfpathlineto{\pgfqpoint{1.666331in}{2.517118in}}%
\pgfpathlineto{\pgfqpoint{1.668849in}{2.577844in}}%
\pgfpathlineto{\pgfqpoint{1.671368in}{2.551790in}}%
\pgfpathlineto{\pgfqpoint{1.673887in}{2.583428in}}%
\pgfpathlineto{\pgfqpoint{1.676406in}{2.588141in}}%
\pgfpathlineto{\pgfqpoint{1.678924in}{2.604424in}}%
\pgfpathlineto{\pgfqpoint{1.681443in}{2.613070in}}%
\pgfpathlineto{\pgfqpoint{1.683962in}{2.654397in}}%
\pgfpathlineto{\pgfqpoint{1.686481in}{2.633931in}}%
\pgfpathlineto{\pgfqpoint{1.688999in}{2.662683in}}%
\pgfpathlineto{\pgfqpoint{1.691518in}{2.686068in}}%
\pgfpathlineto{\pgfqpoint{1.694037in}{2.707866in}}%
\pgfpathlineto{\pgfqpoint{1.696556in}{2.755002in}}%
\pgfpathlineto{\pgfqpoint{1.699074in}{2.702589in}}%
\pgfpathlineto{\pgfqpoint{1.701593in}{2.671314in}}%
\pgfpathlineto{\pgfqpoint{1.704112in}{2.760882in}}%
\pgfpathlineto{\pgfqpoint{1.706631in}{2.768491in}}%
\pgfpathlineto{\pgfqpoint{1.709149in}{2.762926in}}%
\pgfpathlineto{\pgfqpoint{1.711668in}{2.755637in}}%
\pgfpathlineto{\pgfqpoint{1.716706in}{2.809891in}}%
\pgfpathlineto{\pgfqpoint{1.719224in}{2.823165in}}%
\pgfpathlineto{\pgfqpoint{1.721743in}{2.851129in}}%
\pgfpathlineto{\pgfqpoint{1.724262in}{2.856984in}}%
\pgfpathlineto{\pgfqpoint{1.726781in}{2.881634in}}%
\pgfpathlineto{\pgfqpoint{1.729299in}{2.925554in}}%
\pgfpathlineto{\pgfqpoint{1.731818in}{2.911618in}}%
\pgfpathlineto{\pgfqpoint{1.734337in}{2.949000in}}%
\pgfpathlineto{\pgfqpoint{1.736856in}{2.963631in}}%
\pgfpathlineto{\pgfqpoint{1.739374in}{2.922945in}}%
\pgfpathlineto{\pgfqpoint{1.741893in}{2.925150in}}%
\pgfpathlineto{\pgfqpoint{1.744412in}{2.959009in}}%
\pgfpathlineto{\pgfqpoint{1.746931in}{2.949900in}}%
\pgfpathlineto{\pgfqpoint{1.749449in}{2.924256in}}%
\pgfpathlineto{\pgfqpoint{1.751968in}{2.896977in}}%
\pgfpathlineto{\pgfqpoint{1.754487in}{2.921842in}}%
\pgfpathlineto{\pgfqpoint{1.757006in}{2.909070in}}%
\pgfpathlineto{\pgfqpoint{1.759524in}{2.902059in}}%
\pgfpathlineto{\pgfqpoint{1.762043in}{2.926322in}}%
\pgfpathlineto{\pgfqpoint{1.764562in}{2.920948in}}%
\pgfpathlineto{\pgfqpoint{1.767081in}{2.927500in}}%
\pgfpathlineto{\pgfqpoint{1.769599in}{2.922917in}}%
\pgfpathlineto{\pgfqpoint{1.772118in}{2.966487in}}%
\pgfpathlineto{\pgfqpoint{1.774637in}{2.951397in}}%
\pgfpathlineto{\pgfqpoint{1.777156in}{2.921704in}}%
\pgfpathlineto{\pgfqpoint{1.779674in}{2.910314in}}%
\pgfpathlineto{\pgfqpoint{1.782193in}{2.961387in}}%
\pgfpathlineto{\pgfqpoint{1.784712in}{2.945919in}}%
\pgfpathlineto{\pgfqpoint{1.787231in}{2.933727in}}%
\pgfpathlineto{\pgfqpoint{1.789749in}{2.948331in}}%
\pgfpathlineto{\pgfqpoint{1.792268in}{2.957779in}}%
\pgfpathlineto{\pgfqpoint{1.794787in}{2.922222in}}%
\pgfpathlineto{\pgfqpoint{1.797306in}{2.994258in}}%
\pgfpathlineto{\pgfqpoint{1.799824in}{2.959580in}}%
\pgfpathlineto{\pgfqpoint{1.802343in}{2.987670in}}%
\pgfpathlineto{\pgfqpoint{1.804862in}{3.005375in}}%
\pgfpathlineto{\pgfqpoint{1.807381in}{3.024548in}}%
\pgfpathlineto{\pgfqpoint{1.809899in}{3.058949in}}%
\pgfpathlineto{\pgfqpoint{1.812418in}{3.050493in}}%
\pgfpathlineto{\pgfqpoint{1.814937in}{3.081717in}}%
\pgfpathlineto{\pgfqpoint{1.817456in}{3.062870in}}%
\pgfpathlineto{\pgfqpoint{1.819974in}{3.061403in}}%
\pgfpathlineto{\pgfqpoint{1.822493in}{3.046945in}}%
\pgfpathlineto{\pgfqpoint{1.825012in}{3.047431in}}%
\pgfpathlineto{\pgfqpoint{1.827531in}{3.083173in}}%
\pgfpathlineto{\pgfqpoint{1.830049in}{3.111936in}}%
\pgfpathlineto{\pgfqpoint{1.832568in}{3.134441in}}%
\pgfpathlineto{\pgfqpoint{1.835087in}{3.091834in}}%
\pgfpathlineto{\pgfqpoint{1.837606in}{3.068490in}}%
\pgfpathlineto{\pgfqpoint{1.840124in}{3.072971in}}%
\pgfpathlineto{\pgfqpoint{1.842643in}{3.080095in}}%
\pgfpathlineto{\pgfqpoint{1.845162in}{3.057970in}}%
\pgfpathlineto{\pgfqpoint{1.847681in}{3.102389in}}%
\pgfpathlineto{\pgfqpoint{1.850199in}{3.149343in}}%
\pgfpathlineto{\pgfqpoint{1.852718in}{3.120687in}}%
\pgfpathlineto{\pgfqpoint{1.855237in}{3.131505in}}%
\pgfpathlineto{\pgfqpoint{1.857756in}{3.149661in}}%
\pgfpathlineto{\pgfqpoint{1.860274in}{3.141336in}}%
\pgfpathlineto{\pgfqpoint{1.862793in}{3.163759in}}%
\pgfpathlineto{\pgfqpoint{1.865312in}{3.175037in}}%
\pgfpathlineto{\pgfqpoint{1.867831in}{3.160372in}}%
\pgfpathlineto{\pgfqpoint{1.870349in}{3.143906in}}%
\pgfpathlineto{\pgfqpoint{1.872868in}{3.171353in}}%
\pgfpathlineto{\pgfqpoint{1.875387in}{3.168845in}}%
\pgfpathlineto{\pgfqpoint{1.877906in}{3.195149in}}%
\pgfpathlineto{\pgfqpoint{1.880424in}{3.197473in}}%
\pgfpathlineto{\pgfqpoint{1.882943in}{3.202862in}}%
\pgfpathlineto{\pgfqpoint{1.885462in}{3.233660in}}%
\pgfpathlineto{\pgfqpoint{1.887981in}{3.268072in}}%
\pgfpathlineto{\pgfqpoint{1.890499in}{3.280447in}}%
\pgfpathlineto{\pgfqpoint{1.893018in}{3.312245in}}%
\pgfpathlineto{\pgfqpoint{1.895537in}{3.331800in}}%
\pgfpathlineto{\pgfqpoint{1.898056in}{3.331468in}}%
\pgfpathlineto{\pgfqpoint{1.900574in}{3.329674in}}%
\pgfpathlineto{\pgfqpoint{1.903093in}{3.331048in}}%
\pgfpathlineto{\pgfqpoint{1.905612in}{3.317037in}}%
\pgfpathlineto{\pgfqpoint{1.908131in}{3.276156in}}%
\pgfpathlineto{\pgfqpoint{1.910649in}{3.306948in}}%
\pgfpathlineto{\pgfqpoint{1.913168in}{3.329377in}}%
\pgfpathlineto{\pgfqpoint{1.915687in}{3.343500in}}%
\pgfpathlineto{\pgfqpoint{1.918206in}{3.381411in}}%
\pgfpathlineto{\pgfqpoint{1.920724in}{3.379500in}}%
\pgfpathlineto{\pgfqpoint{1.923243in}{3.360087in}}%
\pgfpathlineto{\pgfqpoint{1.925762in}{3.381445in}}%
\pgfpathlineto{\pgfqpoint{1.928281in}{3.405987in}}%
\pgfpathlineto{\pgfqpoint{1.930799in}{3.432953in}}%
\pgfpathlineto{\pgfqpoint{1.933318in}{3.425547in}}%
\pgfpathlineto{\pgfqpoint{1.935837in}{3.392129in}}%
\pgfpathlineto{\pgfqpoint{1.938356in}{3.377250in}}%
\pgfpathlineto{\pgfqpoint{1.940874in}{3.303306in}}%
\pgfpathlineto{\pgfqpoint{1.943393in}{3.304195in}}%
\pgfpathlineto{\pgfqpoint{1.945912in}{3.291785in}}%
\pgfpathlineto{\pgfqpoint{1.948431in}{3.278543in}}%
\pgfpathlineto{\pgfqpoint{1.950949in}{3.294433in}}%
\pgfpathlineto{\pgfqpoint{1.953468in}{3.242621in}}%
\pgfpathlineto{\pgfqpoint{1.955987in}{3.233605in}}%
\pgfpathlineto{\pgfqpoint{1.958506in}{3.261363in}}%
\pgfpathlineto{\pgfqpoint{1.961024in}{3.277884in}}%
\pgfpathlineto{\pgfqpoint{1.963543in}{3.258715in}}%
\pgfpathlineto{\pgfqpoint{1.966062in}{3.222341in}}%
\pgfpathlineto{\pgfqpoint{1.968581in}{3.228920in}}%
\pgfpathlineto{\pgfqpoint{1.971099in}{3.253965in}}%
\pgfpathlineto{\pgfqpoint{1.973618in}{3.271155in}}%
\pgfpathlineto{\pgfqpoint{1.976137in}{3.253854in}}%
\pgfpathlineto{\pgfqpoint{1.978656in}{3.252631in}}%
\pgfpathlineto{\pgfqpoint{1.981174in}{3.252826in}}%
\pgfpathlineto{\pgfqpoint{1.983693in}{3.212467in}}%
\pgfpathlineto{\pgfqpoint{1.986212in}{3.207462in}}%
\pgfpathlineto{\pgfqpoint{1.988731in}{3.163917in}}%
\pgfpathlineto{\pgfqpoint{1.991249in}{3.169557in}}%
\pgfpathlineto{\pgfqpoint{1.993768in}{3.171407in}}%
\pgfpathlineto{\pgfqpoint{1.996287in}{3.115168in}}%
\pgfpathlineto{\pgfqpoint{1.998806in}{3.118849in}}%
\pgfpathlineto{\pgfqpoint{2.001324in}{3.148504in}}%
\pgfpathlineto{\pgfqpoint{2.003843in}{3.176892in}}%
\pgfpathlineto{\pgfqpoint{2.006362in}{3.178726in}}%
\pgfpathlineto{\pgfqpoint{2.008881in}{3.202650in}}%
\pgfpathlineto{\pgfqpoint{2.011399in}{3.247119in}}%
\pgfpathlineto{\pgfqpoint{2.013918in}{3.254737in}}%
\pgfpathlineto{\pgfqpoint{2.016437in}{3.287152in}}%
\pgfpathlineto{\pgfqpoint{2.018956in}{3.237185in}}%
\pgfpathlineto{\pgfqpoint{2.021474in}{3.244854in}}%
\pgfpathlineto{\pgfqpoint{2.023993in}{3.261981in}}%
\pgfpathlineto{\pgfqpoint{2.026512in}{3.234329in}}%
\pgfpathlineto{\pgfqpoint{2.029031in}{3.224277in}}%
\pgfpathlineto{\pgfqpoint{2.031549in}{3.268700in}}%
\pgfpathlineto{\pgfqpoint{2.034068in}{3.234185in}}%
\pgfpathlineto{\pgfqpoint{2.036587in}{3.241507in}}%
\pgfpathlineto{\pgfqpoint{2.039106in}{3.201432in}}%
\pgfpathlineto{\pgfqpoint{2.041624in}{3.223374in}}%
\pgfpathlineto{\pgfqpoint{2.044143in}{3.252013in}}%
\pgfpathlineto{\pgfqpoint{2.046662in}{3.258171in}}%
\pgfpathlineto{\pgfqpoint{2.049181in}{3.262170in}}%
\pgfpathlineto{\pgfqpoint{2.051699in}{3.240249in}}%
\pgfpathlineto{\pgfqpoint{2.054218in}{3.226994in}}%
\pgfpathlineto{\pgfqpoint{2.056737in}{3.249965in}}%
\pgfpathlineto{\pgfqpoint{2.059256in}{3.236026in}}%
\pgfpathlineto{\pgfqpoint{2.061774in}{3.231193in}}%
\pgfpathlineto{\pgfqpoint{2.064293in}{3.279449in}}%
\pgfpathlineto{\pgfqpoint{2.066812in}{3.269999in}}%
\pgfpathlineto{\pgfqpoint{2.069331in}{3.234132in}}%
\pgfpathlineto{\pgfqpoint{2.071849in}{3.226900in}}%
\pgfpathlineto{\pgfqpoint{2.074368in}{3.308491in}}%
\pgfpathlineto{\pgfqpoint{2.076887in}{3.326559in}}%
\pgfpathlineto{\pgfqpoint{2.079406in}{3.306078in}}%
\pgfpathlineto{\pgfqpoint{2.081924in}{3.287253in}}%
\pgfpathlineto{\pgfqpoint{2.084443in}{3.307227in}}%
\pgfpathlineto{\pgfqpoint{2.086962in}{3.317799in}}%
\pgfpathlineto{\pgfqpoint{2.089481in}{3.288333in}}%
\pgfpathlineto{\pgfqpoint{2.091999in}{3.329865in}}%
\pgfpathlineto{\pgfqpoint{2.094518in}{3.311456in}}%
\pgfpathlineto{\pgfqpoint{2.097037in}{3.272063in}}%
\pgfpathlineto{\pgfqpoint{2.099556in}{3.260170in}}%
\pgfpathlineto{\pgfqpoint{2.102074in}{3.247477in}}%
\pgfpathlineto{\pgfqpoint{2.104593in}{3.199291in}}%
\pgfpathlineto{\pgfqpoint{2.107112in}{3.184645in}}%
\pgfpathlineto{\pgfqpoint{2.112149in}{3.142228in}}%
\pgfpathlineto{\pgfqpoint{2.114668in}{3.115329in}}%
\pgfpathlineto{\pgfqpoint{2.117187in}{3.129615in}}%
\pgfpathlineto{\pgfqpoint{2.119706in}{3.118223in}}%
\pgfpathlineto{\pgfqpoint{2.122224in}{3.125953in}}%
\pgfpathlineto{\pgfqpoint{2.124743in}{3.110411in}}%
\pgfpathlineto{\pgfqpoint{2.127262in}{3.111063in}}%
\pgfpathlineto{\pgfqpoint{2.129781in}{3.138832in}}%
\pgfpathlineto{\pgfqpoint{2.132299in}{3.126660in}}%
\pgfpathlineto{\pgfqpoint{2.134818in}{3.179441in}}%
\pgfpathlineto{\pgfqpoint{2.137337in}{3.210701in}}%
\pgfpathlineto{\pgfqpoint{2.139856in}{3.142162in}}%
\pgfpathlineto{\pgfqpoint{2.142374in}{3.121757in}}%
\pgfpathlineto{\pgfqpoint{2.144893in}{3.118552in}}%
\pgfpathlineto{\pgfqpoint{2.147412in}{3.111254in}}%
\pgfpathlineto{\pgfqpoint{2.149931in}{3.115299in}}%
\pgfpathlineto{\pgfqpoint{2.152449in}{3.076951in}}%
\pgfpathlineto{\pgfqpoint{2.154968in}{3.062509in}}%
\pgfpathlineto{\pgfqpoint{2.157487in}{3.064191in}}%
\pgfpathlineto{\pgfqpoint{2.160006in}{3.017938in}}%
\pgfpathlineto{\pgfqpoint{2.162524in}{2.984542in}}%
\pgfpathlineto{\pgfqpoint{2.165043in}{2.974884in}}%
\pgfpathlineto{\pgfqpoint{2.167562in}{3.005337in}}%
\pgfpathlineto{\pgfqpoint{2.172599in}{3.029435in}}%
\pgfpathlineto{\pgfqpoint{2.175118in}{3.056060in}}%
\pgfpathlineto{\pgfqpoint{2.177637in}{3.028235in}}%
\pgfpathlineto{\pgfqpoint{2.180156in}{3.071390in}}%
\pgfpathlineto{\pgfqpoint{2.182674in}{3.089162in}}%
\pgfpathlineto{\pgfqpoint{2.185193in}{3.093437in}}%
\pgfpathlineto{\pgfqpoint{2.187712in}{3.029326in}}%
\pgfpathlineto{\pgfqpoint{2.190231in}{3.044698in}}%
\pgfpathlineto{\pgfqpoint{2.192749in}{3.046763in}}%
\pgfpathlineto{\pgfqpoint{2.195268in}{3.067870in}}%
\pgfpathlineto{\pgfqpoint{2.197787in}{3.076148in}}%
\pgfpathlineto{\pgfqpoint{2.200306in}{3.060259in}}%
\pgfpathlineto{\pgfqpoint{2.202824in}{3.085678in}}%
\pgfpathlineto{\pgfqpoint{2.205343in}{3.084331in}}%
\pgfpathlineto{\pgfqpoint{2.207862in}{3.092363in}}%
\pgfpathlineto{\pgfqpoint{2.210381in}{3.078940in}}%
\pgfpathlineto{\pgfqpoint{2.212899in}{3.093080in}}%
\pgfpathlineto{\pgfqpoint{2.215418in}{3.113622in}}%
\pgfpathlineto{\pgfqpoint{2.217937in}{3.128392in}}%
\pgfpathlineto{\pgfqpoint{2.220456in}{3.146653in}}%
\pgfpathlineto{\pgfqpoint{2.222974in}{3.168049in}}%
\pgfpathlineto{\pgfqpoint{2.225493in}{3.206266in}}%
\pgfpathlineto{\pgfqpoint{2.228012in}{3.198409in}}%
\pgfpathlineto{\pgfqpoint{2.230531in}{3.219608in}}%
\pgfpathlineto{\pgfqpoint{2.233049in}{3.256611in}}%
\pgfpathlineto{\pgfqpoint{2.235568in}{3.253876in}}%
\pgfpathlineto{\pgfqpoint{2.238087in}{3.296984in}}%
\pgfpathlineto{\pgfqpoint{2.240606in}{3.302638in}}%
\pgfpathlineto{\pgfqpoint{2.243124in}{3.254422in}}%
\pgfpathlineto{\pgfqpoint{2.245643in}{3.271660in}}%
\pgfpathlineto{\pgfqpoint{2.248162in}{3.264903in}}%
\pgfpathlineto{\pgfqpoint{2.250681in}{3.249085in}}%
\pgfpathlineto{\pgfqpoint{2.253199in}{3.307337in}}%
\pgfpathlineto{\pgfqpoint{2.255718in}{3.317179in}}%
\pgfpathlineto{\pgfqpoint{2.258237in}{3.322127in}}%
\pgfpathlineto{\pgfqpoint{2.260756in}{3.288029in}}%
\pgfpathlineto{\pgfqpoint{2.263274in}{3.281670in}}%
\pgfpathlineto{\pgfqpoint{2.265793in}{3.317508in}}%
\pgfpathlineto{\pgfqpoint{2.268312in}{3.313273in}}%
\pgfpathlineto{\pgfqpoint{2.270831in}{3.354704in}}%
\pgfpathlineto{\pgfqpoint{2.273349in}{3.371551in}}%
\pgfpathlineto{\pgfqpoint{2.275868in}{3.307572in}}%
\pgfpathlineto{\pgfqpoint{2.278387in}{3.309118in}}%
\pgfpathlineto{\pgfqpoint{2.280906in}{3.339902in}}%
\pgfpathlineto{\pgfqpoint{2.283424in}{3.336748in}}%
\pgfpathlineto{\pgfqpoint{2.285943in}{3.335026in}}%
\pgfpathlineto{\pgfqpoint{2.288462in}{3.339392in}}%
\pgfpathlineto{\pgfqpoint{2.290981in}{3.303045in}}%
\pgfpathlineto{\pgfqpoint{2.293499in}{3.316572in}}%
\pgfpathlineto{\pgfqpoint{2.296018in}{3.327546in}}%
\pgfpathlineto{\pgfqpoint{2.298537in}{3.352452in}}%
\pgfpathlineto{\pgfqpoint{2.301056in}{3.391429in}}%
\pgfpathlineto{\pgfqpoint{2.303574in}{3.387057in}}%
\pgfpathlineto{\pgfqpoint{2.306093in}{3.347811in}}%
\pgfpathlineto{\pgfqpoint{2.308612in}{3.375520in}}%
\pgfpathlineto{\pgfqpoint{2.311131in}{3.386357in}}%
\pgfpathlineto{\pgfqpoint{2.313649in}{3.377375in}}%
\pgfpathlineto{\pgfqpoint{2.316168in}{3.380979in}}%
\pgfpathlineto{\pgfqpoint{2.318687in}{3.367586in}}%
\pgfpathlineto{\pgfqpoint{2.321206in}{3.371727in}}%
\pgfpathlineto{\pgfqpoint{2.323724in}{3.341867in}}%
\pgfpathlineto{\pgfqpoint{2.326243in}{3.392400in}}%
\pgfpathlineto{\pgfqpoint{2.328762in}{3.387079in}}%
\pgfpathlineto{\pgfqpoint{2.331281in}{3.341809in}}%
\pgfpathlineto{\pgfqpoint{2.333799in}{3.345583in}}%
\pgfpathlineto{\pgfqpoint{2.336318in}{3.357464in}}%
\pgfpathlineto{\pgfqpoint{2.338837in}{3.400483in}}%
\pgfpathlineto{\pgfqpoint{2.341356in}{3.426079in}}%
\pgfpathlineto{\pgfqpoint{2.343874in}{3.440996in}}%
\pgfpathlineto{\pgfqpoint{2.346393in}{3.486061in}}%
\pgfpathlineto{\pgfqpoint{2.348912in}{3.471474in}}%
\pgfpathlineto{\pgfqpoint{2.351431in}{3.494553in}}%
\pgfpathlineto{\pgfqpoint{2.353949in}{3.468266in}}%
\pgfpathlineto{\pgfqpoint{2.356468in}{3.461548in}}%
\pgfpathlineto{\pgfqpoint{2.358987in}{3.436962in}}%
\pgfpathlineto{\pgfqpoint{2.361506in}{3.449961in}}%
\pgfpathlineto{\pgfqpoint{2.364024in}{3.454928in}}%
\pgfpathlineto{\pgfqpoint{2.366543in}{3.508361in}}%
\pgfpathlineto{\pgfqpoint{2.369062in}{3.493877in}}%
\pgfpathlineto{\pgfqpoint{2.371581in}{3.538224in}}%
\pgfpathlineto{\pgfqpoint{2.374099in}{3.585036in}}%
\pgfpathlineto{\pgfqpoint{2.376618in}{3.569253in}}%
\pgfpathlineto{\pgfqpoint{2.379137in}{3.545931in}}%
\pgfpathlineto{\pgfqpoint{2.381656in}{3.542477in}}%
\pgfpathlineto{\pgfqpoint{2.384174in}{3.605837in}}%
\pgfpathlineto{\pgfqpoint{2.386693in}{3.549148in}}%
\pgfpathlineto{\pgfqpoint{2.389212in}{3.498313in}}%
\pgfpathlineto{\pgfqpoint{2.391731in}{3.508464in}}%
\pgfpathlineto{\pgfqpoint{2.394249in}{3.527766in}}%
\pgfpathlineto{\pgfqpoint{2.396768in}{3.564518in}}%
\pgfpathlineto{\pgfqpoint{2.399287in}{3.571940in}}%
\pgfpathlineto{\pgfqpoint{2.401806in}{3.584120in}}%
\pgfpathlineto{\pgfqpoint{2.404324in}{3.604423in}}%
\pgfpathlineto{\pgfqpoint{2.406843in}{3.562828in}}%
\pgfpathlineto{\pgfqpoint{2.409362in}{3.603506in}}%
\pgfpathlineto{\pgfqpoint{2.411881in}{3.627644in}}%
\pgfpathlineto{\pgfqpoint{2.414399in}{3.575811in}}%
\pgfpathlineto{\pgfqpoint{2.416918in}{3.578080in}}%
\pgfpathlineto{\pgfqpoint{2.419437in}{3.616392in}}%
\pgfpathlineto{\pgfqpoint{2.421956in}{3.611013in}}%
\pgfpathlineto{\pgfqpoint{2.424474in}{3.595035in}}%
\pgfpathlineto{\pgfqpoint{2.426993in}{3.607742in}}%
\pgfpathlineto{\pgfqpoint{2.429512in}{3.567630in}}%
\pgfpathlineto{\pgfqpoint{2.432031in}{3.638723in}}%
\pgfpathlineto{\pgfqpoint{2.434549in}{3.610959in}}%
\pgfpathlineto{\pgfqpoint{2.437068in}{3.556631in}}%
\pgfpathlineto{\pgfqpoint{2.439587in}{3.582999in}}%
\pgfpathlineto{\pgfqpoint{2.442106in}{3.582011in}}%
\pgfpathlineto{\pgfqpoint{2.444624in}{3.560325in}}%
\pgfpathlineto{\pgfqpoint{2.447143in}{3.561759in}}%
\pgfpathlineto{\pgfqpoint{2.449662in}{3.579707in}}%
\pgfpathlineto{\pgfqpoint{2.452181in}{3.598350in}}%
\pgfpathlineto{\pgfqpoint{2.454699in}{3.651845in}}%
\pgfpathlineto{\pgfqpoint{2.457218in}{3.686903in}}%
\pgfpathlineto{\pgfqpoint{2.462256in}{3.731198in}}%
\pgfpathlineto{\pgfqpoint{2.464774in}{3.783467in}}%
\pgfpathlineto{\pgfqpoint{2.467293in}{3.773474in}}%
\pgfpathlineto{\pgfqpoint{2.469812in}{3.778049in}}%
\pgfpathlineto{\pgfqpoint{2.472331in}{3.732905in}}%
\pgfpathlineto{\pgfqpoint{2.474849in}{3.674556in}}%
\pgfpathlineto{\pgfqpoint{2.477368in}{3.671198in}}%
\pgfpathlineto{\pgfqpoint{2.479887in}{3.668709in}}%
\pgfpathlineto{\pgfqpoint{2.482406in}{3.623357in}}%
\pgfpathlineto{\pgfqpoint{2.484924in}{3.624304in}}%
\pgfpathlineto{\pgfqpoint{2.487443in}{3.604327in}}%
\pgfpathlineto{\pgfqpoint{2.489962in}{3.594755in}}%
\pgfpathlineto{\pgfqpoint{2.492481in}{3.552929in}}%
\pgfpathlineto{\pgfqpoint{2.494999in}{3.589960in}}%
\pgfpathlineto{\pgfqpoint{2.497518in}{3.564941in}}%
\pgfpathlineto{\pgfqpoint{2.500037in}{3.564693in}}%
\pgfpathlineto{\pgfqpoint{2.502556in}{3.590606in}}%
\pgfpathlineto{\pgfqpoint{2.505074in}{3.619122in}}%
\pgfpathlineto{\pgfqpoint{2.507593in}{3.596692in}}%
\pgfpathlineto{\pgfqpoint{2.510112in}{3.625583in}}%
\pgfpathlineto{\pgfqpoint{2.512631in}{3.607026in}}%
\pgfpathlineto{\pgfqpoint{2.515149in}{3.627700in}}%
\pgfpathlineto{\pgfqpoint{2.517668in}{3.639278in}}%
\pgfpathlineto{\pgfqpoint{2.520187in}{3.666444in}}%
\pgfpathlineto{\pgfqpoint{2.522706in}{3.639341in}}%
\pgfpathlineto{\pgfqpoint{2.525224in}{3.592487in}}%
\pgfpathlineto{\pgfqpoint{2.527743in}{3.564982in}}%
\pgfpathlineto{\pgfqpoint{2.530262in}{3.571693in}}%
\pgfpathlineto{\pgfqpoint{2.532781in}{3.598999in}}%
\pgfpathlineto{\pgfqpoint{2.535299in}{3.597197in}}%
\pgfpathlineto{\pgfqpoint{2.537818in}{3.568212in}}%
\pgfpathlineto{\pgfqpoint{2.540337in}{3.607459in}}%
\pgfpathlineto{\pgfqpoint{2.542856in}{3.624723in}}%
\pgfpathlineto{\pgfqpoint{2.545374in}{3.610573in}}%
\pgfpathlineto{\pgfqpoint{2.547893in}{3.624294in}}%
\pgfpathlineto{\pgfqpoint{2.550412in}{3.618000in}}%
\pgfpathlineto{\pgfqpoint{2.552931in}{3.627823in}}%
\pgfpathlineto{\pgfqpoint{2.555449in}{3.623220in}}%
\pgfpathlineto{\pgfqpoint{2.557968in}{3.653351in}}%
\pgfpathlineto{\pgfqpoint{2.560487in}{3.635995in}}%
\pgfpathlineto{\pgfqpoint{2.563006in}{3.600501in}}%
\pgfpathlineto{\pgfqpoint{2.565524in}{3.583786in}}%
\pgfpathlineto{\pgfqpoint{2.568043in}{3.548667in}}%
\pgfpathlineto{\pgfqpoint{2.570562in}{3.587719in}}%
\pgfpathlineto{\pgfqpoint{2.573081in}{3.563790in}}%
\pgfpathlineto{\pgfqpoint{2.575599in}{3.544386in}}%
\pgfpathlineto{\pgfqpoint{2.578118in}{3.550798in}}%
\pgfpathlineto{\pgfqpoint{2.580637in}{3.600574in}}%
\pgfpathlineto{\pgfqpoint{2.583156in}{3.566571in}}%
\pgfpathlineto{\pgfqpoint{2.585674in}{3.570872in}}%
\pgfpathlineto{\pgfqpoint{2.588193in}{3.580203in}}%
\pgfpathlineto{\pgfqpoint{2.590712in}{3.562914in}}%
\pgfpathlineto{\pgfqpoint{2.593231in}{3.496752in}}%
\pgfpathlineto{\pgfqpoint{2.595749in}{3.470033in}}%
\pgfpathlineto{\pgfqpoint{2.598268in}{3.477496in}}%
\pgfpathlineto{\pgfqpoint{2.600787in}{3.464405in}}%
\pgfpathlineto{\pgfqpoint{2.603306in}{3.458489in}}%
\pgfpathlineto{\pgfqpoint{2.605824in}{3.472715in}}%
\pgfpathlineto{\pgfqpoint{2.608343in}{3.483222in}}%
\pgfpathlineto{\pgfqpoint{2.610862in}{3.461304in}}%
\pgfpathlineto{\pgfqpoint{2.613381in}{3.450154in}}%
\pgfpathlineto{\pgfqpoint{2.615899in}{3.434789in}}%
\pgfpathlineto{\pgfqpoint{2.618418in}{3.417530in}}%
\pgfpathlineto{\pgfqpoint{2.620937in}{3.415756in}}%
\pgfpathlineto{\pgfqpoint{2.623456in}{3.387990in}}%
\pgfpathlineto{\pgfqpoint{2.625974in}{3.382726in}}%
\pgfpathlineto{\pgfqpoint{2.628493in}{3.409795in}}%
\pgfpathlineto{\pgfqpoint{2.631012in}{3.443943in}}%
\pgfpathlineto{\pgfqpoint{2.633531in}{3.436397in}}%
\pgfpathlineto{\pgfqpoint{2.636049in}{3.415603in}}%
\pgfpathlineto{\pgfqpoint{2.638568in}{3.407035in}}%
\pgfpathlineto{\pgfqpoint{2.641087in}{3.381570in}}%
\pgfpathlineto{\pgfqpoint{2.643606in}{3.405976in}}%
\pgfpathlineto{\pgfqpoint{2.646124in}{3.434973in}}%
\pgfpathlineto{\pgfqpoint{2.648643in}{3.411908in}}%
\pgfpathlineto{\pgfqpoint{2.651162in}{3.450428in}}%
\pgfpathlineto{\pgfqpoint{2.653681in}{3.479297in}}%
\pgfpathlineto{\pgfqpoint{2.656199in}{3.498729in}}%
\pgfpathlineto{\pgfqpoint{2.658718in}{3.487289in}}%
\pgfpathlineto{\pgfqpoint{2.661237in}{3.482632in}}%
\pgfpathlineto{\pgfqpoint{2.663756in}{3.522922in}}%
\pgfpathlineto{\pgfqpoint{2.666274in}{3.556707in}}%
\pgfpathlineto{\pgfqpoint{2.668793in}{3.588867in}}%
\pgfpathlineto{\pgfqpoint{2.671312in}{3.592215in}}%
\pgfpathlineto{\pgfqpoint{2.673831in}{3.598723in}}%
\pgfpathlineto{\pgfqpoint{2.676349in}{3.586714in}}%
\pgfpathlineto{\pgfqpoint{2.678868in}{3.562032in}}%
\pgfpathlineto{\pgfqpoint{2.683906in}{3.530468in}}%
\pgfpathlineto{\pgfqpoint{2.686424in}{3.562010in}}%
\pgfpathlineto{\pgfqpoint{2.688943in}{3.567495in}}%
\pgfpathlineto{\pgfqpoint{2.691462in}{3.532053in}}%
\pgfpathlineto{\pgfqpoint{2.693981in}{3.520698in}}%
\pgfpathlineto{\pgfqpoint{2.696499in}{3.485845in}}%
\pgfpathlineto{\pgfqpoint{2.699018in}{3.499062in}}%
\pgfpathlineto{\pgfqpoint{2.701537in}{3.518572in}}%
\pgfpathlineto{\pgfqpoint{2.704056in}{3.537250in}}%
\pgfpathlineto{\pgfqpoint{2.706574in}{3.508899in}}%
\pgfpathlineto{\pgfqpoint{2.709093in}{3.483224in}}%
\pgfpathlineto{\pgfqpoint{2.711612in}{3.465418in}}%
\pgfpathlineto{\pgfqpoint{2.714131in}{3.430058in}}%
\pgfpathlineto{\pgfqpoint{2.716649in}{3.396409in}}%
\pgfpathlineto{\pgfqpoint{2.719168in}{3.371249in}}%
\pgfpathlineto{\pgfqpoint{2.721687in}{3.400290in}}%
\pgfpathlineto{\pgfqpoint{2.724206in}{3.386741in}}%
\pgfpathlineto{\pgfqpoint{2.726724in}{3.430954in}}%
\pgfpathlineto{\pgfqpoint{2.729243in}{3.414546in}}%
\pgfpathlineto{\pgfqpoint{2.731762in}{3.409344in}}%
\pgfpathlineto{\pgfqpoint{2.734281in}{3.362330in}}%
\pgfpathlineto{\pgfqpoint{2.736799in}{3.310823in}}%
\pgfpathlineto{\pgfqpoint{2.739318in}{3.332652in}}%
\pgfpathlineto{\pgfqpoint{2.741837in}{3.377962in}}%
\pgfpathlineto{\pgfqpoint{2.744356in}{3.388862in}}%
\pgfpathlineto{\pgfqpoint{2.746874in}{3.386976in}}%
\pgfpathlineto{\pgfqpoint{2.749393in}{3.399684in}}%
\pgfpathlineto{\pgfqpoint{2.751912in}{3.342949in}}%
\pgfpathlineto{\pgfqpoint{2.754431in}{3.371544in}}%
\pgfpathlineto{\pgfqpoint{2.756949in}{3.369438in}}%
\pgfpathlineto{\pgfqpoint{2.759468in}{3.360818in}}%
\pgfpathlineto{\pgfqpoint{2.761987in}{3.380963in}}%
\pgfpathlineto{\pgfqpoint{2.764506in}{3.443493in}}%
\pgfpathlineto{\pgfqpoint{2.767024in}{3.416419in}}%
\pgfpathlineto{\pgfqpoint{2.769543in}{3.404359in}}%
\pgfpathlineto{\pgfqpoint{2.772062in}{3.390782in}}%
\pgfpathlineto{\pgfqpoint{2.774581in}{3.452625in}}%
\pgfpathlineto{\pgfqpoint{2.777099in}{3.500390in}}%
\pgfpathlineto{\pgfqpoint{2.779618in}{3.469847in}}%
\pgfpathlineto{\pgfqpoint{2.782137in}{3.459617in}}%
\pgfpathlineto{\pgfqpoint{2.784656in}{3.478136in}}%
\pgfpathlineto{\pgfqpoint{2.787174in}{3.465612in}}%
\pgfpathlineto{\pgfqpoint{2.789693in}{3.460488in}}%
\pgfpathlineto{\pgfqpoint{2.792212in}{3.500551in}}%
\pgfpathlineto{\pgfqpoint{2.794731in}{3.488840in}}%
\pgfpathlineto{\pgfqpoint{2.797249in}{3.503151in}}%
\pgfpathlineto{\pgfqpoint{2.799768in}{3.469892in}}%
\pgfpathlineto{\pgfqpoint{2.802287in}{3.482669in}}%
\pgfpathlineto{\pgfqpoint{2.804806in}{3.481368in}}%
\pgfpathlineto{\pgfqpoint{2.807324in}{3.495152in}}%
\pgfpathlineto{\pgfqpoint{2.809843in}{3.503162in}}%
\pgfpathlineto{\pgfqpoint{2.812362in}{3.533042in}}%
\pgfpathlineto{\pgfqpoint{2.814881in}{3.476748in}}%
\pgfpathlineto{\pgfqpoint{2.817399in}{3.478895in}}%
\pgfpathlineto{\pgfqpoint{2.819918in}{3.472922in}}%
\pgfpathlineto{\pgfqpoint{2.822437in}{3.478811in}}%
\pgfpathlineto{\pgfqpoint{2.824956in}{3.471818in}}%
\pgfpathlineto{\pgfqpoint{2.827474in}{3.433420in}}%
\pgfpathlineto{\pgfqpoint{2.829993in}{3.399303in}}%
\pgfpathlineto{\pgfqpoint{2.832512in}{3.370121in}}%
\pgfpathlineto{\pgfqpoint{2.835031in}{3.420789in}}%
\pgfpathlineto{\pgfqpoint{2.837549in}{3.419297in}}%
\pgfpathlineto{\pgfqpoint{2.840068in}{3.430530in}}%
\pgfpathlineto{\pgfqpoint{2.842587in}{3.457819in}}%
\pgfpathlineto{\pgfqpoint{2.845106in}{3.470740in}}%
\pgfpathlineto{\pgfqpoint{2.847624in}{3.486943in}}%
\pgfpathlineto{\pgfqpoint{2.850143in}{3.460423in}}%
\pgfpathlineto{\pgfqpoint{2.852662in}{3.425043in}}%
\pgfpathlineto{\pgfqpoint{2.855181in}{3.406846in}}%
\pgfpathlineto{\pgfqpoint{2.857699in}{3.373991in}}%
\pgfpathlineto{\pgfqpoint{2.860218in}{3.362791in}}%
\pgfpathlineto{\pgfqpoint{2.862737in}{3.377115in}}%
\pgfpathlineto{\pgfqpoint{2.865256in}{3.399204in}}%
\pgfpathlineto{\pgfqpoint{2.867774in}{3.425715in}}%
\pgfpathlineto{\pgfqpoint{2.870293in}{3.427473in}}%
\pgfpathlineto{\pgfqpoint{2.872812in}{3.356947in}}%
\pgfpathlineto{\pgfqpoint{2.875331in}{3.399963in}}%
\pgfpathlineto{\pgfqpoint{2.877849in}{3.393005in}}%
\pgfpathlineto{\pgfqpoint{2.880368in}{3.350662in}}%
\pgfpathlineto{\pgfqpoint{2.882887in}{3.318480in}}%
\pgfpathlineto{\pgfqpoint{2.885406in}{3.366686in}}%
\pgfpathlineto{\pgfqpoint{2.887924in}{3.401437in}}%
\pgfpathlineto{\pgfqpoint{2.890443in}{3.369673in}}%
\pgfpathlineto{\pgfqpoint{2.892962in}{3.415831in}}%
\pgfpathlineto{\pgfqpoint{2.895481in}{3.405844in}}%
\pgfpathlineto{\pgfqpoint{2.897999in}{3.409506in}}%
\pgfpathlineto{\pgfqpoint{2.900518in}{3.394287in}}%
\pgfpathlineto{\pgfqpoint{2.903037in}{3.381580in}}%
\pgfpathlineto{\pgfqpoint{2.905556in}{3.380618in}}%
\pgfpathlineto{\pgfqpoint{2.908074in}{3.328531in}}%
\pgfpathlineto{\pgfqpoint{2.910593in}{3.302935in}}%
\pgfpathlineto{\pgfqpoint{2.913112in}{3.313344in}}%
\pgfpathlineto{\pgfqpoint{2.915631in}{3.327905in}}%
\pgfpathlineto{\pgfqpoint{2.918149in}{3.307939in}}%
\pgfpathlineto{\pgfqpoint{2.920668in}{3.309460in}}%
\pgfpathlineto{\pgfqpoint{2.923187in}{3.295772in}}%
\pgfpathlineto{\pgfqpoint{2.925706in}{3.303981in}}%
\pgfpathlineto{\pgfqpoint{2.928224in}{3.246329in}}%
\pgfpathlineto{\pgfqpoint{2.930743in}{3.268370in}}%
\pgfpathlineto{\pgfqpoint{2.933262in}{3.247024in}}%
\pgfpathlineto{\pgfqpoint{2.935781in}{3.262897in}}%
\pgfpathlineto{\pgfqpoint{2.938299in}{3.220601in}}%
\pgfpathlineto{\pgfqpoint{2.940818in}{3.190246in}}%
\pgfpathlineto{\pgfqpoint{2.943337in}{3.164798in}}%
\pgfpathlineto{\pgfqpoint{2.945856in}{3.166869in}}%
\pgfpathlineto{\pgfqpoint{2.948374in}{3.176495in}}%
\pgfpathlineto{\pgfqpoint{2.950893in}{3.124583in}}%
\pgfpathlineto{\pgfqpoint{2.953412in}{3.120193in}}%
\pgfpathlineto{\pgfqpoint{2.955931in}{3.089473in}}%
\pgfpathlineto{\pgfqpoint{2.958449in}{3.070447in}}%
\pgfpathlineto{\pgfqpoint{2.960968in}{3.124536in}}%
\pgfpathlineto{\pgfqpoint{2.963487in}{3.114580in}}%
\pgfpathlineto{\pgfqpoint{2.966006in}{3.061540in}}%
\pgfpathlineto{\pgfqpoint{2.968524in}{3.080372in}}%
\pgfpathlineto{\pgfqpoint{2.971043in}{3.102872in}}%
\pgfpathlineto{\pgfqpoint{2.973562in}{3.127632in}}%
\pgfpathlineto{\pgfqpoint{2.976081in}{3.082514in}}%
\pgfpathlineto{\pgfqpoint{2.978599in}{3.070729in}}%
\pgfpathlineto{\pgfqpoint{2.981118in}{3.064470in}}%
\pgfpathlineto{\pgfqpoint{2.983637in}{3.116129in}}%
\pgfpathlineto{\pgfqpoint{2.986156in}{3.100490in}}%
\pgfpathlineto{\pgfqpoint{2.988674in}{3.097785in}}%
\pgfpathlineto{\pgfqpoint{2.991193in}{3.120477in}}%
\pgfpathlineto{\pgfqpoint{2.993712in}{3.073830in}}%
\pgfpathlineto{\pgfqpoint{2.996231in}{3.108003in}}%
\pgfpathlineto{\pgfqpoint{2.998749in}{3.084031in}}%
\pgfpathlineto{\pgfqpoint{3.001268in}{3.101191in}}%
\pgfpathlineto{\pgfqpoint{3.003787in}{3.083899in}}%
\pgfpathlineto{\pgfqpoint{3.006306in}{3.115350in}}%
\pgfpathlineto{\pgfqpoint{3.008824in}{3.137984in}}%
\pgfpathlineto{\pgfqpoint{3.011343in}{3.177924in}}%
\pgfpathlineto{\pgfqpoint{3.013862in}{3.174935in}}%
\pgfpathlineto{\pgfqpoint{3.016381in}{3.159559in}}%
\pgfpathlineto{\pgfqpoint{3.018899in}{3.192185in}}%
\pgfpathlineto{\pgfqpoint{3.021418in}{3.202671in}}%
\pgfpathlineto{\pgfqpoint{3.023937in}{3.203171in}}%
\pgfpathlineto{\pgfqpoint{3.026456in}{3.228660in}}%
\pgfpathlineto{\pgfqpoint{3.028974in}{3.198996in}}%
\pgfpathlineto{\pgfqpoint{3.031493in}{3.172079in}}%
\pgfpathlineto{\pgfqpoint{3.034012in}{3.147517in}}%
\pgfpathlineto{\pgfqpoint{3.036531in}{3.117644in}}%
\pgfpathlineto{\pgfqpoint{3.041568in}{3.093612in}}%
\pgfpathlineto{\pgfqpoint{3.044087in}{3.109307in}}%
\pgfpathlineto{\pgfqpoint{3.046606in}{3.095625in}}%
\pgfpathlineto{\pgfqpoint{3.049124in}{3.087666in}}%
\pgfpathlineto{\pgfqpoint{3.051643in}{3.114417in}}%
\pgfpathlineto{\pgfqpoint{3.054162in}{3.100015in}}%
\pgfpathlineto{\pgfqpoint{3.056681in}{3.052490in}}%
\pgfpathlineto{\pgfqpoint{3.059199in}{3.063474in}}%
\pgfpathlineto{\pgfqpoint{3.061718in}{3.059228in}}%
\pgfpathlineto{\pgfqpoint{3.064237in}{3.017805in}}%
\pgfpathlineto{\pgfqpoint{3.066756in}{3.021039in}}%
\pgfpathlineto{\pgfqpoint{3.071793in}{3.093253in}}%
\pgfpathlineto{\pgfqpoint{3.074312in}{3.098502in}}%
\pgfpathlineto{\pgfqpoint{3.076831in}{3.080759in}}%
\pgfpathlineto{\pgfqpoint{3.081868in}{3.140044in}}%
\pgfpathlineto{\pgfqpoint{3.084387in}{3.144317in}}%
\pgfpathlineto{\pgfqpoint{3.086906in}{3.146726in}}%
\pgfpathlineto{\pgfqpoint{3.089424in}{3.180821in}}%
\pgfpathlineto{\pgfqpoint{3.091943in}{3.193141in}}%
\pgfpathlineto{\pgfqpoint{3.094462in}{3.199494in}}%
\pgfpathlineto{\pgfqpoint{3.096981in}{3.188086in}}%
\pgfpathlineto{\pgfqpoint{3.099499in}{3.169523in}}%
\pgfpathlineto{\pgfqpoint{3.102018in}{3.142804in}}%
\pgfpathlineto{\pgfqpoint{3.104537in}{3.164080in}}%
\pgfpathlineto{\pgfqpoint{3.107056in}{3.173893in}}%
\pgfpathlineto{\pgfqpoint{3.109574in}{3.176263in}}%
\pgfpathlineto{\pgfqpoint{3.112093in}{3.198097in}}%
\pgfpathlineto{\pgfqpoint{3.114612in}{3.217880in}}%
\pgfpathlineto{\pgfqpoint{3.117131in}{3.245920in}}%
\pgfpathlineto{\pgfqpoint{3.119649in}{3.215274in}}%
\pgfpathlineto{\pgfqpoint{3.122168in}{3.212539in}}%
\pgfpathlineto{\pgfqpoint{3.124687in}{3.281282in}}%
\pgfpathlineto{\pgfqpoint{3.127206in}{3.305910in}}%
\pgfpathlineto{\pgfqpoint{3.129724in}{3.338227in}}%
\pgfpathlineto{\pgfqpoint{3.132243in}{3.340773in}}%
\pgfpathlineto{\pgfqpoint{3.134762in}{3.317880in}}%
\pgfpathlineto{\pgfqpoint{3.137281in}{3.326750in}}%
\pgfpathlineto{\pgfqpoint{3.139799in}{3.339373in}}%
\pgfpathlineto{\pgfqpoint{3.142318in}{3.326840in}}%
\pgfpathlineto{\pgfqpoint{3.144837in}{3.302765in}}%
\pgfpathlineto{\pgfqpoint{3.147356in}{3.334576in}}%
\pgfpathlineto{\pgfqpoint{3.149874in}{3.306940in}}%
\pgfpathlineto{\pgfqpoint{3.152393in}{3.316202in}}%
\pgfpathlineto{\pgfqpoint{3.154912in}{3.274974in}}%
\pgfpathlineto{\pgfqpoint{3.157431in}{3.289220in}}%
\pgfpathlineto{\pgfqpoint{3.159949in}{3.263765in}}%
\pgfpathlineto{\pgfqpoint{3.162468in}{3.234903in}}%
\pgfpathlineto{\pgfqpoint{3.164987in}{3.277708in}}%
\pgfpathlineto{\pgfqpoint{3.167506in}{3.266655in}}%
\pgfpathlineto{\pgfqpoint{3.170024in}{3.260622in}}%
\pgfpathlineto{\pgfqpoint{3.172543in}{3.226888in}}%
\pgfpathlineto{\pgfqpoint{3.175062in}{3.259514in}}%
\pgfpathlineto{\pgfqpoint{3.177581in}{3.241195in}}%
\pgfpathlineto{\pgfqpoint{3.180099in}{3.230814in}}%
\pgfpathlineto{\pgfqpoint{3.182618in}{3.238304in}}%
\pgfpathlineto{\pgfqpoint{3.185137in}{3.223272in}}%
\pgfpathlineto{\pgfqpoint{3.187656in}{3.205125in}}%
\pgfpathlineto{\pgfqpoint{3.190174in}{3.204216in}}%
\pgfpathlineto{\pgfqpoint{3.192693in}{3.170770in}}%
\pgfpathlineto{\pgfqpoint{3.195212in}{3.187087in}}%
\pgfpathlineto{\pgfqpoint{3.197731in}{3.190860in}}%
\pgfpathlineto{\pgfqpoint{3.200249in}{3.146521in}}%
\pgfpathlineto{\pgfqpoint{3.202768in}{3.139193in}}%
\pgfpathlineto{\pgfqpoint{3.205287in}{3.166635in}}%
\pgfpathlineto{\pgfqpoint{3.207806in}{3.187627in}}%
\pgfpathlineto{\pgfqpoint{3.212843in}{3.256959in}}%
\pgfpathlineto{\pgfqpoint{3.215362in}{3.211815in}}%
\pgfpathlineto{\pgfqpoint{3.217881in}{3.248909in}}%
\pgfpathlineto{\pgfqpoint{3.220399in}{3.205841in}}%
\pgfpathlineto{\pgfqpoint{3.222918in}{3.247024in}}%
\pgfpathlineto{\pgfqpoint{3.225437in}{3.176220in}}%
\pgfpathlineto{\pgfqpoint{3.227956in}{3.191329in}}%
\pgfpathlineto{\pgfqpoint{3.230474in}{3.193113in}}%
\pgfpathlineto{\pgfqpoint{3.232993in}{3.192860in}}%
\pgfpathlineto{\pgfqpoint{3.235512in}{3.205250in}}%
\pgfpathlineto{\pgfqpoint{3.238031in}{3.201234in}}%
\pgfpathlineto{\pgfqpoint{3.240549in}{3.203319in}}%
\pgfpathlineto{\pgfqpoint{3.243068in}{3.185255in}}%
\pgfpathlineto{\pgfqpoint{3.245587in}{3.205393in}}%
\pgfpathlineto{\pgfqpoint{3.248106in}{3.170526in}}%
\pgfpathlineto{\pgfqpoint{3.250624in}{3.156561in}}%
\pgfpathlineto{\pgfqpoint{3.255662in}{3.267730in}}%
\pgfpathlineto{\pgfqpoint{3.258181in}{3.255517in}}%
\pgfpathlineto{\pgfqpoint{3.260699in}{3.310925in}}%
\pgfpathlineto{\pgfqpoint{3.263218in}{3.223138in}}%
\pgfpathlineto{\pgfqpoint{3.265737in}{3.150627in}}%
\pgfpathlineto{\pgfqpoint{3.268256in}{3.154256in}}%
\pgfpathlineto{\pgfqpoint{3.270774in}{3.152441in}}%
\pgfpathlineto{\pgfqpoint{3.273293in}{3.164240in}}%
\pgfpathlineto{\pgfqpoint{3.275812in}{3.122322in}}%
\pgfpathlineto{\pgfqpoint{3.278331in}{3.142842in}}%
\pgfpathlineto{\pgfqpoint{3.280849in}{3.137959in}}%
\pgfpathlineto{\pgfqpoint{3.285887in}{3.168600in}}%
\pgfpathlineto{\pgfqpoint{3.288406in}{3.128874in}}%
\pgfpathlineto{\pgfqpoint{3.290924in}{3.116698in}}%
\pgfpathlineto{\pgfqpoint{3.293443in}{3.127503in}}%
\pgfpathlineto{\pgfqpoint{3.295962in}{3.092048in}}%
\pgfpathlineto{\pgfqpoint{3.298481in}{3.109595in}}%
\pgfpathlineto{\pgfqpoint{3.300999in}{3.084794in}}%
\pgfpathlineto{\pgfqpoint{3.303518in}{3.080887in}}%
\pgfpathlineto{\pgfqpoint{3.306037in}{3.061444in}}%
\pgfpathlineto{\pgfqpoint{3.308556in}{3.050956in}}%
\pgfpathlineto{\pgfqpoint{3.311074in}{3.044749in}}%
\pgfpathlineto{\pgfqpoint{3.313593in}{3.034883in}}%
\pgfpathlineto{\pgfqpoint{3.316112in}{3.024492in}}%
\pgfpathlineto{\pgfqpoint{3.318631in}{3.037693in}}%
\pgfpathlineto{\pgfqpoint{3.321149in}{3.073893in}}%
\pgfpathlineto{\pgfqpoint{3.323668in}{3.044052in}}%
\pgfpathlineto{\pgfqpoint{3.326187in}{3.086958in}}%
\pgfpathlineto{\pgfqpoint{3.328706in}{3.055646in}}%
\pgfpathlineto{\pgfqpoint{3.331224in}{3.086086in}}%
\pgfpathlineto{\pgfqpoint{3.333743in}{3.031561in}}%
\pgfpathlineto{\pgfqpoint{3.336262in}{3.038288in}}%
\pgfpathlineto{\pgfqpoint{3.338781in}{3.052536in}}%
\pgfpathlineto{\pgfqpoint{3.341299in}{3.077816in}}%
\pgfpathlineto{\pgfqpoint{3.343818in}{3.044625in}}%
\pgfpathlineto{\pgfqpoint{3.346337in}{3.067565in}}%
\pgfpathlineto{\pgfqpoint{3.348856in}{3.081696in}}%
\pgfpathlineto{\pgfqpoint{3.351374in}{3.084256in}}%
\pgfpathlineto{\pgfqpoint{3.353893in}{3.058052in}}%
\pgfpathlineto{\pgfqpoint{3.356412in}{3.077276in}}%
\pgfpathlineto{\pgfqpoint{3.358931in}{3.083845in}}%
\pgfpathlineto{\pgfqpoint{3.361449in}{3.069393in}}%
\pgfpathlineto{\pgfqpoint{3.363968in}{3.075925in}}%
\pgfpathlineto{\pgfqpoint{3.366487in}{3.069413in}}%
\pgfpathlineto{\pgfqpoint{3.369006in}{3.047948in}}%
\pgfpathlineto{\pgfqpoint{3.371524in}{3.048199in}}%
\pgfpathlineto{\pgfqpoint{3.376562in}{3.008709in}}%
\pgfpathlineto{\pgfqpoint{3.379081in}{3.024617in}}%
\pgfpathlineto{\pgfqpoint{3.381599in}{3.016259in}}%
\pgfpathlineto{\pgfqpoint{3.384118in}{3.023469in}}%
\pgfpathlineto{\pgfqpoint{3.386637in}{3.060195in}}%
\pgfpathlineto{\pgfqpoint{3.389156in}{3.046959in}}%
\pgfpathlineto{\pgfqpoint{3.391674in}{3.044123in}}%
\pgfpathlineto{\pgfqpoint{3.394193in}{3.106061in}}%
\pgfpathlineto{\pgfqpoint{3.396712in}{3.094979in}}%
\pgfpathlineto{\pgfqpoint{3.399231in}{3.124950in}}%
\pgfpathlineto{\pgfqpoint{3.401749in}{3.064311in}}%
\pgfpathlineto{\pgfqpoint{3.404268in}{3.095613in}}%
\pgfpathlineto{\pgfqpoint{3.406787in}{3.100057in}}%
\pgfpathlineto{\pgfqpoint{3.409306in}{3.088622in}}%
\pgfpathlineto{\pgfqpoint{3.411824in}{3.116413in}}%
\pgfpathlineto{\pgfqpoint{3.414343in}{3.048129in}}%
\pgfpathlineto{\pgfqpoint{3.416862in}{3.013026in}}%
\pgfpathlineto{\pgfqpoint{3.419381in}{2.956014in}}%
\pgfpathlineto{\pgfqpoint{3.421899in}{2.996922in}}%
\pgfpathlineto{\pgfqpoint{3.424418in}{2.962286in}}%
\pgfpathlineto{\pgfqpoint{3.426937in}{2.950505in}}%
\pgfpathlineto{\pgfqpoint{3.429456in}{2.887689in}}%
\pgfpathlineto{\pgfqpoint{3.431974in}{2.894212in}}%
\pgfpathlineto{\pgfqpoint{3.434493in}{2.848709in}}%
\pgfpathlineto{\pgfqpoint{3.437012in}{2.856836in}}%
\pgfpathlineto{\pgfqpoint{3.439531in}{2.846766in}}%
\pgfpathlineto{\pgfqpoint{3.442049in}{2.858890in}}%
\pgfpathlineto{\pgfqpoint{3.447087in}{2.893878in}}%
\pgfpathlineto{\pgfqpoint{3.449606in}{2.920984in}}%
\pgfpathlineto{\pgfqpoint{3.452124in}{2.897246in}}%
\pgfpathlineto{\pgfqpoint{3.454643in}{2.901318in}}%
\pgfpathlineto{\pgfqpoint{3.457162in}{2.906581in}}%
\pgfpathlineto{\pgfqpoint{3.459681in}{2.875818in}}%
\pgfpathlineto{\pgfqpoint{3.462199in}{2.853246in}}%
\pgfpathlineto{\pgfqpoint{3.464718in}{2.856489in}}%
\pgfpathlineto{\pgfqpoint{3.467237in}{2.835765in}}%
\pgfpathlineto{\pgfqpoint{3.469756in}{2.849173in}}%
\pgfpathlineto{\pgfqpoint{3.472274in}{2.890094in}}%
\pgfpathlineto{\pgfqpoint{3.474793in}{2.872247in}}%
\pgfpathlineto{\pgfqpoint{3.477312in}{2.847965in}}%
\pgfpathlineto{\pgfqpoint{3.479831in}{2.820999in}}%
\pgfpathlineto{\pgfqpoint{3.482349in}{2.843446in}}%
\pgfpathlineto{\pgfqpoint{3.484868in}{2.809270in}}%
\pgfpathlineto{\pgfqpoint{3.489906in}{2.744462in}}%
\pgfpathlineto{\pgfqpoint{3.492424in}{2.747404in}}%
\pgfpathlineto{\pgfqpoint{3.494943in}{2.688880in}}%
\pgfpathlineto{\pgfqpoint{3.497462in}{2.684273in}}%
\pgfpathlineto{\pgfqpoint{3.499981in}{2.615487in}}%
\pgfpathlineto{\pgfqpoint{3.502499in}{2.598371in}}%
\pgfpathlineto{\pgfqpoint{3.505018in}{2.609217in}}%
\pgfpathlineto{\pgfqpoint{3.507537in}{2.593184in}}%
\pgfpathlineto{\pgfqpoint{3.510056in}{2.633953in}}%
\pgfpathlineto{\pgfqpoint{3.512574in}{2.585066in}}%
\pgfpathlineto{\pgfqpoint{3.515093in}{2.619581in}}%
\pgfpathlineto{\pgfqpoint{3.517612in}{2.623760in}}%
\pgfpathlineto{\pgfqpoint{3.520131in}{2.622027in}}%
\pgfpathlineto{\pgfqpoint{3.522649in}{2.660439in}}%
\pgfpathlineto{\pgfqpoint{3.525168in}{2.724651in}}%
\pgfpathlineto{\pgfqpoint{3.527687in}{2.699284in}}%
\pgfpathlineto{\pgfqpoint{3.530206in}{2.714007in}}%
\pgfpathlineto{\pgfqpoint{3.532724in}{2.686723in}}%
\pgfpathlineto{\pgfqpoint{3.535243in}{2.684987in}}%
\pgfpathlineto{\pgfqpoint{3.537762in}{2.651360in}}%
\pgfpathlineto{\pgfqpoint{3.540281in}{2.684799in}}%
\pgfpathlineto{\pgfqpoint{3.542799in}{2.673797in}}%
\pgfpathlineto{\pgfqpoint{3.545318in}{2.697675in}}%
\pgfpathlineto{\pgfqpoint{3.547837in}{2.633571in}}%
\pgfpathlineto{\pgfqpoint{3.550356in}{2.613550in}}%
\pgfpathlineto{\pgfqpoint{3.552874in}{2.570585in}}%
\pgfpathlineto{\pgfqpoint{3.555393in}{2.582260in}}%
\pgfpathlineto{\pgfqpoint{3.557912in}{2.615708in}}%
\pgfpathlineto{\pgfqpoint{3.560431in}{2.667106in}}%
\pgfpathlineto{\pgfqpoint{3.562949in}{2.651533in}}%
\pgfpathlineto{\pgfqpoint{3.565468in}{2.649461in}}%
\pgfpathlineto{\pgfqpoint{3.567987in}{2.663245in}}%
\pgfpathlineto{\pgfqpoint{3.570506in}{2.669130in}}%
\pgfpathlineto{\pgfqpoint{3.573024in}{2.692452in}}%
\pgfpathlineto{\pgfqpoint{3.575543in}{2.660178in}}%
\pgfpathlineto{\pgfqpoint{3.578062in}{2.634892in}}%
\pgfpathlineto{\pgfqpoint{3.580581in}{2.630991in}}%
\pgfpathlineto{\pgfqpoint{3.583099in}{2.667067in}}%
\pgfpathlineto{\pgfqpoint{3.585618in}{2.663408in}}%
\pgfpathlineto{\pgfqpoint{3.588137in}{2.703341in}}%
\pgfpathlineto{\pgfqpoint{3.590656in}{2.682130in}}%
\pgfpathlineto{\pgfqpoint{3.593174in}{2.678650in}}%
\pgfpathlineto{\pgfqpoint{3.595693in}{2.622410in}}%
\pgfpathlineto{\pgfqpoint{3.598212in}{2.640910in}}%
\pgfpathlineto{\pgfqpoint{3.600731in}{2.660277in}}%
\pgfpathlineto{\pgfqpoint{3.603249in}{2.614127in}}%
\pgfpathlineto{\pgfqpoint{3.605768in}{2.651574in}}%
\pgfpathlineto{\pgfqpoint{3.608287in}{2.647621in}}%
\pgfpathlineto{\pgfqpoint{3.610806in}{2.697333in}}%
\pgfpathlineto{\pgfqpoint{3.613324in}{2.730197in}}%
\pgfpathlineto{\pgfqpoint{3.615843in}{2.696607in}}%
\pgfpathlineto{\pgfqpoint{3.618362in}{2.714777in}}%
\pgfpathlineto{\pgfqpoint{3.620881in}{2.753394in}}%
\pgfpathlineto{\pgfqpoint{3.623399in}{2.768099in}}%
\pgfpathlineto{\pgfqpoint{3.625918in}{2.787349in}}%
\pgfpathlineto{\pgfqpoint{3.628437in}{2.766168in}}%
\pgfpathlineto{\pgfqpoint{3.630956in}{2.771479in}}%
\pgfpathlineto{\pgfqpoint{3.633474in}{2.709032in}}%
\pgfpathlineto{\pgfqpoint{3.635993in}{2.683663in}}%
\pgfpathlineto{\pgfqpoint{3.638512in}{2.684400in}}%
\pgfpathlineto{\pgfqpoint{3.641031in}{2.719102in}}%
\pgfpathlineto{\pgfqpoint{3.643549in}{2.719180in}}%
\pgfpathlineto{\pgfqpoint{3.646068in}{2.712227in}}%
\pgfpathlineto{\pgfqpoint{3.648587in}{2.715060in}}%
\pgfpathlineto{\pgfqpoint{3.651106in}{2.761401in}}%
\pgfpathlineto{\pgfqpoint{3.653624in}{2.740987in}}%
\pgfpathlineto{\pgfqpoint{3.656143in}{2.773715in}}%
\pgfpathlineto{\pgfqpoint{3.658662in}{2.796887in}}%
\pgfpathlineto{\pgfqpoint{3.661181in}{2.799583in}}%
\pgfpathlineto{\pgfqpoint{3.663699in}{2.802121in}}%
\pgfpathlineto{\pgfqpoint{3.666218in}{2.744848in}}%
\pgfpathlineto{\pgfqpoint{3.668737in}{2.773847in}}%
\pgfpathlineto{\pgfqpoint{3.671256in}{2.756708in}}%
\pgfpathlineto{\pgfqpoint{3.673774in}{2.737671in}}%
\pgfpathlineto{\pgfqpoint{3.678812in}{2.820191in}}%
\pgfpathlineto{\pgfqpoint{3.681331in}{2.835349in}}%
\pgfpathlineto{\pgfqpoint{3.683849in}{2.846536in}}%
\pgfpathlineto{\pgfqpoint{3.686368in}{2.825915in}}%
\pgfpathlineto{\pgfqpoint{3.688887in}{2.837305in}}%
\pgfpathlineto{\pgfqpoint{3.691406in}{2.841842in}}%
\pgfpathlineto{\pgfqpoint{3.693924in}{2.849869in}}%
\pgfpathlineto{\pgfqpoint{3.696443in}{2.820034in}}%
\pgfpathlineto{\pgfqpoint{3.698962in}{2.830443in}}%
\pgfpathlineto{\pgfqpoint{3.701481in}{2.824535in}}%
\pgfpathlineto{\pgfqpoint{3.703999in}{2.788934in}}%
\pgfpathlineto{\pgfqpoint{3.706518in}{2.806532in}}%
\pgfpathlineto{\pgfqpoint{3.709037in}{2.833541in}}%
\pgfpathlineto{\pgfqpoint{3.711556in}{2.815035in}}%
\pgfpathlineto{\pgfqpoint{3.714074in}{2.807889in}}%
\pgfpathlineto{\pgfqpoint{3.716593in}{2.859742in}}%
\pgfpathlineto{\pgfqpoint{3.719112in}{2.896433in}}%
\pgfpathlineto{\pgfqpoint{3.721631in}{2.911926in}}%
\pgfpathlineto{\pgfqpoint{3.724149in}{2.896963in}}%
\pgfpathlineto{\pgfqpoint{3.726668in}{2.918014in}}%
\pgfpathlineto{\pgfqpoint{3.729187in}{2.905172in}}%
\pgfpathlineto{\pgfqpoint{3.731706in}{2.879463in}}%
\pgfpathlineto{\pgfqpoint{3.734224in}{2.892058in}}%
\pgfpathlineto{\pgfqpoint{3.736743in}{2.926957in}}%
\pgfpathlineto{\pgfqpoint{3.739262in}{2.916286in}}%
\pgfpathlineto{\pgfqpoint{3.741781in}{2.914994in}}%
\pgfpathlineto{\pgfqpoint{3.744299in}{2.933129in}}%
\pgfpathlineto{\pgfqpoint{3.749337in}{2.845924in}}%
\pgfpathlineto{\pgfqpoint{3.751856in}{2.836638in}}%
\pgfpathlineto{\pgfqpoint{3.754374in}{2.828710in}}%
\pgfpathlineto{\pgfqpoint{3.756893in}{2.826847in}}%
\pgfpathlineto{\pgfqpoint{3.759412in}{2.860784in}}%
\pgfpathlineto{\pgfqpoint{3.761931in}{2.851678in}}%
\pgfpathlineto{\pgfqpoint{3.764449in}{2.852651in}}%
\pgfpathlineto{\pgfqpoint{3.766968in}{2.837330in}}%
\pgfpathlineto{\pgfqpoint{3.769487in}{2.858608in}}%
\pgfpathlineto{\pgfqpoint{3.772006in}{2.895890in}}%
\pgfpathlineto{\pgfqpoint{3.774524in}{2.910109in}}%
\pgfpathlineto{\pgfqpoint{3.777043in}{2.887971in}}%
\pgfpathlineto{\pgfqpoint{3.779562in}{2.884344in}}%
\pgfpathlineto{\pgfqpoint{3.782081in}{2.836745in}}%
\pgfpathlineto{\pgfqpoint{3.784599in}{2.829242in}}%
\pgfpathlineto{\pgfqpoint{3.787118in}{2.851494in}}%
\pgfpathlineto{\pgfqpoint{3.789637in}{2.827372in}}%
\pgfpathlineto{\pgfqpoint{3.792156in}{2.837263in}}%
\pgfpathlineto{\pgfqpoint{3.794674in}{2.792460in}}%
\pgfpathlineto{\pgfqpoint{3.797193in}{2.793638in}}%
\pgfpathlineto{\pgfqpoint{3.799712in}{2.777371in}}%
\pgfpathlineto{\pgfqpoint{3.802231in}{2.755102in}}%
\pgfpathlineto{\pgfqpoint{3.804749in}{2.761292in}}%
\pgfpathlineto{\pgfqpoint{3.807268in}{2.791749in}}%
\pgfpathlineto{\pgfqpoint{3.809787in}{2.851258in}}%
\pgfpathlineto{\pgfqpoint{3.812306in}{2.904117in}}%
\pgfpathlineto{\pgfqpoint{3.814824in}{2.898960in}}%
\pgfpathlineto{\pgfqpoint{3.817343in}{2.931256in}}%
\pgfpathlineto{\pgfqpoint{3.819862in}{2.934198in}}%
\pgfpathlineto{\pgfqpoint{3.822381in}{2.938842in}}%
\pgfpathlineto{\pgfqpoint{3.824899in}{2.991266in}}%
\pgfpathlineto{\pgfqpoint{3.827418in}{2.998189in}}%
\pgfpathlineto{\pgfqpoint{3.829937in}{2.986938in}}%
\pgfpathlineto{\pgfqpoint{3.832456in}{3.026047in}}%
\pgfpathlineto{\pgfqpoint{3.834974in}{3.036745in}}%
\pgfpathlineto{\pgfqpoint{3.837493in}{3.078782in}}%
\pgfpathlineto{\pgfqpoint{3.840012in}{3.013593in}}%
\pgfpathlineto{\pgfqpoint{3.842531in}{2.958534in}}%
\pgfpathlineto{\pgfqpoint{3.845049in}{2.931161in}}%
\pgfpathlineto{\pgfqpoint{3.847568in}{3.000849in}}%
\pgfpathlineto{\pgfqpoint{3.852606in}{3.066491in}}%
\pgfpathlineto{\pgfqpoint{3.855124in}{3.088522in}}%
\pgfpathlineto{\pgfqpoint{3.857643in}{3.096976in}}%
\pgfpathlineto{\pgfqpoint{3.860162in}{3.123929in}}%
\pgfpathlineto{\pgfqpoint{3.862681in}{3.134757in}}%
\pgfpathlineto{\pgfqpoint{3.865199in}{3.127598in}}%
\pgfpathlineto{\pgfqpoint{3.867718in}{3.174128in}}%
\pgfpathlineto{\pgfqpoint{3.870237in}{3.177549in}}%
\pgfpathlineto{\pgfqpoint{3.872756in}{3.164229in}}%
\pgfpathlineto{\pgfqpoint{3.875274in}{3.177629in}}%
\pgfpathlineto{\pgfqpoint{3.877793in}{3.198261in}}%
\pgfpathlineto{\pgfqpoint{3.880312in}{3.179838in}}%
\pgfpathlineto{\pgfqpoint{3.882831in}{3.182781in}}%
\pgfpathlineto{\pgfqpoint{3.885349in}{3.185003in}}%
\pgfpathlineto{\pgfqpoint{3.887868in}{3.175274in}}%
\pgfpathlineto{\pgfqpoint{3.890387in}{3.181597in}}%
\pgfpathlineto{\pgfqpoint{3.892906in}{3.183353in}}%
\pgfpathlineto{\pgfqpoint{3.895424in}{3.160061in}}%
\pgfpathlineto{\pgfqpoint{3.897943in}{3.167335in}}%
\pgfpathlineto{\pgfqpoint{3.900462in}{3.201002in}}%
\pgfpathlineto{\pgfqpoint{3.902981in}{3.210350in}}%
\pgfpathlineto{\pgfqpoint{3.905499in}{3.230148in}}%
\pgfpathlineto{\pgfqpoint{3.908018in}{3.198342in}}%
\pgfpathlineto{\pgfqpoint{3.910537in}{3.187547in}}%
\pgfpathlineto{\pgfqpoint{3.913056in}{3.186585in}}%
\pgfpathlineto{\pgfqpoint{3.915574in}{3.213676in}}%
\pgfpathlineto{\pgfqpoint{3.918093in}{3.224140in}}%
\pgfpathlineto{\pgfqpoint{3.920612in}{3.223420in}}%
\pgfpathlineto{\pgfqpoint{3.923131in}{3.183860in}}%
\pgfpathlineto{\pgfqpoint{3.925649in}{3.137733in}}%
\pgfpathlineto{\pgfqpoint{3.928168in}{3.106352in}}%
\pgfpathlineto{\pgfqpoint{3.930687in}{3.109646in}}%
\pgfpathlineto{\pgfqpoint{3.933206in}{3.107079in}}%
\pgfpathlineto{\pgfqpoint{3.935724in}{3.061668in}}%
\pgfpathlineto{\pgfqpoint{3.938243in}{3.060074in}}%
\pgfpathlineto{\pgfqpoint{3.940762in}{3.014008in}}%
\pgfpathlineto{\pgfqpoint{3.943281in}{2.992226in}}%
\pgfpathlineto{\pgfqpoint{3.945799in}{2.990603in}}%
\pgfpathlineto{\pgfqpoint{3.953356in}{3.039492in}}%
\pgfpathlineto{\pgfqpoint{3.955874in}{3.032637in}}%
\pgfpathlineto{\pgfqpoint{3.958393in}{3.064462in}}%
\pgfpathlineto{\pgfqpoint{3.960912in}{3.042022in}}%
\pgfpathlineto{\pgfqpoint{3.963431in}{3.056672in}}%
\pgfpathlineto{\pgfqpoint{3.965949in}{3.044747in}}%
\pgfpathlineto{\pgfqpoint{3.968468in}{3.047207in}}%
\pgfpathlineto{\pgfqpoint{3.970987in}{3.028407in}}%
\pgfpathlineto{\pgfqpoint{3.973506in}{2.995826in}}%
\pgfpathlineto{\pgfqpoint{3.976024in}{2.996814in}}%
\pgfpathlineto{\pgfqpoint{3.978543in}{3.011252in}}%
\pgfpathlineto{\pgfqpoint{3.981062in}{3.026994in}}%
\pgfpathlineto{\pgfqpoint{3.983581in}{3.047387in}}%
\pgfpathlineto{\pgfqpoint{3.986099in}{3.027903in}}%
\pgfpathlineto{\pgfqpoint{3.988618in}{3.017616in}}%
\pgfpathlineto{\pgfqpoint{3.991137in}{2.962298in}}%
\pgfpathlineto{\pgfqpoint{3.993656in}{2.950780in}}%
\pgfpathlineto{\pgfqpoint{3.996174in}{2.965043in}}%
\pgfpathlineto{\pgfqpoint{3.998693in}{2.990140in}}%
\pgfpathlineto{\pgfqpoint{4.001212in}{3.017874in}}%
\pgfpathlineto{\pgfqpoint{4.003731in}{3.057585in}}%
\pgfpathlineto{\pgfqpoint{4.006249in}{3.033909in}}%
\pgfpathlineto{\pgfqpoint{4.011287in}{2.971608in}}%
\pgfpathlineto{\pgfqpoint{4.013806in}{2.947557in}}%
\pgfpathlineto{\pgfqpoint{4.016324in}{2.981866in}}%
\pgfpathlineto{\pgfqpoint{4.018843in}{2.989120in}}%
\pgfpathlineto{\pgfqpoint{4.021362in}{2.969620in}}%
\pgfpathlineto{\pgfqpoint{4.023881in}{2.931700in}}%
\pgfpathlineto{\pgfqpoint{4.026399in}{2.900145in}}%
\pgfpathlineto{\pgfqpoint{4.028918in}{2.849429in}}%
\pgfpathlineto{\pgfqpoint{4.031437in}{2.851468in}}%
\pgfpathlineto{\pgfqpoint{4.033956in}{2.816824in}}%
\pgfpathlineto{\pgfqpoint{4.036474in}{2.838802in}}%
\pgfpathlineto{\pgfqpoint{4.038993in}{2.835125in}}%
\pgfpathlineto{\pgfqpoint{4.041512in}{2.803776in}}%
\pgfpathlineto{\pgfqpoint{4.044031in}{2.842789in}}%
\pgfpathlineto{\pgfqpoint{4.046549in}{2.834736in}}%
\pgfpathlineto{\pgfqpoint{4.049068in}{2.819885in}}%
\pgfpathlineto{\pgfqpoint{4.051587in}{2.788438in}}%
\pgfpathlineto{\pgfqpoint{4.054106in}{2.773887in}}%
\pgfpathlineto{\pgfqpoint{4.056624in}{2.769398in}}%
\pgfpathlineto{\pgfqpoint{4.059143in}{2.787783in}}%
\pgfpathlineto{\pgfqpoint{4.061662in}{2.766769in}}%
\pgfpathlineto{\pgfqpoint{4.064181in}{2.781867in}}%
\pgfpathlineto{\pgfqpoint{4.066699in}{2.783872in}}%
\pgfpathlineto{\pgfqpoint{4.069218in}{2.802701in}}%
\pgfpathlineto{\pgfqpoint{4.071737in}{2.789818in}}%
\pgfpathlineto{\pgfqpoint{4.074256in}{2.836130in}}%
\pgfpathlineto{\pgfqpoint{4.076774in}{2.835743in}}%
\pgfpathlineto{\pgfqpoint{4.079293in}{2.850367in}}%
\pgfpathlineto{\pgfqpoint{4.081812in}{2.867759in}}%
\pgfpathlineto{\pgfqpoint{4.084331in}{2.911430in}}%
\pgfpathlineto{\pgfqpoint{4.086849in}{2.897584in}}%
\pgfpathlineto{\pgfqpoint{4.089368in}{2.889908in}}%
\pgfpathlineto{\pgfqpoint{4.091887in}{2.926129in}}%
\pgfpathlineto{\pgfqpoint{4.094406in}{2.922476in}}%
\pgfpathlineto{\pgfqpoint{4.096924in}{2.984534in}}%
\pgfpathlineto{\pgfqpoint{4.099443in}{3.030268in}}%
\pgfpathlineto{\pgfqpoint{4.101962in}{3.020751in}}%
\pgfpathlineto{\pgfqpoint{4.104481in}{3.033299in}}%
\pgfpathlineto{\pgfqpoint{4.106999in}{3.067691in}}%
\pgfpathlineto{\pgfqpoint{4.109518in}{3.084256in}}%
\pgfpathlineto{\pgfqpoint{4.112037in}{3.106864in}}%
\pgfpathlineto{\pgfqpoint{4.114556in}{3.094353in}}%
\pgfpathlineto{\pgfqpoint{4.117074in}{3.153083in}}%
\pgfpathlineto{\pgfqpoint{4.119593in}{3.171252in}}%
\pgfpathlineto{\pgfqpoint{4.122112in}{3.145696in}}%
\pgfpathlineto{\pgfqpoint{4.124631in}{3.166859in}}%
\pgfpathlineto{\pgfqpoint{4.127149in}{3.193488in}}%
\pgfpathlineto{\pgfqpoint{4.132187in}{3.167197in}}%
\pgfpathlineto{\pgfqpoint{4.134706in}{3.221828in}}%
\pgfpathlineto{\pgfqpoint{4.137224in}{3.181324in}}%
\pgfpathlineto{\pgfqpoint{4.139743in}{3.197691in}}%
\pgfpathlineto{\pgfqpoint{4.142262in}{3.182781in}}%
\pgfpathlineto{\pgfqpoint{4.144781in}{3.209055in}}%
\pgfpathlineto{\pgfqpoint{4.147299in}{3.211504in}}%
\pgfpathlineto{\pgfqpoint{4.149818in}{3.206819in}}%
\pgfpathlineto{\pgfqpoint{4.152337in}{3.233227in}}%
\pgfpathlineto{\pgfqpoint{4.154856in}{3.234727in}}%
\pgfpathlineto{\pgfqpoint{4.157374in}{3.231345in}}%
\pgfpathlineto{\pgfqpoint{4.159893in}{3.236971in}}%
\pgfpathlineto{\pgfqpoint{4.162412in}{3.291844in}}%
\pgfpathlineto{\pgfqpoint{4.164931in}{3.289121in}}%
\pgfpathlineto{\pgfqpoint{4.167449in}{3.273395in}}%
\pgfpathlineto{\pgfqpoint{4.169968in}{3.267853in}}%
\pgfpathlineto{\pgfqpoint{4.172487in}{3.275412in}}%
\pgfpathlineto{\pgfqpoint{4.175006in}{3.267211in}}%
\pgfpathlineto{\pgfqpoint{4.177524in}{3.289828in}}%
\pgfpathlineto{\pgfqpoint{4.180043in}{3.279759in}}%
\pgfpathlineto{\pgfqpoint{4.182562in}{3.264169in}}%
\pgfpathlineto{\pgfqpoint{4.185081in}{3.277639in}}%
\pgfpathlineto{\pgfqpoint{4.187599in}{3.299862in}}%
\pgfpathlineto{\pgfqpoint{4.190118in}{3.306626in}}%
\pgfpathlineto{\pgfqpoint{4.192637in}{3.324062in}}%
\pgfpathlineto{\pgfqpoint{4.195156in}{3.314350in}}%
\pgfpathlineto{\pgfqpoint{4.197674in}{3.318855in}}%
\pgfpathlineto{\pgfqpoint{4.200193in}{3.305937in}}%
\pgfpathlineto{\pgfqpoint{4.202712in}{3.331370in}}%
\pgfpathlineto{\pgfqpoint{4.205231in}{3.309666in}}%
\pgfpathlineto{\pgfqpoint{4.207749in}{3.310791in}}%
\pgfpathlineto{\pgfqpoint{4.210268in}{3.352838in}}%
\pgfpathlineto{\pgfqpoint{4.212787in}{3.362824in}}%
\pgfpathlineto{\pgfqpoint{4.215306in}{3.365756in}}%
\pgfpathlineto{\pgfqpoint{4.217824in}{3.364175in}}%
\pgfpathlineto{\pgfqpoint{4.220343in}{3.387301in}}%
\pgfpathlineto{\pgfqpoint{4.222862in}{3.377521in}}%
\pgfpathlineto{\pgfqpoint{4.225381in}{3.374600in}}%
\pgfpathlineto{\pgfqpoint{4.227899in}{3.414100in}}%
\pgfpathlineto{\pgfqpoint{4.230418in}{3.443264in}}%
\pgfpathlineto{\pgfqpoint{4.232937in}{3.473520in}}%
\pgfpathlineto{\pgfqpoint{4.235456in}{3.506726in}}%
\pgfpathlineto{\pgfqpoint{4.237974in}{3.491825in}}%
\pgfpathlineto{\pgfqpoint{4.240493in}{3.495809in}}%
\pgfpathlineto{\pgfqpoint{4.243012in}{3.530273in}}%
\pgfpathlineto{\pgfqpoint{4.245531in}{3.548398in}}%
\pgfpathlineto{\pgfqpoint{4.248049in}{3.528733in}}%
\pgfpathlineto{\pgfqpoint{4.250568in}{3.549004in}}%
\pgfpathlineto{\pgfqpoint{4.253087in}{3.543815in}}%
\pgfpathlineto{\pgfqpoint{4.255606in}{3.566134in}}%
\pgfpathlineto{\pgfqpoint{4.258124in}{3.572298in}}%
\pgfpathlineto{\pgfqpoint{4.260643in}{3.590039in}}%
\pgfpathlineto{\pgfqpoint{4.263162in}{3.591842in}}%
\pgfpathlineto{\pgfqpoint{4.265681in}{3.617958in}}%
\pgfpathlineto{\pgfqpoint{4.268199in}{3.629797in}}%
\pgfpathlineto{\pgfqpoint{4.270718in}{3.626523in}}%
\pgfpathlineto{\pgfqpoint{4.273237in}{3.627679in}}%
\pgfpathlineto{\pgfqpoint{4.275756in}{3.627299in}}%
\pgfpathlineto{\pgfqpoint{4.278274in}{3.610999in}}%
\pgfpathlineto{\pgfqpoint{4.280793in}{3.588947in}}%
\pgfpathlineto{\pgfqpoint{4.283312in}{3.596222in}}%
\pgfpathlineto{\pgfqpoint{4.285831in}{3.601886in}}%
\pgfpathlineto{\pgfqpoint{4.288349in}{3.534514in}}%
\pgfpathlineto{\pgfqpoint{4.290868in}{3.507311in}}%
\pgfpathlineto{\pgfqpoint{4.293387in}{3.485604in}}%
\pgfpathlineto{\pgfqpoint{4.295906in}{3.493478in}}%
\pgfpathlineto{\pgfqpoint{4.298424in}{3.493206in}}%
\pgfpathlineto{\pgfqpoint{4.300943in}{3.485484in}}%
\pgfpathlineto{\pgfqpoint{4.303462in}{3.476252in}}%
\pgfpathlineto{\pgfqpoint{4.305981in}{3.485983in}}%
\pgfpathlineto{\pgfqpoint{4.308499in}{3.481661in}}%
\pgfpathlineto{\pgfqpoint{4.311018in}{3.488367in}}%
\pgfpathlineto{\pgfqpoint{4.313537in}{3.513452in}}%
\pgfpathlineto{\pgfqpoint{4.316056in}{3.479520in}}%
\pgfpathlineto{\pgfqpoint{4.318574in}{3.470706in}}%
\pgfpathlineto{\pgfqpoint{4.321093in}{3.464134in}}%
\pgfpathlineto{\pgfqpoint{4.323612in}{3.444470in}}%
\pgfpathlineto{\pgfqpoint{4.326131in}{3.422836in}}%
\pgfpathlineto{\pgfqpoint{4.328649in}{3.362392in}}%
\pgfpathlineto{\pgfqpoint{4.331168in}{3.349203in}}%
\pgfpathlineto{\pgfqpoint{4.333687in}{3.382931in}}%
\pgfpathlineto{\pgfqpoint{4.336206in}{3.356692in}}%
\pgfpathlineto{\pgfqpoint{4.338724in}{3.337920in}}%
\pgfpathlineto{\pgfqpoint{4.341243in}{3.357989in}}%
\pgfpathlineto{\pgfqpoint{4.343762in}{3.303521in}}%
\pgfpathlineto{\pgfqpoint{4.346281in}{3.220439in}}%
\pgfpathlineto{\pgfqpoint{4.348799in}{3.230914in}}%
\pgfpathlineto{\pgfqpoint{4.351318in}{3.195308in}}%
\pgfpathlineto{\pgfqpoint{4.353837in}{3.146538in}}%
\pgfpathlineto{\pgfqpoint{4.356356in}{3.177784in}}%
\pgfpathlineto{\pgfqpoint{4.358874in}{3.206042in}}%
\pgfpathlineto{\pgfqpoint{4.361393in}{3.208629in}}%
\pgfpathlineto{\pgfqpoint{4.363912in}{3.219702in}}%
\pgfpathlineto{\pgfqpoint{4.366431in}{3.185217in}}%
\pgfpathlineto{\pgfqpoint{4.368949in}{3.209615in}}%
\pgfpathlineto{\pgfqpoint{4.371468in}{3.208665in}}%
\pgfpathlineto{\pgfqpoint{4.373987in}{3.165971in}}%
\pgfpathlineto{\pgfqpoint{4.376506in}{3.183179in}}%
\pgfpathlineto{\pgfqpoint{4.379024in}{3.220787in}}%
\pgfpathlineto{\pgfqpoint{4.381543in}{3.229963in}}%
\pgfpathlineto{\pgfqpoint{4.384062in}{3.222462in}}%
\pgfpathlineto{\pgfqpoint{4.386581in}{3.244585in}}%
\pgfpathlineto{\pgfqpoint{4.389099in}{3.240922in}}%
\pgfpathlineto{\pgfqpoint{4.391618in}{3.217398in}}%
\pgfpathlineto{\pgfqpoint{4.394137in}{3.202151in}}%
\pgfpathlineto{\pgfqpoint{4.396656in}{3.192308in}}%
\pgfpathlineto{\pgfqpoint{4.399174in}{3.142867in}}%
\pgfpathlineto{\pgfqpoint{4.401693in}{3.146033in}}%
\pgfpathlineto{\pgfqpoint{4.404212in}{3.140204in}}%
\pgfpathlineto{\pgfqpoint{4.406731in}{3.159549in}}%
\pgfpathlineto{\pgfqpoint{4.409249in}{3.152352in}}%
\pgfpathlineto{\pgfqpoint{4.411768in}{3.102387in}}%
\pgfpathlineto{\pgfqpoint{4.414287in}{3.127618in}}%
\pgfpathlineto{\pgfqpoint{4.416806in}{3.163227in}}%
\pgfpathlineto{\pgfqpoint{4.419324in}{3.156596in}}%
\pgfpathlineto{\pgfqpoint{4.421843in}{3.170337in}}%
\pgfpathlineto{\pgfqpoint{4.424362in}{3.134405in}}%
\pgfpathlineto{\pgfqpoint{4.426881in}{3.083457in}}%
\pgfpathlineto{\pgfqpoint{4.429399in}{3.108168in}}%
\pgfpathlineto{\pgfqpoint{4.431918in}{3.069435in}}%
\pgfpathlineto{\pgfqpoint{4.434437in}{3.090724in}}%
\pgfpathlineto{\pgfqpoint{4.436956in}{3.106903in}}%
\pgfpathlineto{\pgfqpoint{4.439474in}{3.090995in}}%
\pgfpathlineto{\pgfqpoint{4.441993in}{3.116566in}}%
\pgfpathlineto{\pgfqpoint{4.444512in}{3.125458in}}%
\pgfpathlineto{\pgfqpoint{4.447031in}{3.131235in}}%
\pgfpathlineto{\pgfqpoint{4.449549in}{3.143315in}}%
\pgfpathlineto{\pgfqpoint{4.452068in}{3.169382in}}%
\pgfpathlineto{\pgfqpoint{4.454587in}{3.169068in}}%
\pgfpathlineto{\pgfqpoint{4.457106in}{3.175156in}}%
\pgfpathlineto{\pgfqpoint{4.459624in}{3.173038in}}%
\pgfpathlineto{\pgfqpoint{4.462143in}{3.168735in}}%
\pgfpathlineto{\pgfqpoint{4.464662in}{3.148775in}}%
\pgfpathlineto{\pgfqpoint{4.467181in}{3.099474in}}%
\pgfpathlineto{\pgfqpoint{4.469699in}{3.100371in}}%
\pgfpathlineto{\pgfqpoint{4.472218in}{3.121229in}}%
\pgfpathlineto{\pgfqpoint{4.474737in}{3.127998in}}%
\pgfpathlineto{\pgfqpoint{4.477256in}{3.127694in}}%
\pgfpathlineto{\pgfqpoint{4.479774in}{3.114482in}}%
\pgfpathlineto{\pgfqpoint{4.482293in}{3.130099in}}%
\pgfpathlineto{\pgfqpoint{4.484812in}{3.124558in}}%
\pgfpathlineto{\pgfqpoint{4.487331in}{3.153547in}}%
\pgfpathlineto{\pgfqpoint{4.489849in}{3.153976in}}%
\pgfpathlineto{\pgfqpoint{4.492368in}{3.173609in}}%
\pgfpathlineto{\pgfqpoint{4.494887in}{3.191669in}}%
\pgfpathlineto{\pgfqpoint{4.497406in}{3.241269in}}%
\pgfpathlineto{\pgfqpoint{4.499924in}{3.274535in}}%
\pgfpathlineto{\pgfqpoint{4.502443in}{3.303674in}}%
\pgfpathlineto{\pgfqpoint{4.504962in}{3.248515in}}%
\pgfpathlineto{\pgfqpoint{4.507481in}{3.216040in}}%
\pgfpathlineto{\pgfqpoint{4.509999in}{3.201950in}}%
\pgfpathlineto{\pgfqpoint{4.512518in}{3.231016in}}%
\pgfpathlineto{\pgfqpoint{4.515037in}{3.196813in}}%
\pgfpathlineto{\pgfqpoint{4.517556in}{3.195614in}}%
\pgfpathlineto{\pgfqpoint{4.520074in}{3.149897in}}%
\pgfpathlineto{\pgfqpoint{4.522593in}{3.183137in}}%
\pgfpathlineto{\pgfqpoint{4.525112in}{3.137311in}}%
\pgfpathlineto{\pgfqpoint{4.527631in}{3.128900in}}%
\pgfpathlineto{\pgfqpoint{4.530149in}{3.127376in}}%
\pgfpathlineto{\pgfqpoint{4.532668in}{3.150554in}}%
\pgfpathlineto{\pgfqpoint{4.535187in}{3.106056in}}%
\pgfpathlineto{\pgfqpoint{4.537706in}{3.116620in}}%
\pgfpathlineto{\pgfqpoint{4.540224in}{3.176636in}}%
\pgfpathlineto{\pgfqpoint{4.542743in}{3.109935in}}%
\pgfpathlineto{\pgfqpoint{4.545262in}{3.099215in}}%
\pgfpathlineto{\pgfqpoint{4.547781in}{3.145529in}}%
\pgfpathlineto{\pgfqpoint{4.550299in}{3.067671in}}%
\pgfpathlineto{\pgfqpoint{4.552818in}{3.068008in}}%
\pgfpathlineto{\pgfqpoint{4.555337in}{3.103476in}}%
\pgfpathlineto{\pgfqpoint{4.557856in}{3.106682in}}%
\pgfpathlineto{\pgfqpoint{4.560374in}{3.110865in}}%
\pgfpathlineto{\pgfqpoint{4.562893in}{3.160989in}}%
\pgfpathlineto{\pgfqpoint{4.565412in}{3.121821in}}%
\pgfpathlineto{\pgfqpoint{4.570449in}{3.094287in}}%
\pgfpathlineto{\pgfqpoint{4.572968in}{3.145454in}}%
\pgfpathlineto{\pgfqpoint{4.575487in}{3.165936in}}%
\pgfpathlineto{\pgfqpoint{4.578006in}{3.196415in}}%
\pgfpathlineto{\pgfqpoint{4.580524in}{3.173406in}}%
\pgfpathlineto{\pgfqpoint{4.583043in}{3.206736in}}%
\pgfpathlineto{\pgfqpoint{4.585562in}{3.182086in}}%
\pgfpathlineto{\pgfqpoint{4.588081in}{3.205351in}}%
\pgfpathlineto{\pgfqpoint{4.590599in}{3.173059in}}%
\pgfpathlineto{\pgfqpoint{4.593118in}{3.166711in}}%
\pgfpathlineto{\pgfqpoint{4.595637in}{3.137630in}}%
\pgfpathlineto{\pgfqpoint{4.598156in}{3.156497in}}%
\pgfpathlineto{\pgfqpoint{4.600674in}{3.160360in}}%
\pgfpathlineto{\pgfqpoint{4.603193in}{3.187615in}}%
\pgfpathlineto{\pgfqpoint{4.605712in}{3.167178in}}%
\pgfpathlineto{\pgfqpoint{4.608231in}{3.196764in}}%
\pgfpathlineto{\pgfqpoint{4.610749in}{3.196891in}}%
\pgfpathlineto{\pgfqpoint{4.613268in}{3.263666in}}%
\pgfpathlineto{\pgfqpoint{4.615787in}{3.202275in}}%
\pgfpathlineto{\pgfqpoint{4.618306in}{3.191938in}}%
\pgfpathlineto{\pgfqpoint{4.620824in}{3.202896in}}%
\pgfpathlineto{\pgfqpoint{4.623343in}{3.205651in}}%
\pgfpathlineto{\pgfqpoint{4.625862in}{3.197603in}}%
\pgfpathlineto{\pgfqpoint{4.628381in}{3.200411in}}%
\pgfpathlineto{\pgfqpoint{4.630899in}{3.172380in}}%
\pgfpathlineto{\pgfqpoint{4.633418in}{3.158956in}}%
\pgfpathlineto{\pgfqpoint{4.635937in}{3.176977in}}%
\pgfpathlineto{\pgfqpoint{4.638456in}{3.178794in}}%
\pgfpathlineto{\pgfqpoint{4.640974in}{3.176363in}}%
\pgfpathlineto{\pgfqpoint{4.643493in}{3.228701in}}%
\pgfpathlineto{\pgfqpoint{4.646012in}{3.231941in}}%
\pgfpathlineto{\pgfqpoint{4.648531in}{3.211177in}}%
\pgfpathlineto{\pgfqpoint{4.651049in}{3.215028in}}%
\pgfpathlineto{\pgfqpoint{4.653568in}{3.227862in}}%
\pgfpathlineto{\pgfqpoint{4.656087in}{3.199027in}}%
\pgfpathlineto{\pgfqpoint{4.658606in}{3.274759in}}%
\pgfpathlineto{\pgfqpoint{4.661124in}{3.225810in}}%
\pgfpathlineto{\pgfqpoint{4.663643in}{3.230603in}}%
\pgfpathlineto{\pgfqpoint{4.666162in}{3.240680in}}%
\pgfpathlineto{\pgfqpoint{4.668681in}{3.234878in}}%
\pgfpathlineto{\pgfqpoint{4.671199in}{3.243771in}}%
\pgfpathlineto{\pgfqpoint{4.673718in}{3.242318in}}%
\pgfpathlineto{\pgfqpoint{4.676237in}{3.212064in}}%
\pgfpathlineto{\pgfqpoint{4.678756in}{3.224939in}}%
\pgfpathlineto{\pgfqpoint{4.681274in}{3.285020in}}%
\pgfpathlineto{\pgfqpoint{4.683793in}{3.317912in}}%
\pgfpathlineto{\pgfqpoint{4.686312in}{3.344198in}}%
\pgfpathlineto{\pgfqpoint{4.688831in}{3.355289in}}%
\pgfpathlineto{\pgfqpoint{4.691349in}{3.400971in}}%
\pgfpathlineto{\pgfqpoint{4.693868in}{3.431318in}}%
\pgfpathlineto{\pgfqpoint{4.696387in}{3.427410in}}%
\pgfpathlineto{\pgfqpoint{4.698906in}{3.444657in}}%
\pgfpathlineto{\pgfqpoint{4.701424in}{3.441352in}}%
\pgfpathlineto{\pgfqpoint{4.703943in}{3.421320in}}%
\pgfpathlineto{\pgfqpoint{4.706462in}{3.427329in}}%
\pgfpathlineto{\pgfqpoint{4.708981in}{3.450289in}}%
\pgfpathlineto{\pgfqpoint{4.711499in}{3.502546in}}%
\pgfpathlineto{\pgfqpoint{4.714018in}{3.525655in}}%
\pgfpathlineto{\pgfqpoint{4.716537in}{3.569963in}}%
\pgfpathlineto{\pgfqpoint{4.719056in}{3.572386in}}%
\pgfpathlineto{\pgfqpoint{4.721574in}{3.571011in}}%
\pgfpathlineto{\pgfqpoint{4.724093in}{3.533791in}}%
\pgfpathlineto{\pgfqpoint{4.726612in}{3.517061in}}%
\pgfpathlineto{\pgfqpoint{4.729131in}{3.526431in}}%
\pgfpathlineto{\pgfqpoint{4.731649in}{3.539115in}}%
\pgfpathlineto{\pgfqpoint{4.734168in}{3.617323in}}%
\pgfpathlineto{\pgfqpoint{4.736687in}{3.630420in}}%
\pgfpathlineto{\pgfqpoint{4.739206in}{3.638250in}}%
\pgfpathlineto{\pgfqpoint{4.741724in}{3.714133in}}%
\pgfpathlineto{\pgfqpoint{4.744243in}{3.664070in}}%
\pgfpathlineto{\pgfqpoint{4.746762in}{3.637481in}}%
\pgfpathlineto{\pgfqpoint{4.749281in}{3.585001in}}%
\pgfpathlineto{\pgfqpoint{4.751799in}{3.582476in}}%
\pgfpathlineto{\pgfqpoint{4.754318in}{3.598571in}}%
\pgfpathlineto{\pgfqpoint{4.756837in}{3.625626in}}%
\pgfpathlineto{\pgfqpoint{4.759356in}{3.611578in}}%
\pgfpathlineto{\pgfqpoint{4.761874in}{3.620142in}}%
\pgfpathlineto{\pgfqpoint{4.764393in}{3.574257in}}%
\pgfpathlineto{\pgfqpoint{4.766912in}{3.590886in}}%
\pgfpathlineto{\pgfqpoint{4.769431in}{3.559539in}}%
\pgfpathlineto{\pgfqpoint{4.771949in}{3.592172in}}%
\pgfpathlineto{\pgfqpoint{4.774468in}{3.539231in}}%
\pgfpathlineto{\pgfqpoint{4.776987in}{3.550812in}}%
\pgfpathlineto{\pgfqpoint{4.779506in}{3.510465in}}%
\pgfpathlineto{\pgfqpoint{4.782024in}{3.533632in}}%
\pgfpathlineto{\pgfqpoint{4.784543in}{3.529298in}}%
\pgfpathlineto{\pgfqpoint{4.787062in}{3.525524in}}%
\pgfpathlineto{\pgfqpoint{4.789581in}{3.468807in}}%
\pgfpathlineto{\pgfqpoint{4.792099in}{3.510169in}}%
\pgfpathlineto{\pgfqpoint{4.794618in}{3.527988in}}%
\pgfpathlineto{\pgfqpoint{4.797137in}{3.484637in}}%
\pgfpathlineto{\pgfqpoint{4.799656in}{3.482312in}}%
\pgfpathlineto{\pgfqpoint{4.802174in}{3.461318in}}%
\pgfpathlineto{\pgfqpoint{4.804693in}{3.445105in}}%
\pgfpathlineto{\pgfqpoint{4.807212in}{3.421405in}}%
\pgfpathlineto{\pgfqpoint{4.809731in}{3.420429in}}%
\pgfpathlineto{\pgfqpoint{4.812249in}{3.417864in}}%
\pgfpathlineto{\pgfqpoint{4.814768in}{3.448549in}}%
\pgfpathlineto{\pgfqpoint{4.817287in}{3.446196in}}%
\pgfpathlineto{\pgfqpoint{4.819806in}{3.474368in}}%
\pgfpathlineto{\pgfqpoint{4.822324in}{3.493428in}}%
\pgfpathlineto{\pgfqpoint{4.824843in}{3.450189in}}%
\pgfpathlineto{\pgfqpoint{4.827362in}{3.480463in}}%
\pgfpathlineto{\pgfqpoint{4.829881in}{3.494659in}}%
\pgfpathlineto{\pgfqpoint{4.832399in}{3.483724in}}%
\pgfpathlineto{\pgfqpoint{4.834918in}{3.478160in}}%
\pgfpathlineto{\pgfqpoint{4.837437in}{3.521973in}}%
\pgfpathlineto{\pgfqpoint{4.839956in}{3.563410in}}%
\pgfpathlineto{\pgfqpoint{4.842474in}{3.566816in}}%
\pgfpathlineto{\pgfqpoint{4.844993in}{3.550195in}}%
\pgfpathlineto{\pgfqpoint{4.847512in}{3.567360in}}%
\pgfpathlineto{\pgfqpoint{4.850031in}{3.571275in}}%
\pgfpathlineto{\pgfqpoint{4.852549in}{3.587610in}}%
\pgfpathlineto{\pgfqpoint{4.855068in}{3.607323in}}%
\pgfpathlineto{\pgfqpoint{4.857587in}{3.593728in}}%
\pgfpathlineto{\pgfqpoint{4.860106in}{3.598495in}}%
\pgfpathlineto{\pgfqpoint{4.862624in}{3.604632in}}%
\pgfpathlineto{\pgfqpoint{4.865143in}{3.612388in}}%
\pgfpathlineto{\pgfqpoint{4.867662in}{3.602601in}}%
\pgfpathlineto{\pgfqpoint{4.870181in}{3.556884in}}%
\pgfpathlineto{\pgfqpoint{4.872699in}{3.549649in}}%
\pgfpathlineto{\pgfqpoint{4.875218in}{3.559344in}}%
\pgfpathlineto{\pgfqpoint{4.877737in}{3.548591in}}%
\pgfpathlineto{\pgfqpoint{4.880256in}{3.561482in}}%
\pgfpathlineto{\pgfqpoint{4.882774in}{3.565421in}}%
\pgfpathlineto{\pgfqpoint{4.885293in}{3.573834in}}%
\pgfpathlineto{\pgfqpoint{4.887812in}{3.586670in}}%
\pgfpathlineto{\pgfqpoint{4.890331in}{3.589501in}}%
\pgfpathlineto{\pgfqpoint{4.892849in}{3.532182in}}%
\pgfpathlineto{\pgfqpoint{4.895368in}{3.567672in}}%
\pgfpathlineto{\pgfqpoint{4.897887in}{3.521367in}}%
\pgfpathlineto{\pgfqpoint{4.900406in}{3.590969in}}%
\pgfpathlineto{\pgfqpoint{4.902924in}{3.600205in}}%
\pgfpathlineto{\pgfqpoint{4.905443in}{3.605711in}}%
\pgfpathlineto{\pgfqpoint{4.907962in}{3.641621in}}%
\pgfpathlineto{\pgfqpoint{4.910481in}{3.619823in}}%
\pgfpathlineto{\pgfqpoint{4.912999in}{3.629049in}}%
\pgfpathlineto{\pgfqpoint{4.915518in}{3.659714in}}%
\pgfpathlineto{\pgfqpoint{4.918037in}{3.587628in}}%
\pgfpathlineto{\pgfqpoint{4.920556in}{3.570698in}}%
\pgfpathlineto{\pgfqpoint{4.923074in}{3.608689in}}%
\pgfpathlineto{\pgfqpoint{4.925593in}{3.611587in}}%
\pgfpathlineto{\pgfqpoint{4.928112in}{3.631940in}}%
\pgfpathlineto{\pgfqpoint{4.930631in}{3.605088in}}%
\pgfpathlineto{\pgfqpoint{4.933149in}{3.615927in}}%
\pgfpathlineto{\pgfqpoint{4.935668in}{3.589284in}}%
\pgfpathlineto{\pgfqpoint{4.938187in}{3.601480in}}%
\pgfpathlineto{\pgfqpoint{4.940706in}{3.595150in}}%
\pgfpathlineto{\pgfqpoint{4.943224in}{3.603520in}}%
\pgfpathlineto{\pgfqpoint{4.948262in}{3.642277in}}%
\pgfpathlineto{\pgfqpoint{4.950781in}{3.589609in}}%
\pgfpathlineto{\pgfqpoint{4.953299in}{3.647380in}}%
\pgfpathlineto{\pgfqpoint{4.955818in}{3.612615in}}%
\pgfpathlineto{\pgfqpoint{4.958337in}{3.689584in}}%
\pgfpathlineto{\pgfqpoint{4.960856in}{3.679053in}}%
\pgfpathlineto{\pgfqpoint{4.963374in}{3.703012in}}%
\pgfpathlineto{\pgfqpoint{4.965893in}{3.769200in}}%
\pgfpathlineto{\pgfqpoint{4.968412in}{3.741841in}}%
\pgfpathlineto{\pgfqpoint{4.970931in}{3.731215in}}%
\pgfpathlineto{\pgfqpoint{4.973449in}{3.787341in}}%
\pgfpathlineto{\pgfqpoint{4.975968in}{3.798183in}}%
\pgfpathlineto{\pgfqpoint{4.978487in}{3.784409in}}%
\pgfpathlineto{\pgfqpoint{4.981006in}{3.809595in}}%
\pgfpathlineto{\pgfqpoint{4.983524in}{3.867433in}}%
\pgfpathlineto{\pgfqpoint{4.986043in}{3.906410in}}%
\pgfpathlineto{\pgfqpoint{4.988562in}{3.903726in}}%
\pgfpathlineto{\pgfqpoint{4.991081in}{3.897025in}}%
\pgfpathlineto{\pgfqpoint{4.993599in}{3.902093in}}%
\pgfpathlineto{\pgfqpoint{4.996118in}{3.916115in}}%
\pgfpathlineto{\pgfqpoint{4.998637in}{3.876251in}}%
\pgfpathlineto{\pgfqpoint{5.001156in}{3.895184in}}%
\pgfpathlineto{\pgfqpoint{5.003674in}{3.950222in}}%
\pgfpathlineto{\pgfqpoint{5.006193in}{3.960981in}}%
\pgfpathlineto{\pgfqpoint{5.008712in}{3.946404in}}%
\pgfpathlineto{\pgfqpoint{5.011231in}{3.944030in}}%
\pgfpathlineto{\pgfqpoint{5.013749in}{3.925019in}}%
\pgfpathlineto{\pgfqpoint{5.016268in}{3.925938in}}%
\pgfpathlineto{\pgfqpoint{5.018787in}{3.962851in}}%
\pgfpathlineto{\pgfqpoint{5.021306in}{4.001734in}}%
\pgfpathlineto{\pgfqpoint{5.023824in}{4.030333in}}%
\pgfpathlineto{\pgfqpoint{5.026343in}{4.019067in}}%
\pgfpathlineto{\pgfqpoint{5.028862in}{4.031337in}}%
\pgfpathlineto{\pgfqpoint{5.031381in}{4.034485in}}%
\pgfpathlineto{\pgfqpoint{5.033899in}{4.030131in}}%
\pgfpathlineto{\pgfqpoint{5.036418in}{4.013334in}}%
\pgfpathlineto{\pgfqpoint{5.038937in}{4.030642in}}%
\pgfpathlineto{\pgfqpoint{5.041456in}{3.998926in}}%
\pgfpathlineto{\pgfqpoint{5.043974in}{3.985560in}}%
\pgfpathlineto{\pgfqpoint{5.046493in}{3.995180in}}%
\pgfpathlineto{\pgfqpoint{5.049012in}{4.026123in}}%
\pgfpathlineto{\pgfqpoint{5.051531in}{4.014004in}}%
\pgfpathlineto{\pgfqpoint{5.054049in}{4.017300in}}%
\pgfpathlineto{\pgfqpoint{5.056568in}{4.008060in}}%
\pgfpathlineto{\pgfqpoint{5.059087in}{4.033885in}}%
\pgfpathlineto{\pgfqpoint{5.061606in}{4.000827in}}%
\pgfpathlineto{\pgfqpoint{5.064124in}{3.910662in}}%
\pgfpathlineto{\pgfqpoint{5.066643in}{3.900299in}}%
\pgfpathlineto{\pgfqpoint{5.069162in}{3.924141in}}%
\pgfpathlineto{\pgfqpoint{5.071681in}{3.925645in}}%
\pgfpathlineto{\pgfqpoint{5.074199in}{3.943179in}}%
\pgfpathlineto{\pgfqpoint{5.076718in}{3.951517in}}%
\pgfpathlineto{\pgfqpoint{5.079237in}{3.930962in}}%
\pgfpathlineto{\pgfqpoint{5.081756in}{3.886808in}}%
\pgfpathlineto{\pgfqpoint{5.084274in}{3.885456in}}%
\pgfpathlineto{\pgfqpoint{5.086793in}{3.844534in}}%
\pgfpathlineto{\pgfqpoint{5.089312in}{3.837419in}}%
\pgfpathlineto{\pgfqpoint{5.091831in}{3.859621in}}%
\pgfpathlineto{\pgfqpoint{5.094349in}{3.848719in}}%
\pgfpathlineto{\pgfqpoint{5.096868in}{3.835666in}}%
\pgfpathlineto{\pgfqpoint{5.099387in}{3.839565in}}%
\pgfpathlineto{\pgfqpoint{5.101906in}{3.825178in}}%
\pgfpathlineto{\pgfqpoint{5.104424in}{3.873919in}}%
\pgfpathlineto{\pgfqpoint{5.106943in}{3.883420in}}%
\pgfpathlineto{\pgfqpoint{5.109462in}{3.884493in}}%
\pgfpathlineto{\pgfqpoint{5.111981in}{3.884053in}}%
\pgfpathlineto{\pgfqpoint{5.114499in}{3.874399in}}%
\pgfpathlineto{\pgfqpoint{5.117018in}{3.876867in}}%
\pgfpathlineto{\pgfqpoint{5.119537in}{3.876775in}}%
\pgfpathlineto{\pgfqpoint{5.122056in}{3.892764in}}%
\pgfpathlineto{\pgfqpoint{5.124574in}{3.856116in}}%
\pgfpathlineto{\pgfqpoint{5.127093in}{3.902221in}}%
\pgfpathlineto{\pgfqpoint{5.129612in}{3.893701in}}%
\pgfpathlineto{\pgfqpoint{5.132131in}{3.907215in}}%
\pgfpathlineto{\pgfqpoint{5.134649in}{3.898592in}}%
\pgfpathlineto{\pgfqpoint{5.137168in}{3.903786in}}%
\pgfpathlineto{\pgfqpoint{5.139687in}{3.935710in}}%
\pgfpathlineto{\pgfqpoint{5.142206in}{3.941699in}}%
\pgfpathlineto{\pgfqpoint{5.144724in}{3.961961in}}%
\pgfpathlineto{\pgfqpoint{5.147243in}{3.932513in}}%
\pgfpathlineto{\pgfqpoint{5.149762in}{3.960382in}}%
\pgfpathlineto{\pgfqpoint{5.152281in}{3.964814in}}%
\pgfpathlineto{\pgfqpoint{5.154799in}{4.018110in}}%
\pgfpathlineto{\pgfqpoint{5.157318in}{4.016279in}}%
\pgfpathlineto{\pgfqpoint{5.162356in}{4.046992in}}%
\pgfpathlineto{\pgfqpoint{5.164874in}{3.999651in}}%
\pgfpathlineto{\pgfqpoint{5.167393in}{4.039619in}}%
\pgfpathlineto{\pgfqpoint{5.169912in}{4.012948in}}%
\pgfpathlineto{\pgfqpoint{5.172431in}{3.993191in}}%
\pgfpathlineto{\pgfqpoint{5.174949in}{4.030294in}}%
\pgfpathlineto{\pgfqpoint{5.177468in}{4.101904in}}%
\pgfpathlineto{\pgfqpoint{5.179987in}{4.101360in}}%
\pgfpathlineto{\pgfqpoint{5.182506in}{4.140924in}}%
\pgfpathlineto{\pgfqpoint{5.185024in}{4.129103in}}%
\pgfpathlineto{\pgfqpoint{5.187543in}{4.115575in}}%
\pgfpathlineto{\pgfqpoint{5.190062in}{4.123871in}}%
\pgfpathlineto{\pgfqpoint{5.192581in}{4.102353in}}%
\pgfpathlineto{\pgfqpoint{5.195099in}{4.169068in}}%
\pgfpathlineto{\pgfqpoint{5.197618in}{4.161537in}}%
\pgfpathlineto{\pgfqpoint{5.200137in}{4.163356in}}%
\pgfpathlineto{\pgfqpoint{5.202656in}{4.164429in}}%
\pgfpathlineto{\pgfqpoint{5.205174in}{4.208606in}}%
\pgfpathlineto{\pgfqpoint{5.207693in}{4.223080in}}%
\pgfpathlineto{\pgfqpoint{5.210212in}{4.204053in}}%
\pgfpathlineto{\pgfqpoint{5.212731in}{4.172551in}}%
\pgfpathlineto{\pgfqpoint{5.215249in}{4.149558in}}%
\pgfpathlineto{\pgfqpoint{5.217768in}{4.171904in}}%
\pgfpathlineto{\pgfqpoint{5.220287in}{4.114103in}}%
\pgfpathlineto{\pgfqpoint{5.222806in}{4.127708in}}%
\pgfpathlineto{\pgfqpoint{5.225324in}{4.147801in}}%
\pgfpathlineto{\pgfqpoint{5.227843in}{4.146974in}}%
\pgfpathlineto{\pgfqpoint{5.230362in}{4.123717in}}%
\pgfpathlineto{\pgfqpoint{5.232881in}{4.152308in}}%
\pgfpathlineto{\pgfqpoint{5.235399in}{4.159651in}}%
\pgfpathlineto{\pgfqpoint{5.237918in}{4.126728in}}%
\pgfpathlineto{\pgfqpoint{5.240437in}{4.131119in}}%
\pgfpathlineto{\pgfqpoint{5.242956in}{4.144500in}}%
\pgfpathlineto{\pgfqpoint{5.245474in}{4.135726in}}%
\pgfpathlineto{\pgfqpoint{5.247993in}{4.124812in}}%
\pgfpathlineto{\pgfqpoint{5.250512in}{4.163846in}}%
\pgfpathlineto{\pgfqpoint{5.253031in}{4.117310in}}%
\pgfpathlineto{\pgfqpoint{5.255549in}{4.085878in}}%
\pgfpathlineto{\pgfqpoint{5.258068in}{4.071453in}}%
\pgfpathlineto{\pgfqpoint{5.260587in}{4.043565in}}%
\pgfpathlineto{\pgfqpoint{5.263106in}{3.991726in}}%
\pgfpathlineto{\pgfqpoint{5.265624in}{4.024553in}}%
\pgfpathlineto{\pgfqpoint{5.268143in}{4.020419in}}%
\pgfpathlineto{\pgfqpoint{5.270662in}{3.975126in}}%
\pgfpathlineto{\pgfqpoint{5.273181in}{3.986365in}}%
\pgfpathlineto{\pgfqpoint{5.275699in}{3.961501in}}%
\pgfpathlineto{\pgfqpoint{5.278218in}{3.984264in}}%
\pgfpathlineto{\pgfqpoint{5.280737in}{3.968241in}}%
\pgfpathlineto{\pgfqpoint{5.283256in}{3.981615in}}%
\pgfpathlineto{\pgfqpoint{5.285774in}{3.989605in}}%
\pgfpathlineto{\pgfqpoint{5.288293in}{3.974574in}}%
\pgfpathlineto{\pgfqpoint{5.290812in}{3.997257in}}%
\pgfpathlineto{\pgfqpoint{5.293331in}{3.993522in}}%
\pgfpathlineto{\pgfqpoint{5.295849in}{3.940400in}}%
\pgfpathlineto{\pgfqpoint{5.298368in}{3.936921in}}%
\pgfpathlineto{\pgfqpoint{5.300887in}{3.946632in}}%
\pgfpathlineto{\pgfqpoint{5.303406in}{3.997599in}}%
\pgfpathlineto{\pgfqpoint{5.305924in}{4.001863in}}%
\pgfpathlineto{\pgfqpoint{5.308443in}{4.002159in}}%
\pgfpathlineto{\pgfqpoint{5.310962in}{4.030304in}}%
\pgfpathlineto{\pgfqpoint{5.313481in}{3.993510in}}%
\pgfpathlineto{\pgfqpoint{5.315999in}{3.967521in}}%
\pgfpathlineto{\pgfqpoint{5.318518in}{4.023307in}}%
\pgfpathlineto{\pgfqpoint{5.321037in}{4.041260in}}%
\pgfpathlineto{\pgfqpoint{5.323556in}{4.049050in}}%
\pgfpathlineto{\pgfqpoint{5.326074in}{4.087819in}}%
\pgfpathlineto{\pgfqpoint{5.328593in}{4.063408in}}%
\pgfpathlineto{\pgfqpoint{5.331112in}{4.070138in}}%
\pgfpathlineto{\pgfqpoint{5.333631in}{4.075921in}}%
\pgfpathlineto{\pgfqpoint{5.336149in}{4.066127in}}%
\pgfpathlineto{\pgfqpoint{5.338668in}{4.113115in}}%
\pgfpathlineto{\pgfqpoint{5.341187in}{4.080188in}}%
\pgfpathlineto{\pgfqpoint{5.343706in}{4.080808in}}%
\pgfpathlineto{\pgfqpoint{5.346224in}{4.038889in}}%
\pgfpathlineto{\pgfqpoint{5.348743in}{4.052123in}}%
\pgfpathlineto{\pgfqpoint{5.351262in}{4.064165in}}%
\pgfpathlineto{\pgfqpoint{5.353781in}{4.107140in}}%
\pgfpathlineto{\pgfqpoint{5.356299in}{4.054184in}}%
\pgfpathlineto{\pgfqpoint{5.358818in}{4.051365in}}%
\pgfpathlineto{\pgfqpoint{5.361337in}{4.014482in}}%
\pgfpathlineto{\pgfqpoint{5.363856in}{4.011120in}}%
\pgfpathlineto{\pgfqpoint{5.366374in}{4.054047in}}%
\pgfpathlineto{\pgfqpoint{5.368893in}{4.040501in}}%
\pgfpathlineto{\pgfqpoint{5.371412in}{4.037377in}}%
\pgfpathlineto{\pgfqpoint{5.373931in}{4.025228in}}%
\pgfpathlineto{\pgfqpoint{5.376449in}{4.010096in}}%
\pgfpathlineto{\pgfqpoint{5.378968in}{4.013139in}}%
\pgfpathlineto{\pgfqpoint{5.381487in}{4.007589in}}%
\pgfpathlineto{\pgfqpoint{5.384006in}{3.983898in}}%
\pgfpathlineto{\pgfqpoint{5.386524in}{4.018942in}}%
\pgfpathlineto{\pgfqpoint{5.389043in}{3.996549in}}%
\pgfpathlineto{\pgfqpoint{5.391562in}{4.015764in}}%
\pgfpathlineto{\pgfqpoint{5.394081in}{3.962044in}}%
\pgfpathlineto{\pgfqpoint{5.396599in}{3.951390in}}%
\pgfpathlineto{\pgfqpoint{5.399118in}{3.950090in}}%
\pgfpathlineto{\pgfqpoint{5.401637in}{3.938168in}}%
\pgfpathlineto{\pgfqpoint{5.404156in}{3.928267in}}%
\pgfpathlineto{\pgfqpoint{5.406674in}{3.974365in}}%
\pgfpathlineto{\pgfqpoint{5.409193in}{3.950135in}}%
\pgfpathlineto{\pgfqpoint{5.411712in}{3.943451in}}%
\pgfpathlineto{\pgfqpoint{5.414231in}{3.952696in}}%
\pgfpathlineto{\pgfqpoint{5.416749in}{3.975031in}}%
\pgfpathlineto{\pgfqpoint{5.419268in}{4.009420in}}%
\pgfpathlineto{\pgfqpoint{5.421787in}{3.947730in}}%
\pgfpathlineto{\pgfqpoint{5.424306in}{3.982079in}}%
\pgfpathlineto{\pgfqpoint{5.426824in}{4.034853in}}%
\pgfpathlineto{\pgfqpoint{5.429343in}{4.040020in}}%
\pgfpathlineto{\pgfqpoint{5.431862in}{4.035313in}}%
\pgfpathlineto{\pgfqpoint{5.436899in}{4.061882in}}%
\pgfpathlineto{\pgfqpoint{5.439418in}{4.058930in}}%
\pgfpathlineto{\pgfqpoint{5.441937in}{4.019613in}}%
\pgfpathlineto{\pgfqpoint{5.444456in}{3.995062in}}%
\pgfpathlineto{\pgfqpoint{5.446974in}{4.045947in}}%
\pgfpathlineto{\pgfqpoint{5.449493in}{4.051769in}}%
\pgfpathlineto{\pgfqpoint{5.452012in}{4.073681in}}%
\pgfpathlineto{\pgfqpoint{5.454531in}{4.053297in}}%
\pgfpathlineto{\pgfqpoint{5.457049in}{4.099461in}}%
\pgfpathlineto{\pgfqpoint{5.459568in}{4.142579in}}%
\pgfpathlineto{\pgfqpoint{5.462087in}{4.110860in}}%
\pgfpathlineto{\pgfqpoint{5.464606in}{4.115881in}}%
\pgfpathlineto{\pgfqpoint{5.467124in}{4.105701in}}%
\pgfpathlineto{\pgfqpoint{5.469643in}{4.114482in}}%
\pgfpathlineto{\pgfqpoint{5.472162in}{4.098570in}}%
\pgfpathlineto{\pgfqpoint{5.474681in}{4.147199in}}%
\pgfpathlineto{\pgfqpoint{5.477199in}{4.125803in}}%
\pgfpathlineto{\pgfqpoint{5.479718in}{4.106418in}}%
\pgfpathlineto{\pgfqpoint{5.482237in}{4.066594in}}%
\pgfpathlineto{\pgfqpoint{5.484756in}{4.053264in}}%
\pgfpathlineto{\pgfqpoint{5.487274in}{4.045556in}}%
\pgfpathlineto{\pgfqpoint{5.489793in}{4.056559in}}%
\pgfpathlineto{\pgfqpoint{5.492312in}{4.064489in}}%
\pgfpathlineto{\pgfqpoint{5.494831in}{4.054018in}}%
\pgfpathlineto{\pgfqpoint{5.497349in}{4.042470in}}%
\pgfpathlineto{\pgfqpoint{5.502387in}{3.996003in}}%
\pgfpathlineto{\pgfqpoint{5.504906in}{4.006738in}}%
\pgfpathlineto{\pgfqpoint{5.507424in}{4.024454in}}%
\pgfpathlineto{\pgfqpoint{5.509943in}{4.065632in}}%
\pgfpathlineto{\pgfqpoint{5.512462in}{4.078340in}}%
\pgfpathlineto{\pgfqpoint{5.514981in}{4.088960in}}%
\pgfpathlineto{\pgfqpoint{5.517499in}{4.072484in}}%
\pgfpathlineto{\pgfqpoint{5.520018in}{4.082280in}}%
\pgfpathlineto{\pgfqpoint{5.522537in}{4.060182in}}%
\pgfpathlineto{\pgfqpoint{5.525056in}{4.017888in}}%
\pgfpathlineto{\pgfqpoint{5.527574in}{4.050089in}}%
\pgfpathlineto{\pgfqpoint{5.530093in}{4.047842in}}%
\pgfpathlineto{\pgfqpoint{5.532612in}{4.058005in}}%
\pgfpathlineto{\pgfqpoint{5.535131in}{4.068900in}}%
\pgfpathlineto{\pgfqpoint{5.537649in}{4.049674in}}%
\pgfpathlineto{\pgfqpoint{5.540168in}{4.046658in}}%
\pgfpathlineto{\pgfqpoint{5.542687in}{4.071353in}}%
\pgfpathlineto{\pgfqpoint{5.545206in}{4.026662in}}%
\pgfpathlineto{\pgfqpoint{5.547724in}{4.013714in}}%
\pgfpathlineto{\pgfqpoint{5.550243in}{4.049825in}}%
\pgfpathlineto{\pgfqpoint{5.552762in}{4.014841in}}%
\pgfpathlineto{\pgfqpoint{5.555281in}{3.995500in}}%
\pgfpathlineto{\pgfqpoint{5.557799in}{4.045169in}}%
\pgfpathlineto{\pgfqpoint{5.560318in}{4.074454in}}%
\pgfpathlineto{\pgfqpoint{5.562837in}{4.031497in}}%
\pgfpathlineto{\pgfqpoint{5.565356in}{4.037337in}}%
\pgfpathlineto{\pgfqpoint{5.567874in}{3.993275in}}%
\pgfpathlineto{\pgfqpoint{5.570393in}{3.984754in}}%
\pgfpathlineto{\pgfqpoint{5.572912in}{4.010007in}}%
\pgfpathlineto{\pgfqpoint{5.575431in}{4.019461in}}%
\pgfpathlineto{\pgfqpoint{5.577949in}{4.025607in}}%
\pgfpathlineto{\pgfqpoint{5.580468in}{4.021512in}}%
\pgfpathlineto{\pgfqpoint{5.582987in}{4.008063in}}%
\pgfpathlineto{\pgfqpoint{5.585506in}{4.029794in}}%
\pgfpathlineto{\pgfqpoint{5.588024in}{4.039270in}}%
\pgfpathlineto{\pgfqpoint{5.590543in}{4.038703in}}%
\pgfpathlineto{\pgfqpoint{5.593062in}{4.000503in}}%
\pgfpathlineto{\pgfqpoint{5.598099in}{3.971861in}}%
\pgfpathlineto{\pgfqpoint{5.600618in}{3.977269in}}%
\pgfpathlineto{\pgfqpoint{5.600618in}{3.977269in}}%
\pgfusepath{stroke}%
\end{pgfscope}%
\begin{pgfscope}%
\pgfsetrectcap%
\pgfsetmiterjoin%
\pgfsetlinewidth{0.501875pt}%
\definecolor{currentstroke}{rgb}{0.000000,0.000000,0.000000}%
\pgfsetstrokecolor{currentstroke}%
\pgfsetdash{}{0pt}%
\pgfpathmoveto{\pgfqpoint{0.563118in}{0.401080in}}%
\pgfpathlineto{\pgfqpoint{0.563118in}{4.405080in}}%
\pgfusepath{stroke}%
\end{pgfscope}%
\begin{pgfscope}%
\pgfsetrectcap%
\pgfsetmiterjoin%
\pgfsetlinewidth{0.501875pt}%
\definecolor{currentstroke}{rgb}{0.000000,0.000000,0.000000}%
\pgfsetstrokecolor{currentstroke}%
\pgfsetdash{}{0pt}%
\pgfpathmoveto{\pgfqpoint{5.600618in}{0.401080in}}%
\pgfpathlineto{\pgfqpoint{5.600618in}{4.405080in}}%
\pgfusepath{stroke}%
\end{pgfscope}%
\begin{pgfscope}%
\pgfsetrectcap%
\pgfsetmiterjoin%
\pgfsetlinewidth{0.501875pt}%
\definecolor{currentstroke}{rgb}{0.000000,0.000000,0.000000}%
\pgfsetstrokecolor{currentstroke}%
\pgfsetdash{}{0pt}%
\pgfpathmoveto{\pgfqpoint{0.563118in}{0.401080in}}%
\pgfpathlineto{\pgfqpoint{5.600618in}{0.401080in}}%
\pgfusepath{stroke}%
\end{pgfscope}%
\begin{pgfscope}%
\pgfsetrectcap%
\pgfsetmiterjoin%
\pgfsetlinewidth{0.501875pt}%
\definecolor{currentstroke}{rgb}{0.000000,0.000000,0.000000}%
\pgfsetstrokecolor{currentstroke}%
\pgfsetdash{}{0pt}%
\pgfpathmoveto{\pgfqpoint{0.563118in}{4.405080in}}%
\pgfpathlineto{\pgfqpoint{5.600618in}{4.405080in}}%
\pgfusepath{stroke}%
\end{pgfscope}%
\begin{pgfscope}%
\pgfsetrectcap%
\pgfsetroundjoin%
\pgfsetlinewidth{0.803000pt}%
\definecolor{currentstroke}{rgb}{0.000000,0.000000,0.000000}%
\pgfsetstrokecolor{currentstroke}%
\pgfsetdash{}{0pt}%
\pgfpathmoveto{\pgfqpoint{0.688118in}{4.231469in}}%
\pgfpathlineto{\pgfqpoint{0.827007in}{4.231469in}}%
\pgfpathlineto{\pgfqpoint{0.965896in}{4.231469in}}%
\pgfusepath{stroke}%
\end{pgfscope}%
\begin{pgfscope}%
\definecolor{textcolor}{rgb}{0.000000,0.000000,0.000000}%
\pgfsetstrokecolor{textcolor}%
\pgfsetfillcolor{textcolor}%
\pgftext[x=1.077007in,y=4.182858in,left,base]{\color{textcolor}{\rmfamily\fontsize{10.000000}{12.000000}\selectfont\catcode`\^=\active\def^{\ifmmode\sp\else\^{}\fi}\catcode`\%=\active\def%{\%}$\omega_1$}}%
\end{pgfscope}%
\begin{pgfscope}%
\pgfsetbuttcap%
\pgfsetroundjoin%
\pgfsetlinewidth{0.803000pt}%
\definecolor{currentstroke}{rgb}{0.690196,0.188235,0.376471}%
\pgfsetstrokecolor{currentstroke}%
\pgfsetdash{{2.960000pt}{1.280000pt}}{0.000000pt}%
\pgfpathmoveto{\pgfqpoint{0.688118in}{4.037796in}}%
\pgfpathlineto{\pgfqpoint{0.827007in}{4.037796in}}%
\pgfpathlineto{\pgfqpoint{0.965896in}{4.037796in}}%
\pgfusepath{stroke}%
\end{pgfscope}%
\begin{pgfscope}%
\definecolor{textcolor}{rgb}{0.000000,0.000000,0.000000}%
\pgfsetstrokecolor{textcolor}%
\pgfsetfillcolor{textcolor}%
\pgftext[x=1.077007in,y=3.989185in,left,base]{\color{textcolor}{\rmfamily\fontsize{10.000000}{12.000000}\selectfont\catcode`\^=\active\def^{\ifmmode\sp\else\^{}\fi}\catcode`\%=\active\def%{\%}$\omega_2$}}%
\end{pgfscope}%
\end{pgfpicture}%
\makeatother%
\endgroup%

    \caption{Two realisations of the Wiener process on the interval $[0,1]$.}
    \label{fig:twoRealisations.pgf}
\end{figure}

\begin{defn}{Brownian Motion}{}
A stochastic process $(B_t, t \ge 0)$ satisfying
\begin{equation*}
B_t = B_0 + \mu t + \sigma W_t, \quad t \ge 0,
\end{equation*}
where $(W_t)$ is a Wiener process independent of $B_0$, is called a \textit{Brownian motion with drift $\mu$ and diffusion coefficient $\sigma$}.
\end{defn}

A standard Brownian motion is one where $\mu = 0$ and $\sigma = 1$. The only difference between a Wiener process and a Brownian motion is its (random) starting position.

\begin{algorithm}[H]
    \caption{Simulating a Brownian Motion}
    \label{alg:SimulatingBrownianMotion}
    \begin{algorithmic}[1]
        \State Simulate the starting position $B_0$.
        \State Simulate outcomes $W_{t_1}, \dots, W_{t_n}$ of a Wiener process at times $t_1, \dots, t_n$.
        \State Return $B_{t_i} \coloneqq B_0 + \mu t_i + \sigma W_{t_i}, \quad i = 1, \dots, n$ as the outcomes of the Brownian motion at times $t_1, \dots, t_n$.
    \end{algorithmic}
\end{algorithm}

Multidimensional Brownian motions are likewise obtained from an affine transformation $\bB_t = \bB_0 + \bmu t + \bsig \bW_t$ of a multidimensional Wiener process, where $\bsig$ is a matrix.

\begin{defn}{Multidimensional Wiener Process}{}
Let $(W_{t,i}, t \ge 0)$, $i = 1, \dots, d$ be independent Wiener processes and let $\v{W}_t \coloneqq [W_{t,1}, \dots, W_{t,d}]^{\T}$. The process $(\v{W}_t, t \ge 0)$ is called a \textit{$d$-dimensional Wiener process}.
\end{defn}

If $(W_t)$ is a $d$-dimensional Wiener process, then $W_t$ has a $\mathcal{N}(\mathbf{0}, t \mathbf{I}_d)$ multivariate normal distribution, where $\mathbf{I}_d$ denotes the $d$-dimensional identity matrix.

\begin{figure}[H]
    \centering
    %% Creator: Matplotlib, PGF backend
%%
%% To include the figure in your LaTeX document, write
%%   \input{<filename>.pgf}
%%
%% Make sure the required packages are loaded in your preamble
%%   \usepackage{pgf}
%%
%% Also ensure that all the required font packages are loaded; for instance,
%% the lmodern package is sometimes necessary when using math font.
%%   \usepackage{lmodern}
%%
%% Figures using additional raster images can only be included by \input if
%% they are in the same directory as the main LaTeX file. For loading figures
%% from other directories you can use the `import` package
%%   \usepackage{import}
%%
%% and then include the figures with
%%   \import{<path to file>}{<filename>.pgf}
%%
%% Matplotlib used the following preamble
%%   \def\mathdefault#1{#1}
%%   \everymath=\expandafter{\the\everymath\displaystyle}
%%   \IfFileExists{scrextend.sty}{
%%     \usepackage[fontsize=10.000000pt]{scrextend}
%%   }{
%%     \renewcommand{\normalsize}{\fontsize{10.000000}{12.000000}\selectfont}
%%     \normalsize
%%   }
%%   
%%   \makeatletter\@ifpackageloaded{underscore}{}{\usepackage[strings]{underscore}}\makeatother
%%
\begingroup%
\makeatletter%
\begin{pgfpicture}%
\pgfpathrectangle{\pgfpointorigin}{\pgfqpoint{4.374741in}{4.104000in}}%
\pgfusepath{use as bounding box, clip}%
\begin{pgfscope}%
\pgfsetbuttcap%
\pgfsetmiterjoin%
\definecolor{currentfill}{rgb}{1.000000,1.000000,1.000000}%
\pgfsetfillcolor{currentfill}%
\pgfsetlinewidth{0.000000pt}%
\definecolor{currentstroke}{rgb}{1.000000,1.000000,1.000000}%
\pgfsetstrokecolor{currentstroke}%
\pgfsetdash{}{0pt}%
\pgfpathmoveto{\pgfqpoint{0.000000in}{0.000000in}}%
\pgfpathlineto{\pgfqpoint{4.374741in}{0.000000in}}%
\pgfpathlineto{\pgfqpoint{4.374741in}{4.104000in}}%
\pgfpathlineto{\pgfqpoint{0.000000in}{4.104000in}}%
\pgfpathlineto{\pgfqpoint{0.000000in}{0.000000in}}%
\pgfpathclose%
\pgfusepath{fill}%
\end{pgfscope}%
\begin{pgfscope}%
\pgfsetbuttcap%
\pgfsetmiterjoin%
\definecolor{currentfill}{rgb}{1.000000,1.000000,1.000000}%
\pgfsetfillcolor{currentfill}%
\pgfsetlinewidth{0.000000pt}%
\definecolor{currentstroke}{rgb}{0.000000,0.000000,0.000000}%
\pgfsetstrokecolor{currentstroke}%
\pgfsetstrokeopacity{0.000000}%
\pgfsetdash{}{0pt}%
\pgfpathmoveto{\pgfqpoint{0.050000in}{0.050000in}}%
\pgfpathlineto{\pgfqpoint{4.054000in}{0.050000in}}%
\pgfpathlineto{\pgfqpoint{4.054000in}{4.054000in}}%
\pgfpathlineto{\pgfqpoint{0.050000in}{4.054000in}}%
\pgfpathlineto{\pgfqpoint{0.050000in}{0.050000in}}%
\pgfpathclose%
\pgfusepath{fill}%
\end{pgfscope}%
\begin{pgfscope}%
\pgfsetbuttcap%
\pgfsetmiterjoin%
\pgfsetlinewidth{0.000000pt}%
\definecolor{currentstroke}{rgb}{1.000000,1.000000,1.000000}%
\pgfsetstrokecolor{currentstroke}%
\pgfsetstrokeopacity{0.000000}%
\pgfsetdash{}{0pt}%
\pgfpathmoveto{\pgfqpoint{2.106108in}{2.009579in}}%
\pgfpathlineto{\pgfqpoint{0.317931in}{1.246634in}}%
\pgfpathlineto{\pgfqpoint{0.249158in}{2.998637in}}%
\pgfpathlineto{\pgfqpoint{2.106108in}{3.662401in}}%
\pgfusepath{}%
\end{pgfscope}%
\begin{pgfscope}%
\pgfsetbuttcap%
\pgfsetmiterjoin%
\pgfsetlinewidth{0.000000pt}%
\definecolor{currentstroke}{rgb}{1.000000,1.000000,1.000000}%
\pgfsetstrokecolor{currentstroke}%
\pgfsetstrokeopacity{0.000000}%
\pgfsetdash{}{0pt}%
\pgfpathmoveto{\pgfqpoint{3.894286in}{1.246634in}}%
\pgfpathlineto{\pgfqpoint{2.106108in}{2.009579in}}%
\pgfpathlineto{\pgfqpoint{2.106108in}{3.662401in}}%
\pgfpathlineto{\pgfqpoint{3.963058in}{2.998637in}}%
\pgfusepath{}%
\end{pgfscope}%
\begin{pgfscope}%
\pgfsetbuttcap%
\pgfsetmiterjoin%
\pgfsetlinewidth{0.000000pt}%
\definecolor{currentstroke}{rgb}{1.000000,1.000000,1.000000}%
\pgfsetstrokecolor{currentstroke}%
\pgfsetstrokeopacity{0.000000}%
\pgfsetdash{}{0pt}%
\pgfpathmoveto{\pgfqpoint{3.894286in}{1.246634in}}%
\pgfpathlineto{\pgfqpoint{2.106108in}{2.009579in}}%
\pgfpathlineto{\pgfqpoint{0.317931in}{1.246634in}}%
\pgfpathlineto{\pgfqpoint{2.106108in}{0.360184in}}%
\pgfusepath{}%
\end{pgfscope}%
\begin{pgfscope}%
\pgfsetbuttcap%
\pgfsetroundjoin%
\pgfsetlinewidth{0.501875pt}%
\definecolor{currentstroke}{rgb}{0.800000,0.800000,0.800000}%
\pgfsetstrokecolor{currentstroke}%
\pgfsetdash{}{0pt}%
\pgfpathmoveto{\pgfqpoint{1.887012in}{0.468796in}}%
\pgfpathlineto{\pgfqpoint{3.676268in}{1.339654in}}%
\pgfpathlineto{\pgfqpoint{3.736123in}{3.079754in}}%
\pgfusepath{stroke}%
\end{pgfscope}%
\begin{pgfscope}%
\pgfsetbuttcap%
\pgfsetroundjoin%
\pgfsetlinewidth{0.501875pt}%
\definecolor{currentstroke}{rgb}{0.800000,0.800000,0.800000}%
\pgfsetstrokecolor{currentstroke}%
\pgfsetdash{}{0pt}%
\pgfpathmoveto{\pgfqpoint{1.603104in}{0.609537in}}%
\pgfpathlineto{\pgfqpoint{3.393312in}{1.460380in}}%
\pgfpathlineto{\pgfqpoint{3.441814in}{3.184955in}}%
\pgfusepath{stroke}%
\end{pgfscope}%
\begin{pgfscope}%
\pgfsetbuttcap%
\pgfsetroundjoin%
\pgfsetlinewidth{0.501875pt}%
\definecolor{currentstroke}{rgb}{0.800000,0.800000,0.800000}%
\pgfsetstrokecolor{currentstroke}%
\pgfsetdash{}{0pt}%
\pgfpathmoveto{\pgfqpoint{1.325928in}{0.746941in}}%
\pgfpathlineto{\pgfqpoint{3.116578in}{1.578452in}}%
\pgfpathlineto{\pgfqpoint{3.154218in}{3.287755in}}%
\pgfusepath{stroke}%
\end{pgfscope}%
\begin{pgfscope}%
\pgfsetbuttcap%
\pgfsetroundjoin%
\pgfsetlinewidth{0.501875pt}%
\definecolor{currentstroke}{rgb}{0.800000,0.800000,0.800000}%
\pgfsetstrokecolor{currentstroke}%
\pgfsetdash{}{0pt}%
\pgfpathmoveto{\pgfqpoint{1.055248in}{0.881125in}}%
\pgfpathlineto{\pgfqpoint{2.845863in}{1.693955in}}%
\pgfpathlineto{\pgfqpoint{2.873109in}{3.388237in}}%
\pgfusepath{stroke}%
\end{pgfscope}%
\begin{pgfscope}%
\pgfsetbuttcap%
\pgfsetroundjoin%
\pgfsetlinewidth{0.501875pt}%
\definecolor{currentstroke}{rgb}{0.800000,0.800000,0.800000}%
\pgfsetstrokecolor{currentstroke}%
\pgfsetdash{}{0pt}%
\pgfpathmoveto{\pgfqpoint{0.790838in}{1.012201in}}%
\pgfpathlineto{\pgfqpoint{2.580974in}{1.806973in}}%
\pgfpathlineto{\pgfqpoint{2.598268in}{3.486479in}}%
\pgfusepath{stroke}%
\end{pgfscope}%
\begin{pgfscope}%
\pgfsetbuttcap%
\pgfsetroundjoin%
\pgfsetlinewidth{0.501875pt}%
\definecolor{currentstroke}{rgb}{0.800000,0.800000,0.800000}%
\pgfsetstrokecolor{currentstroke}%
\pgfsetdash{}{0pt}%
\pgfpathmoveto{\pgfqpoint{0.532482in}{1.140275in}}%
\pgfpathlineto{\pgfqpoint{2.321724in}{1.917584in}}%
\pgfpathlineto{\pgfqpoint{2.329490in}{3.582553in}}%
\pgfusepath{stroke}%
\end{pgfscope}%
\begin{pgfscope}%
\pgfsetbuttcap%
\pgfsetroundjoin%
\pgfsetlinewidth{0.501875pt}%
\definecolor{currentstroke}{rgb}{0.800000,0.800000,0.800000}%
\pgfsetstrokecolor{currentstroke}%
\pgfsetdash{}{0pt}%
\pgfpathmoveto{\pgfqpoint{2.104510in}{0.360976in}}%
\pgfpathlineto{\pgfqpoint{3.892696in}{1.247313in}}%
\pgfpathlineto{\pgfqpoint{3.961403in}{2.999228in}}%
\pgfusepath{stroke}%
\end{pgfscope}%
\begin{pgfscope}%
\pgfsetbuttcap%
\pgfsetroundjoin%
\pgfsetlinewidth{0.501875pt}%
\definecolor{currentstroke}{rgb}{0.800000,0.800000,0.800000}%
\pgfsetstrokecolor{currentstroke}%
\pgfsetdash{}{0pt}%
\pgfpathmoveto{\pgfqpoint{2.031568in}{0.397136in}}%
\pgfpathlineto{\pgfqpoint{3.820146in}{1.278267in}}%
\pgfpathlineto{\pgfqpoint{3.885869in}{3.026228in}}%
\pgfusepath{stroke}%
\end{pgfscope}%
\begin{pgfscope}%
\pgfsetbuttcap%
\pgfsetroundjoin%
\pgfsetlinewidth{0.501875pt}%
\definecolor{currentstroke}{rgb}{0.800000,0.800000,0.800000}%
\pgfsetstrokecolor{currentstroke}%
\pgfsetdash{}{0pt}%
\pgfpathmoveto{\pgfqpoint{1.959070in}{0.433075in}}%
\pgfpathlineto{\pgfqpoint{3.748004in}{1.309047in}}%
\pgfpathlineto{\pgfqpoint{3.810777in}{3.053069in}}%
\pgfusepath{stroke}%
\end{pgfscope}%
\begin{pgfscope}%
\pgfsetbuttcap%
\pgfsetroundjoin%
\pgfsetlinewidth{0.501875pt}%
\definecolor{currentstroke}{rgb}{0.800000,0.800000,0.800000}%
\pgfsetstrokecolor{currentstroke}%
\pgfsetdash{}{0pt}%
\pgfpathmoveto{\pgfqpoint{1.815390in}{0.504301in}}%
\pgfpathlineto{\pgfqpoint{3.604934in}{1.370089in}}%
\pgfpathlineto{\pgfqpoint{3.661903in}{3.106284in}}%
\pgfusepath{stroke}%
\end{pgfscope}%
\begin{pgfscope}%
\pgfsetbuttcap%
\pgfsetroundjoin%
\pgfsetlinewidth{0.501875pt}%
\definecolor{currentstroke}{rgb}{0.800000,0.800000,0.800000}%
\pgfsetstrokecolor{currentstroke}%
\pgfsetdash{}{0pt}%
\pgfpathmoveto{\pgfqpoint{1.744201in}{0.539591in}}%
\pgfpathlineto{\pgfqpoint{3.533999in}{1.400354in}}%
\pgfpathlineto{\pgfqpoint{3.588115in}{3.132660in}}%
\pgfusepath{stroke}%
\end{pgfscope}%
\begin{pgfscope}%
\pgfsetbuttcap%
\pgfsetroundjoin%
\pgfsetlinewidth{0.501875pt}%
\definecolor{currentstroke}{rgb}{0.800000,0.800000,0.800000}%
\pgfsetstrokecolor{currentstroke}%
\pgfsetdash{}{0pt}%
\pgfpathmoveto{\pgfqpoint{1.673440in}{0.574670in}}%
\pgfpathlineto{\pgfqpoint{3.463459in}{1.430451in}}%
\pgfpathlineto{\pgfqpoint{3.514753in}{3.158883in}}%
\pgfusepath{stroke}%
\end{pgfscope}%
\begin{pgfscope}%
\pgfsetbuttcap%
\pgfsetroundjoin%
\pgfsetlinewidth{0.501875pt}%
\definecolor{currentstroke}{rgb}{0.800000,0.800000,0.800000}%
\pgfsetstrokecolor{currentstroke}%
\pgfsetdash{}{0pt}%
\pgfpathmoveto{\pgfqpoint{1.533188in}{0.644196in}}%
\pgfpathlineto{\pgfqpoint{3.323553in}{1.490143in}}%
\pgfpathlineto{\pgfqpoint{3.369295in}{3.210876in}}%
\pgfusepath{stroke}%
\end{pgfscope}%
\begin{pgfscope}%
\pgfsetbuttcap%
\pgfsetroundjoin%
\pgfsetlinewidth{0.501875pt}%
\definecolor{currentstroke}{rgb}{0.800000,0.800000,0.800000}%
\pgfsetstrokecolor{currentstroke}%
\pgfsetdash{}{0pt}%
\pgfpathmoveto{\pgfqpoint{1.463690in}{0.678649in}}%
\pgfpathlineto{\pgfqpoint{3.254180in}{1.519742in}}%
\pgfpathlineto{\pgfqpoint{3.297191in}{3.236650in}}%
\pgfusepath{stroke}%
\end{pgfscope}%
\begin{pgfscope}%
\pgfsetbuttcap%
\pgfsetroundjoin%
\pgfsetlinewidth{0.501875pt}%
\definecolor{currentstroke}{rgb}{0.800000,0.800000,0.800000}%
\pgfsetstrokecolor{currentstroke}%
\pgfsetdash{}{0pt}%
\pgfpathmoveto{\pgfqpoint{1.394604in}{0.712897in}}%
\pgfpathlineto{\pgfqpoint{3.185189in}{1.549178in}}%
\pgfpathlineto{\pgfqpoint{3.225500in}{3.262276in}}%
\pgfusepath{stroke}%
\end{pgfscope}%
\begin{pgfscope}%
\pgfsetbuttcap%
\pgfsetroundjoin%
\pgfsetlinewidth{0.501875pt}%
\definecolor{currentstroke}{rgb}{0.800000,0.800000,0.800000}%
\pgfsetstrokecolor{currentstroke}%
\pgfsetdash{}{0pt}%
\pgfpathmoveto{\pgfqpoint{1.257658in}{0.780785in}}%
\pgfpathlineto{\pgfqpoint{3.048342in}{1.607565in}}%
\pgfpathlineto{\pgfqpoint{3.083342in}{3.313090in}}%
\pgfusepath{stroke}%
\end{pgfscope}%
\begin{pgfscope}%
\pgfsetbuttcap%
\pgfsetroundjoin%
\pgfsetlinewidth{0.501875pt}%
\definecolor{currentstroke}{rgb}{0.800000,0.800000,0.800000}%
\pgfsetstrokecolor{currentstroke}%
\pgfsetdash{}{0pt}%
\pgfpathmoveto{\pgfqpoint{1.189791in}{0.814429in}}%
\pgfpathlineto{\pgfqpoint{2.980480in}{1.636519in}}%
\pgfpathlineto{\pgfqpoint{3.012867in}{3.338281in}}%
\pgfusepath{stroke}%
\end{pgfscope}%
\begin{pgfscope}%
\pgfsetbuttcap%
\pgfsetroundjoin%
\pgfsetlinewidth{0.501875pt}%
\definecolor{currentstroke}{rgb}{0.800000,0.800000,0.800000}%
\pgfsetstrokecolor{currentstroke}%
\pgfsetdash{}{0pt}%
\pgfpathmoveto{\pgfqpoint{1.122322in}{0.847875in}}%
\pgfpathlineto{\pgfqpoint{2.912988in}{1.665315in}}%
\pgfpathlineto{\pgfqpoint{2.942790in}{3.363330in}}%
\pgfusepath{stroke}%
\end{pgfscope}%
\begin{pgfscope}%
\pgfsetbuttcap%
\pgfsetroundjoin%
\pgfsetlinewidth{0.501875pt}%
\definecolor{currentstroke}{rgb}{0.800000,0.800000,0.800000}%
\pgfsetstrokecolor{currentstroke}%
\pgfsetdash{}{0pt}%
\pgfpathmoveto{\pgfqpoint{0.988566in}{0.914181in}}%
\pgfpathlineto{\pgfqpoint{2.779102in}{1.722439in}}%
\pgfpathlineto{\pgfqpoint{2.803819in}{3.413005in}}%
\pgfusepath{stroke}%
\end{pgfscope}%
\begin{pgfscope}%
\pgfsetbuttcap%
\pgfsetroundjoin%
\pgfsetlinewidth{0.501875pt}%
\definecolor{currentstroke}{rgb}{0.800000,0.800000,0.800000}%
\pgfsetstrokecolor{currentstroke}%
\pgfsetdash{}{0pt}%
\pgfpathmoveto{\pgfqpoint{0.922273in}{0.947045in}}%
\pgfpathlineto{\pgfqpoint{2.712702in}{1.750769in}}%
\pgfpathlineto{\pgfqpoint{2.734918in}{3.437634in}}%
\pgfusepath{stroke}%
\end{pgfscope}%
\begin{pgfscope}%
\pgfsetbuttcap%
\pgfsetroundjoin%
\pgfsetlinewidth{0.501875pt}%
\definecolor{currentstroke}{rgb}{0.800000,0.800000,0.800000}%
\pgfsetstrokecolor{currentstroke}%
\pgfsetdash{}{0pt}%
\pgfpathmoveto{\pgfqpoint{0.856365in}{0.979717in}}%
\pgfpathlineto{\pgfqpoint{2.646660in}{1.778947in}}%
\pgfpathlineto{\pgfqpoint{2.666402in}{3.462124in}}%
\pgfusepath{stroke}%
\end{pgfscope}%
\begin{pgfscope}%
\pgfsetbuttcap%
\pgfsetroundjoin%
\pgfsetlinewidth{0.501875pt}%
\definecolor{currentstroke}{rgb}{0.800000,0.800000,0.800000}%
\pgfsetstrokecolor{currentstroke}%
\pgfsetdash{}{0pt}%
\pgfpathmoveto{\pgfqpoint{0.725690in}{1.044497in}}%
\pgfpathlineto{\pgfqpoint{2.515640in}{1.834848in}}%
\pgfpathlineto{\pgfqpoint{2.530513in}{3.510698in}}%
\pgfusepath{stroke}%
\end{pgfscope}%
\begin{pgfscope}%
\pgfsetbuttcap%
\pgfsetroundjoin%
\pgfsetlinewidth{0.501875pt}%
\definecolor{currentstroke}{rgb}{0.800000,0.800000,0.800000}%
\pgfsetstrokecolor{currentstroke}%
\pgfsetdash{}{0pt}%
\pgfpathmoveto{\pgfqpoint{0.660916in}{1.076607in}}%
\pgfpathlineto{\pgfqpoint{2.450655in}{1.862574in}}%
\pgfpathlineto{\pgfqpoint{2.463134in}{3.534782in}}%
\pgfusepath{stroke}%
\end{pgfscope}%
\begin{pgfscope}%
\pgfsetbuttcap%
\pgfsetroundjoin%
\pgfsetlinewidth{0.501875pt}%
\definecolor{currentstroke}{rgb}{0.800000,0.800000,0.800000}%
\pgfsetstrokecolor{currentstroke}%
\pgfsetdash{}{0pt}%
\pgfpathmoveto{\pgfqpoint{0.596515in}{1.108532in}}%
\pgfpathlineto{\pgfqpoint{2.386017in}{1.890153in}}%
\pgfpathlineto{\pgfqpoint{2.396127in}{3.558734in}}%
\pgfusepath{stroke}%
\end{pgfscope}%
\begin{pgfscope}%
\pgfsetbuttcap%
\pgfsetroundjoin%
\pgfsetlinewidth{0.501875pt}%
\definecolor{currentstroke}{rgb}{0.800000,0.800000,0.800000}%
\pgfsetstrokecolor{currentstroke}%
\pgfsetdash{}{0pt}%
\pgfpathmoveto{\pgfqpoint{0.468815in}{1.171837in}}%
\pgfpathlineto{\pgfqpoint{2.257771in}{1.944870in}}%
\pgfpathlineto{\pgfqpoint{2.263219in}{3.606242in}}%
\pgfusepath{stroke}%
\end{pgfscope}%
\begin{pgfscope}%
\pgfsetbuttcap%
\pgfsetroundjoin%
\pgfsetlinewidth{0.501875pt}%
\definecolor{currentstroke}{rgb}{0.800000,0.800000,0.800000}%
\pgfsetstrokecolor{currentstroke}%
\pgfsetdash{}{0pt}%
\pgfpathmoveto{\pgfqpoint{0.405510in}{1.203219in}}%
\pgfpathlineto{\pgfqpoint{2.194157in}{1.972012in}}%
\pgfpathlineto{\pgfqpoint{2.197311in}{3.629800in}}%
\pgfusepath{stroke}%
\end{pgfscope}%
\begin{pgfscope}%
\pgfsetbuttcap%
\pgfsetroundjoin%
\pgfsetlinewidth{0.501875pt}%
\definecolor{currentstroke}{rgb}{0.800000,0.800000,0.800000}%
\pgfsetstrokecolor{currentstroke}%
\pgfsetdash{}{0pt}%
\pgfpathmoveto{\pgfqpoint{0.342565in}{1.234422in}}%
\pgfpathlineto{\pgfqpoint{2.130879in}{1.999010in}}%
\pgfpathlineto{\pgfqpoint{2.131764in}{3.653230in}}%
\pgfusepath{stroke}%
\end{pgfscope}%
\begin{pgfscope}%
\pgfsetbuttcap%
\pgfsetroundjoin%
\pgfsetlinewidth{0.501875pt}%
\definecolor{currentstroke}{rgb}{0.800000,0.800000,0.800000}%
\pgfsetstrokecolor{currentstroke}%
\pgfsetdash{}{0pt}%
\pgfpathmoveto{\pgfqpoint{2.076020in}{3.651646in}}%
\pgfpathlineto{\pgfqpoint{2.077058in}{1.997184in}}%
\pgfpathlineto{\pgfqpoint{3.865395in}{1.232313in}}%
\pgfusepath{stroke}%
\end{pgfscope}%
\begin{pgfscope}%
\pgfsetbuttcap%
\pgfsetroundjoin%
\pgfsetlinewidth{0.501875pt}%
\definecolor{currentstroke}{rgb}{0.800000,0.800000,0.800000}%
\pgfsetstrokecolor{currentstroke}%
\pgfsetdash{}{0pt}%
\pgfpathmoveto{\pgfqpoint{1.843033in}{3.568365in}}%
\pgfpathlineto{\pgfqpoint{1.852194in}{1.901244in}}%
\pgfpathlineto{\pgfqpoint{3.641594in}{1.121368in}}%
\pgfusepath{stroke}%
\end{pgfscope}%
\begin{pgfscope}%
\pgfsetbuttcap%
\pgfsetroundjoin%
\pgfsetlinewidth{0.501875pt}%
\definecolor{currentstroke}{rgb}{0.800000,0.800000,0.800000}%
\pgfsetstrokecolor{currentstroke}%
\pgfsetdash{}{0pt}%
\pgfpathmoveto{\pgfqpoint{1.605461in}{3.483445in}}%
\pgfpathlineto{\pgfqpoint{1.623060in}{1.803482in}}%
\pgfpathlineto{\pgfqpoint{3.413217in}{1.008155in}}%
\pgfusepath{stroke}%
\end{pgfscope}%
\begin{pgfscope}%
\pgfsetbuttcap%
\pgfsetroundjoin%
\pgfsetlinewidth{0.501875pt}%
\definecolor{currentstroke}{rgb}{0.800000,0.800000,0.800000}%
\pgfsetstrokecolor{currentstroke}%
\pgfsetdash{}{0pt}%
\pgfpathmoveto{\pgfqpoint{1.363168in}{3.396838in}}%
\pgfpathlineto{\pgfqpoint{1.389534in}{1.703845in}}%
\pgfpathlineto{\pgfqpoint{3.180125in}{0.892605in}}%
\pgfusepath{stroke}%
\end{pgfscope}%
\begin{pgfscope}%
\pgfsetbuttcap%
\pgfsetroundjoin%
\pgfsetlinewidth{0.501875pt}%
\definecolor{currentstroke}{rgb}{0.800000,0.800000,0.800000}%
\pgfsetstrokecolor{currentstroke}%
\pgfsetdash{}{0pt}%
\pgfpathmoveto{\pgfqpoint{1.116011in}{3.308492in}}%
\pgfpathlineto{\pgfqpoint{1.151488in}{1.602280in}}%
\pgfpathlineto{\pgfqpoint{2.942168in}{0.774643in}}%
\pgfusepath{stroke}%
\end{pgfscope}%
\begin{pgfscope}%
\pgfsetbuttcap%
\pgfsetroundjoin%
\pgfsetlinewidth{0.501875pt}%
\definecolor{currentstroke}{rgb}{0.800000,0.800000,0.800000}%
\pgfsetstrokecolor{currentstroke}%
\pgfsetdash{}{0pt}%
\pgfpathmoveto{\pgfqpoint{0.863841in}{3.218354in}}%
\pgfpathlineto{\pgfqpoint{0.908790in}{1.498731in}}%
\pgfpathlineto{\pgfqpoint{2.699194in}{0.654193in}}%
\pgfusepath{stroke}%
\end{pgfscope}%
\begin{pgfscope}%
\pgfsetbuttcap%
\pgfsetroundjoin%
\pgfsetlinewidth{0.501875pt}%
\definecolor{currentstroke}{rgb}{0.800000,0.800000,0.800000}%
\pgfsetstrokecolor{currentstroke}%
\pgfsetdash{}{0pt}%
\pgfpathmoveto{\pgfqpoint{0.606506in}{3.126370in}}%
\pgfpathlineto{\pgfqpoint{0.661301in}{1.393137in}}%
\pgfpathlineto{\pgfqpoint{2.451041in}{0.531177in}}%
\pgfusepath{stroke}%
\end{pgfscope}%
\begin{pgfscope}%
\pgfsetbuttcap%
\pgfsetroundjoin%
\pgfsetlinewidth{0.501875pt}%
\definecolor{currentstroke}{rgb}{0.800000,0.800000,0.800000}%
\pgfsetstrokecolor{currentstroke}%
\pgfsetdash{}{0pt}%
\pgfpathmoveto{\pgfqpoint{0.343844in}{3.032482in}}%
\pgfpathlineto{\pgfqpoint{0.408879in}{1.285438in}}%
\pgfpathlineto{\pgfqpoint{2.197543in}{0.405510in}}%
\pgfusepath{stroke}%
\end{pgfscope}%
\begin{pgfscope}%
\pgfsetbuttcap%
\pgfsetroundjoin%
\pgfsetlinewidth{0.501875pt}%
\definecolor{currentstroke}{rgb}{0.800000,0.800000,0.800000}%
\pgfsetstrokecolor{currentstroke}%
\pgfsetdash{}{0pt}%
\pgfpathmoveto{\pgfqpoint{2.018195in}{3.630977in}}%
\pgfpathlineto{\pgfqpoint{2.021235in}{1.973367in}}%
\pgfpathlineto{\pgfqpoint{3.809866in}{1.204785in}}%
\pgfusepath{stroke}%
\end{pgfscope}%
\begin{pgfscope}%
\pgfsetbuttcap%
\pgfsetroundjoin%
\pgfsetlinewidth{0.501875pt}%
\definecolor{currentstroke}{rgb}{0.800000,0.800000,0.800000}%
\pgfsetstrokecolor{currentstroke}%
\pgfsetdash{}{0pt}%
\pgfpathmoveto{\pgfqpoint{1.960091in}{3.610207in}}%
\pgfpathlineto{\pgfqpoint{1.965152in}{1.949438in}}%
\pgfpathlineto{\pgfqpoint{3.754058in}{1.177119in}}%
\pgfusepath{stroke}%
\end{pgfscope}%
\begin{pgfscope}%
\pgfsetbuttcap%
\pgfsetroundjoin%
\pgfsetlinewidth{0.501875pt}%
\definecolor{currentstroke}{rgb}{0.800000,0.800000,0.800000}%
\pgfsetstrokecolor{currentstroke}%
\pgfsetdash{}{0pt}%
\pgfpathmoveto{\pgfqpoint{1.901704in}{3.589337in}}%
\pgfpathlineto{\pgfqpoint{1.908805in}{1.925398in}}%
\pgfpathlineto{\pgfqpoint{3.697967in}{1.149314in}}%
\pgfusepath{stroke}%
\end{pgfscope}%
\begin{pgfscope}%
\pgfsetbuttcap%
\pgfsetroundjoin%
\pgfsetlinewidth{0.501875pt}%
\definecolor{currentstroke}{rgb}{0.800000,0.800000,0.800000}%
\pgfsetstrokecolor{currentstroke}%
\pgfsetdash{}{0pt}%
\pgfpathmoveto{\pgfqpoint{1.784075in}{3.547291in}}%
\pgfpathlineto{\pgfqpoint{1.795315in}{1.876976in}}%
\pgfpathlineto{\pgfqpoint{3.584934in}{1.093280in}}%
\pgfusepath{stroke}%
\end{pgfscope}%
\begin{pgfscope}%
\pgfsetbuttcap%
\pgfsetroundjoin%
\pgfsetlinewidth{0.501875pt}%
\definecolor{currentstroke}{rgb}{0.800000,0.800000,0.800000}%
\pgfsetstrokecolor{currentstroke}%
\pgfsetdash{}{0pt}%
\pgfpathmoveto{\pgfqpoint{1.724829in}{3.526113in}}%
\pgfpathlineto{\pgfqpoint{1.738168in}{1.852594in}}%
\pgfpathlineto{\pgfqpoint{3.527986in}{1.065049in}}%
\pgfusepath{stroke}%
\end{pgfscope}%
\begin{pgfscope}%
\pgfsetbuttcap%
\pgfsetroundjoin%
\pgfsetlinewidth{0.501875pt}%
\definecolor{currentstroke}{rgb}{0.800000,0.800000,0.800000}%
\pgfsetstrokecolor{currentstroke}%
\pgfsetdash{}{0pt}%
\pgfpathmoveto{\pgfqpoint{1.665292in}{3.504832in}}%
\pgfpathlineto{\pgfqpoint{1.680750in}{1.828096in}}%
\pgfpathlineto{\pgfqpoint{3.470748in}{1.036675in}}%
\pgfusepath{stroke}%
\end{pgfscope}%
\begin{pgfscope}%
\pgfsetbuttcap%
\pgfsetroundjoin%
\pgfsetlinewidth{0.501875pt}%
\definecolor{currentstroke}{rgb}{0.800000,0.800000,0.800000}%
\pgfsetstrokecolor{currentstroke}%
\pgfsetdash{}{0pt}%
\pgfpathmoveto{\pgfqpoint{1.545336in}{3.461954in}}%
\pgfpathlineto{\pgfqpoint{1.565095in}{1.778750in}}%
\pgfpathlineto{\pgfqpoint{3.355392in}{0.979490in}}%
\pgfusepath{stroke}%
\end{pgfscope}%
\begin{pgfscope}%
\pgfsetbuttcap%
\pgfsetroundjoin%
\pgfsetlinewidth{0.501875pt}%
\definecolor{currentstroke}{rgb}{0.800000,0.800000,0.800000}%
\pgfsetstrokecolor{currentstroke}%
\pgfsetdash{}{0pt}%
\pgfpathmoveto{\pgfqpoint{1.484914in}{3.440356in}}%
\pgfpathlineto{\pgfqpoint{1.506854in}{1.753901in}}%
\pgfpathlineto{\pgfqpoint{3.297270in}{0.950677in}}%
\pgfusepath{stroke}%
\end{pgfscope}%
\begin{pgfscope}%
\pgfsetbuttcap%
\pgfsetroundjoin%
\pgfsetlinewidth{0.501875pt}%
\definecolor{currentstroke}{rgb}{0.800000,0.800000,0.800000}%
\pgfsetstrokecolor{currentstroke}%
\pgfsetdash{}{0pt}%
\pgfpathmoveto{\pgfqpoint{1.424192in}{3.418651in}}%
\pgfpathlineto{\pgfqpoint{1.448334in}{1.728933in}}%
\pgfpathlineto{\pgfqpoint{3.238848in}{0.921715in}}%
\pgfusepath{stroke}%
\end{pgfscope}%
\begin{pgfscope}%
\pgfsetbuttcap%
\pgfsetroundjoin%
\pgfsetlinewidth{0.501875pt}%
\definecolor{currentstroke}{rgb}{0.800000,0.800000,0.800000}%
\pgfsetstrokecolor{currentstroke}%
\pgfsetdash{}{0pt}%
\pgfpathmoveto{\pgfqpoint{1.301840in}{3.374916in}}%
\pgfpathlineto{\pgfqpoint{1.330452in}{1.678637in}}%
\pgfpathlineto{\pgfqpoint{3.121098in}{0.863343in}}%
\pgfusepath{stroke}%
\end{pgfscope}%
\begin{pgfscope}%
\pgfsetbuttcap%
\pgfsetroundjoin%
\pgfsetlinewidth{0.501875pt}%
\definecolor{currentstroke}{rgb}{0.800000,0.800000,0.800000}%
\pgfsetstrokecolor{currentstroke}%
\pgfsetdash{}{0pt}%
\pgfpathmoveto{\pgfqpoint{1.240206in}{3.352885in}}%
\pgfpathlineto{\pgfqpoint{1.271084in}{1.653307in}}%
\pgfpathlineto{\pgfqpoint{3.061764in}{0.833930in}}%
\pgfusepath{stroke}%
\end{pgfscope}%
\begin{pgfscope}%
\pgfsetbuttcap%
\pgfsetroundjoin%
\pgfsetlinewidth{0.501875pt}%
\definecolor{currentstroke}{rgb}{0.800000,0.800000,0.800000}%
\pgfsetstrokecolor{currentstroke}%
\pgfsetdash{}{0pt}%
\pgfpathmoveto{\pgfqpoint{1.178264in}{3.330744in}}%
\pgfpathlineto{\pgfqpoint{1.211431in}{1.627855in}}%
\pgfpathlineto{\pgfqpoint{3.002122in}{0.804363in}}%
\pgfusepath{stroke}%
\end{pgfscope}%
\begin{pgfscope}%
\pgfsetbuttcap%
\pgfsetroundjoin%
\pgfsetlinewidth{0.501875pt}%
\definecolor{currentstroke}{rgb}{0.800000,0.800000,0.800000}%
\pgfsetstrokecolor{currentstroke}%
\pgfsetdash{}{0pt}%
\pgfpathmoveto{\pgfqpoint{1.053444in}{3.286127in}}%
\pgfpathlineto{\pgfqpoint{1.091255in}{1.576581in}}%
\pgfpathlineto{\pgfqpoint{2.881901in}{0.744767in}}%
\pgfusepath{stroke}%
\end{pgfscope}%
\begin{pgfscope}%
\pgfsetbuttcap%
\pgfsetroundjoin%
\pgfsetlinewidth{0.501875pt}%
\definecolor{currentstroke}{rgb}{0.800000,0.800000,0.800000}%
\pgfsetstrokecolor{currentstroke}%
\pgfsetdash{}{0pt}%
\pgfpathmoveto{\pgfqpoint{0.990562in}{3.263650in}}%
\pgfpathlineto{\pgfqpoint{1.030729in}{1.550757in}}%
\pgfpathlineto{\pgfqpoint{2.821318in}{0.714734in}}%
\pgfusepath{stroke}%
\end{pgfscope}%
\begin{pgfscope}%
\pgfsetbuttcap%
\pgfsetroundjoin%
\pgfsetlinewidth{0.501875pt}%
\definecolor{currentstroke}{rgb}{0.800000,0.800000,0.800000}%
\pgfsetstrokecolor{currentstroke}%
\pgfsetdash{}{0pt}%
\pgfpathmoveto{\pgfqpoint{0.927362in}{3.241060in}}%
\pgfpathlineto{\pgfqpoint{0.969908in}{1.524807in}}%
\pgfpathlineto{\pgfqpoint{2.760416in}{0.684543in}}%
\pgfusepath{stroke}%
\end{pgfscope}%
\begin{pgfscope}%
\pgfsetbuttcap%
\pgfsetroundjoin%
\pgfsetlinewidth{0.501875pt}%
\definecolor{currentstroke}{rgb}{0.800000,0.800000,0.800000}%
\pgfsetstrokecolor{currentstroke}%
\pgfsetdash{}{0pt}%
\pgfpathmoveto{\pgfqpoint{0.799998in}{3.195533in}}%
\pgfpathlineto{\pgfqpoint{0.847372in}{1.472526in}}%
\pgfpathlineto{\pgfqpoint{2.637647in}{0.623683in}}%
\pgfusepath{stroke}%
\end{pgfscope}%
\begin{pgfscope}%
\pgfsetbuttcap%
\pgfsetroundjoin%
\pgfsetlinewidth{0.501875pt}%
\definecolor{currentstroke}{rgb}{0.800000,0.800000,0.800000}%
\pgfsetstrokecolor{currentstroke}%
\pgfsetdash{}{0pt}%
\pgfpathmoveto{\pgfqpoint{0.735829in}{3.172596in}}%
\pgfpathlineto{\pgfqpoint{0.785653in}{1.446193in}}%
\pgfpathlineto{\pgfqpoint{2.575775in}{0.593011in}}%
\pgfusepath{stroke}%
\end{pgfscope}%
\begin{pgfscope}%
\pgfsetbuttcap%
\pgfsetroundjoin%
\pgfsetlinewidth{0.501875pt}%
\definecolor{currentstroke}{rgb}{0.800000,0.800000,0.800000}%
\pgfsetstrokecolor{currentstroke}%
\pgfsetdash{}{0pt}%
\pgfpathmoveto{\pgfqpoint{0.671333in}{3.149542in}}%
\pgfpathlineto{\pgfqpoint{0.723630in}{1.419730in}}%
\pgfpathlineto{\pgfqpoint{2.513574in}{0.562176in}}%
\pgfusepath{stroke}%
\end{pgfscope}%
\begin{pgfscope}%
\pgfsetbuttcap%
\pgfsetroundjoin%
\pgfsetlinewidth{0.501875pt}%
\definecolor{currentstroke}{rgb}{0.800000,0.800000,0.800000}%
\pgfsetstrokecolor{currentstroke}%
\pgfsetdash{}{0pt}%
\pgfpathmoveto{\pgfqpoint{0.541346in}{3.103079in}}%
\pgfpathlineto{\pgfqpoint{0.598664in}{1.366412in}}%
\pgfpathlineto{\pgfqpoint{2.388174in}{0.500012in}}%
\pgfusepath{stroke}%
\end{pgfscope}%
\begin{pgfscope}%
\pgfsetbuttcap%
\pgfsetroundjoin%
\pgfsetlinewidth{0.501875pt}%
\definecolor{currentstroke}{rgb}{0.800000,0.800000,0.800000}%
\pgfsetstrokecolor{currentstroke}%
\pgfsetdash{}{0pt}%
\pgfpathmoveto{\pgfqpoint{0.475851in}{3.079668in}}%
\pgfpathlineto{\pgfqpoint{0.535716in}{1.339555in}}%
\pgfpathlineto{\pgfqpoint{2.324971in}{0.468680in}}%
\pgfusepath{stroke}%
\end{pgfscope}%
\begin{pgfscope}%
\pgfsetbuttcap%
\pgfsetroundjoin%
\pgfsetlinewidth{0.501875pt}%
\definecolor{currentstroke}{rgb}{0.800000,0.800000,0.800000}%
\pgfsetstrokecolor{currentstroke}%
\pgfsetdash{}{0pt}%
\pgfpathmoveto{\pgfqpoint{0.410018in}{3.056136in}}%
\pgfpathlineto{\pgfqpoint{0.472455in}{1.312564in}}%
\pgfpathlineto{\pgfqpoint{2.261428in}{0.437180in}}%
\pgfusepath{stroke}%
\end{pgfscope}%
\begin{pgfscope}%
\pgfsetbuttcap%
\pgfsetroundjoin%
\pgfsetlinewidth{0.501875pt}%
\definecolor{currentstroke}{rgb}{0.800000,0.800000,0.800000}%
\pgfsetstrokecolor{currentstroke}%
\pgfsetdash{}{0pt}%
\pgfpathmoveto{\pgfqpoint{0.277327in}{3.008705in}}%
\pgfpathlineto{\pgfqpoint{0.344985in}{1.258177in}}%
\pgfpathlineto{\pgfqpoint{2.133312in}{0.373670in}}%
\pgfusepath{stroke}%
\end{pgfscope}%
\begin{pgfscope}%
\pgfsetbuttcap%
\pgfsetroundjoin%
\pgfsetlinewidth{0.501875pt}%
\definecolor{currentstroke}{rgb}{0.800000,0.800000,0.800000}%
\pgfsetstrokecolor{currentstroke}%
\pgfsetdash{}{0pt}%
\pgfpathmoveto{\pgfqpoint{3.899661in}{1.383571in}}%
\pgfpathlineto{\pgfqpoint{2.106108in}{2.139082in}}%
\pgfpathlineto{\pgfqpoint{0.312555in}{1.383571in}}%
\pgfusepath{stroke}%
\end{pgfscope}%
\begin{pgfscope}%
\pgfsetbuttcap%
\pgfsetroundjoin%
\pgfsetlinewidth{0.501875pt}%
\definecolor{currentstroke}{rgb}{0.800000,0.800000,0.800000}%
\pgfsetstrokecolor{currentstroke}%
\pgfsetdash{}{0pt}%
\pgfpathmoveto{\pgfqpoint{3.910097in}{1.649440in}}%
\pgfpathlineto{\pgfqpoint{2.106108in}{2.390366in}}%
\pgfpathlineto{\pgfqpoint{0.302119in}{1.649440in}}%
\pgfusepath{stroke}%
\end{pgfscope}%
\begin{pgfscope}%
\pgfsetbuttcap%
\pgfsetroundjoin%
\pgfsetlinewidth{0.501875pt}%
\definecolor{currentstroke}{rgb}{0.800000,0.800000,0.800000}%
\pgfsetstrokecolor{currentstroke}%
\pgfsetdash{}{0pt}%
\pgfpathmoveto{\pgfqpoint{3.920656in}{1.918422in}}%
\pgfpathlineto{\pgfqpoint{2.106108in}{2.644383in}}%
\pgfpathlineto{\pgfqpoint{0.291561in}{1.918422in}}%
\pgfusepath{stroke}%
\end{pgfscope}%
\begin{pgfscope}%
\pgfsetbuttcap%
\pgfsetroundjoin%
\pgfsetlinewidth{0.501875pt}%
\definecolor{currentstroke}{rgb}{0.800000,0.800000,0.800000}%
\pgfsetstrokecolor{currentstroke}%
\pgfsetdash{}{0pt}%
\pgfpathmoveto{\pgfqpoint{3.931338in}{2.190571in}}%
\pgfpathlineto{\pgfqpoint{2.106108in}{2.901178in}}%
\pgfpathlineto{\pgfqpoint{0.280878in}{2.190571in}}%
\pgfusepath{stroke}%
\end{pgfscope}%
\begin{pgfscope}%
\pgfsetbuttcap%
\pgfsetroundjoin%
\pgfsetlinewidth{0.501875pt}%
\definecolor{currentstroke}{rgb}{0.800000,0.800000,0.800000}%
\pgfsetstrokecolor{currentstroke}%
\pgfsetdash{}{0pt}%
\pgfpathmoveto{\pgfqpoint{3.942148in}{2.465943in}}%
\pgfpathlineto{\pgfqpoint{2.106108in}{3.160798in}}%
\pgfpathlineto{\pgfqpoint{0.270069in}{2.465943in}}%
\pgfusepath{stroke}%
\end{pgfscope}%
\begin{pgfscope}%
\pgfsetbuttcap%
\pgfsetroundjoin%
\pgfsetlinewidth{0.501875pt}%
\definecolor{currentstroke}{rgb}{0.800000,0.800000,0.800000}%
\pgfsetstrokecolor{currentstroke}%
\pgfsetdash{}{0pt}%
\pgfpathmoveto{\pgfqpoint{3.953086in}{2.744596in}}%
\pgfpathlineto{\pgfqpoint{2.106108in}{3.423290in}}%
\pgfpathlineto{\pgfqpoint{0.259130in}{2.744596in}}%
\pgfusepath{stroke}%
\end{pgfscope}%
\begin{pgfscope}%
\pgfsetrectcap%
\pgfsetroundjoin%
\pgfsetlinewidth{0.501875pt}%
\definecolor{currentstroke}{rgb}{1.000000,1.000000,1.000000}%
\pgfsetstrokecolor{currentstroke}%
\pgfsetstrokeopacity{0.000000}%
\pgfsetdash{}{0pt}%
\pgfpathmoveto{\pgfqpoint{0.317931in}{1.246634in}}%
\pgfpathlineto{\pgfqpoint{2.106108in}{0.360184in}}%
\pgfusepath{stroke}%
\end{pgfscope}%
\begin{pgfscope}%
\pgfsetrectcap%
\pgfsetroundjoin%
\pgfsetlinewidth{0.501875pt}%
\definecolor{currentstroke}{rgb}{0.000000,0.000000,0.000000}%
\pgfsetstrokecolor{currentstroke}%
\pgfsetdash{}{0pt}%
\pgfpathmoveto{\pgfqpoint{1.902464in}{0.476317in}}%
\pgfpathlineto{\pgfqpoint{1.856048in}{0.453725in}}%
\pgfusepath{stroke}%
\end{pgfscope}%
\begin{pgfscope}%
\definecolor{textcolor}{rgb}{0.000000,0.000000,0.000000}%
\pgfsetstrokecolor{textcolor}%
\pgfsetfillcolor{textcolor}%
\pgftext[x=1.739189in,y=0.260069in,,top]{\color{textcolor}{\rmfamily\fontsize{10.000000}{12.000000}\selectfont\catcode`\^=\active\def^{\ifmmode\sp\else\^{}\fi}\catcode`\%=\active\def%{\%}$\mathdefault{-0.4}$}}%
\end{pgfscope}%
\begin{pgfscope}%
\pgfsetrectcap%
\pgfsetroundjoin%
\pgfsetlinewidth{0.501875pt}%
\definecolor{currentstroke}{rgb}{0.000000,0.000000,0.000000}%
\pgfsetstrokecolor{currentstroke}%
\pgfsetdash{}{0pt}%
\pgfpathmoveto{\pgfqpoint{1.618551in}{0.616879in}}%
\pgfpathlineto{\pgfqpoint{1.572152in}{0.594826in}}%
\pgfusepath{stroke}%
\end{pgfscope}%
\begin{pgfscope}%
\definecolor{textcolor}{rgb}{0.000000,0.000000,0.000000}%
\pgfsetstrokecolor{textcolor}%
\pgfsetfillcolor{textcolor}%
\pgftext[x=1.456187in,y=0.403756in,,top]{\color{textcolor}{\rmfamily\fontsize{10.000000}{12.000000}\selectfont\catcode`\^=\active\def^{\ifmmode\sp\else\^{}\fi}\catcode`\%=\active\def%{\%}$\mathdefault{-0.2}$}}%
\end{pgfscope}%
\begin{pgfscope}%
\pgfsetrectcap%
\pgfsetroundjoin%
\pgfsetlinewidth{0.501875pt}%
\definecolor{currentstroke}{rgb}{0.000000,0.000000,0.000000}%
\pgfsetstrokecolor{currentstroke}%
\pgfsetdash{}{0pt}%
\pgfpathmoveto{\pgfqpoint{1.341365in}{0.754110in}}%
\pgfpathlineto{\pgfqpoint{1.294996in}{0.732577in}}%
\pgfusepath{stroke}%
\end{pgfscope}%
\begin{pgfscope}%
\definecolor{textcolor}{rgb}{0.000000,0.000000,0.000000}%
\pgfsetstrokecolor{textcolor}%
\pgfsetfillcolor{textcolor}%
\pgftext[x=1.179916in,y=0.544024in,,top]{\color{textcolor}{\rmfamily\fontsize{10.000000}{12.000000}\selectfont\catcode`\^=\active\def^{\ifmmode\sp\else\^{}\fi}\catcode`\%=\active\def%{\%}$\mathdefault{0.0}$}}%
\end{pgfscope}%
\begin{pgfscope}%
\pgfsetrectcap%
\pgfsetroundjoin%
\pgfsetlinewidth{0.501875pt}%
\definecolor{currentstroke}{rgb}{0.000000,0.000000,0.000000}%
\pgfsetstrokecolor{currentstroke}%
\pgfsetdash{}{0pt}%
\pgfpathmoveto{\pgfqpoint{1.070672in}{0.888127in}}%
\pgfpathlineto{\pgfqpoint{1.024343in}{0.867096in}}%
\pgfusepath{stroke}%
\end{pgfscope}%
\begin{pgfscope}%
\definecolor{textcolor}{rgb}{0.000000,0.000000,0.000000}%
\pgfsetstrokecolor{textcolor}%
\pgfsetfillcolor{textcolor}%
\pgftext[x=0.910138in,y=0.680996in,,top]{\color{textcolor}{\rmfamily\fontsize{10.000000}{12.000000}\selectfont\catcode`\^=\active\def^{\ifmmode\sp\else\^{}\fi}\catcode`\%=\active\def%{\%}$\mathdefault{0.2}$}}%
\end{pgfscope}%
\begin{pgfscope}%
\pgfsetrectcap%
\pgfsetroundjoin%
\pgfsetlinewidth{0.501875pt}%
\definecolor{currentstroke}{rgb}{0.000000,0.000000,0.000000}%
\pgfsetstrokecolor{currentstroke}%
\pgfsetdash{}{0pt}%
\pgfpathmoveto{\pgfqpoint{0.806245in}{1.019041in}}%
\pgfpathlineto{\pgfqpoint{0.759967in}{0.998495in}}%
\pgfusepath{stroke}%
\end{pgfscope}%
\begin{pgfscope}%
\definecolor{textcolor}{rgb}{0.000000,0.000000,0.000000}%
\pgfsetstrokecolor{textcolor}%
\pgfsetfillcolor{textcolor}%
\pgftext[x=0.646629in,y=0.814786in,,top]{\color{textcolor}{\rmfamily\fontsize{10.000000}{12.000000}\selectfont\catcode`\^=\active\def^{\ifmmode\sp\else\^{}\fi}\catcode`\%=\active\def%{\%}$\mathdefault{0.4}$}}%
\end{pgfscope}%
\begin{pgfscope}%
\pgfsetrectcap%
\pgfsetroundjoin%
\pgfsetlinewidth{0.501875pt}%
\definecolor{currentstroke}{rgb}{0.000000,0.000000,0.000000}%
\pgfsetstrokecolor{currentstroke}%
\pgfsetdash{}{0pt}%
\pgfpathmoveto{\pgfqpoint{0.547869in}{1.146960in}}%
\pgfpathlineto{\pgfqpoint{0.501653in}{1.126882in}}%
\pgfusepath{stroke}%
\end{pgfscope}%
\begin{pgfscope}%
\definecolor{textcolor}{rgb}{0.000000,0.000000,0.000000}%
\pgfsetstrokecolor{textcolor}%
\pgfsetfillcolor{textcolor}%
\pgftext[x=0.389171in,y=0.945503in,,top]{\color{textcolor}{\rmfamily\fontsize{10.000000}{12.000000}\selectfont\catcode`\^=\active\def^{\ifmmode\sp\else\^{}\fi}\catcode`\%=\active\def%{\%}$\mathdefault{0.6}$}}%
\end{pgfscope}%
\begin{pgfscope}%
\pgfsetrectcap%
\pgfsetroundjoin%
\pgfsetlinewidth{0.501875pt}%
\definecolor{currentstroke}{rgb}{0.000000,0.000000,0.000000}%
\pgfsetstrokecolor{currentstroke}%
\pgfsetdash{}{0pt}%
\pgfpathmoveto{\pgfqpoint{2.119963in}{0.368636in}}%
\pgfpathlineto{\pgfqpoint{2.073543in}{0.345627in}}%
\pgfusepath{stroke}%
\end{pgfscope}%
\begin{pgfscope}%
\pgfsetrectcap%
\pgfsetroundjoin%
\pgfsetlinewidth{0.501875pt}%
\definecolor{currentstroke}{rgb}{0.000000,0.000000,0.000000}%
\pgfsetstrokecolor{currentstroke}%
\pgfsetdash{}{0pt}%
\pgfpathmoveto{\pgfqpoint{2.047021in}{0.404749in}}%
\pgfpathlineto{\pgfqpoint{2.000601in}{0.381880in}}%
\pgfusepath{stroke}%
\end{pgfscope}%
\begin{pgfscope}%
\pgfsetrectcap%
\pgfsetroundjoin%
\pgfsetlinewidth{0.501875pt}%
\definecolor{currentstroke}{rgb}{0.000000,0.000000,0.000000}%
\pgfsetstrokecolor{currentstroke}%
\pgfsetdash{}{0pt}%
\pgfpathmoveto{\pgfqpoint{1.974522in}{0.440642in}}%
\pgfpathlineto{\pgfqpoint{1.928104in}{0.417912in}}%
\pgfusepath{stroke}%
\end{pgfscope}%
\begin{pgfscope}%
\pgfsetrectcap%
\pgfsetroundjoin%
\pgfsetlinewidth{0.501875pt}%
\definecolor{currentstroke}{rgb}{0.000000,0.000000,0.000000}%
\pgfsetstrokecolor{currentstroke}%
\pgfsetdash{}{0pt}%
\pgfpathmoveto{\pgfqpoint{1.830841in}{0.511776in}}%
\pgfpathlineto{\pgfqpoint{1.784429in}{0.489321in}}%
\pgfusepath{stroke}%
\end{pgfscope}%
\begin{pgfscope}%
\pgfsetrectcap%
\pgfsetroundjoin%
\pgfsetlinewidth{0.501875pt}%
\definecolor{currentstroke}{rgb}{0.000000,0.000000,0.000000}%
\pgfsetstrokecolor{currentstroke}%
\pgfsetdash{}{0pt}%
\pgfpathmoveto{\pgfqpoint{1.759651in}{0.547022in}}%
\pgfpathlineto{\pgfqpoint{1.713242in}{0.524702in}}%
\pgfusepath{stroke}%
\end{pgfscope}%
\begin{pgfscope}%
\pgfsetrectcap%
\pgfsetroundjoin%
\pgfsetlinewidth{0.501875pt}%
\definecolor{currentstroke}{rgb}{0.000000,0.000000,0.000000}%
\pgfsetstrokecolor{currentstroke}%
\pgfsetdash{}{0pt}%
\pgfpathmoveto{\pgfqpoint{1.688889in}{0.582055in}}%
\pgfpathlineto{\pgfqpoint{1.642484in}{0.559870in}}%
\pgfusepath{stroke}%
\end{pgfscope}%
\begin{pgfscope}%
\pgfsetrectcap%
\pgfsetroundjoin%
\pgfsetlinewidth{0.501875pt}%
\definecolor{currentstroke}{rgb}{0.000000,0.000000,0.000000}%
\pgfsetstrokecolor{currentstroke}%
\pgfsetdash{}{0pt}%
\pgfpathmoveto{\pgfqpoint{1.548633in}{0.651494in}}%
\pgfpathlineto{\pgfqpoint{1.502240in}{0.629574in}}%
\pgfusepath{stroke}%
\end{pgfscope}%
\begin{pgfscope}%
\pgfsetrectcap%
\pgfsetroundjoin%
\pgfsetlinewidth{0.501875pt}%
\definecolor{currentstroke}{rgb}{0.000000,0.000000,0.000000}%
\pgfsetstrokecolor{currentstroke}%
\pgfsetdash{}{0pt}%
\pgfpathmoveto{\pgfqpoint{1.479132in}{0.685903in}}%
\pgfpathlineto{\pgfqpoint{1.432746in}{0.664113in}}%
\pgfusepath{stroke}%
\end{pgfscope}%
\begin{pgfscope}%
\pgfsetrectcap%
\pgfsetroundjoin%
\pgfsetlinewidth{0.501875pt}%
\definecolor{currentstroke}{rgb}{0.000000,0.000000,0.000000}%
\pgfsetstrokecolor{currentstroke}%
\pgfsetdash{}{0pt}%
\pgfpathmoveto{\pgfqpoint{1.410044in}{0.720108in}}%
\pgfpathlineto{\pgfqpoint{1.363666in}{0.698447in}}%
\pgfusepath{stroke}%
\end{pgfscope}%
\begin{pgfscope}%
\pgfsetrectcap%
\pgfsetroundjoin%
\pgfsetlinewidth{0.501875pt}%
\definecolor{currentstroke}{rgb}{0.000000,0.000000,0.000000}%
\pgfsetstrokecolor{currentstroke}%
\pgfsetdash{}{0pt}%
\pgfpathmoveto{\pgfqpoint{1.273092in}{0.787911in}}%
\pgfpathlineto{\pgfqpoint{1.226732in}{0.766506in}}%
\pgfusepath{stroke}%
\end{pgfscope}%
\begin{pgfscope}%
\pgfsetrectcap%
\pgfsetroundjoin%
\pgfsetlinewidth{0.501875pt}%
\definecolor{currentstroke}{rgb}{0.000000,0.000000,0.000000}%
\pgfsetstrokecolor{currentstroke}%
\pgfsetdash{}{0pt}%
\pgfpathmoveto{\pgfqpoint{1.205222in}{0.821513in}}%
\pgfpathlineto{\pgfqpoint{1.158871in}{0.800234in}}%
\pgfusepath{stroke}%
\end{pgfscope}%
\begin{pgfscope}%
\pgfsetrectcap%
\pgfsetroundjoin%
\pgfsetlinewidth{0.501875pt}%
\definecolor{currentstroke}{rgb}{0.000000,0.000000,0.000000}%
\pgfsetstrokecolor{currentstroke}%
\pgfsetdash{}{0pt}%
\pgfpathmoveto{\pgfqpoint{1.137749in}{0.854917in}}%
\pgfpathlineto{\pgfqpoint{1.091409in}{0.833763in}}%
\pgfusepath{stroke}%
\end{pgfscope}%
\begin{pgfscope}%
\pgfsetrectcap%
\pgfsetroundjoin%
\pgfsetlinewidth{0.501875pt}%
\definecolor{currentstroke}{rgb}{0.000000,0.000000,0.000000}%
\pgfsetstrokecolor{currentstroke}%
\pgfsetdash{}{0pt}%
\pgfpathmoveto{\pgfqpoint{1.003986in}{0.921142in}}%
\pgfpathlineto{\pgfqpoint{0.957669in}{0.900234in}}%
\pgfusepath{stroke}%
\end{pgfscope}%
\begin{pgfscope}%
\pgfsetrectcap%
\pgfsetroundjoin%
\pgfsetlinewidth{0.501875pt}%
\definecolor{currentstroke}{rgb}{0.000000,0.000000,0.000000}%
\pgfsetstrokecolor{currentstroke}%
\pgfsetdash{}{0pt}%
\pgfpathmoveto{\pgfqpoint{0.937689in}{0.953965in}}%
\pgfpathlineto{\pgfqpoint{0.891384in}{0.933179in}}%
\pgfusepath{stroke}%
\end{pgfscope}%
\begin{pgfscope}%
\pgfsetrectcap%
\pgfsetroundjoin%
\pgfsetlinewidth{0.501875pt}%
\definecolor{currentstroke}{rgb}{0.000000,0.000000,0.000000}%
\pgfsetstrokecolor{currentstroke}%
\pgfsetdash{}{0pt}%
\pgfpathmoveto{\pgfqpoint{0.871776in}{0.986597in}}%
\pgfpathlineto{\pgfqpoint{0.825485in}{0.965932in}}%
\pgfusepath{stroke}%
\end{pgfscope}%
\begin{pgfscope}%
\pgfsetrectcap%
\pgfsetroundjoin%
\pgfsetlinewidth{0.501875pt}%
\definecolor{currentstroke}{rgb}{0.000000,0.000000,0.000000}%
\pgfsetstrokecolor{currentstroke}%
\pgfsetdash{}{0pt}%
\pgfpathmoveto{\pgfqpoint{0.741092in}{1.051297in}}%
\pgfpathlineto{\pgfqpoint{0.694829in}{1.030870in}}%
\pgfusepath{stroke}%
\end{pgfscope}%
\begin{pgfscope}%
\pgfsetrectcap%
\pgfsetroundjoin%
\pgfsetlinewidth{0.501875pt}%
\definecolor{currentstroke}{rgb}{0.000000,0.000000,0.000000}%
\pgfsetstrokecolor{currentstroke}%
\pgfsetdash{}{0pt}%
\pgfpathmoveto{\pgfqpoint{0.676314in}{1.083368in}}%
\pgfpathlineto{\pgfqpoint{0.630066in}{1.063058in}}%
\pgfusepath{stroke}%
\end{pgfscope}%
\begin{pgfscope}%
\pgfsetrectcap%
\pgfsetroundjoin%
\pgfsetlinewidth{0.501875pt}%
\definecolor{currentstroke}{rgb}{0.000000,0.000000,0.000000}%
\pgfsetstrokecolor{currentstroke}%
\pgfsetdash{}{0pt}%
\pgfpathmoveto{\pgfqpoint{0.611907in}{1.115255in}}%
\pgfpathlineto{\pgfqpoint{0.565675in}{1.095062in}}%
\pgfusepath{stroke}%
\end{pgfscope}%
\begin{pgfscope}%
\pgfsetrectcap%
\pgfsetroundjoin%
\pgfsetlinewidth{0.501875pt}%
\definecolor{currentstroke}{rgb}{0.000000,0.000000,0.000000}%
\pgfsetstrokecolor{currentstroke}%
\pgfsetdash{}{0pt}%
\pgfpathmoveto{\pgfqpoint{0.484197in}{1.178483in}}%
\pgfpathlineto{\pgfqpoint{0.437997in}{1.158520in}}%
\pgfusepath{stroke}%
\end{pgfscope}%
\begin{pgfscope}%
\pgfsetrectcap%
\pgfsetroundjoin%
\pgfsetlinewidth{0.501875pt}%
\definecolor{currentstroke}{rgb}{0.000000,0.000000,0.000000}%
\pgfsetstrokecolor{currentstroke}%
\pgfsetdash{}{0pt}%
\pgfpathmoveto{\pgfqpoint{0.420886in}{1.209827in}}%
\pgfpathlineto{\pgfqpoint{0.374704in}{1.189977in}}%
\pgfusepath{stroke}%
\end{pgfscope}%
\begin{pgfscope}%
\pgfsetrectcap%
\pgfsetroundjoin%
\pgfsetlinewidth{0.501875pt}%
\definecolor{currentstroke}{rgb}{0.000000,0.000000,0.000000}%
\pgfsetstrokecolor{currentstroke}%
\pgfsetdash{}{0pt}%
\pgfpathmoveto{\pgfqpoint{0.357935in}{1.240994in}}%
\pgfpathlineto{\pgfqpoint{0.311770in}{1.221256in}}%
\pgfusepath{stroke}%
\end{pgfscope}%
\begin{pgfscope}%
\definecolor{textcolor}{rgb}{0.000000,0.000000,0.000000}%
\pgfsetstrokecolor{textcolor}%
\pgfsetfillcolor{textcolor}%
\pgftext[x=0.859772in,y=0.381299in,,]{\color{textcolor}{\rmfamily\fontsize{10.000000}{12.000000}\selectfont\catcode`\^=\active\def^{\ifmmode\sp\else\^{}\fi}\catcode`\%=\active\def%{\%}$x$}}%
\end{pgfscope}%
\begin{pgfscope}%
\pgfsetrectcap%
\pgfsetroundjoin%
\pgfsetlinewidth{0.501875pt}%
\definecolor{currentstroke}{rgb}{1.000000,1.000000,1.000000}%
\pgfsetstrokecolor{currentstroke}%
\pgfsetstrokeopacity{0.000000}%
\pgfsetdash{}{0pt}%
\pgfpathmoveto{\pgfqpoint{3.894286in}{1.246634in}}%
\pgfpathlineto{\pgfqpoint{2.106108in}{0.360184in}}%
\pgfusepath{stroke}%
\end{pgfscope}%
\begin{pgfscope}%
\pgfsetrectcap%
\pgfsetroundjoin%
\pgfsetlinewidth{0.501875pt}%
\definecolor{currentstroke}{rgb}{0.000000,0.000000,0.000000}%
\pgfsetstrokecolor{currentstroke}%
\pgfsetdash{}{0pt}%
\pgfpathmoveto{\pgfqpoint{3.850025in}{1.238886in}}%
\pgfpathlineto{\pgfqpoint{3.896191in}{1.219141in}}%
\pgfusepath{stroke}%
\end{pgfscope}%
\begin{pgfscope}%
\definecolor{textcolor}{rgb}{0.000000,0.000000,0.000000}%
\pgfsetstrokecolor{textcolor}%
\pgfsetfillcolor{textcolor}%
\pgftext[x=4.008051in,y=1.039434in,,top]{\color{textcolor}{\rmfamily\fontsize{10.000000}{12.000000}\selectfont\catcode`\^=\active\def^{\ifmmode\sp\else\^{}\fi}\catcode`\%=\active\def%{\%}$\mathdefault{-0.2}$}}%
\end{pgfscope}%
\begin{pgfscope}%
\pgfsetrectcap%
\pgfsetroundjoin%
\pgfsetlinewidth{0.501875pt}%
\definecolor{currentstroke}{rgb}{0.000000,0.000000,0.000000}%
\pgfsetstrokecolor{currentstroke}%
\pgfsetdash{}{0pt}%
\pgfpathmoveto{\pgfqpoint{3.626204in}{1.128075in}}%
\pgfpathlineto{\pgfqpoint{3.672430in}{1.107928in}}%
\pgfusepath{stroke}%
\end{pgfscope}%
\begin{pgfscope}%
\definecolor{textcolor}{rgb}{0.000000,0.000000,0.000000}%
\pgfsetstrokecolor{textcolor}%
\pgfsetfillcolor{textcolor}%
\pgftext[x=3.785039in,y=0.926206in,,top]{\color{textcolor}{\rmfamily\fontsize{10.000000}{12.000000}\selectfont\catcode`\^=\active\def^{\ifmmode\sp\else\^{}\fi}\catcode`\%=\active\def%{\%}$\mathdefault{0.0}$}}%
\end{pgfscope}%
\begin{pgfscope}%
\pgfsetrectcap%
\pgfsetroundjoin%
\pgfsetlinewidth{0.501875pt}%
\definecolor{currentstroke}{rgb}{0.000000,0.000000,0.000000}%
\pgfsetstrokecolor{currentstroke}%
\pgfsetdash{}{0pt}%
\pgfpathmoveto{\pgfqpoint{3.397810in}{1.015000in}}%
\pgfpathlineto{\pgfqpoint{3.444089in}{0.994440in}}%
\pgfusepath{stroke}%
\end{pgfscope}%
\begin{pgfscope}%
\definecolor{textcolor}{rgb}{0.000000,0.000000,0.000000}%
\pgfsetstrokecolor{textcolor}%
\pgfsetfillcolor{textcolor}%
\pgftext[x=3.557455in,y=0.810657in,,top]{\color{textcolor}{\rmfamily\fontsize{10.000000}{12.000000}\selectfont\catcode`\^=\active\def^{\ifmmode\sp\else\^{}\fi}\catcode`\%=\active\def%{\%}$\mathdefault{0.2}$}}%
\end{pgfscope}%
\begin{pgfscope}%
\pgfsetrectcap%
\pgfsetroundjoin%
\pgfsetlinewidth{0.501875pt}%
\definecolor{currentstroke}{rgb}{0.000000,0.000000,0.000000}%
\pgfsetstrokecolor{currentstroke}%
\pgfsetdash{}{0pt}%
\pgfpathmoveto{\pgfqpoint{3.164702in}{0.899592in}}%
\pgfpathlineto{\pgfqpoint{3.211027in}{0.878604in}}%
\pgfusepath{stroke}%
\end{pgfscope}%
\begin{pgfscope}%
\definecolor{textcolor}{rgb}{0.000000,0.000000,0.000000}%
\pgfsetstrokecolor{textcolor}%
\pgfsetfillcolor{textcolor}%
\pgftext[x=3.325156in,y=0.692714in,,top]{\color{textcolor}{\rmfamily\fontsize{10.000000}{12.000000}\selectfont\catcode`\^=\active\def^{\ifmmode\sp\else\^{}\fi}\catcode`\%=\active\def%{\%}$\mathdefault{0.4}$}}%
\end{pgfscope}%
\begin{pgfscope}%
\pgfsetrectcap%
\pgfsetroundjoin%
\pgfsetlinewidth{0.501875pt}%
\definecolor{currentstroke}{rgb}{0.000000,0.000000,0.000000}%
\pgfsetstrokecolor{currentstroke}%
\pgfsetdash{}{0pt}%
\pgfpathmoveto{\pgfqpoint{2.926733in}{0.781776in}}%
\pgfpathlineto{\pgfqpoint{2.973096in}{0.760348in}}%
\pgfusepath{stroke}%
\end{pgfscope}%
\begin{pgfscope}%
\definecolor{textcolor}{rgb}{0.000000,0.000000,0.000000}%
\pgfsetstrokecolor{textcolor}%
\pgfsetfillcolor{textcolor}%
\pgftext[x=3.087996in,y=0.572302in,,top]{\color{textcolor}{\rmfamily\fontsize{10.000000}{12.000000}\selectfont\catcode`\^=\active\def^{\ifmmode\sp\else\^{}\fi}\catcode`\%=\active\def%{\%}$\mathdefault{0.6}$}}%
\end{pgfscope}%
\begin{pgfscope}%
\pgfsetrectcap%
\pgfsetroundjoin%
\pgfsetlinewidth{0.501875pt}%
\definecolor{currentstroke}{rgb}{0.000000,0.000000,0.000000}%
\pgfsetstrokecolor{currentstroke}%
\pgfsetdash{}{0pt}%
\pgfpathmoveto{\pgfqpoint{2.683750in}{0.661478in}}%
\pgfpathlineto{\pgfqpoint{2.730140in}{0.639596in}}%
\pgfusepath{stroke}%
\end{pgfscope}%
\begin{pgfscope}%
\definecolor{textcolor}{rgb}{0.000000,0.000000,0.000000}%
\pgfsetstrokecolor{textcolor}%
\pgfsetfillcolor{textcolor}%
\pgftext[x=2.845819in,y=0.449344in,,top]{\color{textcolor}{\rmfamily\fontsize{10.000000}{12.000000}\selectfont\catcode`\^=\active\def^{\ifmmode\sp\else\^{}\fi}\catcode`\%=\active\def%{\%}$\mathdefault{0.8}$}}%
\end{pgfscope}%
\begin{pgfscope}%
\pgfsetrectcap%
\pgfsetroundjoin%
\pgfsetlinewidth{0.501875pt}%
\definecolor{currentstroke}{rgb}{0.000000,0.000000,0.000000}%
\pgfsetstrokecolor{currentstroke}%
\pgfsetdash{}{0pt}%
\pgfpathmoveto{\pgfqpoint{2.435591in}{0.538618in}}%
\pgfpathlineto{\pgfqpoint{2.482001in}{0.516266in}}%
\pgfusepath{stroke}%
\end{pgfscope}%
\begin{pgfscope}%
\definecolor{textcolor}{rgb}{0.000000,0.000000,0.000000}%
\pgfsetstrokecolor{textcolor}%
\pgfsetfillcolor{textcolor}%
\pgftext[x=2.598465in,y=0.323757in,,top]{\color{textcolor}{\rmfamily\fontsize{10.000000}{12.000000}\selectfont\catcode`\^=\active\def^{\ifmmode\sp\else\^{}\fi}\catcode`\%=\active\def%{\%}$\mathdefault{1.0}$}}%
\end{pgfscope}%
\begin{pgfscope}%
\pgfsetrectcap%
\pgfsetroundjoin%
\pgfsetlinewidth{0.501875pt}%
\definecolor{currentstroke}{rgb}{0.000000,0.000000,0.000000}%
\pgfsetstrokecolor{currentstroke}%
\pgfsetdash{}{0pt}%
\pgfpathmoveto{\pgfqpoint{2.182089in}{0.413113in}}%
\pgfpathlineto{\pgfqpoint{2.228509in}{0.390277in}}%
\pgfusepath{stroke}%
\end{pgfscope}%
\begin{pgfscope}%
\definecolor{textcolor}{rgb}{0.000000,0.000000,0.000000}%
\pgfsetstrokecolor{textcolor}%
\pgfsetfillcolor{textcolor}%
\pgftext[x=2.345766in,y=0.195456in,,top]{\color{textcolor}{\rmfamily\fontsize{10.000000}{12.000000}\selectfont\catcode`\^=\active\def^{\ifmmode\sp\else\^{}\fi}\catcode`\%=\active\def%{\%}$\mathdefault{1.2}$}}%
\end{pgfscope}%
\begin{pgfscope}%
\pgfsetrectcap%
\pgfsetroundjoin%
\pgfsetlinewidth{0.501875pt}%
\definecolor{currentstroke}{rgb}{0.000000,0.000000,0.000000}%
\pgfsetstrokecolor{currentstroke}%
\pgfsetdash{}{0pt}%
\pgfpathmoveto{\pgfqpoint{3.794491in}{1.211392in}}%
\pgfpathlineto{\pgfqpoint{3.840672in}{1.191548in}}%
\pgfusepath{stroke}%
\end{pgfscope}%
\begin{pgfscope}%
\pgfsetrectcap%
\pgfsetroundjoin%
\pgfsetlinewidth{0.501875pt}%
\definecolor{currentstroke}{rgb}{0.000000,0.000000,0.000000}%
\pgfsetstrokecolor{currentstroke}%
\pgfsetdash{}{0pt}%
\pgfpathmoveto{\pgfqpoint{3.738677in}{1.183760in}}%
\pgfpathlineto{\pgfqpoint{3.784874in}{1.163815in}}%
\pgfusepath{stroke}%
\end{pgfscope}%
\begin{pgfscope}%
\pgfsetrectcap%
\pgfsetroundjoin%
\pgfsetlinewidth{0.501875pt}%
\definecolor{currentstroke}{rgb}{0.000000,0.000000,0.000000}%
\pgfsetstrokecolor{currentstroke}%
\pgfsetdash{}{0pt}%
\pgfpathmoveto{\pgfqpoint{3.682582in}{1.155988in}}%
\pgfpathlineto{\pgfqpoint{3.728794in}{1.135942in}}%
\pgfusepath{stroke}%
\end{pgfscope}%
\begin{pgfscope}%
\pgfsetrectcap%
\pgfsetroundjoin%
\pgfsetlinewidth{0.501875pt}%
\definecolor{currentstroke}{rgb}{0.000000,0.000000,0.000000}%
\pgfsetstrokecolor{currentstroke}%
\pgfsetdash{}{0pt}%
\pgfpathmoveto{\pgfqpoint{3.569539in}{1.100021in}}%
\pgfpathlineto{\pgfqpoint{3.615779in}{1.079772in}}%
\pgfusepath{stroke}%
\end{pgfscope}%
\begin{pgfscope}%
\pgfsetrectcap%
\pgfsetroundjoin%
\pgfsetlinewidth{0.501875pt}%
\definecolor{currentstroke}{rgb}{0.000000,0.000000,0.000000}%
\pgfsetstrokecolor{currentstroke}%
\pgfsetdash{}{0pt}%
\pgfpathmoveto{\pgfqpoint{3.512587in}{1.071825in}}%
\pgfpathlineto{\pgfqpoint{3.558841in}{1.051473in}}%
\pgfusepath{stroke}%
\end{pgfscope}%
\begin{pgfscope}%
\pgfsetrectcap%
\pgfsetroundjoin%
\pgfsetlinewidth{0.501875pt}%
\definecolor{currentstroke}{rgb}{0.000000,0.000000,0.000000}%
\pgfsetstrokecolor{currentstroke}%
\pgfsetdash{}{0pt}%
\pgfpathmoveto{\pgfqpoint{3.455345in}{1.043485in}}%
\pgfpathlineto{\pgfqpoint{3.501611in}{1.023029in}}%
\pgfusepath{stroke}%
\end{pgfscope}%
\begin{pgfscope}%
\pgfsetrectcap%
\pgfsetroundjoin%
\pgfsetlinewidth{0.501875pt}%
\definecolor{currentstroke}{rgb}{0.000000,0.000000,0.000000}%
\pgfsetstrokecolor{currentstroke}%
\pgfsetdash{}{0pt}%
\pgfpathmoveto{\pgfqpoint{3.339981in}{0.986370in}}%
\pgfpathlineto{\pgfqpoint{3.386272in}{0.965703in}}%
\pgfusepath{stroke}%
\end{pgfscope}%
\begin{pgfscope}%
\pgfsetrectcap%
\pgfsetroundjoin%
\pgfsetlinewidth{0.501875pt}%
\definecolor{currentstroke}{rgb}{0.000000,0.000000,0.000000}%
\pgfsetstrokecolor{currentstroke}%
\pgfsetdash{}{0pt}%
\pgfpathmoveto{\pgfqpoint{3.281854in}{0.957592in}}%
\pgfpathlineto{\pgfqpoint{3.328157in}{0.936820in}}%
\pgfusepath{stroke}%
\end{pgfscope}%
\begin{pgfscope}%
\pgfsetrectcap%
\pgfsetroundjoin%
\pgfsetlinewidth{0.501875pt}%
\definecolor{currentstroke}{rgb}{0.000000,0.000000,0.000000}%
\pgfsetstrokecolor{currentstroke}%
\pgfsetdash{}{0pt}%
\pgfpathmoveto{\pgfqpoint{3.223429in}{0.928667in}}%
\pgfpathlineto{\pgfqpoint{3.269743in}{0.907787in}}%
\pgfusepath{stroke}%
\end{pgfscope}%
\begin{pgfscope}%
\pgfsetrectcap%
\pgfsetroundjoin%
\pgfsetlinewidth{0.501875pt}%
\definecolor{currentstroke}{rgb}{0.000000,0.000000,0.000000}%
\pgfsetstrokecolor{currentstroke}%
\pgfsetdash{}{0pt}%
\pgfpathmoveto{\pgfqpoint{3.105672in}{0.870367in}}%
\pgfpathlineto{\pgfqpoint{3.152007in}{0.849270in}}%
\pgfusepath{stroke}%
\end{pgfscope}%
\begin{pgfscope}%
\pgfsetrectcap%
\pgfsetroundjoin%
\pgfsetlinewidth{0.501875pt}%
\definecolor{currentstroke}{rgb}{0.000000,0.000000,0.000000}%
\pgfsetstrokecolor{currentstroke}%
\pgfsetdash{}{0pt}%
\pgfpathmoveto{\pgfqpoint{3.046335in}{0.840990in}}%
\pgfpathlineto{\pgfqpoint{3.092679in}{0.819783in}}%
\pgfusepath{stroke}%
\end{pgfscope}%
\begin{pgfscope}%
\pgfsetrectcap%
\pgfsetroundjoin%
\pgfsetlinewidth{0.501875pt}%
\definecolor{currentstroke}{rgb}{0.000000,0.000000,0.000000}%
\pgfsetstrokecolor{currentstroke}%
\pgfsetdash{}{0pt}%
\pgfpathmoveto{\pgfqpoint{2.986690in}{0.811460in}}%
\pgfpathlineto{\pgfqpoint{3.033043in}{0.790143in}}%
\pgfusepath{stroke}%
\end{pgfscope}%
\begin{pgfscope}%
\pgfsetrectcap%
\pgfsetroundjoin%
\pgfsetlinewidth{0.501875pt}%
\definecolor{currentstroke}{rgb}{0.000000,0.000000,0.000000}%
\pgfsetstrokecolor{currentstroke}%
\pgfsetdash{}{0pt}%
\pgfpathmoveto{\pgfqpoint{2.866464in}{0.751938in}}%
\pgfpathlineto{\pgfqpoint{2.912834in}{0.730397in}}%
\pgfusepath{stroke}%
\end{pgfscope}%
\begin{pgfscope}%
\pgfsetrectcap%
\pgfsetroundjoin%
\pgfsetlinewidth{0.501875pt}%
\definecolor{currentstroke}{rgb}{0.000000,0.000000,0.000000}%
\pgfsetstrokecolor{currentstroke}%
\pgfsetdash{}{0pt}%
\pgfpathmoveto{\pgfqpoint{2.805878in}{0.721943in}}%
\pgfpathlineto{\pgfqpoint{2.852256in}{0.700289in}}%
\pgfusepath{stroke}%
\end{pgfscope}%
\begin{pgfscope}%
\pgfsetrectcap%
\pgfsetroundjoin%
\pgfsetlinewidth{0.501875pt}%
\definecolor{currentstroke}{rgb}{0.000000,0.000000,0.000000}%
\pgfsetstrokecolor{currentstroke}%
\pgfsetdash{}{0pt}%
\pgfpathmoveto{\pgfqpoint{2.744974in}{0.691790in}}%
\pgfpathlineto{\pgfqpoint{2.791359in}{0.670022in}}%
\pgfusepath{stroke}%
\end{pgfscope}%
\begin{pgfscope}%
\pgfsetrectcap%
\pgfsetroundjoin%
\pgfsetlinewidth{0.501875pt}%
\definecolor{currentstroke}{rgb}{0.000000,0.000000,0.000000}%
\pgfsetstrokecolor{currentstroke}%
\pgfsetdash{}{0pt}%
\pgfpathmoveto{\pgfqpoint{2.622202in}{0.631006in}}%
\pgfpathlineto{\pgfqpoint{2.668598in}{0.609008in}}%
\pgfusepath{stroke}%
\end{pgfscope}%
\begin{pgfscope}%
\pgfsetrectcap%
\pgfsetroundjoin%
\pgfsetlinewidth{0.501875pt}%
\definecolor{currentstroke}{rgb}{0.000000,0.000000,0.000000}%
\pgfsetstrokecolor{currentstroke}%
\pgfsetdash{}{0pt}%
\pgfpathmoveto{\pgfqpoint{2.560327in}{0.600373in}}%
\pgfpathlineto{\pgfqpoint{2.606729in}{0.578258in}}%
\pgfusepath{stroke}%
\end{pgfscope}%
\begin{pgfscope}%
\pgfsetrectcap%
\pgfsetroundjoin%
\pgfsetlinewidth{0.501875pt}%
\definecolor{currentstroke}{rgb}{0.000000,0.000000,0.000000}%
\pgfsetstrokecolor{currentstroke}%
\pgfsetdash{}{0pt}%
\pgfpathmoveto{\pgfqpoint{2.498125in}{0.569578in}}%
\pgfpathlineto{\pgfqpoint{2.544531in}{0.547345in}}%
\pgfusepath{stroke}%
\end{pgfscope}%
\begin{pgfscope}%
\pgfsetrectcap%
\pgfsetroundjoin%
\pgfsetlinewidth{0.501875pt}%
\definecolor{currentstroke}{rgb}{0.000000,0.000000,0.000000}%
\pgfsetstrokecolor{currentstroke}%
\pgfsetdash{}{0pt}%
\pgfpathmoveto{\pgfqpoint{2.372723in}{0.507493in}}%
\pgfpathlineto{\pgfqpoint{2.419136in}{0.485022in}}%
\pgfusepath{stroke}%
\end{pgfscope}%
\begin{pgfscope}%
\pgfsetrectcap%
\pgfsetroundjoin%
\pgfsetlinewidth{0.501875pt}%
\definecolor{currentstroke}{rgb}{0.000000,0.000000,0.000000}%
\pgfsetstrokecolor{currentstroke}%
\pgfsetdash{}{0pt}%
\pgfpathmoveto{\pgfqpoint{2.309519in}{0.476201in}}%
\pgfpathlineto{\pgfqpoint{2.355934in}{0.453609in}}%
\pgfusepath{stroke}%
\end{pgfscope}%
\begin{pgfscope}%
\pgfsetrectcap%
\pgfsetroundjoin%
\pgfsetlinewidth{0.501875pt}%
\definecolor{currentstroke}{rgb}{0.000000,0.000000,0.000000}%
\pgfsetstrokecolor{currentstroke}%
\pgfsetdash{}{0pt}%
\pgfpathmoveto{\pgfqpoint{2.245975in}{0.444741in}}%
\pgfpathlineto{\pgfqpoint{2.292393in}{0.422028in}}%
\pgfusepath{stroke}%
\end{pgfscope}%
\begin{pgfscope}%
\pgfsetrectcap%
\pgfsetroundjoin%
\pgfsetlinewidth{0.501875pt}%
\definecolor{currentstroke}{rgb}{0.000000,0.000000,0.000000}%
\pgfsetstrokecolor{currentstroke}%
\pgfsetdash{}{0pt}%
\pgfpathmoveto{\pgfqpoint{2.117859in}{0.381313in}}%
\pgfpathlineto{\pgfqpoint{2.164279in}{0.358354in}}%
\pgfusepath{stroke}%
\end{pgfscope}%
\begin{pgfscope}%
\definecolor{textcolor}{rgb}{0.000000,0.000000,0.000000}%
\pgfsetstrokecolor{textcolor}%
\pgfsetfillcolor{textcolor}%
\pgftext[x=3.352445in,y=0.381299in,,]{\color{textcolor}{\rmfamily\fontsize{10.000000}{12.000000}\selectfont\catcode`\^=\active\def^{\ifmmode\sp\else\^{}\fi}\catcode`\%=\active\def%{\%}$y$}}%
\end{pgfscope}%
\begin{pgfscope}%
\pgfsetrectcap%
\pgfsetroundjoin%
\pgfsetlinewidth{0.501875pt}%
\definecolor{currentstroke}{rgb}{1.000000,1.000000,1.000000}%
\pgfsetstrokecolor{currentstroke}%
\pgfsetstrokeopacity{0.000000}%
\pgfsetdash{}{0pt}%
\pgfpathmoveto{\pgfqpoint{3.894286in}{1.246634in}}%
\pgfpathlineto{\pgfqpoint{3.963058in}{2.998637in}}%
\pgfusepath{stroke}%
\end{pgfscope}%
\begin{pgfscope}%
\pgfsetrectcap%
\pgfsetroundjoin%
\pgfsetlinewidth{0.501875pt}%
\definecolor{currentstroke}{rgb}{0.000000,0.000000,0.000000}%
\pgfsetstrokecolor{currentstroke}%
\pgfsetdash{}{0pt}%
\pgfpathmoveto{\pgfqpoint{3.884244in}{1.390065in}}%
\pgfpathlineto{\pgfqpoint{3.930550in}{1.370559in}}%
\pgfusepath{stroke}%
\end{pgfscope}%
\begin{pgfscope}%
\definecolor{textcolor}{rgb}{0.000000,0.000000,0.000000}%
\pgfsetstrokecolor{textcolor}%
\pgfsetfillcolor{textcolor}%
\pgftext[x=4.174397in,y=1.383571in,,top]{\color{textcolor}{\rmfamily\fontsize{10.000000}{12.000000}\selectfont\catcode`\^=\active\def^{\ifmmode\sp\else\^{}\fi}\catcode`\%=\active\def%{\%}$\mathdefault{-0.2}$}}%
\end{pgfscope}%
\begin{pgfscope}%
\pgfsetrectcap%
\pgfsetroundjoin%
\pgfsetlinewidth{0.501875pt}%
\definecolor{currentstroke}{rgb}{0.000000,0.000000,0.000000}%
\pgfsetstrokecolor{currentstroke}%
\pgfsetdash{}{0pt}%
\pgfpathmoveto{\pgfqpoint{3.894584in}{1.655812in}}%
\pgfpathlineto{\pgfqpoint{3.941179in}{1.636674in}}%
\pgfusepath{stroke}%
\end{pgfscope}%
\begin{pgfscope}%
\definecolor{textcolor}{rgb}{0.000000,0.000000,0.000000}%
\pgfsetstrokecolor{textcolor}%
\pgfsetfillcolor{textcolor}%
\pgftext[x=4.186432in,y=1.649440in,,top]{\color{textcolor}{\rmfamily\fontsize{10.000000}{12.000000}\selectfont\catcode`\^=\active\def^{\ifmmode\sp\else\^{}\fi}\catcode`\%=\active\def%{\%}$\mathdefault{0.0}$}}%
\end{pgfscope}%
\begin{pgfscope}%
\pgfsetrectcap%
\pgfsetroundjoin%
\pgfsetlinewidth{0.501875pt}%
\definecolor{currentstroke}{rgb}{0.000000,0.000000,0.000000}%
\pgfsetstrokecolor{currentstroke}%
\pgfsetdash{}{0pt}%
\pgfpathmoveto{\pgfqpoint{3.905046in}{1.924667in}}%
\pgfpathlineto{\pgfqpoint{3.951932in}{1.905909in}}%
\pgfusepath{stroke}%
\end{pgfscope}%
\begin{pgfscope}%
\definecolor{textcolor}{rgb}{0.000000,0.000000,0.000000}%
\pgfsetstrokecolor{textcolor}%
\pgfsetfillcolor{textcolor}%
\pgftext[x=4.198608in,y=1.918422in,,top]{\color{textcolor}{\rmfamily\fontsize{10.000000}{12.000000}\selectfont\catcode`\^=\active\def^{\ifmmode\sp\else\^{}\fi}\catcode`\%=\active\def%{\%}$\mathdefault{0.2}$}}%
\end{pgfscope}%
\begin{pgfscope}%
\pgfsetrectcap%
\pgfsetroundjoin%
\pgfsetlinewidth{0.501875pt}%
\definecolor{currentstroke}{rgb}{0.000000,0.000000,0.000000}%
\pgfsetstrokecolor{currentstroke}%
\pgfsetdash{}{0pt}%
\pgfpathmoveto{\pgfqpoint{3.915630in}{2.196686in}}%
\pgfpathlineto{\pgfqpoint{3.962813in}{2.178317in}}%
\pgfusepath{stroke}%
\end{pgfscope}%
\begin{pgfscope}%
\definecolor{textcolor}{rgb}{0.000000,0.000000,0.000000}%
\pgfsetstrokecolor{textcolor}%
\pgfsetfillcolor{textcolor}%
\pgftext[x=4.210927in,y=2.190571in,,top]{\color{textcolor}{\rmfamily\fontsize{10.000000}{12.000000}\selectfont\catcode`\^=\active\def^{\ifmmode\sp\else\^{}\fi}\catcode`\%=\active\def%{\%}$\mathdefault{0.4}$}}%
\end{pgfscope}%
\begin{pgfscope}%
\pgfsetrectcap%
\pgfsetroundjoin%
\pgfsetlinewidth{0.501875pt}%
\definecolor{currentstroke}{rgb}{0.000000,0.000000,0.000000}%
\pgfsetstrokecolor{currentstroke}%
\pgfsetdash{}{0pt}%
\pgfpathmoveto{\pgfqpoint{3.926340in}{2.471926in}}%
\pgfpathlineto{\pgfqpoint{3.973822in}{2.453956in}}%
\pgfusepath{stroke}%
\end{pgfscope}%
\begin{pgfscope}%
\definecolor{textcolor}{rgb}{0.000000,0.000000,0.000000}%
\pgfsetstrokecolor{textcolor}%
\pgfsetfillcolor{textcolor}%
\pgftext[x=4.223392in,y=2.465943in,,top]{\color{textcolor}{\rmfamily\fontsize{10.000000}{12.000000}\selectfont\catcode`\^=\active\def^{\ifmmode\sp\else\^{}\fi}\catcode`\%=\active\def%{\%}$\mathdefault{0.6}$}}%
\end{pgfscope}%
\begin{pgfscope}%
\pgfsetrectcap%
\pgfsetroundjoin%
\pgfsetlinewidth{0.501875pt}%
\definecolor{currentstroke}{rgb}{0.000000,0.000000,0.000000}%
\pgfsetstrokecolor{currentstroke}%
\pgfsetdash{}{0pt}%
\pgfpathmoveto{\pgfqpoint{3.937177in}{2.750442in}}%
\pgfpathlineto{\pgfqpoint{3.984962in}{2.732883in}}%
\pgfusepath{stroke}%
\end{pgfscope}%
\begin{pgfscope}%
\definecolor{textcolor}{rgb}{0.000000,0.000000,0.000000}%
\pgfsetstrokecolor{textcolor}%
\pgfsetfillcolor{textcolor}%
\pgftext[x=4.236006in,y=2.744596in,,top]{\color{textcolor}{\rmfamily\fontsize{10.000000}{12.000000}\selectfont\catcode`\^=\active\def^{\ifmmode\sp\else\^{}\fi}\catcode`\%=\active\def%{\%}$\mathdefault{0.8}$}}%
\end{pgfscope}%
\begin{pgfscope}%
\definecolor{textcolor}{rgb}{0.000000,0.000000,0.000000}%
\pgfsetstrokecolor{textcolor}%
\pgfsetfillcolor{textcolor}%
\pgftext[x=4.534721in,y=2.106108in,,]{\color{textcolor}{\rmfamily\fontsize{10.000000}{12.000000}\selectfont\catcode`\^=\active\def^{\ifmmode\sp\else\^{}\fi}\catcode`\%=\active\def%{\%}$z$}}%
\end{pgfscope}%
\begin{pgfscope}%
\pgfpathrectangle{\pgfqpoint{0.050000in}{0.050000in}}{\pgfqpoint{4.004000in}{4.004000in}}%
\pgfusepath{clip}%
\pgfsetrectcap%
\pgfsetroundjoin%
\pgfsetlinewidth{0.501875pt}%
\definecolor{currentstroke}{rgb}{0.690196,0.188235,0.376471}%
\pgfsetstrokecolor{currentstroke}%
\pgfsetstrokeopacity{0.900000}%
\pgfsetdash{}{0pt}%
\pgfpathmoveto{\pgfqpoint{2.869564in}{1.857421in}}%
\pgfpathlineto{\pgfqpoint{2.873722in}{1.927706in}}%
\pgfpathlineto{\pgfqpoint{2.961725in}{1.995606in}}%
\pgfpathlineto{\pgfqpoint{2.905157in}{1.962487in}}%
\pgfpathlineto{\pgfqpoint{2.882253in}{1.988913in}}%
\pgfpathlineto{\pgfqpoint{2.817974in}{2.005158in}}%
\pgfpathlineto{\pgfqpoint{2.837468in}{2.048746in}}%
\pgfpathlineto{\pgfqpoint{2.849247in}{2.093846in}}%
\pgfpathlineto{\pgfqpoint{2.816563in}{2.070890in}}%
\pgfpathlineto{\pgfqpoint{2.799159in}{2.085046in}}%
\pgfpathlineto{\pgfqpoint{2.827739in}{2.100010in}}%
\pgfpathlineto{\pgfqpoint{2.836195in}{2.087386in}}%
\pgfpathlineto{\pgfqpoint{2.839887in}{2.155115in}}%
\pgfpathlineto{\pgfqpoint{2.848813in}{2.123859in}}%
\pgfpathlineto{\pgfqpoint{2.843341in}{2.162698in}}%
\pgfpathlineto{\pgfqpoint{2.847395in}{2.160262in}}%
\pgfpathlineto{\pgfqpoint{2.860177in}{2.190911in}}%
\pgfpathlineto{\pgfqpoint{2.886768in}{2.189618in}}%
\pgfpathlineto{\pgfqpoint{2.803232in}{2.200366in}}%
\pgfpathlineto{\pgfqpoint{2.840106in}{2.253930in}}%
\pgfpathlineto{\pgfqpoint{2.848472in}{2.229464in}}%
\pgfpathlineto{\pgfqpoint{2.872014in}{2.263289in}}%
\pgfpathlineto{\pgfqpoint{2.934600in}{2.285698in}}%
\pgfpathlineto{\pgfqpoint{2.944563in}{2.247907in}}%
\pgfpathlineto{\pgfqpoint{2.955662in}{2.310529in}}%
\pgfpathlineto{\pgfqpoint{2.947638in}{2.249091in}}%
\pgfpathlineto{\pgfqpoint{2.975875in}{2.244075in}}%
\pgfpathlineto{\pgfqpoint{3.020320in}{2.255005in}}%
\pgfpathlineto{\pgfqpoint{2.981283in}{2.302236in}}%
\pgfpathlineto{\pgfqpoint{2.910191in}{2.315162in}}%
\pgfpathlineto{\pgfqpoint{2.963852in}{2.346839in}}%
\pgfpathlineto{\pgfqpoint{2.852284in}{2.353032in}}%
\pgfpathlineto{\pgfqpoint{2.810202in}{2.371467in}}%
\pgfpathlineto{\pgfqpoint{2.758977in}{2.379480in}}%
\pgfpathlineto{\pgfqpoint{2.727596in}{2.459221in}}%
\pgfpathlineto{\pgfqpoint{2.742679in}{2.461831in}}%
\pgfpathlineto{\pgfqpoint{2.783785in}{2.427506in}}%
\pgfpathlineto{\pgfqpoint{2.841547in}{2.393907in}}%
\pgfpathlineto{\pgfqpoint{2.781365in}{2.412942in}}%
\pgfpathlineto{\pgfqpoint{2.836897in}{2.437062in}}%
\pgfpathlineto{\pgfqpoint{2.764765in}{2.413145in}}%
\pgfpathlineto{\pgfqpoint{2.788892in}{2.445159in}}%
\pgfpathlineto{\pgfqpoint{2.813310in}{2.369920in}}%
\pgfpathlineto{\pgfqpoint{2.815946in}{2.370415in}}%
\pgfpathlineto{\pgfqpoint{2.804486in}{2.300049in}}%
\pgfpathlineto{\pgfqpoint{2.783830in}{2.311894in}}%
\pgfpathlineto{\pgfqpoint{2.779541in}{2.284000in}}%
\pgfpathlineto{\pgfqpoint{2.791552in}{2.295477in}}%
\pgfpathlineto{\pgfqpoint{2.763986in}{2.263429in}}%
\pgfpathlineto{\pgfqpoint{2.707818in}{2.225724in}}%
\pgfpathlineto{\pgfqpoint{2.725702in}{2.234031in}}%
\pgfpathlineto{\pgfqpoint{2.756364in}{2.183320in}}%
\pgfpathlineto{\pgfqpoint{2.746411in}{2.255313in}}%
\pgfpathlineto{\pgfqpoint{2.810439in}{2.224636in}}%
\pgfpathlineto{\pgfqpoint{2.744710in}{2.260501in}}%
\pgfpathlineto{\pgfqpoint{2.739451in}{2.223050in}}%
\pgfpathlineto{\pgfqpoint{2.692705in}{2.270652in}}%
\pgfpathlineto{\pgfqpoint{2.755319in}{2.277210in}}%
\pgfpathlineto{\pgfqpoint{2.788808in}{2.286823in}}%
\pgfpathlineto{\pgfqpoint{2.755513in}{2.332914in}}%
\pgfpathlineto{\pgfqpoint{2.749292in}{2.323116in}}%
\pgfpathlineto{\pgfqpoint{2.788597in}{2.273464in}}%
\pgfpathlineto{\pgfqpoint{2.803466in}{2.362455in}}%
\pgfpathlineto{\pgfqpoint{2.897710in}{2.357302in}}%
\pgfpathlineto{\pgfqpoint{2.932952in}{2.316934in}}%
\pgfpathlineto{\pgfqpoint{2.895230in}{2.302126in}}%
\pgfpathlineto{\pgfqpoint{2.919565in}{2.280498in}}%
\pgfpathlineto{\pgfqpoint{2.814693in}{2.304350in}}%
\pgfpathlineto{\pgfqpoint{2.796447in}{2.319304in}}%
\pgfpathlineto{\pgfqpoint{2.788283in}{2.309728in}}%
\pgfpathlineto{\pgfqpoint{2.756039in}{2.268014in}}%
\pgfpathlineto{\pgfqpoint{2.724905in}{2.309346in}}%
\pgfpathlineto{\pgfqpoint{2.722570in}{2.265819in}}%
\pgfpathlineto{\pgfqpoint{2.743719in}{2.239006in}}%
\pgfpathlineto{\pgfqpoint{2.730655in}{2.185316in}}%
\pgfpathlineto{\pgfqpoint{2.733127in}{2.135440in}}%
\pgfpathlineto{\pgfqpoint{2.710479in}{2.112073in}}%
\pgfpathlineto{\pgfqpoint{2.642541in}{2.157495in}}%
\pgfpathlineto{\pgfqpoint{2.601455in}{2.116301in}}%
\pgfpathlineto{\pgfqpoint{2.602126in}{2.099583in}}%
\pgfpathlineto{\pgfqpoint{2.553355in}{2.106693in}}%
\pgfpathlineto{\pgfqpoint{2.671665in}{2.079423in}}%
\pgfpathlineto{\pgfqpoint{2.675157in}{2.037990in}}%
\pgfpathlineto{\pgfqpoint{2.678248in}{2.060264in}}%
\pgfpathlineto{\pgfqpoint{2.717350in}{2.124364in}}%
\pgfpathlineto{\pgfqpoint{2.684716in}{2.151318in}}%
\pgfpathlineto{\pgfqpoint{2.648643in}{2.141196in}}%
\pgfpathlineto{\pgfqpoint{2.620210in}{2.163480in}}%
\pgfpathlineto{\pgfqpoint{2.563566in}{2.152313in}}%
\pgfpathlineto{\pgfqpoint{2.594168in}{2.168744in}}%
\pgfpathlineto{\pgfqpoint{2.643641in}{2.167107in}}%
\pgfpathlineto{\pgfqpoint{2.643205in}{2.135948in}}%
\pgfpathlineto{\pgfqpoint{2.635192in}{2.149572in}}%
\pgfpathlineto{\pgfqpoint{2.594591in}{2.168426in}}%
\pgfpathlineto{\pgfqpoint{2.616099in}{2.215363in}}%
\pgfpathlineto{\pgfqpoint{2.625023in}{2.203431in}}%
\pgfpathlineto{\pgfqpoint{2.643413in}{2.198147in}}%
\pgfpathlineto{\pgfqpoint{2.684562in}{2.183316in}}%
\pgfpathlineto{\pgfqpoint{2.657840in}{2.204298in}}%
\pgfpathlineto{\pgfqpoint{2.649327in}{2.163622in}}%
\pgfpathlineto{\pgfqpoint{2.637628in}{2.091012in}}%
\pgfpathlineto{\pgfqpoint{2.640278in}{2.071504in}}%
\pgfpathlineto{\pgfqpoint{2.637721in}{2.028447in}}%
\pgfpathlineto{\pgfqpoint{2.598408in}{1.962927in}}%
\pgfpathlineto{\pgfqpoint{2.652819in}{1.946615in}}%
\pgfpathlineto{\pgfqpoint{2.589936in}{1.971990in}}%
\pgfpathlineto{\pgfqpoint{2.528645in}{1.901734in}}%
\pgfpathlineto{\pgfqpoint{2.592332in}{1.937524in}}%
\pgfpathlineto{\pgfqpoint{2.588308in}{1.961185in}}%
\pgfpathlineto{\pgfqpoint{2.638865in}{1.969211in}}%
\pgfpathlineto{\pgfqpoint{2.644885in}{1.924826in}}%
\pgfpathlineto{\pgfqpoint{2.697152in}{1.934376in}}%
\pgfpathlineto{\pgfqpoint{2.741109in}{1.893290in}}%
\pgfpathlineto{\pgfqpoint{2.769710in}{1.949789in}}%
\pgfpathlineto{\pgfqpoint{2.722120in}{1.902428in}}%
\pgfpathlineto{\pgfqpoint{2.688073in}{1.908429in}}%
\pgfpathlineto{\pgfqpoint{2.764175in}{1.892543in}}%
\pgfpathlineto{\pgfqpoint{2.740561in}{1.918221in}}%
\pgfpathlineto{\pgfqpoint{2.688971in}{1.849710in}}%
\pgfpathlineto{\pgfqpoint{2.744302in}{1.814290in}}%
\pgfpathlineto{\pgfqpoint{2.681835in}{1.782859in}}%
\pgfpathlineto{\pgfqpoint{2.670068in}{1.771380in}}%
\pgfpathlineto{\pgfqpoint{2.671205in}{1.819488in}}%
\pgfpathlineto{\pgfqpoint{2.699910in}{1.775004in}}%
\pgfpathlineto{\pgfqpoint{2.724043in}{1.723951in}}%
\pgfpathlineto{\pgfqpoint{2.744297in}{1.821872in}}%
\pgfpathlineto{\pgfqpoint{2.647918in}{1.788494in}}%
\pgfpathlineto{\pgfqpoint{2.630552in}{1.763995in}}%
\pgfpathlineto{\pgfqpoint{2.613645in}{1.736309in}}%
\pgfpathlineto{\pgfqpoint{2.599782in}{1.748178in}}%
\pgfpathlineto{\pgfqpoint{2.679634in}{1.756968in}}%
\pgfpathlineto{\pgfqpoint{2.657563in}{1.757525in}}%
\pgfpathlineto{\pgfqpoint{2.661381in}{1.792263in}}%
\pgfpathlineto{\pgfqpoint{2.699849in}{1.818950in}}%
\pgfpathlineto{\pgfqpoint{2.713895in}{1.805400in}}%
\pgfpathlineto{\pgfqpoint{2.730445in}{1.733099in}}%
\pgfpathlineto{\pgfqpoint{2.743015in}{1.731207in}}%
\pgfpathlineto{\pgfqpoint{2.781011in}{1.709647in}}%
\pgfpathlineto{\pgfqpoint{2.744421in}{1.700214in}}%
\pgfpathlineto{\pgfqpoint{2.666729in}{1.748538in}}%
\pgfpathlineto{\pgfqpoint{2.680469in}{1.667767in}}%
\pgfpathlineto{\pgfqpoint{2.723356in}{1.695835in}}%
\pgfpathlineto{\pgfqpoint{2.640813in}{1.644581in}}%
\pgfpathlineto{\pgfqpoint{2.624790in}{1.653291in}}%
\pgfpathlineto{\pgfqpoint{2.612497in}{1.610622in}}%
\pgfpathlineto{\pgfqpoint{2.591242in}{1.598931in}}%
\pgfpathlineto{\pgfqpoint{2.662269in}{1.617725in}}%
\pgfpathlineto{\pgfqpoint{2.605282in}{1.653695in}}%
\pgfpathlineto{\pgfqpoint{2.613008in}{1.652140in}}%
\pgfpathlineto{\pgfqpoint{2.636188in}{1.633055in}}%
\pgfpathlineto{\pgfqpoint{2.674142in}{1.630866in}}%
\pgfpathlineto{\pgfqpoint{2.675835in}{1.641057in}}%
\pgfpathlineto{\pgfqpoint{2.684423in}{1.683443in}}%
\pgfpathlineto{\pgfqpoint{2.649548in}{1.684448in}}%
\pgfpathlineto{\pgfqpoint{2.664081in}{1.665930in}}%
\pgfpathlineto{\pgfqpoint{2.679236in}{1.694966in}}%
\pgfpathlineto{\pgfqpoint{2.712781in}{1.740854in}}%
\pgfpathlineto{\pgfqpoint{2.729643in}{1.700168in}}%
\pgfpathlineto{\pgfqpoint{2.707363in}{1.789823in}}%
\pgfpathlineto{\pgfqpoint{2.728396in}{1.838725in}}%
\pgfpathlineto{\pgfqpoint{2.836905in}{1.817330in}}%
\pgfpathlineto{\pgfqpoint{2.755222in}{1.810681in}}%
\pgfpathlineto{\pgfqpoint{2.704630in}{1.811876in}}%
\pgfpathlineto{\pgfqpoint{2.683984in}{1.782345in}}%
\pgfpathlineto{\pgfqpoint{2.663557in}{1.793557in}}%
\pgfpathlineto{\pgfqpoint{2.647797in}{1.748442in}}%
\pgfpathlineto{\pgfqpoint{2.724578in}{1.730810in}}%
\pgfpathlineto{\pgfqpoint{2.673455in}{1.702176in}}%
\pgfpathlineto{\pgfqpoint{2.730278in}{1.657821in}}%
\pgfpathlineto{\pgfqpoint{2.696646in}{1.664551in}}%
\pgfpathlineto{\pgfqpoint{2.713735in}{1.666941in}}%
\pgfpathlineto{\pgfqpoint{2.574931in}{1.613269in}}%
\pgfpathlineto{\pgfqpoint{2.621656in}{1.651954in}}%
\pgfpathlineto{\pgfqpoint{2.590222in}{1.638337in}}%
\pgfpathlineto{\pgfqpoint{2.592871in}{1.722002in}}%
\pgfpathlineto{\pgfqpoint{2.536119in}{1.698068in}}%
\pgfpathlineto{\pgfqpoint{2.528876in}{1.708461in}}%
\pgfpathlineto{\pgfqpoint{2.530004in}{1.736554in}}%
\pgfpathlineto{\pgfqpoint{2.526749in}{1.800819in}}%
\pgfpathlineto{\pgfqpoint{2.540710in}{1.801581in}}%
\pgfpathlineto{\pgfqpoint{2.531564in}{1.776689in}}%
\pgfpathlineto{\pgfqpoint{2.488175in}{1.759331in}}%
\pgfpathlineto{\pgfqpoint{2.427499in}{1.766925in}}%
\pgfpathlineto{\pgfqpoint{2.406381in}{1.768025in}}%
\pgfpathlineto{\pgfqpoint{2.438965in}{1.828755in}}%
\pgfpathlineto{\pgfqpoint{2.498267in}{1.801286in}}%
\pgfpathlineto{\pgfqpoint{2.530592in}{1.749160in}}%
\pgfpathlineto{\pgfqpoint{2.609663in}{1.719333in}}%
\pgfpathlineto{\pgfqpoint{2.533394in}{1.655351in}}%
\pgfpathlineto{\pgfqpoint{2.553832in}{1.632418in}}%
\pgfpathlineto{\pgfqpoint{2.522115in}{1.558215in}}%
\pgfpathlineto{\pgfqpoint{2.533445in}{1.553849in}}%
\pgfpathlineto{\pgfqpoint{2.530434in}{1.569077in}}%
\pgfpathlineto{\pgfqpoint{2.553291in}{1.529509in}}%
\pgfpathlineto{\pgfqpoint{2.540523in}{1.512957in}}%
\pgfpathlineto{\pgfqpoint{2.580871in}{1.589687in}}%
\pgfpathlineto{\pgfqpoint{2.567247in}{1.577856in}}%
\pgfpathlineto{\pgfqpoint{2.579221in}{1.537599in}}%
\pgfpathlineto{\pgfqpoint{2.579403in}{1.570398in}}%
\pgfpathlineto{\pgfqpoint{2.576026in}{1.567600in}}%
\pgfpathlineto{\pgfqpoint{2.584333in}{1.543674in}}%
\pgfpathlineto{\pgfqpoint{2.645746in}{1.548136in}}%
\pgfpathlineto{\pgfqpoint{2.622154in}{1.557994in}}%
\pgfpathlineto{\pgfqpoint{2.598521in}{1.562537in}}%
\pgfpathlineto{\pgfqpoint{2.616412in}{1.530401in}}%
\pgfpathlineto{\pgfqpoint{2.589028in}{1.563701in}}%
\pgfpathlineto{\pgfqpoint{2.612868in}{1.614944in}}%
\pgfpathlineto{\pgfqpoint{2.575180in}{1.628209in}}%
\pgfpathlineto{\pgfqpoint{2.525201in}{1.603680in}}%
\pgfpathlineto{\pgfqpoint{2.618317in}{1.687315in}}%
\pgfpathlineto{\pgfqpoint{2.604183in}{1.700983in}}%
\pgfpathlineto{\pgfqpoint{2.615847in}{1.733208in}}%
\pgfpathlineto{\pgfqpoint{2.590985in}{1.748990in}}%
\pgfpathlineto{\pgfqpoint{2.608371in}{1.754692in}}%
\pgfpathlineto{\pgfqpoint{2.685042in}{1.777709in}}%
\pgfpathlineto{\pgfqpoint{2.670500in}{1.830226in}}%
\pgfpathlineto{\pgfqpoint{2.677819in}{1.839907in}}%
\pgfpathlineto{\pgfqpoint{2.661404in}{1.803951in}}%
\pgfpathlineto{\pgfqpoint{2.628048in}{1.778112in}}%
\pgfpathlineto{\pgfqpoint{2.680073in}{1.799991in}}%
\pgfpathlineto{\pgfqpoint{2.646601in}{1.800082in}}%
\pgfpathlineto{\pgfqpoint{2.579460in}{1.739894in}}%
\pgfpathlineto{\pgfqpoint{2.621395in}{1.745956in}}%
\pgfpathlineto{\pgfqpoint{2.652682in}{1.745474in}}%
\pgfpathlineto{\pgfqpoint{2.654999in}{1.743570in}}%
\pgfpathlineto{\pgfqpoint{2.636377in}{1.731140in}}%
\pgfpathlineto{\pgfqpoint{2.668432in}{1.751272in}}%
\pgfpathlineto{\pgfqpoint{2.689157in}{1.696623in}}%
\pgfpathlineto{\pgfqpoint{2.689626in}{1.678759in}}%
\pgfpathlineto{\pgfqpoint{2.668836in}{1.671562in}}%
\pgfpathlineto{\pgfqpoint{2.697891in}{1.676990in}}%
\pgfpathlineto{\pgfqpoint{2.670587in}{1.689504in}}%
\pgfpathlineto{\pgfqpoint{2.667449in}{1.703140in}}%
\pgfpathlineto{\pgfqpoint{2.637868in}{1.750061in}}%
\pgfpathlineto{\pgfqpoint{2.638469in}{1.699720in}}%
\pgfpathlineto{\pgfqpoint{2.604802in}{1.744053in}}%
\pgfpathlineto{\pgfqpoint{2.601176in}{1.727898in}}%
\pgfpathlineto{\pgfqpoint{2.555317in}{1.814746in}}%
\pgfpathlineto{\pgfqpoint{2.561560in}{1.801028in}}%
\pgfpathlineto{\pgfqpoint{2.601454in}{1.776400in}}%
\pgfpathlineto{\pgfqpoint{2.648400in}{1.770647in}}%
\pgfpathlineto{\pgfqpoint{2.695633in}{1.790402in}}%
\pgfpathlineto{\pgfqpoint{2.620048in}{1.801748in}}%
\pgfpathlineto{\pgfqpoint{2.679119in}{1.783108in}}%
\pgfpathlineto{\pgfqpoint{2.702349in}{1.808218in}}%
\pgfpathlineto{\pgfqpoint{2.730216in}{1.800360in}}%
\pgfpathlineto{\pgfqpoint{2.732712in}{1.863490in}}%
\pgfpathlineto{\pgfqpoint{2.697684in}{1.905966in}}%
\pgfpathlineto{\pgfqpoint{2.699139in}{1.915029in}}%
\pgfpathlineto{\pgfqpoint{2.695418in}{1.903868in}}%
\pgfpathlineto{\pgfqpoint{2.676833in}{1.874135in}}%
\pgfpathlineto{\pgfqpoint{2.720809in}{1.874224in}}%
\pgfpathlineto{\pgfqpoint{2.701360in}{1.861456in}}%
\pgfpathlineto{\pgfqpoint{2.659401in}{1.838603in}}%
\pgfpathlineto{\pgfqpoint{2.679597in}{1.882311in}}%
\pgfpathlineto{\pgfqpoint{2.730472in}{1.868210in}}%
\pgfpathlineto{\pgfqpoint{2.681381in}{1.841970in}}%
\pgfpathlineto{\pgfqpoint{2.657704in}{1.894748in}}%
\pgfpathlineto{\pgfqpoint{2.663781in}{1.946653in}}%
\pgfpathlineto{\pgfqpoint{2.635334in}{1.993541in}}%
\pgfpathlineto{\pgfqpoint{2.624425in}{1.933490in}}%
\pgfpathlineto{\pgfqpoint{2.592791in}{1.933559in}}%
\pgfpathlineto{\pgfqpoint{2.555689in}{1.945480in}}%
\pgfpathlineto{\pgfqpoint{2.546219in}{1.944201in}}%
\pgfpathlineto{\pgfqpoint{2.482366in}{1.903119in}}%
\pgfpathlineto{\pgfqpoint{2.469758in}{1.836927in}}%
\pgfpathlineto{\pgfqpoint{2.489414in}{1.853993in}}%
\pgfpathlineto{\pgfqpoint{2.572818in}{1.839940in}}%
\pgfpathlineto{\pgfqpoint{2.502519in}{1.829035in}}%
\pgfpathlineto{\pgfqpoint{2.554918in}{1.844646in}}%
\pgfpathlineto{\pgfqpoint{2.547004in}{1.853054in}}%
\pgfpathlineto{\pgfqpoint{2.558719in}{1.833813in}}%
\pgfpathlineto{\pgfqpoint{2.600946in}{1.847647in}}%
\pgfpathlineto{\pgfqpoint{2.588271in}{1.786694in}}%
\pgfpathlineto{\pgfqpoint{2.550487in}{1.790939in}}%
\pgfpathlineto{\pgfqpoint{2.521531in}{1.802771in}}%
\pgfpathlineto{\pgfqpoint{2.549348in}{1.813689in}}%
\pgfpathlineto{\pgfqpoint{2.565835in}{1.772812in}}%
\pgfpathlineto{\pgfqpoint{2.599832in}{1.769274in}}%
\pgfpathlineto{\pgfqpoint{2.595004in}{1.802751in}}%
\pgfpathlineto{\pgfqpoint{2.581180in}{1.856144in}}%
\pgfpathlineto{\pgfqpoint{2.593510in}{1.875720in}}%
\pgfpathlineto{\pgfqpoint{2.557516in}{1.876062in}}%
\pgfpathlineto{\pgfqpoint{2.565074in}{1.877075in}}%
\pgfpathlineto{\pgfqpoint{2.555785in}{1.854711in}}%
\pgfpathlineto{\pgfqpoint{2.525587in}{1.881308in}}%
\pgfpathlineto{\pgfqpoint{2.501000in}{1.859195in}}%
\pgfpathlineto{\pgfqpoint{2.542049in}{1.867005in}}%
\pgfpathlineto{\pgfqpoint{2.578784in}{1.859946in}}%
\pgfpathlineto{\pgfqpoint{2.609205in}{1.902318in}}%
\pgfpathlineto{\pgfqpoint{2.650274in}{1.920149in}}%
\pgfpathlineto{\pgfqpoint{2.595578in}{1.937701in}}%
\pgfpathlineto{\pgfqpoint{2.579704in}{1.880093in}}%
\pgfpathlineto{\pgfqpoint{2.603298in}{1.857315in}}%
\pgfpathlineto{\pgfqpoint{2.617791in}{1.834261in}}%
\pgfpathlineto{\pgfqpoint{2.666663in}{1.782173in}}%
\pgfpathlineto{\pgfqpoint{2.670405in}{1.818938in}}%
\pgfpathlineto{\pgfqpoint{2.626766in}{1.794359in}}%
\pgfpathlineto{\pgfqpoint{2.539427in}{1.819211in}}%
\pgfpathlineto{\pgfqpoint{2.479110in}{1.865909in}}%
\pgfpathlineto{\pgfqpoint{2.445018in}{1.890294in}}%
\pgfpathlineto{\pgfqpoint{2.536376in}{1.875135in}}%
\pgfpathlineto{\pgfqpoint{2.570379in}{1.864876in}}%
\pgfpathlineto{\pgfqpoint{2.482877in}{1.843712in}}%
\pgfpathlineto{\pgfqpoint{2.492096in}{1.805820in}}%
\pgfpathlineto{\pgfqpoint{2.498275in}{1.784425in}}%
\pgfpathlineto{\pgfqpoint{2.441183in}{1.803807in}}%
\pgfpathlineto{\pgfqpoint{2.451553in}{1.839090in}}%
\pgfpathlineto{\pgfqpoint{2.514599in}{1.849834in}}%
\pgfpathlineto{\pgfqpoint{2.572036in}{1.834932in}}%
\pgfpathlineto{\pgfqpoint{2.546891in}{1.784851in}}%
\pgfpathlineto{\pgfqpoint{2.598264in}{1.728017in}}%
\pgfpathlineto{\pgfqpoint{2.519167in}{1.754094in}}%
\pgfpathlineto{\pgfqpoint{2.519871in}{1.757117in}}%
\pgfpathlineto{\pgfqpoint{2.474357in}{1.727004in}}%
\pgfpathlineto{\pgfqpoint{2.498057in}{1.716772in}}%
\pgfpathlineto{\pgfqpoint{2.507923in}{1.760059in}}%
\pgfpathlineto{\pgfqpoint{2.484136in}{1.767553in}}%
\pgfpathlineto{\pgfqpoint{2.466529in}{1.724987in}}%
\pgfpathlineto{\pgfqpoint{2.393114in}{1.753636in}}%
\pgfpathlineto{\pgfqpoint{2.335795in}{1.790912in}}%
\pgfpathlineto{\pgfqpoint{2.321231in}{1.781216in}}%
\pgfpathlineto{\pgfqpoint{2.306075in}{1.808453in}}%
\pgfpathlineto{\pgfqpoint{2.267557in}{1.806831in}}%
\pgfpathlineto{\pgfqpoint{2.252980in}{1.821395in}}%
\pgfpathlineto{\pgfqpoint{2.283782in}{1.853545in}}%
\pgfpathlineto{\pgfqpoint{2.225075in}{1.819087in}}%
\pgfpathlineto{\pgfqpoint{2.245563in}{1.814864in}}%
\pgfpathlineto{\pgfqpoint{2.233289in}{1.765469in}}%
\pgfpathlineto{\pgfqpoint{2.269221in}{1.730900in}}%
\pgfpathlineto{\pgfqpoint{2.279258in}{1.701862in}}%
\pgfpathlineto{\pgfqpoint{2.333103in}{1.739608in}}%
\pgfpathlineto{\pgfqpoint{2.385544in}{1.781919in}}%
\pgfpathlineto{\pgfqpoint{2.410065in}{1.740911in}}%
\pgfpathlineto{\pgfqpoint{2.387900in}{1.736052in}}%
\pgfpathlineto{\pgfqpoint{2.451126in}{1.740962in}}%
\pgfpathlineto{\pgfqpoint{2.454502in}{1.737478in}}%
\pgfpathlineto{\pgfqpoint{2.444004in}{1.737555in}}%
\pgfpathlineto{\pgfqpoint{2.409623in}{1.727812in}}%
\pgfpathlineto{\pgfqpoint{2.487344in}{1.713503in}}%
\pgfpathlineto{\pgfqpoint{2.495078in}{1.723779in}}%
\pgfpathlineto{\pgfqpoint{2.438807in}{1.748248in}}%
\pgfpathlineto{\pgfqpoint{2.374780in}{1.799208in}}%
\pgfpathlineto{\pgfqpoint{2.367801in}{1.823056in}}%
\pgfpathlineto{\pgfqpoint{2.331382in}{1.799236in}}%
\pgfpathlineto{\pgfqpoint{2.304583in}{1.822687in}}%
\pgfpathlineto{\pgfqpoint{2.341801in}{1.825586in}}%
\pgfpathlineto{\pgfqpoint{2.324712in}{1.816811in}}%
\pgfpathlineto{\pgfqpoint{2.341178in}{1.783075in}}%
\pgfpathlineto{\pgfqpoint{2.268369in}{1.829805in}}%
\pgfpathlineto{\pgfqpoint{2.288177in}{1.870742in}}%
\pgfpathlineto{\pgfqpoint{2.260701in}{1.899533in}}%
\pgfpathlineto{\pgfqpoint{2.238350in}{1.900784in}}%
\pgfpathlineto{\pgfqpoint{2.153783in}{1.899897in}}%
\pgfpathlineto{\pgfqpoint{2.152463in}{1.895775in}}%
\pgfpathlineto{\pgfqpoint{2.165080in}{1.887163in}}%
\pgfpathlineto{\pgfqpoint{2.182198in}{1.852236in}}%
\pgfpathlineto{\pgfqpoint{2.157998in}{1.857700in}}%
\pgfpathlineto{\pgfqpoint{2.176798in}{1.845066in}}%
\pgfpathlineto{\pgfqpoint{2.104984in}{1.871780in}}%
\pgfpathlineto{\pgfqpoint{2.139874in}{1.805660in}}%
\pgfpathlineto{\pgfqpoint{2.160032in}{1.803635in}}%
\pgfpathlineto{\pgfqpoint{2.118457in}{1.825688in}}%
\pgfpathlineto{\pgfqpoint{2.099875in}{1.817558in}}%
\pgfpathlineto{\pgfqpoint{2.169940in}{1.830974in}}%
\pgfpathlineto{\pgfqpoint{2.165345in}{1.852959in}}%
\pgfpathlineto{\pgfqpoint{2.171517in}{1.854161in}}%
\pgfpathlineto{\pgfqpoint{2.049060in}{1.860125in}}%
\pgfpathlineto{\pgfqpoint{1.997457in}{1.853156in}}%
\pgfpathlineto{\pgfqpoint{2.049872in}{1.845446in}}%
\pgfpathlineto{\pgfqpoint{2.088623in}{1.862780in}}%
\pgfpathlineto{\pgfqpoint{2.074827in}{1.878495in}}%
\pgfpathlineto{\pgfqpoint{2.045590in}{1.899437in}}%
\pgfpathlineto{\pgfqpoint{2.057055in}{1.879729in}}%
\pgfpathlineto{\pgfqpoint{2.036763in}{1.850904in}}%
\pgfpathlineto{\pgfqpoint{2.116381in}{1.841592in}}%
\pgfpathlineto{\pgfqpoint{2.159054in}{1.883099in}}%
\pgfpathlineto{\pgfqpoint{2.065292in}{1.912610in}}%
\pgfpathlineto{\pgfqpoint{2.127218in}{1.996951in}}%
\pgfpathlineto{\pgfqpoint{2.123906in}{1.992801in}}%
\pgfpathlineto{\pgfqpoint{2.185111in}{1.959843in}}%
\pgfpathlineto{\pgfqpoint{2.236890in}{2.007896in}}%
\pgfpathlineto{\pgfqpoint{2.159000in}{1.976916in}}%
\pgfpathlineto{\pgfqpoint{2.140542in}{1.974708in}}%
\pgfpathlineto{\pgfqpoint{2.199421in}{2.017277in}}%
\pgfpathlineto{\pgfqpoint{2.202286in}{1.928289in}}%
\pgfpathlineto{\pgfqpoint{2.186076in}{1.860312in}}%
\pgfpathlineto{\pgfqpoint{2.173029in}{1.904367in}}%
\pgfpathlineto{\pgfqpoint{2.147321in}{1.891653in}}%
\pgfpathlineto{\pgfqpoint{2.176090in}{1.871518in}}%
\pgfpathlineto{\pgfqpoint{2.182566in}{1.838685in}}%
\pgfpathlineto{\pgfqpoint{2.185141in}{1.859754in}}%
\pgfpathlineto{\pgfqpoint{2.185718in}{1.834028in}}%
\pgfpathlineto{\pgfqpoint{2.133002in}{1.840930in}}%
\pgfpathlineto{\pgfqpoint{2.123355in}{1.798816in}}%
\pgfpathlineto{\pgfqpoint{2.104597in}{1.790901in}}%
\pgfpathlineto{\pgfqpoint{2.079916in}{1.814692in}}%
\pgfpathlineto{\pgfqpoint{2.073086in}{1.831537in}}%
\pgfpathlineto{\pgfqpoint{2.051732in}{1.815968in}}%
\pgfpathlineto{\pgfqpoint{2.061060in}{1.799109in}}%
\pgfpathlineto{\pgfqpoint{2.034430in}{1.759967in}}%
\pgfpathlineto{\pgfqpoint{1.991242in}{1.793188in}}%
\pgfpathlineto{\pgfqpoint{2.048643in}{1.752599in}}%
\pgfpathlineto{\pgfqpoint{2.083171in}{1.730892in}}%
\pgfpathlineto{\pgfqpoint{2.064222in}{1.705455in}}%
\pgfpathlineto{\pgfqpoint{2.052231in}{1.710428in}}%
\pgfpathlineto{\pgfqpoint{2.027151in}{1.693944in}}%
\pgfpathlineto{\pgfqpoint{2.002361in}{1.679960in}}%
\pgfpathlineto{\pgfqpoint{1.955382in}{1.699329in}}%
\pgfpathlineto{\pgfqpoint{1.905309in}{1.691440in}}%
\pgfpathlineto{\pgfqpoint{1.913062in}{1.712868in}}%
\pgfpathlineto{\pgfqpoint{1.912749in}{1.661862in}}%
\pgfpathlineto{\pgfqpoint{1.948612in}{1.716202in}}%
\pgfpathlineto{\pgfqpoint{1.971095in}{1.736067in}}%
\pgfpathlineto{\pgfqpoint{1.906211in}{1.763366in}}%
\pgfpathlineto{\pgfqpoint{1.887292in}{1.773752in}}%
\pgfpathlineto{\pgfqpoint{1.885663in}{1.798466in}}%
\pgfpathlineto{\pgfqpoint{1.841283in}{1.792940in}}%
\pgfpathlineto{\pgfqpoint{1.783516in}{1.824962in}}%
\pgfpathlineto{\pgfqpoint{1.774257in}{1.905495in}}%
\pgfpathlineto{\pgfqpoint{1.758506in}{1.902958in}}%
\pgfpathlineto{\pgfqpoint{1.698019in}{1.948946in}}%
\pgfpathlineto{\pgfqpoint{1.718749in}{1.990478in}}%
\pgfpathlineto{\pgfqpoint{1.725122in}{1.987805in}}%
\pgfpathlineto{\pgfqpoint{1.735457in}{2.025959in}}%
\pgfpathlineto{\pgfqpoint{1.708138in}{2.031000in}}%
\pgfpathlineto{\pgfqpoint{1.708828in}{2.061912in}}%
\pgfpathlineto{\pgfqpoint{1.646468in}{2.107377in}}%
\pgfpathlineto{\pgfqpoint{1.714172in}{2.067164in}}%
\pgfpathlineto{\pgfqpoint{1.717043in}{2.095608in}}%
\pgfpathlineto{\pgfqpoint{1.663583in}{2.067760in}}%
\pgfpathlineto{\pgfqpoint{1.707264in}{2.097655in}}%
\pgfpathlineto{\pgfqpoint{1.757777in}{2.107053in}}%
\pgfpathlineto{\pgfqpoint{1.810847in}{2.171170in}}%
\pgfpathlineto{\pgfqpoint{1.784620in}{2.169029in}}%
\pgfpathlineto{\pgfqpoint{1.782629in}{2.172794in}}%
\pgfpathlineto{\pgfqpoint{1.827105in}{2.154557in}}%
\pgfpathlineto{\pgfqpoint{1.836311in}{2.201969in}}%
\pgfpathlineto{\pgfqpoint{1.781112in}{2.218902in}}%
\pgfpathlineto{\pgfqpoint{1.796908in}{2.212814in}}%
\pgfpathlineto{\pgfqpoint{1.805265in}{2.205814in}}%
\pgfpathlineto{\pgfqpoint{1.844319in}{2.123889in}}%
\pgfpathlineto{\pgfqpoint{1.915530in}{2.151863in}}%
\pgfpathlineto{\pgfqpoint{1.953195in}{2.167436in}}%
\pgfpathlineto{\pgfqpoint{1.935956in}{2.182351in}}%
\pgfpathlineto{\pgfqpoint{1.995712in}{2.092453in}}%
\pgfpathlineto{\pgfqpoint{2.024352in}{2.094308in}}%
\pgfpathlineto{\pgfqpoint{2.023263in}{2.156332in}}%
\pgfpathlineto{\pgfqpoint{1.966554in}{2.140822in}}%
\pgfpathlineto{\pgfqpoint{1.930991in}{2.120382in}}%
\pgfpathlineto{\pgfqpoint{1.931339in}{2.110669in}}%
\pgfpathlineto{\pgfqpoint{1.855350in}{2.127725in}}%
\pgfpathlineto{\pgfqpoint{1.819368in}{2.121406in}}%
\pgfpathlineto{\pgfqpoint{1.784583in}{2.152943in}}%
\pgfpathlineto{\pgfqpoint{1.784009in}{2.186383in}}%
\pgfpathlineto{\pgfqpoint{1.797256in}{2.192183in}}%
\pgfpathlineto{\pgfqpoint{1.796614in}{2.141403in}}%
\pgfpathlineto{\pgfqpoint{1.741155in}{2.125870in}}%
\pgfpathlineto{\pgfqpoint{1.769437in}{2.149834in}}%
\pgfpathlineto{\pgfqpoint{1.763306in}{2.207744in}}%
\pgfpathlineto{\pgfqpoint{1.754307in}{2.247844in}}%
\pgfpathlineto{\pgfqpoint{1.727283in}{2.264023in}}%
\pgfpathlineto{\pgfqpoint{1.762287in}{2.232786in}}%
\pgfpathlineto{\pgfqpoint{1.796508in}{2.272158in}}%
\pgfpathlineto{\pgfqpoint{1.844162in}{2.296959in}}%
\pgfpathlineto{\pgfqpoint{1.858523in}{2.282030in}}%
\pgfpathlineto{\pgfqpoint{1.866064in}{2.266441in}}%
\pgfpathlineto{\pgfqpoint{1.782851in}{2.267490in}}%
\pgfpathlineto{\pgfqpoint{1.806845in}{2.211399in}}%
\pgfpathlineto{\pgfqpoint{1.820304in}{2.176201in}}%
\pgfpathlineto{\pgfqpoint{1.797604in}{2.191823in}}%
\pgfpathlineto{\pgfqpoint{1.817592in}{2.182732in}}%
\pgfpathlineto{\pgfqpoint{1.853460in}{2.251312in}}%
\pgfpathlineto{\pgfqpoint{1.836457in}{2.272114in}}%
\pgfpathlineto{\pgfqpoint{1.892930in}{2.290094in}}%
\pgfpathlineto{\pgfqpoint{1.874309in}{2.297986in}}%
\pgfpathlineto{\pgfqpoint{1.897211in}{2.334714in}}%
\pgfpathlineto{\pgfqpoint{1.887509in}{2.296253in}}%
\pgfpathlineto{\pgfqpoint{1.878134in}{2.356930in}}%
\pgfpathlineto{\pgfqpoint{1.901299in}{2.358374in}}%
\pgfpathlineto{\pgfqpoint{1.899068in}{2.358790in}}%
\pgfpathlineto{\pgfqpoint{1.893470in}{2.404277in}}%
\pgfpathlineto{\pgfqpoint{1.891786in}{2.385141in}}%
\pgfpathlineto{\pgfqpoint{1.893992in}{2.374116in}}%
\pgfpathlineto{\pgfqpoint{1.934592in}{2.339704in}}%
\pgfpathlineto{\pgfqpoint{1.890770in}{2.347702in}}%
\pgfpathlineto{\pgfqpoint{1.858882in}{2.317213in}}%
\pgfpathlineto{\pgfqpoint{1.926168in}{2.290415in}}%
\pgfpathlineto{\pgfqpoint{1.914906in}{2.262656in}}%
\pgfpathlineto{\pgfqpoint{1.924273in}{2.253569in}}%
\pgfpathlineto{\pgfqpoint{1.923518in}{2.238608in}}%
\pgfpathlineto{\pgfqpoint{1.915288in}{2.233259in}}%
\pgfpathlineto{\pgfqpoint{1.894993in}{2.206934in}}%
\pgfpathlineto{\pgfqpoint{1.959651in}{2.130468in}}%
\pgfpathlineto{\pgfqpoint{1.979566in}{2.141135in}}%
\pgfpathlineto{\pgfqpoint{2.019987in}{2.201654in}}%
\pgfpathlineto{\pgfqpoint{2.077636in}{2.150558in}}%
\pgfpathlineto{\pgfqpoint{2.156209in}{2.158246in}}%
\pgfpathlineto{\pgfqpoint{2.081613in}{2.126039in}}%
\pgfpathlineto{\pgfqpoint{2.057222in}{2.143993in}}%
\pgfpathlineto{\pgfqpoint{2.104443in}{2.189943in}}%
\pgfpathlineto{\pgfqpoint{2.138214in}{2.163473in}}%
\pgfpathlineto{\pgfqpoint{2.177751in}{2.144727in}}%
\pgfpathlineto{\pgfqpoint{2.158411in}{2.157141in}}%
\pgfpathlineto{\pgfqpoint{2.066741in}{2.185761in}}%
\pgfpathlineto{\pgfqpoint{2.077186in}{2.216767in}}%
\pgfpathlineto{\pgfqpoint{2.049719in}{2.203976in}}%
\pgfpathlineto{\pgfqpoint{2.008569in}{2.214214in}}%
\pgfpathlineto{\pgfqpoint{2.048296in}{2.188871in}}%
\pgfpathlineto{\pgfqpoint{2.004414in}{2.197764in}}%
\pgfpathlineto{\pgfqpoint{2.025001in}{2.197689in}}%
\pgfpathlineto{\pgfqpoint{2.012384in}{2.173087in}}%
\pgfpathlineto{\pgfqpoint{1.956585in}{2.185117in}}%
\pgfpathlineto{\pgfqpoint{1.957795in}{2.215267in}}%
\pgfpathlineto{\pgfqpoint{1.894600in}{2.199045in}}%
\pgfpathlineto{\pgfqpoint{1.896306in}{2.196711in}}%
\pgfpathlineto{\pgfqpoint{1.903586in}{2.214882in}}%
\pgfpathlineto{\pgfqpoint{1.877634in}{2.302320in}}%
\pgfpathlineto{\pgfqpoint{1.848910in}{2.273001in}}%
\pgfpathlineto{\pgfqpoint{1.793283in}{2.232861in}}%
\pgfpathlineto{\pgfqpoint{1.750008in}{2.191327in}}%
\pgfpathlineto{\pgfqpoint{1.789773in}{2.144837in}}%
\pgfpathlineto{\pgfqpoint{1.812999in}{2.142205in}}%
\pgfpathlineto{\pgfqpoint{1.801495in}{2.156197in}}%
\pgfpathlineto{\pgfqpoint{1.778393in}{2.141012in}}%
\pgfpathlineto{\pgfqpoint{1.761250in}{2.129068in}}%
\pgfpathlineto{\pgfqpoint{1.839374in}{2.129954in}}%
\pgfpathlineto{\pgfqpoint{1.839899in}{2.102956in}}%
\pgfpathlineto{\pgfqpoint{1.764304in}{2.065118in}}%
\pgfpathlineto{\pgfqpoint{1.760558in}{2.051695in}}%
\pgfpathlineto{\pgfqpoint{1.725738in}{1.983649in}}%
\pgfpathlineto{\pgfqpoint{1.772907in}{1.978559in}}%
\pgfpathlineto{\pgfqpoint{1.803948in}{1.991666in}}%
\pgfpathlineto{\pgfqpoint{1.824539in}{1.963890in}}%
\pgfpathlineto{\pgfqpoint{1.806372in}{2.001374in}}%
\pgfpathlineto{\pgfqpoint{1.849657in}{2.052899in}}%
\pgfpathlineto{\pgfqpoint{1.795959in}{2.127560in}}%
\pgfpathlineto{\pgfqpoint{1.758811in}{2.126161in}}%
\pgfpathlineto{\pgfqpoint{1.771563in}{2.106431in}}%
\pgfpathlineto{\pgfqpoint{1.792392in}{2.068163in}}%
\pgfpathlineto{\pgfqpoint{1.752401in}{2.054490in}}%
\pgfpathlineto{\pgfqpoint{1.689599in}{2.046482in}}%
\pgfpathlineto{\pgfqpoint{1.726972in}{2.102740in}}%
\pgfpathlineto{\pgfqpoint{1.699112in}{2.046038in}}%
\pgfpathlineto{\pgfqpoint{1.684834in}{2.156919in}}%
\pgfpathlineto{\pgfqpoint{1.669697in}{2.125070in}}%
\pgfpathlineto{\pgfqpoint{1.589098in}{2.147459in}}%
\pgfpathlineto{\pgfqpoint{1.576598in}{2.141086in}}%
\pgfpathlineto{\pgfqpoint{1.585828in}{2.128093in}}%
\pgfpathlineto{\pgfqpoint{1.570867in}{2.106325in}}%
\pgfpathlineto{\pgfqpoint{1.577041in}{2.041236in}}%
\pgfpathlineto{\pgfqpoint{1.546890in}{1.979062in}}%
\pgfpathlineto{\pgfqpoint{1.514933in}{1.965088in}}%
\pgfpathlineto{\pgfqpoint{1.514146in}{1.933153in}}%
\pgfpathlineto{\pgfqpoint{1.524079in}{1.910534in}}%
\pgfpathlineto{\pgfqpoint{1.466333in}{1.963581in}}%
\pgfpathlineto{\pgfqpoint{1.458599in}{1.952727in}}%
\pgfpathlineto{\pgfqpoint{1.495727in}{1.939938in}}%
\pgfpathlineto{\pgfqpoint{1.482049in}{1.882893in}}%
\pgfpathlineto{\pgfqpoint{1.489314in}{1.815459in}}%
\pgfpathlineto{\pgfqpoint{1.499417in}{1.774257in}}%
\pgfpathlineto{\pgfqpoint{1.439445in}{1.725315in}}%
\pgfpathlineto{\pgfqpoint{1.399544in}{1.737919in}}%
\pgfpathlineto{\pgfqpoint{1.485350in}{1.861575in}}%
\pgfpathlineto{\pgfqpoint{1.481855in}{1.887778in}}%
\pgfpathlineto{\pgfqpoint{1.435378in}{1.871953in}}%
\pgfpathlineto{\pgfqpoint{1.421650in}{1.884990in}}%
\pgfpathlineto{\pgfqpoint{1.403848in}{1.933581in}}%
\pgfpathlineto{\pgfqpoint{1.421391in}{1.925075in}}%
\pgfpathlineto{\pgfqpoint{1.406240in}{1.934026in}}%
\pgfpathlineto{\pgfqpoint{1.371533in}{1.932852in}}%
\pgfpathlineto{\pgfqpoint{1.362838in}{1.918269in}}%
\pgfpathlineto{\pgfqpoint{1.375220in}{1.909666in}}%
\pgfpathlineto{\pgfqpoint{1.356135in}{1.895209in}}%
\pgfpathlineto{\pgfqpoint{1.422762in}{1.870556in}}%
\pgfpathlineto{\pgfqpoint{1.452435in}{1.944350in}}%
\pgfpathlineto{\pgfqpoint{1.422896in}{1.953838in}}%
\pgfpathlineto{\pgfqpoint{1.446946in}{1.999518in}}%
\pgfpathlineto{\pgfqpoint{1.443359in}{2.019892in}}%
\pgfpathlineto{\pgfqpoint{1.488279in}{1.978403in}}%
\pgfpathlineto{\pgfqpoint{1.534319in}{1.967143in}}%
\pgfpathlineto{\pgfqpoint{1.494775in}{1.934223in}}%
\pgfpathlineto{\pgfqpoint{1.527673in}{1.997370in}}%
\pgfpathlineto{\pgfqpoint{1.529176in}{2.026887in}}%
\pgfpathlineto{\pgfqpoint{1.566586in}{2.046829in}}%
\pgfpathlineto{\pgfqpoint{1.603104in}{2.012255in}}%
\pgfpathlineto{\pgfqpoint{1.634849in}{2.007746in}}%
\pgfpathlineto{\pgfqpoint{1.693760in}{1.973381in}}%
\pgfpathlineto{\pgfqpoint{1.656929in}{1.957931in}}%
\pgfpathlineto{\pgfqpoint{1.686736in}{1.945654in}}%
\pgfpathlineto{\pgfqpoint{1.652132in}{1.948303in}}%
\pgfpathlineto{\pgfqpoint{1.625462in}{1.988105in}}%
\pgfpathlineto{\pgfqpoint{1.599507in}{2.009314in}}%
\pgfpathlineto{\pgfqpoint{1.620983in}{2.006537in}}%
\pgfpathlineto{\pgfqpoint{1.724818in}{2.048679in}}%
\pgfpathlineto{\pgfqpoint{1.726818in}{1.998757in}}%
\pgfpathlineto{\pgfqpoint{1.719485in}{1.982480in}}%
\pgfpathlineto{\pgfqpoint{1.764913in}{1.959409in}}%
\pgfpathlineto{\pgfqpoint{1.766393in}{1.990277in}}%
\pgfpathlineto{\pgfqpoint{1.771572in}{2.010157in}}%
\pgfpathlineto{\pgfqpoint{1.775825in}{1.929300in}}%
\pgfpathlineto{\pgfqpoint{1.709033in}{1.898173in}}%
\pgfpathlineto{\pgfqpoint{1.679545in}{1.926745in}}%
\pgfpathlineto{\pgfqpoint{1.659743in}{1.917569in}}%
\pgfpathlineto{\pgfqpoint{1.597531in}{1.928817in}}%
\pgfpathlineto{\pgfqpoint{1.599489in}{1.922371in}}%
\pgfpathlineto{\pgfqpoint{1.521971in}{1.904516in}}%
\pgfpathlineto{\pgfqpoint{1.476320in}{1.867967in}}%
\pgfpathlineto{\pgfqpoint{1.460586in}{1.906439in}}%
\pgfpathlineto{\pgfqpoint{1.477004in}{1.913332in}}%
\pgfpathlineto{\pgfqpoint{1.433760in}{1.912291in}}%
\pgfpathlineto{\pgfqpoint{1.434183in}{1.907301in}}%
\pgfpathlineto{\pgfqpoint{1.452806in}{1.848089in}}%
\pgfpathlineto{\pgfqpoint{1.544947in}{1.803126in}}%
\pgfpathlineto{\pgfqpoint{1.525363in}{1.782878in}}%
\pgfpathlineto{\pgfqpoint{1.470001in}{1.787589in}}%
\pgfpathlineto{\pgfqpoint{1.472957in}{1.744426in}}%
\pgfpathlineto{\pgfqpoint{1.436019in}{1.790839in}}%
\pgfpathlineto{\pgfqpoint{1.481479in}{1.756748in}}%
\pgfpathlineto{\pgfqpoint{1.526342in}{1.751652in}}%
\pgfpathlineto{\pgfqpoint{1.523591in}{1.790330in}}%
\pgfpathlineto{\pgfqpoint{1.541623in}{1.700882in}}%
\pgfpathlineto{\pgfqpoint{1.522646in}{1.706695in}}%
\pgfpathlineto{\pgfqpoint{1.537002in}{1.687441in}}%
\pgfpathlineto{\pgfqpoint{1.487975in}{1.670857in}}%
\pgfpathlineto{\pgfqpoint{1.468285in}{1.661099in}}%
\pgfpathlineto{\pgfqpoint{1.457196in}{1.667604in}}%
\pgfpathlineto{\pgfqpoint{1.449517in}{1.675240in}}%
\pgfpathlineto{\pgfqpoint{1.519171in}{1.641298in}}%
\pgfpathlineto{\pgfqpoint{1.536030in}{1.602396in}}%
\pgfpathlineto{\pgfqpoint{1.548881in}{1.589179in}}%
\pgfpathlineto{\pgfqpoint{1.572800in}{1.601036in}}%
\pgfpathlineto{\pgfqpoint{1.556194in}{1.685981in}}%
\pgfpathlineto{\pgfqpoint{1.601714in}{1.639075in}}%
\pgfpathlineto{\pgfqpoint{1.540403in}{1.627611in}}%
\pgfpathlineto{\pgfqpoint{1.614977in}{1.572412in}}%
\pgfpathlineto{\pgfqpoint{1.671775in}{1.569527in}}%
\pgfpathlineto{\pgfqpoint{1.651598in}{1.576593in}}%
\pgfpathlineto{\pgfqpoint{1.686982in}{1.601043in}}%
\pgfpathlineto{\pgfqpoint{1.684710in}{1.576707in}}%
\pgfpathlineto{\pgfqpoint{1.677882in}{1.559142in}}%
\pgfpathlineto{\pgfqpoint{1.635713in}{1.559281in}}%
\pgfpathlineto{\pgfqpoint{1.682184in}{1.537459in}}%
\pgfpathlineto{\pgfqpoint{1.659837in}{1.560976in}}%
\pgfpathlineto{\pgfqpoint{1.636563in}{1.578893in}}%
\pgfpathlineto{\pgfqpoint{1.577292in}{1.595943in}}%
\pgfpathlineto{\pgfqpoint{1.536341in}{1.606670in}}%
\pgfpathlineto{\pgfqpoint{1.582961in}{1.569779in}}%
\pgfpathlineto{\pgfqpoint{1.525892in}{1.509779in}}%
\pgfpathlineto{\pgfqpoint{1.538394in}{1.483755in}}%
\pgfpathlineto{\pgfqpoint{1.524414in}{1.456716in}}%
\pgfpathlineto{\pgfqpoint{1.465754in}{1.484475in}}%
\pgfpathlineto{\pgfqpoint{1.451658in}{1.444938in}}%
\pgfpathlineto{\pgfqpoint{1.402983in}{1.446551in}}%
\pgfpathlineto{\pgfqpoint{1.407698in}{1.437386in}}%
\pgfpathlineto{\pgfqpoint{1.379074in}{1.464747in}}%
\pgfpathlineto{\pgfqpoint{1.353960in}{1.481715in}}%
\pgfpathlineto{\pgfqpoint{1.329963in}{1.508020in}}%
\pgfpathlineto{\pgfqpoint{1.370599in}{1.499748in}}%
\pgfpathlineto{\pgfqpoint{1.420279in}{1.477473in}}%
\pgfpathlineto{\pgfqpoint{1.425512in}{1.535170in}}%
\pgfpathlineto{\pgfqpoint{1.471241in}{1.579364in}}%
\pgfpathlineto{\pgfqpoint{1.424815in}{1.612658in}}%
\pgfpathlineto{\pgfqpoint{1.360990in}{1.605874in}}%
\pgfpathlineto{\pgfqpoint{1.338987in}{1.634957in}}%
\pgfpathlineto{\pgfqpoint{1.276213in}{1.574517in}}%
\pgfpathlineto{\pgfqpoint{1.267178in}{1.573254in}}%
\pgfpathlineto{\pgfqpoint{1.285272in}{1.618368in}}%
\pgfpathlineto{\pgfqpoint{1.349122in}{1.560211in}}%
\pgfpathlineto{\pgfqpoint{1.340159in}{1.514319in}}%
\pgfpathlineto{\pgfqpoint{1.298467in}{1.538339in}}%
\pgfpathlineto{\pgfqpoint{1.280715in}{1.579933in}}%
\pgfpathlineto{\pgfqpoint{1.345186in}{1.582819in}}%
\pgfpathlineto{\pgfqpoint{1.408581in}{1.557280in}}%
\pgfpathlineto{\pgfqpoint{1.462532in}{1.560921in}}%
\pgfpathlineto{\pgfqpoint{1.384888in}{1.589131in}}%
\pgfpathlineto{\pgfqpoint{1.357441in}{1.569479in}}%
\pgfpathlineto{\pgfqpoint{1.323229in}{1.579880in}}%
\pgfpathlineto{\pgfqpoint{1.280636in}{1.535020in}}%
\pgfpathlineto{\pgfqpoint{1.297673in}{1.495654in}}%
\pgfpathlineto{\pgfqpoint{1.319608in}{1.584895in}}%
\pgfpathlineto{\pgfqpoint{1.270603in}{1.524187in}}%
\pgfpathlineto{\pgfqpoint{1.273485in}{1.565844in}}%
\pgfpathlineto{\pgfqpoint{1.233080in}{1.605337in}}%
\pgfpathlineto{\pgfqpoint{1.272361in}{1.637887in}}%
\pgfpathlineto{\pgfqpoint{1.294511in}{1.659154in}}%
\pgfpathlineto{\pgfqpoint{1.304666in}{1.701212in}}%
\pgfpathlineto{\pgfqpoint{1.290310in}{1.688436in}}%
\pgfpathlineto{\pgfqpoint{1.407656in}{1.671163in}}%
\pgfpathlineto{\pgfqpoint{1.445105in}{1.698673in}}%
\pgfpathlineto{\pgfqpoint{1.412401in}{1.695164in}}%
\pgfpathlineto{\pgfqpoint{1.418053in}{1.663015in}}%
\pgfpathlineto{\pgfqpoint{1.392381in}{1.676379in}}%
\pgfpathlineto{\pgfqpoint{1.449635in}{1.654610in}}%
\pgfpathlineto{\pgfqpoint{1.413671in}{1.591871in}}%
\pgfpathlineto{\pgfqpoint{1.347794in}{1.601326in}}%
\pgfpathlineto{\pgfqpoint{1.318834in}{1.613277in}}%
\pgfpathlineto{\pgfqpoint{1.270864in}{1.608874in}}%
\pgfpathlineto{\pgfqpoint{1.265887in}{1.616660in}}%
\pgfpathlineto{\pgfqpoint{1.280150in}{1.622140in}}%
\pgfpathlineto{\pgfqpoint{1.265658in}{1.633602in}}%
\pgfpathlineto{\pgfqpoint{1.280570in}{1.605752in}}%
\pgfpathlineto{\pgfqpoint{1.227954in}{1.611236in}}%
\pgfpathlineto{\pgfqpoint{1.187076in}{1.610965in}}%
\pgfpathlineto{\pgfqpoint{1.139625in}{1.578943in}}%
\pgfpathlineto{\pgfqpoint{1.180642in}{1.560442in}}%
\pgfpathlineto{\pgfqpoint{1.246249in}{1.533782in}}%
\pgfpathlineto{\pgfqpoint{1.253186in}{1.631560in}}%
\pgfpathlineto{\pgfqpoint{1.316734in}{1.567853in}}%
\pgfpathlineto{\pgfqpoint{1.344260in}{1.589259in}}%
\pgfpathlineto{\pgfqpoint{1.370880in}{1.601923in}}%
\pgfpathlineto{\pgfqpoint{1.377835in}{1.613470in}}%
\pgfpathlineto{\pgfqpoint{1.373127in}{1.576033in}}%
\pgfpathlineto{\pgfqpoint{1.418317in}{1.611289in}}%
\pgfpathlineto{\pgfqpoint{1.469761in}{1.660261in}}%
\pgfpathlineto{\pgfqpoint{1.483954in}{1.664299in}}%
\pgfpathlineto{\pgfqpoint{1.542339in}{1.672148in}}%
\pgfpathlineto{\pgfqpoint{1.571609in}{1.695463in}}%
\pgfpathlineto{\pgfqpoint{1.563031in}{1.663884in}}%
\pgfpathlineto{\pgfqpoint{1.601235in}{1.665503in}}%
\pgfpathlineto{\pgfqpoint{1.639483in}{1.660191in}}%
\pgfpathlineto{\pgfqpoint{1.646740in}{1.667203in}}%
\pgfpathlineto{\pgfqpoint{1.622595in}{1.700907in}}%
\pgfpathlineto{\pgfqpoint{1.603633in}{1.724805in}}%
\pgfpathlineto{\pgfqpoint{1.634214in}{1.776941in}}%
\pgfpathlineto{\pgfqpoint{1.640745in}{1.755852in}}%
\pgfpathlineto{\pgfqpoint{1.623605in}{1.781904in}}%
\pgfpathlineto{\pgfqpoint{1.631808in}{1.798759in}}%
\pgfpathlineto{\pgfqpoint{1.619518in}{1.794780in}}%
\pgfpathlineto{\pgfqpoint{1.600642in}{1.784976in}}%
\pgfpathlineto{\pgfqpoint{1.551117in}{1.781676in}}%
\pgfpathlineto{\pgfqpoint{1.549322in}{1.839349in}}%
\pgfpathlineto{\pgfqpoint{1.556168in}{1.853218in}}%
\pgfpathlineto{\pgfqpoint{1.560253in}{1.887338in}}%
\pgfpathlineto{\pgfqpoint{1.620056in}{1.896251in}}%
\pgfpathlineto{\pgfqpoint{1.596242in}{1.908282in}}%
\pgfpathlineto{\pgfqpoint{1.596000in}{1.916222in}}%
\pgfpathlineto{\pgfqpoint{1.558305in}{1.962560in}}%
\pgfpathlineto{\pgfqpoint{1.605589in}{1.935842in}}%
\pgfpathlineto{\pgfqpoint{1.601806in}{1.879225in}}%
\pgfpathlineto{\pgfqpoint{1.648517in}{1.859524in}}%
\pgfpathlineto{\pgfqpoint{1.677005in}{1.890826in}}%
\pgfpathlineto{\pgfqpoint{1.610108in}{1.844208in}}%
\pgfpathlineto{\pgfqpoint{1.605329in}{1.869140in}}%
\pgfpathlineto{\pgfqpoint{1.594573in}{1.875266in}}%
\pgfpathlineto{\pgfqpoint{1.684389in}{1.863923in}}%
\pgfpathlineto{\pgfqpoint{1.610881in}{1.854251in}}%
\pgfpathlineto{\pgfqpoint{1.618361in}{1.832544in}}%
\pgfpathlineto{\pgfqpoint{1.658181in}{1.832124in}}%
\pgfpathlineto{\pgfqpoint{1.607967in}{1.837428in}}%
\pgfpathlineto{\pgfqpoint{1.671607in}{1.860289in}}%
\pgfpathlineto{\pgfqpoint{1.650044in}{1.904943in}}%
\pgfpathlineto{\pgfqpoint{1.609331in}{1.929645in}}%
\pgfpathlineto{\pgfqpoint{1.608478in}{1.947659in}}%
\pgfpathlineto{\pgfqpoint{1.657991in}{1.980290in}}%
\pgfpathlineto{\pgfqpoint{1.722942in}{2.026455in}}%
\pgfpathlineto{\pgfqpoint{1.689883in}{1.957197in}}%
\pgfpathlineto{\pgfqpoint{1.628159in}{1.966633in}}%
\pgfpathlineto{\pgfqpoint{1.631570in}{2.003907in}}%
\pgfpathlineto{\pgfqpoint{1.619367in}{2.029174in}}%
\pgfpathlineto{\pgfqpoint{1.609823in}{2.033194in}}%
\pgfpathlineto{\pgfqpoint{1.606873in}{2.020981in}}%
\pgfpathlineto{\pgfqpoint{1.630282in}{2.071377in}}%
\pgfpathlineto{\pgfqpoint{1.656872in}{2.108481in}}%
\pgfpathlineto{\pgfqpoint{1.650886in}{2.123213in}}%
\pgfpathlineto{\pgfqpoint{1.698054in}{2.102961in}}%
\pgfpathlineto{\pgfqpoint{1.734463in}{2.110590in}}%
\pgfpathlineto{\pgfqpoint{1.669328in}{2.099396in}}%
\pgfpathlineto{\pgfqpoint{1.713853in}{2.080075in}}%
\pgfpathlineto{\pgfqpoint{1.708187in}{2.063189in}}%
\pgfpathlineto{\pgfqpoint{1.717500in}{2.023994in}}%
\pgfpathlineto{\pgfqpoint{1.710869in}{1.998430in}}%
\pgfpathlineto{\pgfqpoint{1.685987in}{1.950250in}}%
\pgfpathlineto{\pgfqpoint{1.700074in}{1.954302in}}%
\pgfpathlineto{\pgfqpoint{1.675674in}{1.965032in}}%
\pgfpathlineto{\pgfqpoint{1.698887in}{1.971305in}}%
\pgfpathlineto{\pgfqpoint{1.728642in}{2.014517in}}%
\pgfpathlineto{\pgfqpoint{1.671492in}{2.085302in}}%
\pgfpathlineto{\pgfqpoint{1.649140in}{2.071105in}}%
\pgfpathlineto{\pgfqpoint{1.621023in}{2.039394in}}%
\pgfpathlineto{\pgfqpoint{1.608150in}{2.034140in}}%
\pgfpathlineto{\pgfqpoint{1.678185in}{2.024948in}}%
\pgfpathlineto{\pgfqpoint{1.685856in}{2.031285in}}%
\pgfpathlineto{\pgfqpoint{1.686763in}{2.026901in}}%
\pgfpathlineto{\pgfqpoint{1.669766in}{2.006850in}}%
\pgfpathlineto{\pgfqpoint{1.665651in}{1.989058in}}%
\pgfpathlineto{\pgfqpoint{1.661852in}{2.017633in}}%
\pgfpathlineto{\pgfqpoint{1.659567in}{2.033138in}}%
\pgfpathlineto{\pgfqpoint{1.755280in}{2.052810in}}%
\pgfpathlineto{\pgfqpoint{1.858109in}{2.056780in}}%
\pgfpathlineto{\pgfqpoint{1.837207in}{1.989478in}}%
\pgfpathlineto{\pgfqpoint{1.804641in}{1.986353in}}%
\pgfpathlineto{\pgfqpoint{1.797913in}{2.043934in}}%
\pgfpathlineto{\pgfqpoint{1.769856in}{2.058741in}}%
\pgfpathlineto{\pgfqpoint{1.767719in}{2.042843in}}%
\pgfpathlineto{\pgfqpoint{1.680726in}{2.072079in}}%
\pgfpathlineto{\pgfqpoint{1.643148in}{2.034533in}}%
\pgfpathlineto{\pgfqpoint{1.634419in}{2.078914in}}%
\pgfpathlineto{\pgfqpoint{1.612205in}{2.065503in}}%
\pgfpathlineto{\pgfqpoint{1.619798in}{2.075961in}}%
\pgfpathlineto{\pgfqpoint{1.617653in}{2.060858in}}%
\pgfpathlineto{\pgfqpoint{1.530729in}{2.064311in}}%
\pgfpathlineto{\pgfqpoint{1.491051in}{1.998771in}}%
\pgfpathlineto{\pgfqpoint{1.506417in}{2.007057in}}%
\pgfpathlineto{\pgfqpoint{1.516623in}{2.037376in}}%
\pgfpathlineto{\pgfqpoint{1.569461in}{2.019893in}}%
\pgfpathlineto{\pgfqpoint{1.548411in}{2.065094in}}%
\pgfpathlineto{\pgfqpoint{1.607923in}{2.096387in}}%
\pgfpathlineto{\pgfqpoint{1.600809in}{2.082305in}}%
\pgfpathlineto{\pgfqpoint{1.569712in}{2.051134in}}%
\pgfpathlineto{\pgfqpoint{1.545596in}{2.033340in}}%
\pgfpathlineto{\pgfqpoint{1.536916in}{2.031324in}}%
\pgfpathlineto{\pgfqpoint{1.515500in}{2.013755in}}%
\pgfpathlineto{\pgfqpoint{1.545266in}{2.008839in}}%
\pgfpathlineto{\pgfqpoint{1.570142in}{2.021675in}}%
\pgfpathlineto{\pgfqpoint{1.623027in}{2.002498in}}%
\pgfpathlineto{\pgfqpoint{1.671318in}{2.033829in}}%
\pgfpathlineto{\pgfqpoint{1.626156in}{2.011446in}}%
\pgfpathlineto{\pgfqpoint{1.676627in}{1.988828in}}%
\pgfpathlineto{\pgfqpoint{1.700488in}{2.025297in}}%
\pgfpathlineto{\pgfqpoint{1.768817in}{1.972347in}}%
\pgfpathlineto{\pgfqpoint{1.822123in}{1.962974in}}%
\pgfpathlineto{\pgfqpoint{1.827822in}{1.954459in}}%
\pgfpathlineto{\pgfqpoint{1.868640in}{1.961145in}}%
\pgfpathlineto{\pgfqpoint{1.854168in}{2.014956in}}%
\pgfpathlineto{\pgfqpoint{1.882316in}{2.080451in}}%
\pgfpathlineto{\pgfqpoint{1.876594in}{2.087341in}}%
\pgfpathlineto{\pgfqpoint{1.846055in}{2.080070in}}%
\pgfpathlineto{\pgfqpoint{1.894216in}{2.047191in}}%
\pgfpathlineto{\pgfqpoint{1.894305in}{2.028302in}}%
\pgfpathlineto{\pgfqpoint{1.908741in}{1.984096in}}%
\pgfpathlineto{\pgfqpoint{1.907004in}{2.000057in}}%
\pgfpathlineto{\pgfqpoint{1.939211in}{2.003907in}}%
\pgfpathlineto{\pgfqpoint{1.924298in}{2.005981in}}%
\pgfpathlineto{\pgfqpoint{1.992162in}{1.992252in}}%
\pgfpathlineto{\pgfqpoint{2.003988in}{2.052114in}}%
\pgfpathlineto{\pgfqpoint{1.922851in}{2.045644in}}%
\pgfpathlineto{\pgfqpoint{1.876935in}{2.017634in}}%
\pgfpathlineto{\pgfqpoint{1.821729in}{2.011555in}}%
\pgfpathlineto{\pgfqpoint{1.804217in}{2.019921in}}%
\pgfpathlineto{\pgfqpoint{1.787045in}{2.008383in}}%
\pgfpathlineto{\pgfqpoint{1.790820in}{2.014406in}}%
\pgfpathlineto{\pgfqpoint{1.827850in}{1.946627in}}%
\pgfpathlineto{\pgfqpoint{1.864511in}{1.919928in}}%
\pgfpathlineto{\pgfqpoint{1.911455in}{1.860542in}}%
\pgfpathlineto{\pgfqpoint{1.877696in}{1.853660in}}%
\pgfpathlineto{\pgfqpoint{1.866946in}{1.876002in}}%
\pgfpathlineto{\pgfqpoint{1.865996in}{1.875602in}}%
\pgfpathlineto{\pgfqpoint{1.835639in}{1.896214in}}%
\pgfpathlineto{\pgfqpoint{1.838812in}{1.847545in}}%
\pgfpathlineto{\pgfqpoint{1.865425in}{1.867716in}}%
\pgfpathlineto{\pgfqpoint{1.843391in}{1.831175in}}%
\pgfpathlineto{\pgfqpoint{1.841524in}{1.857795in}}%
\pgfpathlineto{\pgfqpoint{1.807083in}{1.852355in}}%
\pgfpathlineto{\pgfqpoint{1.792023in}{1.900359in}}%
\pgfpathlineto{\pgfqpoint{1.755779in}{1.969927in}}%
\pgfpathlineto{\pgfqpoint{1.805765in}{1.964218in}}%
\pgfpathlineto{\pgfqpoint{1.852598in}{1.949641in}}%
\pgfpathlineto{\pgfqpoint{1.856256in}{1.943167in}}%
\pgfpathlineto{\pgfqpoint{1.863640in}{2.000913in}}%
\pgfpathlineto{\pgfqpoint{1.880186in}{1.979097in}}%
\pgfpathlineto{\pgfqpoint{1.848809in}{1.964847in}}%
\pgfpathlineto{\pgfqpoint{1.798895in}{1.932076in}}%
\pgfpathlineto{\pgfqpoint{1.741492in}{1.975317in}}%
\pgfpathlineto{\pgfqpoint{1.672870in}{1.975399in}}%
\pgfpathlineto{\pgfqpoint{1.725803in}{1.950612in}}%
\pgfpathlineto{\pgfqpoint{1.756227in}{1.988237in}}%
\pgfpathlineto{\pgfqpoint{1.721022in}{1.993816in}}%
\pgfpathlineto{\pgfqpoint{1.791642in}{1.984703in}}%
\pgfpathlineto{\pgfqpoint{1.827195in}{1.986553in}}%
\pgfpathlineto{\pgfqpoint{1.825769in}{1.993755in}}%
\pgfpathlineto{\pgfqpoint{1.860200in}{2.003006in}}%
\pgfpathlineto{\pgfqpoint{1.821310in}{2.030303in}}%
\pgfpathlineto{\pgfqpoint{1.794153in}{1.980707in}}%
\pgfpathlineto{\pgfqpoint{1.803395in}{1.993797in}}%
\pgfpathlineto{\pgfqpoint{1.709174in}{2.015093in}}%
\pgfpathlineto{\pgfqpoint{1.725300in}{1.990013in}}%
\pgfpathlineto{\pgfqpoint{1.704632in}{2.026983in}}%
\pgfpathlineto{\pgfqpoint{1.645372in}{2.021347in}}%
\pgfpathlineto{\pgfqpoint{1.629345in}{2.000795in}}%
\pgfpathlineto{\pgfqpoint{1.563898in}{1.966553in}}%
\pgfpathlineto{\pgfqpoint{1.554778in}{1.960975in}}%
\pgfpathlineto{\pgfqpoint{1.544642in}{1.971468in}}%
\pgfpathlineto{\pgfqpoint{1.523598in}{1.984528in}}%
\pgfpathlineto{\pgfqpoint{1.515840in}{1.994884in}}%
\pgfpathlineto{\pgfqpoint{1.531470in}{2.024920in}}%
\pgfpathlineto{\pgfqpoint{1.517771in}{2.013910in}}%
\pgfpathlineto{\pgfqpoint{1.547070in}{1.985746in}}%
\pgfpathlineto{\pgfqpoint{1.565882in}{1.971027in}}%
\pgfpathlineto{\pgfqpoint{1.577259in}{1.991816in}}%
\pgfpathlineto{\pgfqpoint{1.475568in}{2.013250in}}%
\pgfpathlineto{\pgfqpoint{1.487532in}{2.075888in}}%
\pgfpathlineto{\pgfqpoint{1.397869in}{2.108612in}}%
\pgfpathlineto{\pgfqpoint{1.361776in}{2.151178in}}%
\pgfpathlineto{\pgfqpoint{1.310982in}{2.130963in}}%
\pgfpathlineto{\pgfqpoint{1.198886in}{2.139109in}}%
\pgfpathlineto{\pgfqpoint{1.189091in}{2.125198in}}%
\pgfpathlineto{\pgfqpoint{1.156888in}{2.149092in}}%
\pgfpathlineto{\pgfqpoint{1.104421in}{2.141754in}}%
\pgfpathlineto{\pgfqpoint{1.158501in}{2.113274in}}%
\pgfpathlineto{\pgfqpoint{1.215830in}{2.111900in}}%
\pgfpathlineto{\pgfqpoint{1.306221in}{2.120059in}}%
\pgfpathlineto{\pgfqpoint{1.289023in}{2.126077in}}%
\pgfpathlineto{\pgfqpoint{1.310915in}{2.138527in}}%
\pgfpathlineto{\pgfqpoint{1.239024in}{2.118345in}}%
\pgfpathlineto{\pgfqpoint{1.266312in}{2.169619in}}%
\pgfpathlineto{\pgfqpoint{1.129391in}{2.205015in}}%
\pgfpathlineto{\pgfqpoint{1.188215in}{2.238346in}}%
\pgfpathlineto{\pgfqpoint{1.178336in}{2.240559in}}%
\pgfpathlineto{\pgfqpoint{1.178451in}{2.237355in}}%
\pgfpathlineto{\pgfqpoint{1.158520in}{2.225265in}}%
\pgfpathlineto{\pgfqpoint{1.155216in}{2.256460in}}%
\pgfpathlineto{\pgfqpoint{1.174700in}{2.210957in}}%
\pgfpathlineto{\pgfqpoint{1.254699in}{2.204277in}}%
\pgfpathlineto{\pgfqpoint{1.251375in}{2.202343in}}%
\pgfpathlineto{\pgfqpoint{1.337136in}{2.215483in}}%
\pgfpathlineto{\pgfqpoint{1.396893in}{2.221649in}}%
\pgfpathlineto{\pgfqpoint{1.421849in}{2.239625in}}%
\pgfpathlineto{\pgfqpoint{1.416232in}{2.267881in}}%
\pgfpathlineto{\pgfqpoint{1.399130in}{2.255527in}}%
\pgfpathlineto{\pgfqpoint{1.428926in}{2.254877in}}%
\pgfpathlineto{\pgfqpoint{1.468854in}{2.259035in}}%
\pgfpathlineto{\pgfqpoint{1.492358in}{2.191022in}}%
\pgfpathlineto{\pgfqpoint{1.507847in}{2.190817in}}%
\pgfpathlineto{\pgfqpoint{1.545356in}{2.182156in}}%
\pgfpathlineto{\pgfqpoint{1.603058in}{2.100294in}}%
\pgfpathlineto{\pgfqpoint{1.596281in}{2.158650in}}%
\pgfpathlineto{\pgfqpoint{1.571026in}{2.124684in}}%
\pgfpathlineto{\pgfqpoint{1.545635in}{2.167787in}}%
\pgfpathlineto{\pgfqpoint{1.613211in}{2.174691in}}%
\pgfpathlineto{\pgfqpoint{1.645449in}{2.166197in}}%
\pgfpathlineto{\pgfqpoint{1.647536in}{2.223731in}}%
\pgfpathlineto{\pgfqpoint{1.635739in}{2.227898in}}%
\pgfpathlineto{\pgfqpoint{1.608823in}{2.240471in}}%
\pgfpathlineto{\pgfqpoint{1.668814in}{2.286088in}}%
\pgfpathlineto{\pgfqpoint{1.743970in}{2.245928in}}%
\pgfpathlineto{\pgfqpoint{1.759835in}{2.321856in}}%
\pgfpathlineto{\pgfqpoint{1.765283in}{2.313834in}}%
\pgfpathlineto{\pgfqpoint{1.694438in}{2.347668in}}%
\pgfpathlineto{\pgfqpoint{1.722903in}{2.358710in}}%
\pgfpathlineto{\pgfqpoint{1.749548in}{2.318602in}}%
\pgfpathlineto{\pgfqpoint{1.826004in}{2.278882in}}%
\pgfpathlineto{\pgfqpoint{1.815039in}{2.284168in}}%
\pgfpathlineto{\pgfqpoint{1.792081in}{2.290916in}}%
\pgfpathlineto{\pgfqpoint{1.860791in}{2.290363in}}%
\pgfpathlineto{\pgfqpoint{1.905920in}{2.281755in}}%
\pgfpathlineto{\pgfqpoint{1.965860in}{2.285647in}}%
\pgfpathlineto{\pgfqpoint{1.964960in}{2.318005in}}%
\pgfpathlineto{\pgfqpoint{1.965273in}{2.291630in}}%
\pgfpathlineto{\pgfqpoint{1.972516in}{2.299056in}}%
\pgfpathlineto{\pgfqpoint{1.920974in}{2.345260in}}%
\pgfpathlineto{\pgfqpoint{1.884447in}{2.331362in}}%
\pgfpathlineto{\pgfqpoint{1.918400in}{2.322182in}}%
\pgfpathlineto{\pgfqpoint{1.983812in}{2.298769in}}%
\pgfpathlineto{\pgfqpoint{1.975724in}{2.280630in}}%
\pgfpathlineto{\pgfqpoint{1.955031in}{2.294732in}}%
\pgfpathlineto{\pgfqpoint{2.011131in}{2.289336in}}%
\pgfpathlineto{\pgfqpoint{2.009129in}{2.245047in}}%
\pgfpathlineto{\pgfqpoint{1.955530in}{2.206970in}}%
\pgfpathlineto{\pgfqpoint{1.931203in}{2.249078in}}%
\pgfpathlineto{\pgfqpoint{1.926275in}{2.239191in}}%
\pgfpathlineto{\pgfqpoint{1.888076in}{2.252242in}}%
\pgfpathlineto{\pgfqpoint{1.887379in}{2.229763in}}%
\pgfpathlineto{\pgfqpoint{1.908921in}{2.228952in}}%
\pgfpathlineto{\pgfqpoint{1.912853in}{2.294889in}}%
\pgfpathlineto{\pgfqpoint{1.874049in}{2.280715in}}%
\pgfpathlineto{\pgfqpoint{1.937329in}{2.296713in}}%
\pgfpathlineto{\pgfqpoint{1.929185in}{2.218509in}}%
\pgfpathlineto{\pgfqpoint{1.929248in}{2.229865in}}%
\pgfpathlineto{\pgfqpoint{1.879088in}{2.253695in}}%
\pgfpathlineto{\pgfqpoint{1.915399in}{2.240988in}}%
\pgfpathlineto{\pgfqpoint{1.842818in}{2.234101in}}%
\pgfpathlineto{\pgfqpoint{1.795592in}{2.193050in}}%
\pgfpathlineto{\pgfqpoint{1.768029in}{2.169865in}}%
\pgfpathlineto{\pgfqpoint{1.753079in}{2.129149in}}%
\pgfpathlineto{\pgfqpoint{1.801303in}{2.102546in}}%
\pgfpathlineto{\pgfqpoint{1.886151in}{2.127412in}}%
\pgfpathlineto{\pgfqpoint{1.896998in}{2.167098in}}%
\pgfpathlineto{\pgfqpoint{1.912685in}{2.192196in}}%
\pgfpathlineto{\pgfqpoint{1.865194in}{2.145203in}}%
\pgfpathlineto{\pgfqpoint{1.895336in}{2.100711in}}%
\pgfpathlineto{\pgfqpoint{1.960235in}{2.069141in}}%
\pgfpathlineto{\pgfqpoint{1.950928in}{2.055860in}}%
\pgfpathlineto{\pgfqpoint{1.978914in}{2.088509in}}%
\pgfpathlineto{\pgfqpoint{1.986878in}{2.089987in}}%
\pgfpathlineto{\pgfqpoint{2.020933in}{2.076593in}}%
\pgfpathlineto{\pgfqpoint{1.903507in}{2.086169in}}%
\pgfpathlineto{\pgfqpoint{1.887291in}{2.039598in}}%
\pgfpathlineto{\pgfqpoint{1.854371in}{2.067012in}}%
\pgfpathlineto{\pgfqpoint{1.863856in}{2.148786in}}%
\pgfpathlineto{\pgfqpoint{1.836561in}{2.110670in}}%
\pgfpathlineto{\pgfqpoint{1.811793in}{2.089325in}}%
\pgfpathlineto{\pgfqpoint{1.845839in}{2.073125in}}%
\pgfpathlineto{\pgfqpoint{1.823933in}{2.102749in}}%
\pgfpathlineto{\pgfqpoint{1.860132in}{2.091500in}}%
\pgfpathlineto{\pgfqpoint{1.848635in}{2.040288in}}%
\pgfpathlineto{\pgfqpoint{1.854697in}{2.071085in}}%
\pgfpathlineto{\pgfqpoint{1.823140in}{2.037310in}}%
\pgfpathlineto{\pgfqpoint{1.835824in}{2.041915in}}%
\pgfpathlineto{\pgfqpoint{1.834023in}{2.026516in}}%
\pgfpathlineto{\pgfqpoint{1.827241in}{1.982395in}}%
\pgfpathlineto{\pgfqpoint{1.834158in}{2.008065in}}%
\pgfpathlineto{\pgfqpoint{1.855999in}{1.965815in}}%
\pgfpathlineto{\pgfqpoint{1.841938in}{1.984393in}}%
\pgfpathlineto{\pgfqpoint{1.785899in}{1.990181in}}%
\pgfpathlineto{\pgfqpoint{1.781770in}{1.984496in}}%
\pgfpathlineto{\pgfqpoint{1.717325in}{2.010205in}}%
\pgfpathlineto{\pgfqpoint{1.718031in}{2.051538in}}%
\pgfpathlineto{\pgfqpoint{1.708618in}{2.047677in}}%
\pgfpathlineto{\pgfqpoint{1.720332in}{2.084664in}}%
\pgfpathlineto{\pgfqpoint{1.752635in}{2.056813in}}%
\pgfpathlineto{\pgfqpoint{1.737210in}{2.036560in}}%
\pgfpathlineto{\pgfqpoint{1.627637in}{2.037536in}}%
\pgfpathlineto{\pgfqpoint{1.649112in}{2.008053in}}%
\pgfpathlineto{\pgfqpoint{1.626243in}{1.979533in}}%
\pgfpathlineto{\pgfqpoint{1.597037in}{2.023086in}}%
\pgfpathlineto{\pgfqpoint{1.554638in}{2.053774in}}%
\pgfpathlineto{\pgfqpoint{1.509942in}{1.991432in}}%
\pgfpathlineto{\pgfqpoint{1.493471in}{2.011356in}}%
\pgfpathlineto{\pgfqpoint{1.550764in}{2.036847in}}%
\pgfpathlineto{\pgfqpoint{1.560343in}{1.953859in}}%
\pgfpathlineto{\pgfqpoint{1.567142in}{1.873130in}}%
\pgfpathlineto{\pgfqpoint{1.571406in}{1.902206in}}%
\pgfpathlineto{\pgfqpoint{1.561544in}{1.904137in}}%
\pgfpathlineto{\pgfqpoint{1.588584in}{1.882299in}}%
\pgfpathlineto{\pgfqpoint{1.577125in}{1.918962in}}%
\pgfpathlineto{\pgfqpoint{1.602286in}{1.908327in}}%
\pgfpathlineto{\pgfqpoint{1.533164in}{1.918202in}}%
\pgfpathlineto{\pgfqpoint{1.553189in}{1.898175in}}%
\pgfpathlineto{\pgfqpoint{1.493730in}{1.860074in}}%
\pgfpathlineto{\pgfqpoint{1.484221in}{1.866237in}}%
\pgfpathlineto{\pgfqpoint{1.529265in}{1.853731in}}%
\pgfpathlineto{\pgfqpoint{1.527504in}{1.787593in}}%
\pgfpathlineto{\pgfqpoint{1.575052in}{1.783547in}}%
\pgfpathlineto{\pgfqpoint{1.598073in}{1.747518in}}%
\pgfpathlineto{\pgfqpoint{1.586233in}{1.747208in}}%
\pgfpathlineto{\pgfqpoint{1.651910in}{1.711369in}}%
\pgfpathlineto{\pgfqpoint{1.671295in}{1.664347in}}%
\pgfpathlineto{\pgfqpoint{1.710681in}{1.676112in}}%
\pgfpathlineto{\pgfqpoint{1.700375in}{1.620651in}}%
\pgfpathlineto{\pgfqpoint{1.664533in}{1.541495in}}%
\pgfpathlineto{\pgfqpoint{1.629255in}{1.552140in}}%
\pgfpathlineto{\pgfqpoint{1.660807in}{1.556806in}}%
\pgfpathlineto{\pgfqpoint{1.660561in}{1.554543in}}%
\pgfpathlineto{\pgfqpoint{1.685885in}{1.564849in}}%
\pgfpathlineto{\pgfqpoint{1.719154in}{1.552928in}}%
\pgfpathlineto{\pgfqpoint{1.678288in}{1.631507in}}%
\pgfpathlineto{\pgfqpoint{1.712873in}{1.683219in}}%
\pgfpathlineto{\pgfqpoint{1.771319in}{1.627509in}}%
\pgfpathlineto{\pgfqpoint{1.837728in}{1.604300in}}%
\pgfpathlineto{\pgfqpoint{1.758719in}{1.592280in}}%
\pgfpathlineto{\pgfqpoint{1.804994in}{1.616600in}}%
\pgfpathlineto{\pgfqpoint{1.794754in}{1.652855in}}%
\pgfpathlineto{\pgfqpoint{1.802599in}{1.639507in}}%
\pgfpathlineto{\pgfqpoint{1.762002in}{1.638974in}}%
\pgfpathlineto{\pgfqpoint{1.709536in}{1.663958in}}%
\pgfpathlineto{\pgfqpoint{1.661657in}{1.659022in}}%
\pgfpathlineto{\pgfqpoint{1.728225in}{1.686758in}}%
\pgfpathlineto{\pgfqpoint{1.698318in}{1.715944in}}%
\pgfpathlineto{\pgfqpoint{1.729330in}{1.691063in}}%
\pgfpathlineto{\pgfqpoint{1.703603in}{1.654801in}}%
\pgfpathlineto{\pgfqpoint{1.705466in}{1.667019in}}%
\pgfpathlineto{\pgfqpoint{1.689758in}{1.632450in}}%
\pgfpathlineto{\pgfqpoint{1.721538in}{1.631105in}}%
\pgfpathlineto{\pgfqpoint{1.668804in}{1.630574in}}%
\pgfpathlineto{\pgfqpoint{1.744582in}{1.635348in}}%
\pgfpathlineto{\pgfqpoint{1.762685in}{1.606796in}}%
\pgfpathlineto{\pgfqpoint{1.667829in}{1.610985in}}%
\pgfpathlineto{\pgfqpoint{1.637590in}{1.554816in}}%
\pgfpathlineto{\pgfqpoint{1.606884in}{1.529531in}}%
\pgfpathlineto{\pgfqpoint{1.660994in}{1.578786in}}%
\pgfpathlineto{\pgfqpoint{1.614512in}{1.597560in}}%
\pgfpathlineto{\pgfqpoint{1.569321in}{1.578726in}}%
\pgfpathlineto{\pgfqpoint{1.613431in}{1.575009in}}%
\pgfpathlineto{\pgfqpoint{1.595502in}{1.602662in}}%
\pgfpathlineto{\pgfqpoint{1.605361in}{1.580444in}}%
\pgfpathlineto{\pgfqpoint{1.549687in}{1.526518in}}%
\pgfpathlineto{\pgfqpoint{1.602288in}{1.500855in}}%
\pgfpathlineto{\pgfqpoint{1.577289in}{1.537713in}}%
\pgfpathlineto{\pgfqpoint{1.506006in}{1.501706in}}%
\pgfpathlineto{\pgfqpoint{1.510582in}{1.450226in}}%
\pgfpathlineto{\pgfqpoint{1.489221in}{1.497525in}}%
\pgfpathlineto{\pgfqpoint{1.550541in}{1.498552in}}%
\pgfpathlineto{\pgfqpoint{1.492339in}{1.431019in}}%
\pgfpathlineto{\pgfqpoint{1.443438in}{1.457958in}}%
\pgfpathlineto{\pgfqpoint{1.463121in}{1.406800in}}%
\pgfpathlineto{\pgfqpoint{1.424743in}{1.414442in}}%
\pgfpathlineto{\pgfqpoint{1.390015in}{1.455608in}}%
\pgfpathlineto{\pgfqpoint{1.367086in}{1.493459in}}%
\pgfpathlineto{\pgfqpoint{1.311968in}{1.494100in}}%
\pgfpathlineto{\pgfqpoint{1.329996in}{1.471265in}}%
\pgfpathlineto{\pgfqpoint{1.294422in}{1.476649in}}%
\pgfpathlineto{\pgfqpoint{1.291269in}{1.436826in}}%
\pgfpathlineto{\pgfqpoint{1.213482in}{1.367126in}}%
\pgfpathlineto{\pgfqpoint{1.238499in}{1.419454in}}%
\pgfpathlineto{\pgfqpoint{1.194678in}{1.442404in}}%
\pgfpathlineto{\pgfqpoint{1.192027in}{1.458108in}}%
\pgfpathlineto{\pgfqpoint{1.238689in}{1.478060in}}%
\pgfpathlineto{\pgfqpoint{1.201097in}{1.519982in}}%
\pgfpathlineto{\pgfqpoint{1.189269in}{1.528860in}}%
\pgfpathlineto{\pgfqpoint{1.203467in}{1.557509in}}%
\pgfpathlineto{\pgfqpoint{1.214198in}{1.522031in}}%
\pgfpathlineto{\pgfqpoint{1.183839in}{1.527466in}}%
\pgfpathlineto{\pgfqpoint{1.142561in}{1.572095in}}%
\pgfpathlineto{\pgfqpoint{1.138346in}{1.568563in}}%
\pgfpathlineto{\pgfqpoint{1.188813in}{1.600519in}}%
\pgfpathlineto{\pgfqpoint{1.180241in}{1.570859in}}%
\pgfpathlineto{\pgfqpoint{1.156693in}{1.540472in}}%
\pgfpathlineto{\pgfqpoint{1.200591in}{1.561608in}}%
\pgfpathlineto{\pgfqpoint{1.196082in}{1.540976in}}%
\pgfpathlineto{\pgfqpoint{1.215579in}{1.522110in}}%
\pgfpathlineto{\pgfqpoint{1.127374in}{1.536572in}}%
\pgfpathlineto{\pgfqpoint{1.070601in}{1.553250in}}%
\pgfpathlineto{\pgfqpoint{1.048598in}{1.530372in}}%
\pgfpathlineto{\pgfqpoint{1.063101in}{1.562155in}}%
\pgfpathlineto{\pgfqpoint{1.058707in}{1.512754in}}%
\pgfpathlineto{\pgfqpoint{1.096215in}{1.502214in}}%
\pgfpathlineto{\pgfqpoint{1.115826in}{1.453445in}}%
\pgfpathlineto{\pgfqpoint{1.154602in}{1.386563in}}%
\pgfpathlineto{\pgfqpoint{1.289996in}{1.351495in}}%
\pgfpathlineto{\pgfqpoint{1.332901in}{1.408981in}}%
\pgfpathlineto{\pgfqpoint{1.349443in}{1.376431in}}%
\pgfpathlineto{\pgfqpoint{1.287062in}{1.362762in}}%
\pgfpathlineto{\pgfqpoint{1.276746in}{1.366375in}}%
\pgfpathlineto{\pgfqpoint{1.330627in}{1.366836in}}%
\pgfpathlineto{\pgfqpoint{1.406881in}{1.397508in}}%
\pgfpathlineto{\pgfqpoint{1.454794in}{1.419631in}}%
\pgfpathlineto{\pgfqpoint{1.467730in}{1.376475in}}%
\pgfpathlineto{\pgfqpoint{1.463833in}{1.390029in}}%
\pgfpathlineto{\pgfqpoint{1.426132in}{1.400880in}}%
\pgfpathlineto{\pgfqpoint{1.406944in}{1.451913in}}%
\pgfpathlineto{\pgfqpoint{1.429198in}{1.439011in}}%
\pgfpathlineto{\pgfqpoint{1.393037in}{1.440808in}}%
\pgfpathlineto{\pgfqpoint{1.447093in}{1.449926in}}%
\pgfpathlineto{\pgfqpoint{1.422116in}{1.476150in}}%
\pgfpathlineto{\pgfqpoint{1.387919in}{1.489742in}}%
\pgfpathlineto{\pgfqpoint{1.380405in}{1.491271in}}%
\pgfpathlineto{\pgfqpoint{1.410574in}{1.508994in}}%
\pgfpathlineto{\pgfqpoint{1.392682in}{1.545007in}}%
\pgfpathlineto{\pgfqpoint{1.295406in}{1.550252in}}%
\pgfpathlineto{\pgfqpoint{1.309510in}{1.654480in}}%
\pgfpathlineto{\pgfqpoint{1.273221in}{1.669109in}}%
\pgfpathlineto{\pgfqpoint{1.337747in}{1.741810in}}%
\pgfpathlineto{\pgfqpoint{1.362524in}{1.774916in}}%
\pgfpathlineto{\pgfqpoint{1.383217in}{1.846304in}}%
\pgfpathlineto{\pgfqpoint{1.390053in}{1.903487in}}%
\pgfpathlineto{\pgfqpoint{1.488673in}{1.930281in}}%
\pgfpathlineto{\pgfqpoint{1.486320in}{1.917583in}}%
\pgfpathlineto{\pgfqpoint{1.504457in}{1.940042in}}%
\pgfpathlineto{\pgfqpoint{1.474428in}{1.936524in}}%
\pgfpathlineto{\pgfqpoint{1.518750in}{1.968447in}}%
\pgfpathlineto{\pgfqpoint{1.557450in}{2.028286in}}%
\pgfpathlineto{\pgfqpoint{1.610526in}{1.973002in}}%
\pgfpathlineto{\pgfqpoint{1.598682in}{1.967606in}}%
\pgfpathlineto{\pgfqpoint{1.619444in}{1.954155in}}%
\pgfpathlineto{\pgfqpoint{1.576476in}{1.947063in}}%
\pgfpathlineto{\pgfqpoint{1.593869in}{1.950078in}}%
\pgfpathlineto{\pgfqpoint{1.594049in}{1.924347in}}%
\pgfpathlineto{\pgfqpoint{1.601288in}{1.902593in}}%
\pgfpathlineto{\pgfqpoint{1.548309in}{1.885002in}}%
\pgfpathlineto{\pgfqpoint{1.589388in}{1.914549in}}%
\pgfpathlineto{\pgfqpoint{1.537169in}{1.932303in}}%
\pgfpathlineto{\pgfqpoint{1.568442in}{1.906697in}}%
\pgfpathlineto{\pgfqpoint{1.573082in}{1.984342in}}%
\pgfpathlineto{\pgfqpoint{1.562736in}{1.965695in}}%
\pgfpathlineto{\pgfqpoint{1.571078in}{2.029856in}}%
\pgfpathlineto{\pgfqpoint{1.562096in}{2.035484in}}%
\pgfpathlineto{\pgfqpoint{1.515317in}{2.058786in}}%
\pgfpathlineto{\pgfqpoint{1.519675in}{2.009904in}}%
\pgfpathlineto{\pgfqpoint{1.561282in}{2.054146in}}%
\pgfpathlineto{\pgfqpoint{1.547582in}{2.015503in}}%
\pgfpathlineto{\pgfqpoint{1.623256in}{2.017335in}}%
\pgfpathlineto{\pgfqpoint{1.623471in}{2.036924in}}%
\pgfpathlineto{\pgfqpoint{1.666355in}{2.070893in}}%
\pgfpathlineto{\pgfqpoint{1.614921in}{2.091129in}}%
\pgfpathlineto{\pgfqpoint{1.557521in}{2.055249in}}%
\pgfpathlineto{\pgfqpoint{1.516668in}{2.066357in}}%
\pgfpathlineto{\pgfqpoint{1.537820in}{2.060267in}}%
\pgfpathlineto{\pgfqpoint{1.521193in}{2.047321in}}%
\pgfpathlineto{\pgfqpoint{1.489651in}{2.031855in}}%
\pgfpathlineto{\pgfqpoint{1.460220in}{2.030037in}}%
\pgfpathlineto{\pgfqpoint{1.431013in}{1.976596in}}%
\pgfpathlineto{\pgfqpoint{1.461898in}{2.002015in}}%
\pgfpathlineto{\pgfqpoint{1.474962in}{1.989645in}}%
\pgfpathlineto{\pgfqpoint{1.534222in}{1.986656in}}%
\pgfpathlineto{\pgfqpoint{1.593836in}{1.963768in}}%
\pgfpathlineto{\pgfqpoint{1.644524in}{1.901518in}}%
\pgfpathlineto{\pgfqpoint{1.641479in}{1.928199in}}%
\pgfpathlineto{\pgfqpoint{1.610669in}{1.945197in}}%
\pgfpathlineto{\pgfqpoint{1.579054in}{1.924635in}}%
\pgfpathlineto{\pgfqpoint{1.560624in}{1.946676in}}%
\pgfpathlineto{\pgfqpoint{1.578463in}{1.913567in}}%
\pgfpathlineto{\pgfqpoint{1.634371in}{1.927519in}}%
\pgfpathlineto{\pgfqpoint{1.630938in}{1.935543in}}%
\pgfpathlineto{\pgfqpoint{1.652204in}{1.995725in}}%
\pgfpathlineto{\pgfqpoint{1.706382in}{2.035196in}}%
\pgfpathlineto{\pgfqpoint{1.763860in}{2.005252in}}%
\pgfpathlineto{\pgfqpoint{1.742387in}{2.017388in}}%
\pgfpathlineto{\pgfqpoint{1.776530in}{2.029759in}}%
\pgfpathlineto{\pgfqpoint{1.798313in}{1.991572in}}%
\pgfpathlineto{\pgfqpoint{1.884194in}{2.062654in}}%
\pgfpathlineto{\pgfqpoint{1.955097in}{2.090282in}}%
\pgfpathlineto{\pgfqpoint{1.935082in}{2.056773in}}%
\pgfpathlineto{\pgfqpoint{1.935007in}{2.131961in}}%
\pgfpathlineto{\pgfqpoint{1.923731in}{2.180432in}}%
\pgfpathlineto{\pgfqpoint{1.952943in}{2.194584in}}%
\pgfpathlineto{\pgfqpoint{1.960506in}{2.193204in}}%
\pgfpathlineto{\pgfqpoint{1.999623in}{2.212695in}}%
\pgfpathlineto{\pgfqpoint{2.027339in}{2.186447in}}%
\pgfpathlineto{\pgfqpoint{2.057982in}{2.215848in}}%
\pgfpathlineto{\pgfqpoint{2.026007in}{2.174091in}}%
\pgfpathlineto{\pgfqpoint{2.036735in}{2.254949in}}%
\pgfpathlineto{\pgfqpoint{2.114017in}{2.225452in}}%
\pgfpathlineto{\pgfqpoint{2.168535in}{2.204099in}}%
\pgfpathlineto{\pgfqpoint{2.087940in}{2.233926in}}%
\pgfpathlineto{\pgfqpoint{2.033579in}{2.235801in}}%
\pgfpathlineto{\pgfqpoint{2.077973in}{2.167882in}}%
\pgfpathlineto{\pgfqpoint{2.103507in}{2.231627in}}%
\pgfpathlineto{\pgfqpoint{2.088588in}{2.221837in}}%
\pgfpathlineto{\pgfqpoint{2.128615in}{2.208979in}}%
\pgfpathlineto{\pgfqpoint{2.107851in}{2.168717in}}%
\pgfpathlineto{\pgfqpoint{2.045414in}{2.136643in}}%
\pgfpathlineto{\pgfqpoint{2.108813in}{2.121760in}}%
\pgfpathlineto{\pgfqpoint{2.004873in}{2.096981in}}%
\pgfpathlineto{\pgfqpoint{1.967233in}{2.185142in}}%
\pgfpathlineto{\pgfqpoint{2.012545in}{2.220260in}}%
\pgfpathlineto{\pgfqpoint{2.053644in}{2.210002in}}%
\pgfpathlineto{\pgfqpoint{2.044568in}{2.236341in}}%
\pgfpathlineto{\pgfqpoint{1.985372in}{2.216162in}}%
\pgfpathlineto{\pgfqpoint{1.985386in}{2.235141in}}%
\pgfpathlineto{\pgfqpoint{1.970808in}{2.278600in}}%
\pgfpathlineto{\pgfqpoint{1.941458in}{2.232625in}}%
\pgfpathlineto{\pgfqpoint{1.969936in}{2.164028in}}%
\pgfpathlineto{\pgfqpoint{1.970733in}{2.153716in}}%
\pgfpathlineto{\pgfqpoint{1.909448in}{2.139535in}}%
\pgfpathlineto{\pgfqpoint{1.923347in}{2.161648in}}%
\pgfpathlineto{\pgfqpoint{1.981958in}{2.136593in}}%
\pgfpathlineto{\pgfqpoint{1.942444in}{2.120323in}}%
\pgfpathlineto{\pgfqpoint{1.907143in}{2.141211in}}%
\pgfpathlineto{\pgfqpoint{1.948180in}{2.130556in}}%
\pgfpathlineto{\pgfqpoint{1.923829in}{2.146917in}}%
\pgfpathlineto{\pgfqpoint{1.888450in}{2.146024in}}%
\pgfpathlineto{\pgfqpoint{1.875099in}{2.190317in}}%
\pgfpathlineto{\pgfqpoint{1.927757in}{2.131871in}}%
\pgfpathlineto{\pgfqpoint{1.937121in}{2.134696in}}%
\pgfpathlineto{\pgfqpoint{1.969950in}{2.182189in}}%
\pgfpathlineto{\pgfqpoint{1.956713in}{2.217895in}}%
\pgfpathlineto{\pgfqpoint{1.984598in}{2.194649in}}%
\pgfpathlineto{\pgfqpoint{1.935361in}{2.178828in}}%
\pgfpathlineto{\pgfqpoint{1.897748in}{2.127645in}}%
\pgfpathlineto{\pgfqpoint{1.861723in}{2.128373in}}%
\pgfpathlineto{\pgfqpoint{1.899933in}{2.141948in}}%
\pgfpathlineto{\pgfqpoint{1.874255in}{2.097572in}}%
\pgfpathlineto{\pgfqpoint{1.846116in}{2.125703in}}%
\pgfpathlineto{\pgfqpoint{1.849973in}{2.068510in}}%
\pgfpathlineto{\pgfqpoint{1.877824in}{2.030518in}}%
\pgfpathlineto{\pgfqpoint{1.898263in}{1.973169in}}%
\pgfpathlineto{\pgfqpoint{1.902918in}{1.999221in}}%
\pgfpathlineto{\pgfqpoint{1.911176in}{2.022046in}}%
\pgfpathlineto{\pgfqpoint{1.942885in}{2.067337in}}%
\pgfpathlineto{\pgfqpoint{1.943073in}{2.078438in}}%
\pgfpathlineto{\pgfqpoint{1.958026in}{2.026344in}}%
\pgfpathlineto{\pgfqpoint{1.911137in}{2.011583in}}%
\pgfpathlineto{\pgfqpoint{1.916985in}{2.002827in}}%
\pgfpathlineto{\pgfqpoint{1.921270in}{1.991020in}}%
\pgfpathlineto{\pgfqpoint{2.026045in}{1.964694in}}%
\pgfpathlineto{\pgfqpoint{1.954979in}{1.993794in}}%
\pgfpathlineto{\pgfqpoint{1.930607in}{1.949773in}}%
\pgfpathlineto{\pgfqpoint{1.909508in}{2.011892in}}%
\pgfpathlineto{\pgfqpoint{1.943446in}{2.036284in}}%
\pgfpathlineto{\pgfqpoint{1.906564in}{2.031943in}}%
\pgfpathlineto{\pgfqpoint{1.848378in}{2.030691in}}%
\pgfpathlineto{\pgfqpoint{1.713817in}{2.084627in}}%
\pgfpathlineto{\pgfqpoint{1.717653in}{2.138882in}}%
\pgfpathlineto{\pgfqpoint{1.766090in}{2.107831in}}%
\pgfpathlineto{\pgfqpoint{1.707946in}{2.144084in}}%
\pgfpathlineto{\pgfqpoint{1.709815in}{2.148167in}}%
\pgfpathlineto{\pgfqpoint{1.750679in}{2.181679in}}%
\pgfpathlineto{\pgfqpoint{1.804346in}{2.199488in}}%
\pgfpathlineto{\pgfqpoint{1.857044in}{2.204487in}}%
\pgfpathlineto{\pgfqpoint{1.920409in}{2.182739in}}%
\pgfpathlineto{\pgfqpoint{1.860921in}{2.168824in}}%
\pgfpathlineto{\pgfqpoint{1.827374in}{2.116172in}}%
\pgfpathlineto{\pgfqpoint{1.847678in}{2.146937in}}%
\pgfpathlineto{\pgfqpoint{1.832730in}{2.183480in}}%
\pgfpathlineto{\pgfqpoint{1.863224in}{2.278738in}}%
\pgfpathlineto{\pgfqpoint{1.989774in}{2.291288in}}%
\pgfpathlineto{\pgfqpoint{1.982108in}{2.218632in}}%
\pgfpathlineto{\pgfqpoint{1.975049in}{2.211579in}}%
\pgfpathlineto{\pgfqpoint{2.002485in}{2.163231in}}%
\pgfpathlineto{\pgfqpoint{2.032436in}{2.105078in}}%
\pgfpathlineto{\pgfqpoint{1.969210in}{2.069078in}}%
\pgfpathlineto{\pgfqpoint{2.023075in}{2.092533in}}%
\pgfpathlineto{\pgfqpoint{1.993182in}{2.092932in}}%
\pgfpathlineto{\pgfqpoint{1.975931in}{2.052043in}}%
\pgfpathlineto{\pgfqpoint{1.911806in}{2.094483in}}%
\pgfpathlineto{\pgfqpoint{1.830849in}{2.077897in}}%
\pgfpathlineto{\pgfqpoint{1.824422in}{2.073377in}}%
\pgfpathlineto{\pgfqpoint{1.839479in}{2.029137in}}%
\pgfpathlineto{\pgfqpoint{1.912745in}{2.013010in}}%
\pgfpathlineto{\pgfqpoint{2.006877in}{1.942549in}}%
\pgfpathlineto{\pgfqpoint{1.972101in}{1.889565in}}%
\pgfpathlineto{\pgfqpoint{2.024878in}{1.882373in}}%
\pgfpathlineto{\pgfqpoint{2.038189in}{1.927916in}}%
\pgfpathlineto{\pgfqpoint{1.998831in}{1.931394in}}%
\pgfpathlineto{\pgfqpoint{2.042983in}{1.959682in}}%
\pgfpathlineto{\pgfqpoint{2.058463in}{2.022539in}}%
\pgfpathlineto{\pgfqpoint{2.077448in}{1.984399in}}%
\pgfpathlineto{\pgfqpoint{2.037202in}{2.002871in}}%
\pgfpathlineto{\pgfqpoint{2.064948in}{2.026385in}}%
\pgfpathlineto{\pgfqpoint{2.034035in}{2.053536in}}%
\pgfpathlineto{\pgfqpoint{2.046094in}{2.028041in}}%
\pgfpathlineto{\pgfqpoint{2.018426in}{2.048868in}}%
\pgfpathlineto{\pgfqpoint{2.009106in}{2.053163in}}%
\pgfpathlineto{\pgfqpoint{2.025037in}{2.007529in}}%
\pgfpathlineto{\pgfqpoint{1.997391in}{2.048670in}}%
\pgfpathlineto{\pgfqpoint{1.962525in}{1.996764in}}%
\pgfpathlineto{\pgfqpoint{1.957381in}{2.126444in}}%
\pgfpathlineto{\pgfqpoint{1.900958in}{2.072808in}}%
\pgfpathlineto{\pgfqpoint{1.919235in}{2.055904in}}%
\pgfpathlineto{\pgfqpoint{1.829299in}{2.089329in}}%
\pgfpathlineto{\pgfqpoint{1.834781in}{2.111028in}}%
\pgfpathlineto{\pgfqpoint{1.907309in}{2.110971in}}%
\pgfpathlineto{\pgfqpoint{1.896312in}{2.156231in}}%
\pgfpathlineto{\pgfqpoint{1.918364in}{2.168490in}}%
\pgfpathlineto{\pgfqpoint{1.872836in}{2.234076in}}%
\pgfpathlineto{\pgfqpoint{1.897953in}{2.183175in}}%
\pgfpathlineto{\pgfqpoint{1.963779in}{2.181949in}}%
\pgfpathlineto{\pgfqpoint{2.067095in}{2.158460in}}%
\pgfpathlineto{\pgfqpoint{2.047071in}{2.131611in}}%
\pgfpathlineto{\pgfqpoint{1.994304in}{2.158617in}}%
\pgfpathlineto{\pgfqpoint{2.010276in}{2.225981in}}%
\pgfpathlineto{\pgfqpoint{2.019660in}{2.201662in}}%
\pgfpathlineto{\pgfqpoint{1.990070in}{2.238874in}}%
\pgfpathlineto{\pgfqpoint{1.947585in}{2.226565in}}%
\pgfpathlineto{\pgfqpoint{1.946322in}{2.202571in}}%
\pgfpathlineto{\pgfqpoint{2.002056in}{2.221519in}}%
\pgfpathlineto{\pgfqpoint{2.026561in}{2.159881in}}%
\pgfpathlineto{\pgfqpoint{2.033632in}{2.058984in}}%
\pgfpathlineto{\pgfqpoint{2.069819in}{2.071982in}}%
\pgfpathlineto{\pgfqpoint{2.006289in}{2.056936in}}%
\pgfpathlineto{\pgfqpoint{2.043443in}{2.065364in}}%
\pgfpathlineto{\pgfqpoint{2.065923in}{2.063449in}}%
\pgfpathlineto{\pgfqpoint{2.095994in}{2.046203in}}%
\pgfpathlineto{\pgfqpoint{2.053101in}{2.082716in}}%
\pgfpathlineto{\pgfqpoint{2.094497in}{2.083390in}}%
\pgfpathlineto{\pgfqpoint{2.129733in}{2.130254in}}%
\pgfpathlineto{\pgfqpoint{2.203579in}{2.163766in}}%
\pgfpathlineto{\pgfqpoint{2.279059in}{2.150961in}}%
\pgfpathlineto{\pgfqpoint{2.286209in}{2.174777in}}%
\pgfpathlineto{\pgfqpoint{2.278278in}{2.192063in}}%
\pgfpathlineto{\pgfqpoint{2.266039in}{2.146091in}}%
\pgfpathlineto{\pgfqpoint{2.277793in}{2.096440in}}%
\pgfpathlineto{\pgfqpoint{2.324156in}{2.081421in}}%
\pgfpathlineto{\pgfqpoint{2.306185in}{2.155492in}}%
\pgfpathlineto{\pgfqpoint{2.236192in}{2.189455in}}%
\pgfpathlineto{\pgfqpoint{2.238011in}{2.219543in}}%
\pgfpathlineto{\pgfqpoint{2.285197in}{2.221666in}}%
\pgfpathlineto{\pgfqpoint{2.256508in}{2.214319in}}%
\pgfpathlineto{\pgfqpoint{2.276035in}{2.189775in}}%
\pgfpathlineto{\pgfqpoint{2.324049in}{2.199714in}}%
\pgfpathlineto{\pgfqpoint{2.410574in}{2.190874in}}%
\pgfpathlineto{\pgfqpoint{2.504484in}{2.247884in}}%
\pgfpathlineto{\pgfqpoint{2.445908in}{2.307274in}}%
\pgfpathlineto{\pgfqpoint{2.417298in}{2.330129in}}%
\pgfpathlineto{\pgfqpoint{2.364465in}{2.334584in}}%
\pgfpathlineto{\pgfqpoint{2.303409in}{2.330940in}}%
\pgfpathlineto{\pgfqpoint{2.279315in}{2.273826in}}%
\pgfpathlineto{\pgfqpoint{2.245519in}{2.337768in}}%
\pgfpathlineto{\pgfqpoint{2.304408in}{2.338430in}}%
\pgfpathlineto{\pgfqpoint{2.377560in}{2.352158in}}%
\pgfpathlineto{\pgfqpoint{2.361396in}{2.356954in}}%
\pgfpathlineto{\pgfqpoint{2.412537in}{2.299018in}}%
\pgfpathlineto{\pgfqpoint{2.412734in}{2.242297in}}%
\pgfpathlineto{\pgfqpoint{2.567983in}{2.143322in}}%
\pgfpathlineto{\pgfqpoint{2.576555in}{2.195489in}}%
\pgfpathlineto{\pgfqpoint{2.577926in}{2.239069in}}%
\pgfpathlineto{\pgfqpoint{2.622574in}{2.156014in}}%
\pgfpathlineto{\pgfqpoint{2.522278in}{2.154106in}}%
\pgfpathlineto{\pgfqpoint{2.538228in}{2.163871in}}%
\pgfpathlineto{\pgfqpoint{2.517607in}{2.198377in}}%
\pgfpathlineto{\pgfqpoint{2.578113in}{2.251276in}}%
\pgfpathlineto{\pgfqpoint{2.555615in}{2.244909in}}%
\pgfpathlineto{\pgfqpoint{2.527378in}{2.254258in}}%
\pgfpathlineto{\pgfqpoint{2.585220in}{2.258449in}}%
\pgfpathlineto{\pgfqpoint{2.578122in}{2.265941in}}%
\pgfpathlineto{\pgfqpoint{2.585807in}{2.246752in}}%
\pgfpathlineto{\pgfqpoint{2.611415in}{2.266550in}}%
\pgfpathlineto{\pgfqpoint{2.614372in}{2.253560in}}%
\pgfpathlineto{\pgfqpoint{2.556847in}{2.291866in}}%
\pgfpathlineto{\pgfqpoint{2.579139in}{2.317642in}}%
\pgfpathlineto{\pgfqpoint{2.530421in}{2.348363in}}%
\pgfpathlineto{\pgfqpoint{2.532973in}{2.389083in}}%
\pgfpathlineto{\pgfqpoint{2.501117in}{2.436786in}}%
\pgfpathlineto{\pgfqpoint{2.478524in}{2.476121in}}%
\pgfpathlineto{\pgfqpoint{2.515943in}{2.443771in}}%
\pgfpathlineto{\pgfqpoint{2.536149in}{2.402423in}}%
\pgfpathlineto{\pgfqpoint{2.522091in}{2.415298in}}%
\pgfpathlineto{\pgfqpoint{2.482677in}{2.466454in}}%
\pgfpathlineto{\pgfqpoint{2.425012in}{2.493095in}}%
\pgfpathlineto{\pgfqpoint{2.365234in}{2.540384in}}%
\pgfpathlineto{\pgfqpoint{2.386245in}{2.481790in}}%
\pgfpathlineto{\pgfqpoint{2.438154in}{2.477006in}}%
\pgfpathlineto{\pgfqpoint{2.522236in}{2.461978in}}%
\pgfpathlineto{\pgfqpoint{2.460211in}{2.450549in}}%
\pgfpathlineto{\pgfqpoint{2.480744in}{2.420597in}}%
\pgfpathlineto{\pgfqpoint{2.447685in}{2.404406in}}%
\pgfpathlineto{\pgfqpoint{2.405916in}{2.429148in}}%
\pgfpathlineto{\pgfqpoint{2.424663in}{2.407368in}}%
\pgfpathlineto{\pgfqpoint{2.428649in}{2.396970in}}%
\pgfpathlineto{\pgfqpoint{2.366975in}{2.419876in}}%
\pgfpathlineto{\pgfqpoint{2.376219in}{2.404907in}}%
\pgfpathlineto{\pgfqpoint{2.410223in}{2.355158in}}%
\pgfpathlineto{\pgfqpoint{2.444239in}{2.315019in}}%
\pgfpathlineto{\pgfqpoint{2.474626in}{2.322588in}}%
\pgfpathlineto{\pgfqpoint{2.506170in}{2.329910in}}%
\pgfpathlineto{\pgfqpoint{2.542475in}{2.374287in}}%
\pgfpathlineto{\pgfqpoint{2.573874in}{2.367488in}}%
\pgfpathlineto{\pgfqpoint{2.517376in}{2.347539in}}%
\pgfpathlineto{\pgfqpoint{2.518696in}{2.368810in}}%
\pgfpathlineto{\pgfqpoint{2.449617in}{2.377977in}}%
\pgfpathlineto{\pgfqpoint{2.458842in}{2.379913in}}%
\pgfpathlineto{\pgfqpoint{2.445224in}{2.400581in}}%
\pgfpathlineto{\pgfqpoint{2.409937in}{2.396691in}}%
\pgfpathlineto{\pgfqpoint{2.421389in}{2.388040in}}%
\pgfpathlineto{\pgfqpoint{2.408876in}{2.427679in}}%
\pgfpathlineto{\pgfqpoint{2.422361in}{2.382605in}}%
\pgfpathlineto{\pgfqpoint{2.508887in}{2.399458in}}%
\pgfpathlineto{\pgfqpoint{2.547056in}{2.399139in}}%
\pgfpathlineto{\pgfqpoint{2.496487in}{2.380725in}}%
\pgfpathlineto{\pgfqpoint{2.530585in}{2.390653in}}%
\pgfpathlineto{\pgfqpoint{2.553132in}{2.372134in}}%
\pgfpathlineto{\pgfqpoint{2.570971in}{2.438416in}}%
\pgfpathlineto{\pgfqpoint{2.598860in}{2.418177in}}%
\pgfpathlineto{\pgfqpoint{2.590090in}{2.422313in}}%
\pgfpathlineto{\pgfqpoint{2.613241in}{2.447474in}}%
\pgfpathlineto{\pgfqpoint{2.646117in}{2.417519in}}%
\pgfpathlineto{\pgfqpoint{2.631502in}{2.478563in}}%
\pgfpathlineto{\pgfqpoint{2.665184in}{2.465157in}}%
\pgfpathlineto{\pgfqpoint{2.581813in}{2.552583in}}%
\pgfpathlineto{\pgfqpoint{2.574427in}{2.448903in}}%
\pgfpathlineto{\pgfqpoint{2.543293in}{2.486058in}}%
\pgfpathlineto{\pgfqpoint{2.550875in}{2.467957in}}%
\pgfpathlineto{\pgfqpoint{2.594006in}{2.411637in}}%
\pgfpathlineto{\pgfqpoint{2.645337in}{2.433791in}}%
\pgfpathlineto{\pgfqpoint{2.640689in}{2.428915in}}%
\pgfpathlineto{\pgfqpoint{2.668941in}{2.475994in}}%
\pgfpathlineto{\pgfqpoint{2.680407in}{2.466845in}}%
\pgfpathlineto{\pgfqpoint{2.625187in}{2.507968in}}%
\pgfpathlineto{\pgfqpoint{2.594356in}{2.551083in}}%
\pgfpathlineto{\pgfqpoint{2.484809in}{2.533222in}}%
\pgfpathlineto{\pgfqpoint{2.469465in}{2.555714in}}%
\pgfpathlineto{\pgfqpoint{2.385457in}{2.550247in}}%
\pgfpathlineto{\pgfqpoint{2.354422in}{2.647257in}}%
\pgfpathlineto{\pgfqpoint{2.323496in}{2.640476in}}%
\pgfpathlineto{\pgfqpoint{2.298301in}{2.658234in}}%
\pgfpathlineto{\pgfqpoint{2.194212in}{2.707131in}}%
\pgfpathlineto{\pgfqpoint{2.174591in}{2.704062in}}%
\pgfpathlineto{\pgfqpoint{2.161666in}{2.674118in}}%
\pgfpathlineto{\pgfqpoint{2.177952in}{2.705217in}}%
\pgfpathlineto{\pgfqpoint{2.198860in}{2.723796in}}%
\pgfpathlineto{\pgfqpoint{2.222383in}{2.619161in}}%
\pgfpathlineto{\pgfqpoint{2.223693in}{2.630158in}}%
\pgfpathlineto{\pgfqpoint{2.207662in}{2.616535in}}%
\pgfpathlineto{\pgfqpoint{2.210487in}{2.643925in}}%
\pgfpathlineto{\pgfqpoint{2.307897in}{2.687875in}}%
\pgfpathlineto{\pgfqpoint{2.278293in}{2.634655in}}%
\pgfpathlineto{\pgfqpoint{2.276594in}{2.583147in}}%
\pgfpathlineto{\pgfqpoint{2.325169in}{2.595572in}}%
\pgfpathlineto{\pgfqpoint{2.274493in}{2.599172in}}%
\pgfpathlineto{\pgfqpoint{2.387958in}{2.669853in}}%
\pgfpathlineto{\pgfqpoint{2.336609in}{2.645146in}}%
\pgfpathlineto{\pgfqpoint{2.329066in}{2.657802in}}%
\pgfpathlineto{\pgfqpoint{2.347442in}{2.602388in}}%
\pgfpathlineto{\pgfqpoint{2.341455in}{2.628341in}}%
\pgfpathlineto{\pgfqpoint{2.297473in}{2.607086in}}%
\pgfpathlineto{\pgfqpoint{2.301801in}{2.656529in}}%
\pgfpathlineto{\pgfqpoint{2.299779in}{2.658630in}}%
\pgfpathlineto{\pgfqpoint{2.322737in}{2.678564in}}%
\pgfpathlineto{\pgfqpoint{2.296140in}{2.662164in}}%
\pgfpathlineto{\pgfqpoint{2.345357in}{2.650380in}}%
\pgfpathlineto{\pgfqpoint{2.353093in}{2.707786in}}%
\pgfpathlineto{\pgfqpoint{2.392617in}{2.738035in}}%
\pgfpathlineto{\pgfqpoint{2.369573in}{2.737126in}}%
\pgfpathlineto{\pgfqpoint{2.354558in}{2.708147in}}%
\pgfpathlineto{\pgfqpoint{2.326354in}{2.701675in}}%
\pgfpathlineto{\pgfqpoint{2.281282in}{2.726496in}}%
\pgfpathlineto{\pgfqpoint{2.216754in}{2.773778in}}%
\pgfpathlineto{\pgfqpoint{2.164156in}{2.790771in}}%
\pgfpathlineto{\pgfqpoint{2.216016in}{2.700379in}}%
\pgfpathlineto{\pgfqpoint{2.225124in}{2.612507in}}%
\pgfpathlineto{\pgfqpoint{2.244696in}{2.598384in}}%
\pgfpathlineto{\pgfqpoint{2.301904in}{2.564112in}}%
\pgfpathlineto{\pgfqpoint{2.279297in}{2.560573in}}%
\pgfpathlineto{\pgfqpoint{2.228383in}{2.630687in}}%
\pgfpathlineto{\pgfqpoint{2.194686in}{2.586713in}}%
\pgfpathlineto{\pgfqpoint{2.100253in}{2.640956in}}%
\pgfpathlineto{\pgfqpoint{2.122579in}{2.613940in}}%
\pgfpathlineto{\pgfqpoint{2.122311in}{2.590073in}}%
\pgfpathlineto{\pgfqpoint{2.135641in}{2.554040in}}%
\pgfpathlineto{\pgfqpoint{2.194095in}{2.579367in}}%
\pgfpathlineto{\pgfqpoint{2.121913in}{2.581686in}}%
\pgfpathlineto{\pgfqpoint{2.075314in}{2.618421in}}%
\pgfpathlineto{\pgfqpoint{2.047986in}{2.615789in}}%
\pgfpathlineto{\pgfqpoint{2.052234in}{2.600642in}}%
\pgfpathlineto{\pgfqpoint{2.019626in}{2.643937in}}%
\pgfpathlineto{\pgfqpoint{1.992024in}{2.670612in}}%
\pgfpathlineto{\pgfqpoint{2.014595in}{2.656429in}}%
\pgfpathlineto{\pgfqpoint{2.050573in}{2.681263in}}%
\pgfpathlineto{\pgfqpoint{2.079298in}{2.691920in}}%
\pgfpathlineto{\pgfqpoint{2.035299in}{2.675067in}}%
\pgfpathlineto{\pgfqpoint{2.047852in}{2.619665in}}%
\pgfpathlineto{\pgfqpoint{2.096605in}{2.552811in}}%
\pgfpathlineto{\pgfqpoint{2.104701in}{2.551852in}}%
\pgfpathlineto{\pgfqpoint{2.063349in}{2.562076in}}%
\pgfpathlineto{\pgfqpoint{2.046896in}{2.561383in}}%
\pgfpathlineto{\pgfqpoint{2.017456in}{2.604277in}}%
\pgfpathlineto{\pgfqpoint{2.013656in}{2.576060in}}%
\pgfpathlineto{\pgfqpoint{2.052720in}{2.507426in}}%
\pgfpathlineto{\pgfqpoint{2.083918in}{2.471827in}}%
\pgfpathlineto{\pgfqpoint{2.060544in}{2.454508in}}%
\pgfpathlineto{\pgfqpoint{2.090559in}{2.487399in}}%
\pgfpathlineto{\pgfqpoint{2.071147in}{2.434473in}}%
\pgfpathlineto{\pgfqpoint{2.093427in}{2.412255in}}%
\pgfpathlineto{\pgfqpoint{2.054166in}{2.440771in}}%
\pgfpathlineto{\pgfqpoint{2.045770in}{2.380211in}}%
\pgfpathlineto{\pgfqpoint{2.039015in}{2.348085in}}%
\pgfpathlineto{\pgfqpoint{2.037226in}{2.333316in}}%
\pgfpathlineto{\pgfqpoint{1.985408in}{2.405871in}}%
\pgfpathlineto{\pgfqpoint{1.946474in}{2.379004in}}%
\pgfpathlineto{\pgfqpoint{2.024441in}{2.368984in}}%
\pgfpathlineto{\pgfqpoint{2.022471in}{2.358379in}}%
\pgfpathlineto{\pgfqpoint{2.023842in}{2.330217in}}%
\pgfpathlineto{\pgfqpoint{1.988289in}{2.338288in}}%
\pgfpathlineto{\pgfqpoint{2.008231in}{2.303803in}}%
\pgfpathlineto{\pgfqpoint{2.007520in}{2.324732in}}%
\pgfpathlineto{\pgfqpoint{2.047318in}{2.334106in}}%
\pgfpathlineto{\pgfqpoint{2.005744in}{2.337555in}}%
\pgfpathlineto{\pgfqpoint{2.016380in}{2.321021in}}%
\pgfpathlineto{\pgfqpoint{1.969153in}{2.356272in}}%
\pgfpathlineto{\pgfqpoint{1.968374in}{2.385648in}}%
\pgfpathlineto{\pgfqpoint{1.977281in}{2.382100in}}%
\pgfpathlineto{\pgfqpoint{1.943755in}{2.366062in}}%
\pgfpathlineto{\pgfqpoint{1.924631in}{2.370179in}}%
\pgfpathlineto{\pgfqpoint{1.872521in}{2.334919in}}%
\pgfpathlineto{\pgfqpoint{1.960547in}{2.312625in}}%
\pgfpathlineto{\pgfqpoint{2.012739in}{2.309513in}}%
\pgfpathlineto{\pgfqpoint{2.088005in}{2.345697in}}%
\pgfpathlineto{\pgfqpoint{2.083281in}{2.387806in}}%
\pgfpathlineto{\pgfqpoint{2.058128in}{2.452167in}}%
\pgfpathlineto{\pgfqpoint{1.983086in}{2.449665in}}%
\pgfpathlineto{\pgfqpoint{1.949356in}{2.405544in}}%
\pgfpathlineto{\pgfqpoint{1.902707in}{2.437202in}}%
\pgfpathlineto{\pgfqpoint{1.802390in}{2.389064in}}%
\pgfpathlineto{\pgfqpoint{1.786186in}{2.340474in}}%
\pgfpathlineto{\pgfqpoint{1.781873in}{2.344577in}}%
\pgfpathlineto{\pgfqpoint{1.773828in}{2.334934in}}%
\pgfpathlineto{\pgfqpoint{1.843742in}{2.306689in}}%
\pgfpathlineto{\pgfqpoint{1.828096in}{2.259061in}}%
\pgfpathlineto{\pgfqpoint{1.830368in}{2.272695in}}%
\pgfpathlineto{\pgfqpoint{1.796507in}{2.306004in}}%
\pgfpathlineto{\pgfqpoint{1.818619in}{2.285170in}}%
\pgfpathlineto{\pgfqpoint{1.797271in}{2.333949in}}%
\pgfpathlineto{\pgfqpoint{1.823053in}{2.316229in}}%
\pgfpathlineto{\pgfqpoint{1.790541in}{2.366037in}}%
\pgfpathlineto{\pgfqpoint{1.795290in}{2.350090in}}%
\pgfpathlineto{\pgfqpoint{1.760159in}{2.299544in}}%
\pgfpathlineto{\pgfqpoint{1.745394in}{2.287100in}}%
\pgfpathlineto{\pgfqpoint{1.775641in}{2.246941in}}%
\pgfpathlineto{\pgfqpoint{1.705594in}{2.252757in}}%
\pgfpathlineto{\pgfqpoint{1.634734in}{2.173543in}}%
\pgfpathlineto{\pgfqpoint{1.608919in}{2.152347in}}%
\pgfpathlineto{\pgfqpoint{1.657100in}{2.162977in}}%
\pgfpathlineto{\pgfqpoint{1.614781in}{2.194004in}}%
\pgfpathlineto{\pgfqpoint{1.576047in}{2.187871in}}%
\pgfpathlineto{\pgfqpoint{1.565924in}{2.263244in}}%
\pgfpathlineto{\pgfqpoint{1.538417in}{2.201641in}}%
\pgfpathlineto{\pgfqpoint{1.569173in}{2.184716in}}%
\pgfpathlineto{\pgfqpoint{1.529318in}{2.194310in}}%
\pgfpathlineto{\pgfqpoint{1.518095in}{2.163648in}}%
\pgfpathlineto{\pgfqpoint{1.511526in}{2.145122in}}%
\pgfpathlineto{\pgfqpoint{1.587386in}{2.131237in}}%
\pgfpathlineto{\pgfqpoint{1.653937in}{2.153177in}}%
\pgfpathlineto{\pgfqpoint{1.660585in}{2.158471in}}%
\pgfpathlineto{\pgfqpoint{1.701309in}{2.139244in}}%
\pgfpathlineto{\pgfqpoint{1.758139in}{2.128149in}}%
\pgfpathlineto{\pgfqpoint{1.733462in}{2.118976in}}%
\pgfpathlineto{\pgfqpoint{1.750835in}{2.140063in}}%
\pgfpathlineto{\pgfqpoint{1.730719in}{2.130437in}}%
\pgfpathlineto{\pgfqpoint{1.697809in}{2.113467in}}%
\pgfpathlineto{\pgfqpoint{1.760721in}{2.081728in}}%
\pgfpathlineto{\pgfqpoint{1.774842in}{2.116334in}}%
\pgfpathlineto{\pgfqpoint{1.732488in}{2.176634in}}%
\pgfpathlineto{\pgfqpoint{1.661737in}{2.160779in}}%
\pgfpathlineto{\pgfqpoint{1.610921in}{2.126084in}}%
\pgfpathlineto{\pgfqpoint{1.611931in}{2.191624in}}%
\pgfpathlineto{\pgfqpoint{1.555802in}{2.226542in}}%
\pgfpathlineto{\pgfqpoint{1.535206in}{2.263651in}}%
\pgfpathlineto{\pgfqpoint{1.592037in}{2.244778in}}%
\pgfpathlineto{\pgfqpoint{1.583636in}{2.205705in}}%
\pgfpathlineto{\pgfqpoint{1.535591in}{2.150818in}}%
\pgfpathlineto{\pgfqpoint{1.462722in}{2.151715in}}%
\pgfpathlineto{\pgfqpoint{1.508549in}{2.124902in}}%
\pgfpathlineto{\pgfqpoint{1.442128in}{2.140025in}}%
\pgfpathlineto{\pgfqpoint{1.560051in}{2.198945in}}%
\pgfpathlineto{\pgfqpoint{1.525466in}{2.167195in}}%
\pgfpathlineto{\pgfqpoint{1.533938in}{2.140794in}}%
\pgfpathlineto{\pgfqpoint{1.490719in}{2.130953in}}%
\pgfpathlineto{\pgfqpoint{1.462905in}{2.118320in}}%
\pgfpathlineto{\pgfqpoint{1.499281in}{2.082231in}}%
\pgfpathlineto{\pgfqpoint{1.464724in}{2.138396in}}%
\pgfpathlineto{\pgfqpoint{1.429155in}{2.208491in}}%
\pgfpathlineto{\pgfqpoint{1.429870in}{2.188228in}}%
\pgfpathlineto{\pgfqpoint{1.455166in}{2.133904in}}%
\pgfpathlineto{\pgfqpoint{1.499519in}{2.151775in}}%
\pgfpathlineto{\pgfqpoint{1.508844in}{2.172701in}}%
\pgfpathlineto{\pgfqpoint{1.507187in}{2.176356in}}%
\pgfpathlineto{\pgfqpoint{1.490532in}{2.219810in}}%
\pgfpathlineto{\pgfqpoint{1.500005in}{2.191171in}}%
\pgfpathlineto{\pgfqpoint{1.582578in}{2.206981in}}%
\pgfpathlineto{\pgfqpoint{1.555280in}{2.247966in}}%
\pgfpathlineto{\pgfqpoint{1.554119in}{2.230815in}}%
\pgfpathlineto{\pgfqpoint{1.585841in}{2.180415in}}%
\pgfpathlineto{\pgfqpoint{1.523314in}{2.222924in}}%
\pgfpathlineto{\pgfqpoint{1.454595in}{2.220988in}}%
\pgfpathlineto{\pgfqpoint{1.395862in}{2.193446in}}%
\pgfpathlineto{\pgfqpoint{1.411367in}{2.139317in}}%
\pgfpathlineto{\pgfqpoint{1.550076in}{2.132882in}}%
\pgfpathlineto{\pgfqpoint{1.562962in}{2.096413in}}%
\pgfpathlineto{\pgfqpoint{1.516099in}{2.129155in}}%
\pgfpathlineto{\pgfqpoint{1.571714in}{2.119965in}}%
\pgfpathlineto{\pgfqpoint{1.536069in}{2.120199in}}%
\pgfpathlineto{\pgfqpoint{1.516413in}{2.093560in}}%
\pgfpathlineto{\pgfqpoint{1.467992in}{2.161418in}}%
\pgfpathlineto{\pgfqpoint{1.480487in}{2.134995in}}%
\pgfpathlineto{\pgfqpoint{1.499564in}{2.077095in}}%
\pgfpathlineto{\pgfqpoint{1.540248in}{2.022464in}}%
\pgfpathlineto{\pgfqpoint{1.624962in}{1.994775in}}%
\pgfpathlineto{\pgfqpoint{1.647708in}{1.994129in}}%
\pgfpathlineto{\pgfqpoint{1.648354in}{2.026305in}}%
\pgfpathlineto{\pgfqpoint{1.749411in}{1.983606in}}%
\pgfpathlineto{\pgfqpoint{1.847726in}{1.998013in}}%
\pgfpathlineto{\pgfqpoint{1.877941in}{1.974529in}}%
\pgfpathlineto{\pgfqpoint{1.900639in}{1.981506in}}%
\pgfpathlineto{\pgfqpoint{1.871138in}{1.970984in}}%
\pgfpathlineto{\pgfqpoint{1.847175in}{1.980176in}}%
\pgfpathlineto{\pgfqpoint{1.835576in}{1.955284in}}%
\pgfpathlineto{\pgfqpoint{1.870076in}{1.936881in}}%
\pgfpathlineto{\pgfqpoint{1.872712in}{1.912055in}}%
\pgfpathlineto{\pgfqpoint{1.865768in}{1.918864in}}%
\pgfpathlineto{\pgfqpoint{1.922062in}{1.847011in}}%
\pgfpathlineto{\pgfqpoint{1.839692in}{1.882574in}}%
\pgfpathlineto{\pgfqpoint{1.766881in}{1.849876in}}%
\pgfpathlineto{\pgfqpoint{1.766698in}{1.834790in}}%
\pgfpathlineto{\pgfqpoint{1.830421in}{1.838274in}}%
\pgfpathlineto{\pgfqpoint{1.876947in}{1.893320in}}%
\pgfpathlineto{\pgfqpoint{1.792462in}{1.917467in}}%
\pgfpathlineto{\pgfqpoint{1.817309in}{1.882163in}}%
\pgfpathlineto{\pgfqpoint{1.787367in}{1.846434in}}%
\pgfpathlineto{\pgfqpoint{1.755584in}{1.828411in}}%
\pgfpathlineto{\pgfqpoint{1.790928in}{1.796245in}}%
\pgfpathlineto{\pgfqpoint{1.820856in}{1.792990in}}%
\pgfpathlineto{\pgfqpoint{1.841766in}{1.772935in}}%
\pgfpathlineto{\pgfqpoint{1.907145in}{1.763427in}}%
\pgfpathlineto{\pgfqpoint{1.874496in}{1.758221in}}%
\pgfpathlineto{\pgfqpoint{1.841869in}{1.773492in}}%
\pgfpathlineto{\pgfqpoint{1.895437in}{1.775527in}}%
\pgfpathlineto{\pgfqpoint{1.904282in}{1.787074in}}%
\pgfpathlineto{\pgfqpoint{1.893424in}{1.838924in}}%
\pgfpathlineto{\pgfqpoint{1.937407in}{1.869749in}}%
\pgfpathlineto{\pgfqpoint{1.941524in}{1.908149in}}%
\pgfpathlineto{\pgfqpoint{1.947579in}{1.919826in}}%
\pgfpathlineto{\pgfqpoint{1.863761in}{1.892206in}}%
\pgfpathlineto{\pgfqpoint{1.790478in}{1.914892in}}%
\pgfpathlineto{\pgfqpoint{1.771683in}{1.899928in}}%
\pgfpathlineto{\pgfqpoint{1.842515in}{1.959182in}}%
\pgfpathlineto{\pgfqpoint{1.924698in}{1.956144in}}%
\pgfpathlineto{\pgfqpoint{1.926065in}{1.928866in}}%
\pgfpathlineto{\pgfqpoint{1.960716in}{1.896986in}}%
\pgfpathlineto{\pgfqpoint{1.942680in}{1.878362in}}%
\pgfpathlineto{\pgfqpoint{1.965116in}{1.922782in}}%
\pgfpathlineto{\pgfqpoint{1.922484in}{1.895135in}}%
\pgfpathlineto{\pgfqpoint{1.900862in}{1.912027in}}%
\pgfpathlineto{\pgfqpoint{1.965486in}{1.875524in}}%
\pgfpathlineto{\pgfqpoint{1.941708in}{1.949536in}}%
\pgfpathlineto{\pgfqpoint{1.980136in}{1.937107in}}%
\pgfpathlineto{\pgfqpoint{1.974007in}{1.975181in}}%
\pgfpathlineto{\pgfqpoint{1.943929in}{1.979960in}}%
\pgfpathlineto{\pgfqpoint{1.886917in}{2.040423in}}%
\pgfpathlineto{\pgfqpoint{1.893888in}{1.984393in}}%
\pgfpathlineto{\pgfqpoint{1.905374in}{1.988303in}}%
\pgfpathlineto{\pgfqpoint{1.879631in}{2.022640in}}%
\pgfpathlineto{\pgfqpoint{1.945945in}{2.051857in}}%
\pgfpathlineto{\pgfqpoint{1.905002in}{2.077812in}}%
\pgfpathlineto{\pgfqpoint{1.882231in}{2.072861in}}%
\pgfpathlineto{\pgfqpoint{1.947528in}{2.068988in}}%
\pgfpathlineto{\pgfqpoint{1.928731in}{2.025302in}}%
\pgfpathlineto{\pgfqpoint{1.912447in}{2.041892in}}%
\pgfpathlineto{\pgfqpoint{1.902289in}{2.032542in}}%
\pgfpathlineto{\pgfqpoint{1.860970in}{2.058943in}}%
\pgfpathlineto{\pgfqpoint{1.889554in}{2.060087in}}%
\pgfpathlineto{\pgfqpoint{1.917463in}{2.068784in}}%
\pgfpathlineto{\pgfqpoint{1.902322in}{2.058853in}}%
\pgfpathlineto{\pgfqpoint{1.910664in}{2.016712in}}%
\pgfpathlineto{\pgfqpoint{1.966866in}{1.973465in}}%
\pgfpathlineto{\pgfqpoint{1.970445in}{2.021551in}}%
\pgfpathlineto{\pgfqpoint{1.922235in}{1.976815in}}%
\pgfpathlineto{\pgfqpoint{1.943384in}{1.934014in}}%
\pgfpathlineto{\pgfqpoint{2.026939in}{1.936733in}}%
\pgfpathlineto{\pgfqpoint{2.013172in}{1.852775in}}%
\pgfpathlineto{\pgfqpoint{2.010660in}{1.902743in}}%
\pgfpathlineto{\pgfqpoint{2.076393in}{1.901762in}}%
\pgfpathlineto{\pgfqpoint{2.081630in}{1.906126in}}%
\pgfpathlineto{\pgfqpoint{2.108448in}{1.904776in}}%
\pgfpathlineto{\pgfqpoint{2.195331in}{1.905587in}}%
\pgfpathlineto{\pgfqpoint{2.269842in}{1.960520in}}%
\pgfpathlineto{\pgfqpoint{2.246242in}{1.896703in}}%
\pgfpathlineto{\pgfqpoint{2.257817in}{1.935010in}}%
\pgfpathlineto{\pgfqpoint{2.297425in}{1.901368in}}%
\pgfpathlineto{\pgfqpoint{2.307950in}{1.874639in}}%
\pgfpathlineto{\pgfqpoint{2.273140in}{1.973611in}}%
\pgfpathlineto{\pgfqpoint{2.331164in}{1.963551in}}%
\pgfpathlineto{\pgfqpoint{2.327560in}{1.969423in}}%
\pgfpathlineto{\pgfqpoint{2.263730in}{1.965102in}}%
\pgfpathlineto{\pgfqpoint{2.285213in}{2.040130in}}%
\pgfpathlineto{\pgfqpoint{2.346806in}{2.020486in}}%
\pgfpathlineto{\pgfqpoint{2.306055in}{1.977085in}}%
\pgfpathlineto{\pgfqpoint{2.286615in}{1.926006in}}%
\pgfpathlineto{\pgfqpoint{2.248222in}{1.932416in}}%
\pgfpathlineto{\pgfqpoint{2.290654in}{2.038593in}}%
\pgfpathlineto{\pgfqpoint{2.282853in}{2.042469in}}%
\pgfpathlineto{\pgfqpoint{2.302369in}{2.039807in}}%
\pgfpathlineto{\pgfqpoint{2.225134in}{2.104573in}}%
\pgfpathlineto{\pgfqpoint{2.249127in}{2.050695in}}%
\pgfpathlineto{\pgfqpoint{2.205831in}{1.994895in}}%
\pgfpathlineto{\pgfqpoint{2.254278in}{2.000079in}}%
\pgfpathlineto{\pgfqpoint{2.260165in}{2.073109in}}%
\pgfpathlineto{\pgfqpoint{2.228205in}{2.056080in}}%
\pgfpathlineto{\pgfqpoint{2.316209in}{2.086858in}}%
\pgfpathlineto{\pgfqpoint{2.326473in}{2.076338in}}%
\pgfpathlineto{\pgfqpoint{2.263461in}{2.159738in}}%
\pgfpathlineto{\pgfqpoint{2.330570in}{2.172182in}}%
\pgfpathlineto{\pgfqpoint{2.338212in}{2.219332in}}%
\pgfpathlineto{\pgfqpoint{2.353360in}{2.214900in}}%
\pgfpathlineto{\pgfqpoint{2.321264in}{2.251652in}}%
\pgfpathlineto{\pgfqpoint{2.313752in}{2.309081in}}%
\pgfpathlineto{\pgfqpoint{2.278813in}{2.275456in}}%
\pgfpathlineto{\pgfqpoint{2.274646in}{2.287633in}}%
\pgfpathlineto{\pgfqpoint{2.239659in}{2.320398in}}%
\pgfpathlineto{\pgfqpoint{2.170056in}{2.298076in}}%
\pgfpathlineto{\pgfqpoint{2.180973in}{2.380954in}}%
\pgfpathlineto{\pgfqpoint{2.164255in}{2.335772in}}%
\pgfpathlineto{\pgfqpoint{2.088896in}{2.345266in}}%
\pgfpathlineto{\pgfqpoint{1.975649in}{2.305182in}}%
\pgfpathlineto{\pgfqpoint{1.920650in}{2.260644in}}%
\pgfpathlineto{\pgfqpoint{1.964268in}{2.262595in}}%
\pgfpathlineto{\pgfqpoint{1.935742in}{2.313529in}}%
\pgfpathlineto{\pgfqpoint{1.913338in}{2.329504in}}%
\pgfpathlineto{\pgfqpoint{1.931390in}{2.319109in}}%
\pgfpathlineto{\pgfqpoint{1.969575in}{2.292255in}}%
\pgfpathlineto{\pgfqpoint{1.948544in}{2.275312in}}%
\pgfpathlineto{\pgfqpoint{1.920398in}{2.304616in}}%
\pgfpathlineto{\pgfqpoint{1.920695in}{2.340109in}}%
\pgfpathlineto{\pgfqpoint{1.889398in}{2.319732in}}%
\pgfpathlineto{\pgfqpoint{1.907891in}{2.323190in}}%
\pgfpathlineto{\pgfqpoint{1.906577in}{2.287904in}}%
\pgfpathlineto{\pgfqpoint{1.883725in}{2.287297in}}%
\pgfpathlineto{\pgfqpoint{1.862447in}{2.293631in}}%
\pgfpathlineto{\pgfqpoint{1.860362in}{2.256539in}}%
\pgfpathlineto{\pgfqpoint{1.904908in}{2.251971in}}%
\pgfpathlineto{\pgfqpoint{1.906456in}{2.192988in}}%
\pgfpathlineto{\pgfqpoint{1.825003in}{2.240596in}}%
\pgfpathlineto{\pgfqpoint{1.741110in}{2.221000in}}%
\pgfpathlineto{\pgfqpoint{1.710276in}{2.182280in}}%
\pgfpathlineto{\pgfqpoint{1.752567in}{2.129916in}}%
\pgfpathlineto{\pgfqpoint{1.740237in}{2.108167in}}%
\pgfpathlineto{\pgfqpoint{1.648249in}{2.104019in}}%
\pgfpathlineto{\pgfqpoint{1.713691in}{2.164100in}}%
\pgfpathlineto{\pgfqpoint{1.749669in}{2.168341in}}%
\pgfpathlineto{\pgfqpoint{1.689217in}{2.172099in}}%
\pgfpathlineto{\pgfqpoint{1.701030in}{2.133406in}}%
\pgfpathlineto{\pgfqpoint{1.711805in}{2.092711in}}%
\pgfpathlineto{\pgfqpoint{1.626840in}{2.187418in}}%
\pgfpathlineto{\pgfqpoint{1.686042in}{2.177818in}}%
\pgfpathlineto{\pgfqpoint{1.707872in}{2.244253in}}%
\pgfpathlineto{\pgfqpoint{1.774664in}{2.236397in}}%
\pgfpathlineto{\pgfqpoint{1.825881in}{2.214783in}}%
\pgfpathlineto{\pgfqpoint{1.937558in}{2.163928in}}%
\pgfpathlineto{\pgfqpoint{1.939222in}{2.116748in}}%
\pgfpathlineto{\pgfqpoint{1.912507in}{2.095998in}}%
\pgfpathlineto{\pgfqpoint{1.920266in}{2.054698in}}%
\pgfpathlineto{\pgfqpoint{1.863802in}{2.050165in}}%
\pgfpathlineto{\pgfqpoint{1.874460in}{2.081605in}}%
\pgfpathlineto{\pgfqpoint{1.856123in}{2.069409in}}%
\pgfpathlineto{\pgfqpoint{1.833764in}{2.027505in}}%
\pgfpathlineto{\pgfqpoint{1.801575in}{1.997629in}}%
\pgfpathlineto{\pgfqpoint{1.868050in}{1.959408in}}%
\pgfpathlineto{\pgfqpoint{1.864220in}{1.972497in}}%
\pgfpathlineto{\pgfqpoint{1.906423in}{1.964717in}}%
\pgfpathlineto{\pgfqpoint{1.927716in}{1.941840in}}%
\pgfpathlineto{\pgfqpoint{1.951831in}{1.965266in}}%
\pgfpathlineto{\pgfqpoint{1.954329in}{1.931805in}}%
\pgfpathlineto{\pgfqpoint{1.898278in}{1.905880in}}%
\pgfpathlineto{\pgfqpoint{1.909822in}{1.933772in}}%
\pgfpathlineto{\pgfqpoint{1.869354in}{1.985526in}}%
\pgfpathlineto{\pgfqpoint{1.816000in}{1.989359in}}%
\pgfpathlineto{\pgfqpoint{1.824075in}{1.960432in}}%
\pgfpathlineto{\pgfqpoint{1.763406in}{1.931180in}}%
\pgfpathlineto{\pgfqpoint{1.761285in}{1.886802in}}%
\pgfpathlineto{\pgfqpoint{1.770013in}{1.914630in}}%
\pgfpathlineto{\pgfqpoint{1.792949in}{1.947926in}}%
\pgfpathlineto{\pgfqpoint{1.783186in}{1.908597in}}%
\pgfpathlineto{\pgfqpoint{1.801340in}{1.810503in}}%
\pgfpathlineto{\pgfqpoint{1.847626in}{1.829208in}}%
\pgfpathlineto{\pgfqpoint{1.788416in}{1.799875in}}%
\pgfpathlineto{\pgfqpoint{1.842535in}{1.834185in}}%
\pgfpathlineto{\pgfqpoint{1.903276in}{1.885943in}}%
\pgfpathlineto{\pgfqpoint{1.894355in}{1.795931in}}%
\pgfpathlineto{\pgfqpoint{1.979940in}{1.760532in}}%
\pgfpathlineto{\pgfqpoint{2.019120in}{1.754406in}}%
\pgfpathlineto{\pgfqpoint{2.014257in}{1.751582in}}%
\pgfpathlineto{\pgfqpoint{2.040853in}{1.714761in}}%
\pgfpathlineto{\pgfqpoint{2.013402in}{1.760166in}}%
\pgfpathlineto{\pgfqpoint{2.037370in}{1.716746in}}%
\pgfpathlineto{\pgfqpoint{2.015819in}{1.722824in}}%
\pgfpathlineto{\pgfqpoint{2.067050in}{1.789866in}}%
\pgfpathlineto{\pgfqpoint{2.083009in}{1.779554in}}%
\pgfpathlineto{\pgfqpoint{2.072407in}{1.780044in}}%
\pgfpathlineto{\pgfqpoint{2.027861in}{1.836096in}}%
\pgfpathlineto{\pgfqpoint{2.022205in}{1.840876in}}%
\pgfpathlineto{\pgfqpoint{1.987309in}{1.882258in}}%
\pgfpathlineto{\pgfqpoint{1.996403in}{1.891426in}}%
\pgfpathlineto{\pgfqpoint{2.014990in}{1.888931in}}%
\pgfpathlineto{\pgfqpoint{2.024122in}{1.910933in}}%
\pgfpathlineto{\pgfqpoint{2.001497in}{1.927467in}}%
\pgfpathlineto{\pgfqpoint{1.944599in}{1.911947in}}%
\pgfpathlineto{\pgfqpoint{1.895548in}{1.936744in}}%
\pgfpathlineto{\pgfqpoint{1.916382in}{1.945272in}}%
\pgfpathlineto{\pgfqpoint{1.960686in}{1.991776in}}%
\pgfpathlineto{\pgfqpoint{1.972094in}{2.047506in}}%
\pgfpathlineto{\pgfqpoint{1.929381in}{2.086649in}}%
\pgfpathlineto{\pgfqpoint{1.867387in}{2.111807in}}%
\pgfpathlineto{\pgfqpoint{1.781880in}{2.125381in}}%
\pgfpathlineto{\pgfqpoint{1.761117in}{2.101637in}}%
\pgfpathlineto{\pgfqpoint{1.789416in}{2.216097in}}%
\pgfpathlineto{\pgfqpoint{1.694349in}{2.242377in}}%
\pgfpathlineto{\pgfqpoint{1.732441in}{2.278608in}}%
\pgfpathlineto{\pgfqpoint{1.735242in}{2.296037in}}%
\pgfpathlineto{\pgfqpoint{1.798029in}{2.260911in}}%
\pgfpathlineto{\pgfqpoint{1.872467in}{2.269113in}}%
\pgfpathlineto{\pgfqpoint{1.913597in}{2.295857in}}%
\pgfpathlineto{\pgfqpoint{1.887957in}{2.290558in}}%
\pgfpathlineto{\pgfqpoint{1.936814in}{2.285910in}}%
\pgfpathlineto{\pgfqpoint{1.989874in}{2.286477in}}%
\pgfpathlineto{\pgfqpoint{1.994921in}{2.288620in}}%
\pgfpathlineto{\pgfqpoint{2.067277in}{2.302105in}}%
\pgfpathlineto{\pgfqpoint{2.105215in}{2.246163in}}%
\pgfpathlineto{\pgfqpoint{2.113981in}{2.221695in}}%
\pgfpathlineto{\pgfqpoint{2.037572in}{2.213870in}}%
\pgfpathlineto{\pgfqpoint{2.061393in}{2.270110in}}%
\pgfpathlineto{\pgfqpoint{2.014686in}{2.264680in}}%
\pgfpathlineto{\pgfqpoint{2.071823in}{2.297257in}}%
\pgfpathlineto{\pgfqpoint{2.067244in}{2.323757in}}%
\pgfpathlineto{\pgfqpoint{1.967705in}{2.346346in}}%
\pgfpathlineto{\pgfqpoint{1.966554in}{2.386310in}}%
\pgfpathlineto{\pgfqpoint{1.989838in}{2.283220in}}%
\pgfpathlineto{\pgfqpoint{2.048953in}{2.302999in}}%
\pgfpathlineto{\pgfqpoint{2.021139in}{2.298980in}}%
\pgfpathlineto{\pgfqpoint{2.052187in}{2.303343in}}%
\pgfpathlineto{\pgfqpoint{2.063637in}{2.247220in}}%
\pgfpathlineto{\pgfqpoint{2.084780in}{2.271757in}}%
\pgfpathlineto{\pgfqpoint{2.088345in}{2.301225in}}%
\pgfpathlineto{\pgfqpoint{2.057877in}{2.302446in}}%
\pgfpathlineto{\pgfqpoint{2.077876in}{2.273442in}}%
\pgfpathlineto{\pgfqpoint{2.082338in}{2.283631in}}%
\pgfpathlineto{\pgfqpoint{2.026852in}{2.234986in}}%
\pgfpathlineto{\pgfqpoint{2.048726in}{2.256576in}}%
\pgfpathlineto{\pgfqpoint{2.019285in}{2.237304in}}%
\pgfpathlineto{\pgfqpoint{2.029544in}{2.251282in}}%
\pgfpathlineto{\pgfqpoint{1.998029in}{2.297712in}}%
\pgfpathlineto{\pgfqpoint{2.015629in}{2.321228in}}%
\pgfpathlineto{\pgfqpoint{2.005582in}{2.315332in}}%
\pgfpathlineto{\pgfqpoint{2.102690in}{2.326965in}}%
\pgfpathlineto{\pgfqpoint{2.109936in}{2.324054in}}%
\pgfpathlineto{\pgfqpoint{2.105331in}{2.367841in}}%
\pgfpathlineto{\pgfqpoint{2.112202in}{2.381016in}}%
\pgfpathlineto{\pgfqpoint{2.050276in}{2.435263in}}%
\pgfpathlineto{\pgfqpoint{2.053058in}{2.444049in}}%
\pgfpathlineto{\pgfqpoint{2.068672in}{2.490439in}}%
\pgfpathlineto{\pgfqpoint{2.119185in}{2.480435in}}%
\pgfpathlineto{\pgfqpoint{2.151475in}{2.480082in}}%
\pgfpathlineto{\pgfqpoint{2.114597in}{2.451982in}}%
\pgfpathlineto{\pgfqpoint{2.202856in}{2.412431in}}%
\pgfpathlineto{\pgfqpoint{2.211110in}{2.446311in}}%
\pgfpathlineto{\pgfqpoint{2.220191in}{2.365372in}}%
\pgfpathlineto{\pgfqpoint{2.152446in}{2.333124in}}%
\pgfpathlineto{\pgfqpoint{2.249921in}{2.304482in}}%
\pgfpathlineto{\pgfqpoint{2.299108in}{2.371004in}}%
\pgfpathlineto{\pgfqpoint{2.232676in}{2.388341in}}%
\pgfpathlineto{\pgfqpoint{2.167988in}{2.317581in}}%
\pgfpathlineto{\pgfqpoint{2.206705in}{2.264872in}}%
\pgfpathlineto{\pgfqpoint{2.182966in}{2.285753in}}%
\pgfpathlineto{\pgfqpoint{2.197547in}{2.259291in}}%
\pgfpathlineto{\pgfqpoint{2.167185in}{2.244320in}}%
\pgfpathlineto{\pgfqpoint{2.112462in}{2.258865in}}%
\pgfpathlineto{\pgfqpoint{2.096054in}{2.280275in}}%
\pgfpathlineto{\pgfqpoint{2.158653in}{2.280338in}}%
\pgfpathlineto{\pgfqpoint{2.164871in}{2.287313in}}%
\pgfpathlineto{\pgfqpoint{2.096957in}{2.232328in}}%
\pgfpathlineto{\pgfqpoint{2.118128in}{2.207399in}}%
\pgfpathlineto{\pgfqpoint{2.069266in}{2.218127in}}%
\pgfpathlineto{\pgfqpoint{2.020203in}{2.229064in}}%
\pgfpathlineto{\pgfqpoint{1.948877in}{2.184824in}}%
\pgfpathlineto{\pgfqpoint{1.938108in}{2.132929in}}%
\pgfpathlineto{\pgfqpoint{1.878210in}{2.182524in}}%
\pgfpathlineto{\pgfqpoint{1.854203in}{2.132554in}}%
\pgfpathlineto{\pgfqpoint{1.820923in}{2.148749in}}%
\pgfpathlineto{\pgfqpoint{1.891049in}{2.113453in}}%
\pgfpathlineto{\pgfqpoint{1.895409in}{2.092145in}}%
\pgfpathlineto{\pgfqpoint{1.835793in}{2.152120in}}%
\pgfpathlineto{\pgfqpoint{1.747945in}{2.139590in}}%
\pgfpathlineto{\pgfqpoint{1.745881in}{2.123873in}}%
\pgfpathlineto{\pgfqpoint{1.699237in}{2.131903in}}%
\pgfpathlineto{\pgfqpoint{1.735983in}{2.093802in}}%
\pgfpathlineto{\pgfqpoint{1.757577in}{2.032462in}}%
\pgfpathlineto{\pgfqpoint{1.740255in}{2.055893in}}%
\pgfpathlineto{\pgfqpoint{1.773509in}{2.056585in}}%
\pgfpathlineto{\pgfqpoint{1.790208in}{1.964778in}}%
\pgfpathlineto{\pgfqpoint{1.852366in}{1.978786in}}%
\pgfpathlineto{\pgfqpoint{1.871524in}{1.974095in}}%
\pgfpathlineto{\pgfqpoint{1.892550in}{2.043189in}}%
\pgfpathlineto{\pgfqpoint{1.874104in}{2.082802in}}%
\pgfpathlineto{\pgfqpoint{1.872257in}{2.089242in}}%
\pgfpathlineto{\pgfqpoint{1.873583in}{2.074391in}}%
\pgfpathlineto{\pgfqpoint{1.893878in}{2.112354in}}%
\pgfpathlineto{\pgfqpoint{1.843088in}{2.182851in}}%
\pgfpathlineto{\pgfqpoint{1.826750in}{2.163473in}}%
\pgfpathlineto{\pgfqpoint{1.821943in}{2.059427in}}%
\pgfpathlineto{\pgfqpoint{1.845491in}{2.054871in}}%
\pgfpathlineto{\pgfqpoint{1.826836in}{2.089609in}}%
\pgfpathlineto{\pgfqpoint{1.839619in}{2.083362in}}%
\pgfpathlineto{\pgfqpoint{1.783260in}{2.043804in}}%
\pgfpathlineto{\pgfqpoint{1.810044in}{2.018756in}}%
\pgfpathlineto{\pgfqpoint{1.761979in}{1.998149in}}%
\pgfpathlineto{\pgfqpoint{1.755433in}{1.914354in}}%
\pgfpathlineto{\pgfqpoint{1.719131in}{1.873553in}}%
\pgfpathlineto{\pgfqpoint{1.738491in}{1.860441in}}%
\pgfpathlineto{\pgfqpoint{1.668179in}{1.890396in}}%
\pgfpathlineto{\pgfqpoint{1.647122in}{1.893731in}}%
\pgfpathlineto{\pgfqpoint{1.687877in}{1.891891in}}%
\pgfpathlineto{\pgfqpoint{1.735681in}{1.922419in}}%
\pgfpathlineto{\pgfqpoint{1.678236in}{1.926445in}}%
\pgfpathlineto{\pgfqpoint{1.700121in}{1.935088in}}%
\pgfpathlineto{\pgfqpoint{1.686829in}{1.927396in}}%
\pgfpathlineto{\pgfqpoint{1.645878in}{1.942350in}}%
\pgfpathlineto{\pgfqpoint{1.620640in}{1.910649in}}%
\pgfpathlineto{\pgfqpoint{1.666470in}{1.846846in}}%
\pgfpathlineto{\pgfqpoint{1.755274in}{1.869356in}}%
\pgfpathlineto{\pgfqpoint{1.746684in}{1.871510in}}%
\pgfpathlineto{\pgfqpoint{1.743088in}{1.939957in}}%
\pgfpathlineto{\pgfqpoint{1.715987in}{2.022147in}}%
\pgfpathlineto{\pgfqpoint{1.787280in}{1.954443in}}%
\pgfpathlineto{\pgfqpoint{1.758108in}{1.948511in}}%
\pgfpathlineto{\pgfqpoint{1.768289in}{2.032825in}}%
\pgfpathlineto{\pgfqpoint{1.764293in}{2.003014in}}%
\pgfpathlineto{\pgfqpoint{1.737464in}{1.991344in}}%
\pgfpathlineto{\pgfqpoint{1.790347in}{1.995528in}}%
\pgfpathlineto{\pgfqpoint{1.858561in}{2.017670in}}%
\pgfpathlineto{\pgfqpoint{1.887430in}{2.020789in}}%
\pgfpathlineto{\pgfqpoint{1.911576in}{2.036591in}}%
\pgfpathlineto{\pgfqpoint{1.847711in}{2.063958in}}%
\pgfpathlineto{\pgfqpoint{1.820520in}{2.086287in}}%
\pgfpathlineto{\pgfqpoint{1.800373in}{2.099155in}}%
\pgfpathlineto{\pgfqpoint{1.808721in}{2.099290in}}%
\pgfpathlineto{\pgfqpoint{1.861039in}{2.022808in}}%
\pgfpathlineto{\pgfqpoint{1.796077in}{2.030259in}}%
\pgfpathlineto{\pgfqpoint{1.829895in}{2.034129in}}%
\pgfpathlineto{\pgfqpoint{1.843583in}{1.962033in}}%
\pgfpathlineto{\pgfqpoint{1.818403in}{1.981999in}}%
\pgfpathlineto{\pgfqpoint{1.768250in}{2.035569in}}%
\pgfpathlineto{\pgfqpoint{1.762934in}{2.074456in}}%
\pgfpathlineto{\pgfqpoint{1.722080in}{2.093823in}}%
\pgfpathlineto{\pgfqpoint{1.658892in}{2.051671in}}%
\pgfpathlineto{\pgfqpoint{1.697400in}{2.027503in}}%
\pgfpathlineto{\pgfqpoint{1.719780in}{1.999319in}}%
\pgfpathlineto{\pgfqpoint{1.764868in}{1.958766in}}%
\pgfpathlineto{\pgfqpoint{1.742307in}{2.080292in}}%
\pgfpathlineto{\pgfqpoint{1.720040in}{2.078063in}}%
\pgfpathlineto{\pgfqpoint{1.748220in}{2.013826in}}%
\pgfpathlineto{\pgfqpoint{1.765562in}{2.005533in}}%
\pgfpathlineto{\pgfqpoint{1.713576in}{2.011212in}}%
\pgfpathlineto{\pgfqpoint{1.630386in}{2.007531in}}%
\pgfpathlineto{\pgfqpoint{1.640414in}{2.037260in}}%
\pgfpathlineto{\pgfqpoint{1.637136in}{2.061439in}}%
\pgfpathlineto{\pgfqpoint{1.608398in}{2.053412in}}%
\pgfpathlineto{\pgfqpoint{1.598924in}{2.033657in}}%
\pgfpathlineto{\pgfqpoint{1.610531in}{2.095946in}}%
\pgfpathlineto{\pgfqpoint{1.669444in}{2.135884in}}%
\pgfpathlineto{\pgfqpoint{1.659617in}{2.116929in}}%
\pgfpathlineto{\pgfqpoint{1.559544in}{2.146907in}}%
\pgfpathlineto{\pgfqpoint{1.511911in}{2.143890in}}%
\pgfpathlineto{\pgfqpoint{1.511911in}{2.143890in}}%
\pgfusepath{stroke}%
\end{pgfscope}%
\end{pgfpicture}%
\makeatother%
\endgroup%

    \caption{Three-dimensional Wiener process.}
    \label{fig:threeDimensionalBrownianMotionRealisations}
\end{figure}

\begin{thm}{Time-reversal and Time-shifting}{}
If $(W_s, s \in [0, t])$ is a Wiener process on $[0, t]$, then so is the time-reversed process $(\widehat{W}_s, s \in [0, t])$ defined by $\widehat{W}_s \coloneqq W_{t-s} - W_t$. Similarly, if $(W_s, s \ge 0)$ is a Wiener process, then for any $t \ge 0$ the process $(\widetilde{W}_s, s \ge 0)$ defined by $\widetilde{W}_s \coloneqq W_{t+s} - W_t$ is also a Wiener process.
\end{thm}

\begin{thm}{Scaling}{}
If $(W_t)$ is a Wiener process, then so is $(X_t)$, with $X_t \coloneqq W_{at}/\sqrt{a}$, $t \ge 0$ for any $a > 0$.
\end{thm}


\begin{lem}{Probability Bound on the Supremum}{}
For a Wiener process $(W_t)$, it holds for every $\varepsilon > 0$ that
\begin{equation*}
\Prob \left( \sup_{0 \le t \le 1} |W_t| > \varepsilon \right) \le 2 \e^{-\varepsilon^2/2}.
\end{equation*}
\end{lem}


\begin{thm}{Reciprocal Time}{}
If $(W_t)$ is a Wiener process, then so is $(X_t)$, with $X_t \coloneqq t W_{1/t}, \, t > 0, \, X_0 \coloneqq 0$.
\end{thm}

\begin{thm}{Existence of the Wiener Process}{}
There exists a probability space $(\Omega, \cH, \Prob)$ and a zero-mean Gaussian process $W \coloneqq (W_t, t \in [0, 1])$ with covariance function $\E W_s W_t = s \wedge t$ and uniformly continuous sample paths.
\end{thm}

\begin{thm}{Right-continuous Filtration}{}
For the Wiener process, the augmentation $(\xbar{\cF}_t)$ of the natural filtration is right-continuous; that is,
\begin{equation*}
\xbar{\cF}_t = \xbar{\cF}_t^+ \coloneqq \bigcap_{\varepsilon > 0} \xbar{\cF}_{t+\varepsilon}.
\end{equation*}
\end{thm}

\begin{thm}{Wiener Process is Markov w.r.t $\cF^+$}{}
Let $W$ be a Wiener process. For every $t \ge 0$, the process $(W_{t+u} - W_t, u \ge 0)$ is independent of $\cF_t^+$. Consequently, $W$ is Markovian with respect to $\cF^+$.
\end{thm}

\begin{thm}{Strong Markov Property}{}
For any finite stopping time $T$ with respect to $\cF^+$, the process $(W_{T+u} - W_T, u \ge 0)$ is a Wiener process independent of $\cF_T^+$.
\end{thm}

\begin{thm}{Reflection Principle}{}
Let $T$ be a stopping time. Then, $(\tilde{W}_t, t \ge 0)$, defined by
\begin{equation}
    \tilde{W}_t \coloneqq W_t \I_{\{t \le T\}} + (2W_T - W_t) \I_{\{t > T\}}, \quad t \ge 0,
    \label{eq:ReflectionProcess}
\end{equation}
is a Wiener process.
\end{thm}

\begin{thm}{Transition Kernel of the Wiener Process}{}
The Wiener process $W$ admits the transition function
\begin{equation*}
    P_{t,u}(x, A) = \Prob(W_u \in A \mid W_t = x) = \int_A \di y \, p_{u-t}(y \mid x),
\end{equation*}
where the \textit{transition density} $p_t(y \mid x)$ is given by
\begin{equation}
    p_t(y \mid x) \coloneqq \frac{1}{\sqrt{2 \pi t}} \e^{-\frac{(y-x)^2}{2t}}, \quad t \ge 0, \quad x,y \in \R.
    \label{eq:WienerTransitionDensity}
\end{equation}
\end{thm}

\begin{thm}{Exponential Martingale}{}
A continuous process $X \coloneqq (X_t, t \ge 0)$ is a Wiener process if and only if for each $r \in \R$ the process $S \coloneqq (S_t)$ defined by
\begin{equation}
    S_t \coloneqq \e^{r X_t - \frac{1}{2}r^2 t}, \quad t \ge 0,
    \label{eq:ExponentialMartingaleProcess}
\end{equation}
is a martingale.
\end{thm}

\begin{thm}{Martingales from Functions of a Brownian Motion}{}
Let $f : \R^d \to \R$ be twice continuously differentiable and let $\v{B}$ be $d$-dimensional standard Brownian motion. If for all $t > 0$ and $\v{x} \in \R^d$,
\begin{equation*}
    \E^{\v{x}} |f(\v{B}_t)| < \infty \quad \text{and} \quad \E^{\v{x}} \int_0^t |\Delta f(\v{B}_s)| \di s < \infty,
\end{equation*}
then the process $(X_t, t \ge 0)$ defined by
\begin{equation}
    X_t \coloneqq f(\v{B}_t) - \frac{1}{2} \int_0^t \Delta f(\v{B}_s) \di s
    \label{eq:MartingaleFunctionProcess}
\end{equation}
is a martingale.
\end{thm}